\ifdefined\ReaderVersion
  \documentclass[pdflatex,sn-vancouver-num,iicol]{sn-jnl}
\else
  \documentclass[pdflatex,sn-vancouver-num]{sn-jnl}
\fi

\usepackage{amsmath,amssymb,mathtools}
\usepackage{array,booktabs,multirow,tabularx}
\usepackage{graphicx,microtype,balance}
\usepackage{tikz,pgfplots}
\usepgfplotslibrary{groupplots}
\pgfplotsset{compat=1.18}
\usetikzlibrary{arrows.meta}
\usepackage[scaled=0.92]{helvet}
\usepackage{sansmath}
\tikzset{every picture/.append style={execute at begin picture={\sansmath}}}
\newcommand{\KL}{\operatorname{KL}}
\usepackage{needspace}
\usepackage{xurl}
\ifdefined\ReaderVersion\else
  \usepackage[switch]{lineno}
  \usepackage{setspace}
\fi
\graphicspath{{figures/}}
\renewcommand{\dbltopfraction}{0.9}
\renewcommand{\topfraction}{0.9}
\renewcommand{\textfraction}{0.08}
\renewcommand{\dblfloatpagefraction}{0.85}
\renewcommand{\floatpagefraction}{0.85}
\hypersetup{
  pdftitle={Using unpaired measurements to reduce the number of paired cells needed to estimate dependence between modalities},
  pdfauthor={}
}
% Machine-checked values from versioned result artifacts.
\newcommand{\SimReps}{128}
\newcommand{\SimHetSameHier}{0.01493}
\newcommand{\SimHetSameCommon}{0.01521}
\newcommand{\SimHetSameCommonReduction}{1.85\%}
\newcommand{\SimHetSameCommonReductionLo}{1.52\%}
\newcommand{\SimHetSameCommonReductionHi}{2.18\%}
\newcommand{\SimHetShiftHier}{0.01159}
\newcommand{\SimHetShiftCommon}{0.01171}
\newcommand{\SimHetShiftResidual}{0.03359}
\newcommand{\SimHetShiftResidualReduction}{65.49\%}
\newcommand{\SimHetShiftResidualReductionLo}{63.02\%}
\newcommand{\SimHetShiftResidualReductionHi}{67.72\%}
\newcommand{\SimHomSameHier}{0.008752}
\newcommand{\SimHomSameCommon}{0.008660}
\newcommand{\SimHomSameHierPenalty}{1.06\%}
\newcommand{\SimHomShiftHier}{0.008674}
\newcommand{\SimHomShiftCommon}{0.008640}
\newcommand{\SimHomShiftResidual}{0.02943}
\newcommand{\SimHomShiftResidualReduction}{70.52\%}
\newcommand{\SimHomShiftResidualReductionLo}{69.21\%}
\newcommand{\SimHomShiftResidualReductionHi}{71.74\%}
\newcommand{\RiskStressConditions}{48}
\newcommand{\RiskStressEnvelopePass}{48/48}
\newcommand{\RiskStressCoverage}{48/48}
\newcommand{\RiskStressIntegralError}{1.91\times10^{-14}}
\newcommand{\RiskStressMaxZ}{1.833}
\newcommand{\PSciN}{85}
\newcommand{\PSciR}{0.525}
\newcommand{\PSciRLo}{0.483}
\newcommand{\PSciRHi}{0.566}
\newcommand{\PSciRMSE}{0.844}
\newcommand{\PSciRMSELo}{0.819}
\newcommand{\PSciRMSEHi}{0.868}
\newcommand{\PSciDirectR}{0.429}
\newcommand{\PSciDestroyedR}{0.074}
\newcommand{\PSciFourStateR}{0.640}
\newcommand{\PSciNeighbor}{0.064}
\newcommand{\PSciNeighborLo}{0.050}
\newcommand{\PSciNeighborHi}{0.079}
\newcommand{\PSciDestroyedNeighbor}{0.051}
\newcommand{\PSciDestroyedNeighborLo}{0.041}
\newcommand{\PSciDestroyedNeighborHi}{0.063}
\newcommand{\PSciNeighborDiff}{0.013}
\newcommand{\PSciNeighborDiffLo}{-0.003}
\newcommand{\PSciNeighborDiffHi}{0.030}
\newcommand{\PSciNeighborP}{0.283}
\newcommand{\FrangiehN}{179}
\newcommand{\FrangiehR}{0.040}
\newcommand{\FrangiehRLo}{0.013}
\newcommand{\FrangiehRHi}{0.066}
\newcommand{\FrangiehRMSE}{1.009}
\newcommand{\PapalexiN}{24}
\newcommand{\PapalexiR}{-0.025}
\newcommand{\PapalexiRLo}{-0.150}
\newcommand{\PapalexiRHi}{0.096}
\newcommand{\PapalexiRMSE}{1.037}
\newcommand{\MultiN}{35}
\newcommand{\MultiR}{-0.224}
\newcommand{\MultiRLo}{-0.330}
\newcommand{\MultiRHi}{-0.124}
\newcommand{\MultiRMSE}{1.047}
\newcommand{\PerturbFateN}{49}
\newcommand{\PerturbFateR}{-0.092}
\newcommand{\PerturbFateRLo}{-0.152}
\newcommand{\PerturbFateRHi}{-0.035}
\newcommand{\PerturbFateRMSE}{1.023}
\newcommand{\ReSisArmP}{1/65}
\newcommand{\ReSisArmQ}{0.015}
\newcommand{\ReSisNKCos}{0.952}
\newcommand{\ReSisNKLo}{-0.801}
\newcommand{\ReSisNKHi}{0.968}
\newcommand{\ArceN}{28}
\newcommand{\ArceR}{0.412}
\newcommand{\ArceRLo}{0.331}
\newcommand{\ArceRHi}{0.494}
\newcommand{\ArceRMSE}{1.093}
\newcommand{\ArceRMSELo}{1.012}
\newcommand{\ArceRMSEHi}{1.185}
\newcommand{\ArceDecision}{refuse}
\newcommand{\StressNull}{3.0\%}
\newcommand{\StressPower}{96.5\%}
\newcommand{\StressHeavyNull}{5.5\%}
\newcommand{\StressGuidePower}{95.5\%}
\newcommand{\StressObsPower}{90.5\%}
\newcommand{\StephDevN}{36}
\newcommand{\StephHeldN}{56}
\newcommand{\StephPrimaryLoss}{0.01220}
\newcommand{\StephResidualLoss}{0.01478}
\newcommand{\StephResidualReduction}{17.46\%}
\newcommand{\StephResidualCILo}{-0.00413}
\newcommand{\StephResidualCIHi}{-0.00080}
\newcommand{\StephResidualFav}{50/56}
\newcommand{\StephSignP}{5.09\times10^{-10}}
\newcommand{\StephDestroyedLoss}{0.02274}
\newcommand{\StephDestroyedReduction}{46.36\%}
\newcommand{\StephDestroyedCILo}{-0.01323}
\newcommand{\StephDestroyedCIHi}{-0.00749}
\newcommand{\StephDestroyedFav}{51/56}
\newcommand{\StephRepeatedDestroyedN}{29}
\newcommand{\StephRepeatedDestroyedReduction}{47.30\%}
\newcommand{\StephRepeatedDestroyedReductionLo}{37.56\%}
\newcommand{\StephRepeatedDestroyedReductionHi}{54.59\%}
\newcommand{\StephRepeatedDestroyedDiffLo}{-0.01367}
\newcommand{\StephRepeatedDestroyedDiffHi}{-0.00795}
\newcommand{\StephRepeatedDestroyedFav}{51/56}
\newcommand{\StephRepeatedDestroyedP}{5.85\times10^{-11}}
\newcommand{\StephTLineageReduction}{-3.41\%}
\newcommand{\StephBLineageReduction}{-0.78\%}
\newcommand{\StephMyeloidReduction}{15.51\%}
\newcommand{\StephMyeloidReductionLo}{5.51\%}
\newcommand{\StephMyeloidReductionHi}{24.61\%}
\newcommand{\StephMyeloidFav}{21/46}
\newcommand{\StephMyeloidSignP}{0.769}
\newcommand{\StephCommonExactLoss}{0.01300}
\newcommand{\StephCommonExactReduction}{6.18\%}
\newcommand{\StephCommonExactReductionLo}{3.84\%}
\newcommand{\StephCommonExactReductionHi}{8.47\%}
\newcommand{\StephCommonExactCILo}{-0.00121}
\newcommand{\StephCommonExactCIHi}{-0.00046}
\newcommand{\StephCommonExactFav}{46/56}
\newcommand{\StephCommonExactSignP}{6.23\times10^{-7}}
\newcommand{\StephPooledPoissonLoss}{0.01263}
\newcommand{\StephPooledPoissonReduction}{3.39\%}
\newcommand{\StephPooledPoissonReductionLo}{1.42\%}
\newcommand{\StephPooledPoissonReductionHi}{5.37\%}
\newcommand{\StephPooledPoissonCILo}{-0.00075}
\newcommand{\StephPooledPoissonCIHi}{-0.00017}
\newcommand{\StephPooledPoissonFav}{38/56}
\newcommand{\StephPooledPoissonSignP}{0.00523}
\newcommand{\GSEDevN}{15}
\newcommand{\GSEHeldN}{9}
\newcommand{\GSEPrimaryLoss}{0.00851}
\newcommand{\GSEResidualLoss}{0.01413}
\newcommand{\GSEResidualReduction}{39.81\%}
\newcommand{\GSEResidualCILo}{-0.00710}
\newcommand{\GSEResidualCIHi}{-0.00423}
\newcommand{\GSEResidualFav}{9/9}
\newcommand{\GSESignP}{0.00195}
\newcommand{\GSEDestroyedLoss}{0.04029}
\newcommand{\GSEDestroyedReduction}{78.89\%}
\newcommand{\GSEDestroyedCILo}{-0.03567}
\newcommand{\GSEDestroyedCIHi}{-0.02816}
\newcommand{\GSEDestroyedFav}{9/9}
\newcommand{\GSECommonExactLoss}{0.00832}
\newcommand{\GSECommonExactRelative}{-2.20\%}
\newcommand{\GSECommonExactAdvantage}{2.20\%}
\newcommand{\GSECommonExactRelativeLo}{-3.03\%}
\newcommand{\GSECommonExactRelativeHi}{-1.34\%}
\newcommand{\GSECommonExactCILo}{0.00010}
\newcommand{\GSECommonExactCIHi}{0.00028}
\newcommand{\GSECommonExactFav}{0/9}
\newcommand{\GSECommonExactSignP}{1.000}
\newcommand{\GSEPooledPoissonLoss}{0.00996}
\newcommand{\GSEPooledPoissonReduction}{14.58\%}
\newcommand{\GSEPooledPoissonReductionLo}{11.94\%}
\newcommand{\GSEPooledPoissonReductionHi}{17.16\%}
\newcommand{\GSEPooledPoissonCILo}{-0.00181}
\newcommand{\GSEPooledPoissonCIHi}{-0.00115}
\newcommand{\GSEPooledPoissonFav}{9/9}
\newcommand{\GSEPooledPoissonSignP}{0.00195}
\newcommand{\GSEStandardPoissonLoss}{0.00998}
\newcommand{\GSEStandardPoissonReduction}{14.79\%}
\newcommand{\GSEStandardPoissonReductionLo}{12.18\%}
\newcommand{\GSEStandardPoissonReductionHi}{17.37\%}
\newcommand{\GSEStandardPoissonCILo}{-0.00184}
\newcommand{\GSEStandardPoissonCIHi}{-0.00117}
\newcommand{\GSEStandardPoissonFav}{9/9}
\newcommand{\GSEStandardPoissonSignP}{0.00195}
\newcommand{\CombatPrimaryLoss}{0.01156}
\newcommand{\CombatResidualLoss}{0.01561}
\newcommand{\CombatReduction}{25.94\%}
\newcommand{\CombatCILo}{-0.00597}
\newcommand{\CombatCIHi}{-0.00225}
\newcommand{\CombatFav}{17/24}
\newcommand{\BMMCPrimaryLoss}{0.01085}
\newcommand{\BMMCResidualLoss}{0.01315}
\newcommand{\BMMCReduction}{17.47\%}
\newcommand{\BMMCCILo}{-0.00504}
\newcommand{\BMMCCIHi}{-0.00038}
\newcommand{\BMMCFav}{4/4}
\newcommand{\GSEThreeOneFourPilotN}{20}
\newcommand{\GSEThreeOneFourPrimaryLoss}{0.001927}
\newcommand{\GSEThreeOneFourResidualLoss}{0.001936}
\newcommand{\GSEThreeOneFourResidualReduction}{0.49\%}
\newcommand{\GSEThreeOneFourResidualCILo}{-0.0000182}
\newcommand{\GSEThreeOneFourResidualCIHi}{-0.0000011}
\newcommand{\GSEThreeOneFourResidualFav}{13/20}
\newcommand{\GSEThreeOneFourDestroyedLoss}{0.001936}
\newcommand{\GSEThreeOneFourDestroyedReduction}{0.45\%}
\newcommand{\GSEThreeOneFourDestroyedCILo}{-0.0000390}
\newcommand{\GSEThreeOneFourDestroyedCIHi}{0.0000223}
\newcommand{\GSEThreeOneFourDestroyedFav}{11/20}
\newcommand{\GSEThreeOneFourBroadReduction}{2.54\%}
\newcommand{\GSEThreeOneSevenCalibrationN}{7}
\newcommand{\GSEThreeOneSevenVisitN}{28}
\newcommand{\GSEThreeOneSevenPrimaryLoss}{0.006530}
\newcommand{\GSEThreeOneSevenPoissonLoss}{0.006609}
\newcommand{\GSEThreeOneSevenPoissonReduction}{1.19\%}
\newcommand{\GSEThreeOneSevenPoissonFav}{6/7}
\newcommand{\GSEThreeOneSevenTimeFav}{4/4}
\newcommand{\GSEThreeOneSevenDestroyedLoss}{0.007904}
\newcommand{\GSEThreeOneSevenDestroyedReduction}{17.39\%}
\newcommand{\GSEThreeOneSevenDestroyedFav}{7/7}

% Generated by extension/identified_set/make_manuscript_assets.py. Do not edit.
\newcommand{\IdRecipientQueries}{8{,}586}
\newcommand{\IdUnidentifiedCount}{8{,}586}
\newcommand{\IdUnidentifiedPct}{100.0}
\newcommand{\IdUnidCognatePct}{100.0}
\newcommand{\IdCertifiedEndpoints}{17{,}172}
\newcommand{\IdCertMaxDiff}{4.1\times10^{-11}}
\newcommand{\IdScalingFailures}{0}
\newcommand{\IdUnidCam}{100.0}
\newcommand{\IdUnidCognateCam}{100.0}
\newcommand{\IdWidthCam}{14.8}
\newcommand{\IdWidthLoCam}{8.9}
\newcommand{\IdWidthHiCam}{22.9}
\newcommand{\IdWidthMinCam}{5.0}
\newcommand{\IdRiskCam}{0.266}
\newcommand{\IdHeldFreqCam}{0.0223}
\newcommand{\IdHeldFullCam}{0.0186}
\newcommand{\IdSpearRiskCam}{0.03}
\newcommand{\IdSpearProdCam}{0.03}
\newcommand{\IdRevCam}{374}
\newcommand{\IdRevPctCam}{19.2}
\newcommand{\IdRevTiesCam}{8}
\newcommand{\IdAgreeCam}{59.6}
\newcommand{\IdAgreeNCam}{218/366}
\newcommand{\IdAgreeLoCam}{55.0}
\newcommand{\IdAgreeHiCam}{64.2}
\newcommand{\IdAgreeAdjLoCam}{53.8}
\newcommand{\IdAgreeAdjHiCam}{65.5}
\newcommand{\IdNonemptyCam}{59.5}
\newcommand{\IdNonemptyAdjLoCam}{54.1}
\newcommand{\IdRevGainShareCam}{18.8}
\newcommand{\IdUnidNew}{100.0}
\newcommand{\IdUnidCognateNew}{100.0}
\newcommand{\IdWidthNew}{15.1}
\newcommand{\IdWidthLoNew}{11.1}
\newcommand{\IdWidthHiNew}{19.7}
\newcommand{\IdWidthMinNew}{5.3}
\newcommand{\IdRiskNew}{0.319}
\newcommand{\IdHeldFreqNew}{0.0323}
\newcommand{\IdHeldFullNew}{0.0229}
\newcommand{\IdSpearRiskNew}{0.08}
\newcommand{\IdSpearProdNew}{0.09}
\newcommand{\IdRevNew}{1{,}029}
\newcommand{\IdRevPctNew}{22.7}
\newcommand{\IdRevTiesNew}{65}
\newcommand{\IdAgreeNew}{75.3}
\newcommand{\IdAgreeNNew}{726/964}
\newcommand{\IdAgreeLoNew}{71.8}
\newcommand{\IdAgreeHiNew}{78.7}
\newcommand{\IdAgreeAdjLoNew}{70.9}
\newcommand{\IdAgreeAdjHiNew}{79.7}
\newcommand{\IdNonemptyNew}{75.9}
\newcommand{\IdNonemptyAdjLoNew}{71.6}
\newcommand{\IdRevGainShareNew}{37.3}
\newcommand{\IdUnidTall}{100.0}
\newcommand{\IdUnidCognateTall}{100.0}
\newcommand{\IdWidthTall}{16.4}
\newcommand{\IdWidthLoTall}{12.2}
\newcommand{\IdWidthHiTall}{22.2}
\newcommand{\IdWidthMinTall}{5.6}
\newcommand{\IdRiskTall}{0.370}
\newcommand{\IdHeldFreqTall}{0.0585}
\newcommand{\IdHeldFullTall}{0.0260}
\newcommand{\IdSpearRiskTall}{0.04}
\newcommand{\IdSpearProdTall}{0.03}
\newcommand{\IdRevTall}{344}
\newcommand{\IdRevPctTall}{28.3}
\newcommand{\IdRevTiesTall}{3}
\newcommand{\IdAgreeTall}{84.5}
\newcommand{\IdAgreeNTall}{288/341}
\newcommand{\IdAgreeLoTall}{80.6}
\newcommand{\IdAgreeHiTall}{88.1}
\newcommand{\IdAgreeAdjLoTall}{79.5}
\newcommand{\IdAgreeAdjHiTall}{89.0}
\newcommand{\IdNonemptyTall}{83.3}
\newcommand{\IdNonemptyAdjLoTall}{77.7}
\newcommand{\IdRevGainShareTall}{37.5}
\newcommand{\IdUnidColon}{100.0}
\newcommand{\IdUnidCognateColon}{100.0}
\newcommand{\IdWidthColon}{12.6}
\newcommand{\IdWidthLoColon}{9.6}
\newcommand{\IdWidthHiColon}{19.5}
\newcommand{\IdWidthMinColon}{5.6}
\newcommand{\IdRiskColon}{0.265}
\newcommand{\IdHeldFreqColon}{0.0259}
\newcommand{\IdHeldFullColon}{0.0218}
\newcommand{\IdSpearRiskColon}{0.11}
\newcommand{\IdSpearProdColon}{0.10}
\newcommand{\IdRevColon}{201}
\newcommand{\IdRevPctColon}{22.6}
\newcommand{\IdRevTiesColon}{4}
\newcommand{\IdAgreeColon}{68.0}
\newcommand{\IdAgreeNColon}{134/197}
\newcommand{\IdAgreeLoColon}{59.1}
\newcommand{\IdAgreeHiColon}{76.3}
\newcommand{\IdAgreeAdjLoColon}{56.7}
\newcommand{\IdAgreeAdjHiColon}{78.7}
\newcommand{\IdNonemptyColon}{69.2}
\newcommand{\IdNonemptyAdjLoColon}{57.8}
\newcommand{\IdRevGainShareColon}{67.3}
\newcommand{\IdAgreeRetroMin}{60}
\newcommand{\IdAgreeRetroMax}{84}
\newcommand{\IdRevPctMin}{19}
\newcommand{\IdRevPctMax}{28}
\newcommand{\IdAgreeColonRound}{68}
\newcommand{\IdFreeExtLowMedian}{0.89}
\newcommand{\IdFreeExtHighMedian}{1.16}
\newcommand{\IdFreeExtMax}{5.40}
\newcommand{\IdFreeStarWidth}{10.6}
\newcommand{\IdFreeWidth}{12.5}

% Generated by extension/spectral_transfer/make_scale_assets.py. Do not edit.
\newcommand{\SoPairCam}{91.1}
\newcommand{\SoPairLoCam}{88.9}
\newcommand{\SoPairHiCam}{93.0}
\newcommand{\SoClosedCam}{75.3}
\newcommand{\SoClosedLoCam}{59.7}
\newcommand{\SoClosedHiCam}{86.4}
\newcommand{\SoFirstCam}{\ensuremath{-}202}
\newcommand{\SoPairNew}{98.0}
\newcommand{\SoPairLoNew}{97.1}
\newcommand{\SoPairHiNew}{98.8}
\newcommand{\SoClosedNew}{72.2}
\newcommand{\SoClosedLoNew}{62.9}
\newcommand{\SoClosedHiNew}{80.0}
\newcommand{\SoFirstNew}{\ensuremath{-}314}
\newcommand{\SoPairTall}{96.2}
\newcommand{\SoPairLoTall}{94.6}
\newcommand{\SoPairHiTall}{97.4}
\newcommand{\SoClosedTall}{88.8}
\newcommand{\SoClosedLoTall}{83.6}
\newcommand{\SoClosedHiTall}{92.8}
\newcommand{\SoFirstTall}{\ensuremath{-}246}
\newcommand{\SoPairColon}{98.0}
\newcommand{\SoPairLoColon}{96.2}
\newcommand{\SoPairHiColon}{99.8}
\newcommand{\SoClosedColon}{66.4}
\newcommand{\SoClosedLoColon}{47.2}
\newcommand{\SoClosedHiColon}{83.2}
\newcommand{\SoFirstColon}{\ensuremath{-}15}
\newcommand{\SoClipped}{116}
\newcommand{\SoPairMin}{91}
\newcommand{\SoPairMax}{98}
\newcommand{\SoClosedMin}{66}
\newcommand{\SoClosedMax}{89}
\newcommand{\ScDonors}{12}
\newcommand{\ScGenes}{209}
\newcommand{\ScProteins}{177}
\newcommand{\ScPairs}{36{,}993}
\newcommand{\ScErrIndep}{1.000}
\newcommand{\ScCorrIndep}{0.000}
\newcommand{\ScRsqIndep}{0.000}
\newcommand{\ScErrVar}{0.936}
\newcommand{\ScCorrVar}{0.438}
\newcommand{\ScRsqVar}{\ensuremath{-}0.393}
\newcommand{\ScErrRefCorr}{0.477}
\newcommand{\ScCorrRefCorr}{0.882}
\newcommand{\ScRsqRefCorr}{\ensuremath{-}6.523}
\newcommand{\ScErrRefCov}{0.669}
\newcommand{\ScCorrRefCov}{0.822}
\newcommand{\ScRsqRefCov}{\ensuremath{-}180}
\newcommand{\ScErrReg}{0.380}
\newcommand{\ScCorrReg}{0.922}
\newcommand{\ScRsqReg}{0.283}
\newcommand{\ScErrCells}{0.371}
\newcommand{\ScCorrCells}{0.927}
\newcommand{\ScRsqCells}{0.160}
\newcommand{\ScErrClosed}{0.341}
\newcommand{\ScCorrClosed}{0.934}
\newcommand{\ScRsqClosed}{0.285}
\newcommand{\ScErrFirst}{678}
\newcommand{\ScCorrFirst}{0.708}
\newcommand{\ScRsqFirst}{\ensuremath{-}125{,}313}
\newcommand{\ScErrOracle}{0.367}
\newcommand{\ScCorrOracle}{0.924}
\newcommand{\ScRsqOracle}{0.018}
\newcommand{\ScErrPfMeans}{0.536}
\newcommand{\ScCorrPfMeans}{0.840}
\newcommand{\ScRsqPfMeans}{0.089}
\newcommand{\ScErrPfLow}{0.460}
\newcommand{\ScCorrPfLow}{0.881}
\newcommand{\ScRsqPfLow}{0.115}
\newcommand{\ScErrPfMeansReg}{0.740}
\newcommand{\ScCorrPfMeansReg}{0.731}
\newcommand{\ScRsqPfMeansReg}{0.014}
\newcommand{\ScErrPfLowReg}{0.515}
\newcommand{\ScCorrPfLowReg}{0.855}
\newcommand{\ScRsqPfLowReg}{0.116}
\newcommand{\ScDiffVar}{\ensuremath{-}0.595}
\newcommand{\ScDiffLoVar}{\ensuremath{-}0.670}
\newcommand{\ScDiffHiVar}{\ensuremath{-}0.497}
\newcommand{\ScWinsVar}{12/12}
\newcommand{\ScDiffReg}{\ensuremath{-}0.039}
\newcommand{\ScDiffLoReg}{\ensuremath{-}0.056}
\newcommand{\ScDiffHiReg}{\ensuremath{-}0.024}
\newcommand{\ScWinsReg}{12/12}
\newcommand{\ScDiffRefCorr}{\ensuremath{-}0.136}
\newcommand{\ScDiffLoRefCorr}{\ensuremath{-}0.186}
\newcommand{\ScDiffHiRefCorr}{\ensuremath{-}0.090}
\newcommand{\ScWinsRefCorr}{12/12}
\newcommand{\ScDiffRefCov}{\ensuremath{-}0.328}
\newcommand{\ScDiffLoRefCov}{\ensuremath{-}0.401}
\newcommand{\ScDiffHiRefCov}{\ensuremath{-}0.256}
\newcommand{\ScWinsRefCov}{12/12}
\newcommand{\ScDiffCells}{0.030}
\newcommand{\ScDiffLoCells}{0.020}
\newcommand{\ScDiffHiCells}{0.040}
\newcommand{\PfDiffClosed}{0.119}
\newcommand{\PfDiffLoClosed}{0.091}
\newcommand{\PfDiffHiClosed}{0.155}
\newcommand{\PfDiffReg}{0.080}
\newcommand{\PfDiffLoReg}{0.042}
\newcommand{\PfDiffHiReg}{0.115}
\newcommand{\PfDiffVar}{\ensuremath{-}0.476}
\newcommand{\PfDiffLoVar}{\ensuremath{-}0.561}
\newcommand{\PfDiffHiVar}{\ensuremath{-}0.381}
\newcommand{\PfDiffIndep}{\ensuremath{-}0.540}
\newcommand{\PfDiffLoIndep}{\ensuremath{-}0.618}
\newcommand{\PfDiffHiIndep}{\ensuremath{-}0.448}
\newcommand{\PfShareOfPaired}{82}
\newcommand{\PfShareOfPairedMeans}{70}
\newcommand{\HcErrIndep}{1.000}
\newcommand{\HcRsqIndep}{0.000}
\newcommand{\HcErrVar}{0.943}
\newcommand{\HcRsqVar}{\ensuremath{-}0.799}
\newcommand{\HcErrRefCorr}{0.510}
\newcommand{\HcRsqRefCorr}{\ensuremath{-}9.618}
\newcommand{\HcErrRefCov}{0.705}
\newcommand{\HcRsqRefCov}{\ensuremath{-}258}
\newcommand{\HcErrReg}{0.424}
\newcommand{\HcRsqReg}{0.259}
\newcommand{\HcErrCells}{0.397}
\newcommand{\HcRsqCells}{0.102}
\newcommand{\HcErrClosed}{0.367}
\newcommand{\HcRsqClosed}{0.263}
\newcommand{\HcErrFirst}{1{,}041}
\newcommand{\HcRsqFirst}{\ensuremath{-}351{,}681}
\newcommand{\HcErrOracle}{0.384}
\newcommand{\HcRsqOracle}{\ensuremath{-}0.028}
\newcommand{\HcDiffVar}{\ensuremath{-}0.576}
\newcommand{\HcDiffLoVar}{\ensuremath{-}0.670}
\newcommand{\HcDiffHiVar}{\ensuremath{-}0.444}
\newcommand{\HcDiffReg}{\ensuremath{-}0.057}
\newcommand{\HcDiffLoReg}{\ensuremath{-}0.082}
\newcommand{\HcDiffHiReg}{\ensuremath{-}0.031}
\newcommand{\HcRecipients}{8}
\newcommand{\PcObs}{0.536}
\newcommand{\PcNull}{1.255}
\newcommand{\PcPerms}{200}
\newcommand{\PcP}{1/201}
\newcommand{\PcDonors}{12}
\newcommand{\LcErrFive}{0.723}
\newcommand{\LcErrTen}{0.610}
\newcommand{\LcErrFifteen}{0.553}
\newcommand{\LcErrAll}{0.536}
\newcommand{\PfRidgeEdge}{11 of 12}
\newcommand{\PfSamples}{19--20}

% Generated by extension/cross_study/make_assets.py. Do not edit.
\newcommand{\XsPatients}{118}
\newcommand{\XsTrainCells}{223{,}640}
\newcommand{\XsHalfCells}{111{,}816}
\newcommand{\XsGenes}{206}
\newcommand{\XsProteins}{122}
\newcommand{\XsPfScale}{0.01}
\newcommand{\XsPfRank}{32}
\newcommand{\XsPfMomentIter}{9{,}858}
\newcommand{\XsPfMomentGrad}{2.4\times10^{-4}}
\newcommand{\XsPairedRidge}{0.01}
\newcommand{\XsRegLambda}{0.1}
\newcommand{\XsHaoCells}{161{,}764}
\newcommand{\XsHaoScoring}{80{,}884}
\newcommand{\XsColonScoring}{7{,}979}
\newcommand{\XsHaoDonors}{8}
\newcommand{\XsHaoErrPfMeans}{0.292}
\newcommand{\XsHaoRhoPfMeans}{0.964}
\newcommand{\XsHaoErrPfMoment}{0.389}
\newcommand{\XsHaoRhoPfMoment}{0.957}
\newcommand{\XsHaoErrPaired}{0.134}
\newcommand{\XsHaoRhoPaired}{0.995}
\newcommand{\XsHaoErrReg}{0.522}
\newcommand{\XsHaoRhoReg}{0.901}
\newcommand{\XsHaoErrCorr}{0.630}
\newcommand{\XsHaoRhoCorr}{0.846}
\newcommand{\XsHaoErrCov}{0.951}
\newcommand{\XsHaoRhoCov}{0.570}
\newcommand{\XsHaoErrVar}{0.976}
\newcommand{\XsHaoRhoVar}{0.394}
\newcommand{\XsHaoErrBench}{0.120}
\newcommand{\XsHaoRhoBench}{0.998}
\newcommand{\XsHaoErrIndep}{1.000}
\newcommand{\XsHaoRhoIndep}{0.000}
\newcommand{\XsHaoNullErr}{0.949}
\newcommand{\XsHaoNullMin}{0.639}
\newcommand{\XsHaoPermP}{0.005}
\newcommand{\XsHaoPermMaxDonorP}{0.005}
\newcommand{\XsHaoShare}{82}
\newcommand{\XsHaoShareLo}{80}
\newcommand{\XsHaoShareHi}{84}
\newcommand{\XsHaoShareMoment}{71}
\newcommand{\XsHaoShareMomentLo}{69}
\newcommand{\XsHaoShareMomentHi}{72}
\newcommand{\XsHaoWErrPfMeans}{1.082}
\newcommand{\XsHaoWErrPfMeansLo}{1.014}
\newcommand{\XsHaoWErrPfMeansHi}{1.162}
\newcommand{\XsHaoWBelowPfMeans}{2}
\newcommand{\XsHaoWErrPfMoment}{1.038}
\newcommand{\XsHaoWErrPfMomentLo}{0.981}
\newcommand{\XsHaoWErrPfMomentHi}{1.104}
\newcommand{\XsHaoWBelowPfMoment}{4}
\newcommand{\XsHaoWErrPaired}{0.528}
\newcommand{\XsHaoWErrPairedLo}{0.502}
\newcommand{\XsHaoWErrPairedHi}{0.565}
\newcommand{\XsHaoWBelowPaired}{8}
\newcommand{\XsHaoWErrBench}{0.348}
\newcommand{\XsHaoWErrBenchLo}{0.323}
\newcommand{\XsHaoWErrBenchHi}{0.378}
\newcommand{\XsHaoWBelowBench}{8}
\newcommand{\XsHaoWShare}{\ensuremath{-}20}
\newcommand{\XsHaoWShareLo}{\ensuremath{-}41}
\newcommand{\XsHaoWShareHi}{\ensuremath{-}3}
\newcommand{\XsHaoDiffReg}{0.388}
\newcommand{\XsHaoDiffRegLo}{0.364}
\newcommand{\XsHaoDiffRegHi}{0.415}
\newcommand{\XsHaoFavReg}{8}
\newcommand{\XsHaoDiffCorr}{0.495}
\newcommand{\XsHaoDiffCorrLo}{0.466}
\newcommand{\XsHaoDiffCorrHi}{0.531}
\newcommand{\XsHaoFavCorr}{8}
\newcommand{\XsHaoDiffCov}{0.817}
\newcommand{\XsHaoDiffCovLo}{0.615}
\newcommand{\XsHaoDiffCovHi}{1.078}
\newcommand{\XsHaoFavCov}{8}
\newcommand{\XsHaoDiffVar}{0.842}
\newcommand{\XsHaoDiffVarLo}{0.834}
\newcommand{\XsHaoDiffVarHi}{0.852}
\newcommand{\XsHaoFavVar}{8}
\newcommand{\XsHaoMomentMinusMeans}{0.097}
\newcommand{\XsHaoMomentMinusMeansLo}{0.078}
\newcommand{\XsHaoMomentMinusMeansHi}{0.119}
\newcommand{\XsHaoMomentBetter}{0}
\newcommand{\XsHaoPfMeansBetterThanPaired}{0}
\newcommand{\XsHaoPfMomentBetterThanPaired}{0}
\newcommand{\XsHaoPfMeansBetterThanBench}{0}
\newcommand{\XsHaoPairedBetterThanBench}{1}
\newcommand{\XsColonDonors}{12}
\newcommand{\XsColonErrPfMeans}{0.988}
\newcommand{\XsColonRhoPfMeans}{0.397}
\newcommand{\XsColonErrPfMoment}{0.907}
\newcommand{\XsColonRhoPfMoment}{0.463}
\newcommand{\XsColonErrPaired}{0.719}
\newcommand{\XsColonRhoPaired}{0.689}
\newcommand{\XsColonErrReg}{0.859}
\newcommand{\XsColonRhoReg}{0.516}
\newcommand{\XsColonErrCorr}{1.132}
\newcommand{\XsColonRhoCorr}{0.416}
\newcommand{\XsColonErrCov}{4.651}
\newcommand{\XsColonRhoCov}{0.230}
\newcommand{\XsColonErrVar}{0.971}
\newcommand{\XsColonRhoVar}{0.278}
\newcommand{\XsColonErrBench}{0.539}
\newcommand{\XsColonRhoBench}{0.842}
\newcommand{\XsColonErrIndep}{1.000}
\newcommand{\XsColonRhoIndep}{0.000}
\newcommand{\XsColonNullErr}{1.168}
\newcommand{\XsColonNullMin}{1.019}
\newcommand{\XsColonPermP}{0.005}
\newcommand{\XsColonPermMaxDonorP}{0.090}
\newcommand{\XsColonShare}{\ensuremath{-}12}
\newcommand{\XsColonShareLo}{\ensuremath{-}49}
\newcommand{\XsColonShareHi}{15}
\newcommand{\XsColonShareMoment}{25}
\newcommand{\XsColonShareMomentLo}{5}
\newcommand{\XsColonShareMomentHi}{42}
\newcommand{\XsColonWErrPfMeans}{1.274}
\newcommand{\XsColonWErrPfMeansLo}{1.239}
\newcommand{\XsColonWErrPfMeansHi}{1.307}
\newcommand{\XsColonWBelowPfMeans}{0}
\newcommand{\XsColonWErrPfMoment}{1.227}
\newcommand{\XsColonWErrPfMomentLo}{1.200}
\newcommand{\XsColonWErrPfMomentHi}{1.253}
\newcommand{\XsColonWBelowPfMoment}{0}
\newcommand{\XsColonWErrPaired}{1.000}
\newcommand{\XsColonWErrPairedLo}{0.987}
\newcommand{\XsColonWErrPairedHi}{1.013}
\newcommand{\XsColonWBelowPaired}{7}
\newcommand{\XsColonWErrBench}{0.936}
\newcommand{\XsColonWErrBenchLo}{0.922}
\newcommand{\XsColonWErrBenchHi}{0.953}
\newcommand{\XsColonWBelowBench}{11}
\newcommand{\XsColonWShare}{\ensuremath{-}4126}
\newcommand{\XsColonWShareLo}{\ensuremath{-}12228}
\newcommand{\XsColonWShareHi}{891}
\newcommand{\XsColonDiffReg}{0.140}
\newcommand{\XsColonDiffRegLo}{0.099}
\newcommand{\XsColonDiffRegHi}{0.180}
\newcommand{\XsColonFavReg}{12}
\newcommand{\XsColonDiffCorr}{0.413}
\newcommand{\XsColonDiffCorrLo}{0.348}
\newcommand{\XsColonDiffCorrHi}{0.500}
\newcommand{\XsColonFavCorr}{12}
\newcommand{\XsColonDiffCov}{3.932}
\newcommand{\XsColonDiffCovLo}{2.510}
\newcommand{\XsColonDiffCovHi}{5.443}
\newcommand{\XsColonFavCov}{12}
\newcommand{\XsColonDiffVar}{0.252}
\newcommand{\XsColonDiffVarLo}{0.182}
\newcommand{\XsColonDiffVarHi}{0.324}
\newcommand{\XsColonFavVar}{12}
\newcommand{\XsColonMomentMinusMeans}{\ensuremath{-}0.081}
\newcommand{\XsColonMomentMinusMeansLo}{\ensuremath{-}0.106}
\newcommand{\XsColonMomentMinusMeansHi}{\ensuremath{-}0.053}
\newcommand{\XsColonMomentBetter}{11}
\newcommand{\XsColonPfMeansBetterThanPaired}{0}
\newcommand{\XsColonPfMomentBetterThanPaired}{0}
\newcommand{\XsColonPfMeansBetterThanBench}{0}
\newcommand{\XsColonPairedBetterThanBench}{0}
\newcommand{\XsHaoWRhoPfMeans}{0.43}
\newcommand{\XsHaoWRhoPaired}{0.85}
\newcommand{\XsHaoWRhoBench}{0.96}
\newcommand{\XsHaoLTwoErrPfMeans}{1.484}
\newcommand{\XsHaoLTwoErrPfMoment}{1.435}
\newcommand{\XsHaoLTwoErrPaired}{0.868}
\newcommand{\XsHaoLTwoErrBench}{0.541}
\newcommand{\XsDimM}{11}
\newcommand{\XsEffDim}{16.9}
\newcommand{\XsHaoIdRho}{38}
\newcommand{\XsHaoIdRhoMin}{35}
\newcommand{\XsHaoIdRhoMax}{41}
\newcommand{\XsHaoIdFloor}{0.16}
\newcommand{\XsHaoIdFloorMin}{0.13}
\newcommand{\XsHaoIdFloorMax}{0.20}
\newcommand{\XsHaoIdRadius}{0.88}
\newcommand{\XsHaoIdRadiusMin}{0.86}
\newcommand{\XsHaoIdRadiusMax}{0.90}
\newcommand{\XsHaoWIdRho}{12}
\newcommand{\XsHaoWIdRhoMin}{11}
\newcommand{\XsHaoWIdRhoMax}{13}
\newcommand{\XsHaoWIdFloor}{0.72}
\newcommand{\XsHaoWIdFloorMin}{0.61}
\newcommand{\XsHaoWIdFloorMax}{0.82}
\newcommand{\XsHaoWIdRadius}{1.84}
\newcommand{\XsHaoWIdRadiusMin}{1.70}
\newcommand{\XsHaoWIdRadiusMax}{2.23}
\newcommand{\XsColonIdRho}{19}
\newcommand{\XsColonIdRhoMin}{15}
\newcommand{\XsColonIdRhoMax}{24}
\newcommand{\XsColonIdFloor}{0.71}
\newcommand{\XsColonIdFloorMin}{0.59}
\newcommand{\XsColonIdFloorMax}{0.87}
\newcommand{\XsColonIdRadius}{1.32}
\newcommand{\XsColonIdRadiusMin}{1.05}
\newcommand{\XsColonIdRadiusMax}{1.51}
\newcommand{\XsColonWIdRho}{9}
\newcommand{\XsColonWIdRhoMin}{8}
\newcommand{\XsColonWIdRhoMax}{10}
\newcommand{\XsColonWIdFloor}{0.92}
\newcommand{\XsColonWIdFloorMin}{0.87}
\newcommand{\XsColonWIdFloorMax}{0.98}
\newcommand{\XsColonWIdRadius}{1.89}
\newcommand{\XsColonWIdRadiusMin}{1.52}
\newcommand{\XsColonWIdRadiusMax}{2.26}
\newcommand{\XsIdUnits}{48}
\newcommand{\XsFloorErrCorr}{0.89}
\newcommand{\XsRhoErrCorr}{\ensuremath{-}0.94}
\newcommand{\XsMaxFMin}{0.44}
\newcommand{\XsMaxFMax}{4.29}
\newcommand{\XsMaxFAboveOne}{46}
\newcommand{\XsMaxFTotalMin}{1.02}
\newcommand{\XsGeneSpearmanUnidHao}{0.86}
\newcommand{\XsGeneSpearmanAxisHao}{0.17}
\newcommand{\XsGeneSpearmanUnidColon}{0.72}
\newcommand{\XsGeneSpearmanAxisColon}{\ensuremath{-}0.01}
\newcommand{\XsRestartBChange}{9\times10^{-7}}
\newcommand{\XsRestartEChange}{2\times10^{-6}}
\newcommand{\XsLcHaoEight}{0.637}
\newcommand{\XsLcColonEight}{1.024}
\newcommand{\XsLcDimEight}{3.7}
\newcommand{\XsLcHaoSixteen}{0.493}
\newcommand{\XsLcColonSixteen}{0.992}
\newcommand{\XsLcDimSixteen}{5.1}
\newcommand{\XsLcHaoThirtyTwo}{0.395}
\newcommand{\XsLcColonThirtyTwo}{0.931}
\newcommand{\XsLcDimThirtyTwo}{4.2}
\newcommand{\XsLcHaoSixtyFour}{0.336}
\newcommand{\XsLcColonSixtyFour}{0.996}
\newcommand{\XsLcDimSixtyFour}{11.4}
\newcommand{\XsLcHaoAll}{0.292}
\newcommand{\XsLcColonAll}{0.988}
\newcommand{\XsLcDimAll}{11}
\newcommand{\XsRigErrMeans}{0.566}
\newcommand{\XsRigErrMoment}{0.680}
\newcommand{\XsRigNullMeans}{1.086}
\newcommand{\XsRigNullMoment}{1.048}
\newcommand{\XsRigMaxPMeans}{0.005}
\newcommand{\XsRigMaxPMoment}{0.143}
\newcommand{\XsRigDonorsSigMeans}{12}
\newcommand{\XsRigDonorsSigMoment}{10}
\newcommand{\XsRigEdge}{0}
\newcommand{\XsRigExtended}{0}
\newcommand{\XsRigRankMin}{8}
\newcommand{\XsRigRankMax}{8}
\newcommand{\XsRigConverged}{9}
\newcommand{\XsRigDimMin}{11}
\newcommand{\XsRigDimMax}{12}
\newcommand{\XsRigBiopsiesMin}{19}
\newcommand{\XsRigBiopsiesMax}{20}
\newcommand{\XsRigIdRho}{37}
\newcommand{\XsRigIdFloor}{0.59}
\newcommand{\XsRigIdRadius}{1.13}
\newcommand{\XsRigMaxF}{1.46}
\newcommand{\XsRigShareMeans}{65}
\newcommand{\XsRigShareMoment}{50}
\newcommand{\XsPoneUnits}{512}
\newcommand{\XsPoneDim}{86}
\newcommand{\XsPoneErr}{1.287}
\newcommand{\XsPoneErrLTwo}{2.163}
\newcommand{\XsPoneErrTotal}{0.660}
\newcommand{\XsPoneRho}{50}
\newcommand{\XsPoneFloor}{0.28}
\newcommand{\XsMixHaoBetweenShare}{97}
\newcommand{\XsMixHaoBetweenShareMin}{96}
\newcommand{\XsMixHaoBetweenShareMax}{98}
\newcommand{\XsMixHaoWithinShare}{11}
\newcommand{\XsMixHaoWithinShareMin}{10}
\newcommand{\XsMixHaoWithinShareMax}{14}
\newcommand{\XsMixHaoErrBetween}{0.71}
\newcommand{\XsMixHaoErrBetweenMin}{0.65}
\newcommand{\XsMixHaoErrBetweenMax}{0.76}
\newcommand{\XsMixHaoErrBetweenBelow}{8}
\newcommand{\XsMixHaoErrWithin}{1.92}
\newcommand{\XsMixHaoErrWithinMin}{1.55}
\newcommand{\XsMixHaoErrWithinMax}{2.56}
\newcommand{\XsMixHaoErrWithinBelow}{0}
\newcommand{\XsMixHaoErrDirect}{0.75}
\newcommand{\XsMixHaoErrDirectMin}{0.69}
\newcommand{\XsMixHaoErrDirectMax}{0.81}
\newcommand{\XsMixHaoErrDirectBelow}{8}
\newcommand{\XsMixHaoRBetween}{0.75}
\newcommand{\XsMixHaoRWithin}{0.15}
\newcommand{\XsMixHaoNormWithin}{1.8}
\newcommand{\XsMixHaoNormBetween}{0.99}
\newcommand{\XsMixHaoRhoBetween}{88}
\newcommand{\XsMixHaoRhoBetweenMin}{87}
\newcommand{\XsMixHaoRhoBetweenMax}{89}
\newcommand{\XsMixHaoRhoWithin}{12}
\newcommand{\XsMixHaoRhoWithinMin}{10}
\newcommand{\XsMixHaoRhoWithinMax}{13}
\newcommand{\XsMixHaoVarBetween}{39}
\newcommand{\XsMixHaoVarBetweenMin}{34}
\newcommand{\XsMixHaoVarBetweenMax}{42}
\newcommand{\XsMixColonBetweenShare}{81}
\newcommand{\XsMixColonBetweenShareMin}{45}
\newcommand{\XsMixColonBetweenShareMax}{94}
\newcommand{\XsMixColonWithinShare}{43}
\newcommand{\XsMixColonWithinShareMin}{27}
\newcommand{\XsMixColonWithinShareMax}{84}
\newcommand{\XsMixColonErrBetween}{1.14}
\newcommand{\XsMixColonErrBetweenMin}{0.85}
\newcommand{\XsMixColonErrBetweenMax}{1.29}
\newcommand{\XsMixColonErrBetweenBelow}{1}
\newcommand{\XsMixColonErrWithin}{1.38}
\newcommand{\XsMixColonErrWithinMin}{1.21}
\newcommand{\XsMixColonErrWithinMax}{1.53}
\newcommand{\XsMixColonErrWithinBelow}{0}
\newcommand{\XsMixColonErrDirect}{1.20}
\newcommand{\XsMixColonErrDirectMin}{0.91}
\newcommand{\XsMixColonErrDirectMax}{1.34}
\newcommand{\XsMixColonErrDirectBelow}{1}
\newcommand{\XsMixColonRBetween}{0.35}
\newcommand{\XsMixColonRWithin}{0.07}
\newcommand{\XsMixColonNormWithin}{1.0}
\newcommand{\XsMixColonNormBetween}{1.00}
\newcommand{\XsMixColonRhoBetween}{57}
\newcommand{\XsMixColonRhoBetweenMin}{50}
\newcommand{\XsMixColonRhoBetweenMax}{66}
\newcommand{\XsMixColonRhoWithin}{9}
\newcommand{\XsMixColonRhoWithinMin}{8}
\newcommand{\XsMixColonRhoWithinMax}{10}
\newcommand{\XsMixColonVarBetween}{20}
\newcommand{\XsMixColonVarBetweenMin}{9}
\newcommand{\XsMixColonVarBetweenMax}{28}
\newcommand{\XsRigContinued}{3}
\newcommand{\XsRigContinueConverged}{3}
\newcommand{\XsRigContinueMaxChange}{0.0015}
\newcommand{\XsRigContinueBound}{0.002}
\newcommand{\XsRigContinueMaxWChange}{0.031}
\newcommand{\XsPairFcgraTotTruth}{0.75}
\newcommand{\XsPairFcgraTotPfMeans}{0.29}
\newcommand{\XsPairFcgraTotPaired}{0.61}
\newcommand{\XsPairFcgraTotBench}{0.67}
\newcommand{\XsPairFcgraWTruth}{0.63}
\newcommand{\XsPairFcgraWPfMeans}{0.12}
\newcommand{\XsPairFcgraWPaired}{0.52}
\newcommand{\XsPairFcgraWBench}{0.46}
\newcommand{\XsPairKlrbTotTruth}{0.69}
\newcommand{\XsPairKlrbTotPfMeans}{0.26}
\newcommand{\XsPairKlrbTotPaired}{0.44}
\newcommand{\XsPairKlrbTotBench}{0.60}
\newcommand{\XsPairKlrbWTruth}{0.51}
\newcommand{\XsPairKlrbWPfMeans}{0.07}
\newcommand{\XsPairKlrbWPaired}{0.30}
\newcommand{\XsPairKlrbWBench}{0.28}
\newcommand{\XsPairIghdTotTruth}{0.76}
\newcommand{\XsPairIghdTotPfMeans}{0.42}
\newcommand{\XsPairIghdTotPaired}{0.60}
\newcommand{\XsPairIghdTotBench}{0.67}
\newcommand{\XsPairIghdWTruth}{0.51}
\newcommand{\XsPairIghdWPfMeans}{0.10}
\newcommand{\XsPairIghdWPaired}{0.36}
\newcommand{\XsPairIghdWBench}{0.35}
\newcommand{\XsPairIlrTotTruth}{0.82}
\newcommand{\XsPairIlrTotPfMeans}{0.57}
\newcommand{\XsPairIlrTotPaired}{0.68}
\newcommand{\XsPairIlrTotBench}{0.73}
\newcommand{\XsPairIlrWTruth}{0.48}
\newcommand{\XsPairIlrWPfMeans}{0.20}
\newcommand{\XsPairIlrWPaired}{0.39}
\newcommand{\XsPairIlrWBench}{0.29}
\newcommand{\XsPairLyzTotTruth}{0.89}
\newcommand{\XsPairLyzTotPfMeans}{0.77}
\newcommand{\XsPairLyzTotPaired}{0.81}
\newcommand{\XsPairLyzTotBench}{0.80}
\newcommand{\XsPairLyzWTruth}{0.21}
\newcommand{\XsPairLyzWPfMeans}{0.22}
\newcommand{\XsPairLyzWPaired}{0.32}
\newcommand{\XsPairLyzWBench}{0.17}
\newcommand{\XsPairCdaTotTruth}{0.87}
\newcommand{\XsPairCdaTotPfMeans}{0.61}
\newcommand{\XsPairCdaTotPaired}{0.77}
\newcommand{\XsPairCdaTotBench}{0.77}
\newcommand{\XsPairCdaWTruth}{0.05}
\newcommand{\XsPairCdaWPfMeans}{0.09}
\newcommand{\XsPairCdaWPaired}{0.21}
\newcommand{\XsPairCdaWBench}{0.04}

% Generated by extension/recoverability/make_assets.py. Do not edit.
\newcommand{\RcDevFReported}{2.78}
\newcommand{\RcDevFClr}{2.69}
\newcommand{\RcDevFRaw}{2.75}
\newcommand{\RcDevTotalAnalyses}{20}
\newcommand{\RcDevFMeasTotal}{2.78}
\newcommand{\RcDevFLatTotal}{0.28}
\newcommand{\RcDevFMeasAbove}{20}
\newcommand{\RcDevFMeasLowerAbove}{18}
\newcommand{\RcDevFLatBloodMin}{0.85}
\newcommand{\RcDevFLatBloodMax}{1.90}
\newcommand{\RcDevFLatBloodAbove}{6}
\newcommand{\RcDevFLatBloodLowerAbove}{6}
\newcommand{\RcDevFLatColonMax}{0.31}
\newcommand{\RcDevWithinAnalyses}{28}
\newcommand{\RcDevFMeasWithin}{3.23}
\newcommand{\RcDevFLatWithin}{0.51}
\newcommand{\RcDevFLatWithinLowerAbove}{4}
\newcommand{\RcDevPOneAnalyses}{48}
\newcommand{\RcDevPOneFLat}{3.02}
\newcommand{\RcDevPOneLowerAbove}{45}
\newcommand{\RcDevDeparture}{0.82}
\newcommand{\RcDevDepartureMin}{0.70}
\newcommand{\RcDevDepartureMax}{1.24}
\newcommand{\RcDevBoundsNonempty}{37}
\newcommand{\RcDevBoundsAnalyses}{48}
\newcommand{\RcDevBoundsSigned}{0}
\newcommand{\RcDevBoundsInsideMin}{0.928}
\newcommand{\RcDevLabelBlood}{0.12}
\newcommand{\RcDevLabelBloodMin}{0.10}
\newcommand{\RcDevLabelBloodMax}{0.16}
\newcommand{\RcDevPfBlood}{0.29}
\newcommand{\RcDevPairedBlood}{0.13}
\newcommand{\RcDevLabelColon}{0.57}
\newcommand{\RcDevLabelColonMin}{0.38}
\newcommand{\RcDevLabelColonMax}{0.95}
\newcommand{\RcDevPfColon}{0.99}
\newcommand{\RcDevPairedColon}{0.72}
\newcommand{\RcDevLabelBeatsPaired}{16}
\newcommand{\RcDevLabelUnits}{20}
\newcommand{\RcDevPairs}{1,776}
\newcommand{\RcBmSamples}{12}
\newcommand{\RcBmDonors}{9}
\newcommand{\RcBmCells}{90,261}
\newcommand{\RcBmGenes}{201}
\newcommand{\RcBmProteins}{88}
\newcommand{\RcBmFineMin}{24}
\newcommand{\RcBmFineMax}{36}
\newcommand{\RcBmFrozenTot}{0.52}
\newcommand{\RcBmFrozenMeasTot}{0.52}
\newcommand{\RcBmFrozenBelowTot}{12}
\newcommand{\RcBmPOneTot}{0.68}
\newcommand{\RcBmColonTot}{0.49}
\newcommand{\RcBmColonBelowTot}{12}
\newcommand{\RcBmBenchTot}{0.17}
\newcommand{\RcBmFrozenCoarse}{0.99}
\newcommand{\RcBmFrozenMeasCoarse}{1.02}
\newcommand{\RcBmFrozenBelowCoarse}{6}
\newcommand{\RcBmPOneCoarse}{1.01}
\newcommand{\RcBmColonCoarse}{0.93}
\newcommand{\RcBmColonBelowCoarse}{10}
\newcommand{\RcBmBenchCoarse}{0.41}
\newcommand{\RcBmFrozenFine}{1.30}
\newcommand{\RcBmFrozenMeasFine}{1.43}
\newcommand{\RcBmFrozenBelowFine}{0}
\newcommand{\RcBmPOneFine}{1.44}
\newcommand{\RcBmColonFine}{1.10}
\newcommand{\RcBmColonBelowFine}{2}
\newcommand{\RcBmBenchFine}{0.66}
\newcommand{\RcBmLabelTot}{0.21}
\newcommand{\RcBmFMeasTot}{2.80}
\newcommand{\RcBmFLatTot}{0.52}
\newcommand{\RcBmFLatMax}{0.94}
\newcommand{\RcBmFMeasMin}{1.85}
\newcommand{\RcBmFMeasLowerAbove}{36}
\newcommand{\RcBmColonFLatMax}{0.73}
\newcommand{\RcBmPOneFLatMin}{1.00}
\newcommand{\RcBmPOneFLatMax}{3.65}
\newcommand{\RcBmPOneLowerAbove}{35}
\newcommand{\RcBmAnalyses}{36}
\newcommand{\RcBmPairs}{1,476}
\newcommand{\RcBmFailures}{838}
\newcommand{\RcBmPermAbstain}{95}
\newcommand{\RcBmLcAbstain}{7.5}
\newcommand{\RcBmAucPrimary}{0.87}
\newcommand{\RcBmAucLabelFree}{0.84}
\newcommand{\RcBmAucCoverage}{0.84}
\newcommand{\RcBmAucComposition}{0.84}
\newcommand{\RcBmAucCells}{0.46}
\newcommand{\RcBmAucNTrain}{0.54}
\newcommand{\RcBmAucShift}{0.83}
\newcommand{\RcBmAucCompat}{0.29}
\newcommand{\RcBmAucTypeDis}{0.59}
\newcommand{\RcBmDiffPrimaryCoverage}{0.030}
\newcommand{\RcBmDiffPrimaryCoverageLo}{\ensuremath{-}0.014}
\newcommand{\RcBmDiffPrimaryCoverageHi}{0.059}
\newcommand{\RcBmDiffPrimaryComposition}{0.034}
\newcommand{\RcBmDiffPrimaryCompositionLo}{0.024}
\newcommand{\RcBmDiffPrimaryCompositionHi}{0.048}
\newcommand{\RcBmDiffPrimaryCells}{0.417}
\newcommand{\RcBmDiffPrimaryCellsLo}{0.257}
\newcommand{\RcBmDiffPrimaryCellsHi}{0.451}
\newcommand{\RcBmDiffPrimaryNTrain}{0.334}
\newcommand{\RcBmDiffPrimaryNTrainLo}{0.261}
\newcommand{\RcBmDiffPrimaryNTrainHi}{0.392}
\newcommand{\RcBmDiffPrimaryShift}{0.045}
\newcommand{\RcBmDiffPrimaryShiftLo}{\ensuremath{-}0.012}
\newcommand{\RcBmDiffPrimaryShiftHi}{0.086}
\newcommand{\RcBmDiffLabelFreeCoverage}{0.001}
\newcommand{\RcBmDiffLabelFreeCoverageLo}{\ensuremath{-}0.014}
\newcommand{\RcBmDiffLabelFreeCoverageHi}{0.020}
\newcommand{\RcBmHiPairs}{251}
\newcommand{\RcBmHiFailures}{19}
\newcommand{\RcBmHiPOne}{36}
\newcommand{\RcBmHiFailuresPOne}{18}
\newcommand{\RcBmHiAucPrimary}{0.99}
\newcommand{\RcBmHiAucPrimaryLo}{0.98}
\newcommand{\RcBmHiAucPrimaryHi}{1.00}
\newcommand{\RcBmHiAucTypeDis}{0.93}
\newcommand{\RcBmHiAucTypeDisLo}{0.90}
\newcommand{\RcBmHiAucTypeDisHi}{0.95}
\newcommand{\RcBmHiAucCompat}{0.89}
\newcommand{\RcBmHiAucCompatLo}{0.85}
\newcommand{\RcBmHiAucCompatHi}{0.92}
\newcommand{\RcBmHiAucLabelFree}{0.22}
\newcommand{\RcBmHiAucLabelFreeLo}{0.14}
\newcommand{\RcBmHiAucLabelFreeHi}{0.37}
\newcommand{\RcBmRankPrimary}{0.69}
\newcommand{\RcBmRankCoverage}{0.48}
\newcommand{\RcBmRankShift}{0.46}
\newcommand{\RcBmRefRandomFifty}{1.01}
\newcommand{\RcBmRefBalancedFifty}{1.04}
\newcommand{\RcBmRefGuidedFifty}{0.96}
\newcommand{\RcBmRefHybridFifty}{0.21}
\newcommand{\RcBmRefPairedFifty}{0.64}
\newcommand{\RcBmRefRandomHundred}{0.91}
\newcommand{\RcBmRefBalancedHundred}{0.93}
\newcommand{\RcBmRefGuidedHundred}{0.87}
\newcommand{\RcBmRefHybridHundred}{0.19}
\newcommand{\RcBmRefPairedHundred}{0.51}
\newcommand{\RcBmRefRandomTwoHundred}{0.80}
\newcommand{\RcBmRefBalancedTwoHundred}{0.84}
\newcommand{\RcBmRefGuidedTwoHundred}{0.77}
\newcommand{\RcBmRefHybridTwoHundred}{0.18}
\newcommand{\RcBmRefPairedTwoHundred}{0.38}
\newcommand{\RcBmRefRandomFourHundred}{0.69}
\newcommand{\RcBmRefBalancedFourHundred}{0.74}
\newcommand{\RcBmRefGuidedFourHundred}{0.67}
\newcommand{\RcBmRefHybridFourHundred}{0.16}
\newcommand{\RcBmRefPairedFourHundred}{0.29}
\newcommand{\RcBmRefRandomEightHundred}{0.58}
\newcommand{\RcBmRefBalancedEightHundred}{0.62}
\newcommand{\RcBmRefGuidedEightHundred}{0.56}
\newcommand{\RcBmRefHybridEightHundred}{0.14}
\newcommand{\RcBmRefPairedEightHundred}{0.21}
\newcommand{\RcBmRefBetweenOnly}{0.21}
\newcommand{\RcBmRefGainRandom}{0.034}
\newcommand{\RcBmRefGainRandomLo}{0.026}
\newcommand{\RcBmRefGainRandomHi}{0.041}
\newcommand{\RcBmRefGainBalanced}{0.068}
\newcommand{\RcBmRefGainBalancedLo}{0.056}
\newcommand{\RcBmRefGainBalancedHi}{0.087}
\newcommand{\RcBmRefAllSamples}{all}
\newcommand{\RcBmRefEquivTwoHundred}{243}
\newcommand{\RcBmRefEquivEightHundred}{above 800}
\newcommand{\RcBmBoundsSigned}{0}
\newcommand{\RcBmBoundsAnalyses}{36}
\newcommand{\RcBmBoundsHalfwidth}{0.74}
\newcommand{\RcBmBoundsTarget}{0.14}
\newcommand{\RcBmHOne}{not met}
\newcommand{\RcBmHThree}{met}
\newcommand{\RcBmHSix}{met}
\newcommand{\RcFmPairs}{1,776}
\newcommand{\RcFmFailures}{1,235}
\newcommand{\RcFmPrimaryAuc}{0.95}
\newcommand{\RcFmPrimaryCoverage}{\ensuremath{-}0.12}
\newcommand{\RcFmPrimaryLogF}{0.57}
\newcommand{\RcFmPrimaryAgreement}{\ensuremath{-}1.25}
\newcommand{\RcFmPrimaryComposition}{\ensuremath{-}2.97}
\newcommand{\RcFmLabelFreeAuc}{0.89}
\newcommand{\RcFmLabelFreeCoverage}{\ensuremath{-}2.23}
\newcommand{\RcFmLabelFreeLogF}{0.86}
\newcommand{\RcPhMeansOnly}{0.21}
\newcommand{\RcPhMeansOnlyLo}{0.19}
\newcommand{\RcPhMeansOnlyHi}{0.24}
\newcommand{\RcPhTotFiftyRandom}{0.22}
\newcommand{\RcPhTotFiftyRandomDiff}{0.009}
\newcommand{\RcPhTotFiftyRandomDiffLo}{0.000}
\newcommand{\RcPhTotFiftyRandomDiffHi}{0.014}
\newcommand{\RcPhTotFiftyGuided}{0.21}
\newcommand{\RcPhTotFiftyGuidedDiff}{\ensuremath{-}0.002}
\newcommand{\RcPhTotFiftyGuidedDiffLo}{\ensuremath{-}0.009}
\newcommand{\RcPhTotFiftyGuidedDiffHi}{0.003}
\newcommand{\RcPhTotFiftyAlone}{0.64}
\newcommand{\RcPhTotFiftyGuidedMinusRandom}{\ensuremath{-}0.010}
\newcommand{\RcPhTotFiftyGuidedMinusRandomLo}{\ensuremath{-}0.014}
\newcommand{\RcPhTotFiftyGuidedMinusRandomHi}{\ensuremath{-}0.006}
\newcommand{\RcPhTotHundredRandom}{0.20}
\newcommand{\RcPhTotHundredRandomDiff}{\ensuremath{-}0.011}
\newcommand{\RcPhTotHundredRandomDiffLo}{\ensuremath{-}0.022}
\newcommand{\RcPhTotHundredRandomDiffHi}{\ensuremath{-}0.004}
\newcommand{\RcPhTotHundredGuided}{0.19}
\newcommand{\RcPhTotHundredGuidedDiff}{\ensuremath{-}0.016}
\newcommand{\RcPhTotHundredGuidedDiffLo}{\ensuremath{-}0.029}
\newcommand{\RcPhTotHundredGuidedDiffHi}{\ensuremath{-}0.008}
\newcommand{\RcPhTotHundredAlone}{0.51}
\newcommand{\RcPhTotHundredGuidedMinusRandom}{\ensuremath{-}0.006}
\newcommand{\RcPhTotHundredGuidedMinusRandomLo}{\ensuremath{-}0.008}
\newcommand{\RcPhTotHundredGuidedMinusRandomHi}{\ensuremath{-}0.004}
\newcommand{\RcPhTotTwoHundredRandom}{0.18}
\newcommand{\RcPhTotTwoHundredRandomDiff}{\ensuremath{-}0.028}
\newcommand{\RcPhTotTwoHundredRandomDiffLo}{\ensuremath{-}0.045}
\newcommand{\RcPhTotTwoHundredRandomDiffHi}{\ensuremath{-}0.017}
\newcommand{\RcPhTotTwoHundredGuided}{0.18}
\newcommand{\RcPhTotTwoHundredGuidedDiff}{\ensuremath{-}0.032}
\newcommand{\RcPhTotTwoHundredGuidedDiffLo}{\ensuremath{-}0.051}
\newcommand{\RcPhTotTwoHundredGuidedDiffHi}{\ensuremath{-}0.020}
\newcommand{\RcPhTotTwoHundredAlone}{0.38}
\newcommand{\RcPhTotTwoHundredGuidedMinusRandom}{\ensuremath{-}0.004}
\newcommand{\RcPhTotTwoHundredGuidedMinusRandomLo}{\ensuremath{-}0.006}
\newcommand{\RcPhTotTwoHundredGuidedMinusRandomHi}{\ensuremath{-}0.003}
\newcommand{\RcPhTotFourHundredRandom}{0.16}
\newcommand{\RcPhTotFourHundredRandomDiff}{\ensuremath{-}0.046}
\newcommand{\RcPhTotFourHundredRandomDiffLo}{\ensuremath{-}0.068}
\newcommand{\RcPhTotFourHundredRandomDiffHi}{\ensuremath{-}0.032}
\newcommand{\RcPhTotFourHundredGuided}{0.16}
\newcommand{\RcPhTotFourHundredGuidedDiff}{\ensuremath{-}0.049}
\newcommand{\RcPhTotFourHundredGuidedDiffLo}{\ensuremath{-}0.072}
\newcommand{\RcPhTotFourHundredGuidedDiffHi}{\ensuremath{-}0.035}
\newcommand{\RcPhTotFourHundredAlone}{0.29}
\newcommand{\RcPhTotFourHundredGuidedMinusRandom}{\ensuremath{-}0.003}
\newcommand{\RcPhTotFourHundredGuidedMinusRandomLo}{\ensuremath{-}0.005}
\newcommand{\RcPhTotFourHundredGuidedMinusRandomHi}{\ensuremath{-}0.002}
\newcommand{\RcPhTotEightHundredRandom}{0.15}
\newcommand{\RcPhTotEightHundredRandomDiff}{\ensuremath{-}0.064}
\newcommand{\RcPhTotEightHundredRandomDiffLo}{\ensuremath{-}0.090}
\newcommand{\RcPhTotEightHundredRandomDiffHi}{\ensuremath{-}0.048}
\newcommand{\RcPhTotEightHundredGuided}{0.14}
\newcommand{\RcPhTotEightHundredGuidedDiff}{\ensuremath{-}0.067}
\newcommand{\RcPhTotEightHundredGuidedDiffLo}{\ensuremath{-}0.093}
\newcommand{\RcPhTotEightHundredGuidedDiffHi}{\ensuremath{-}0.051}
\newcommand{\RcPhTotEightHundredAlone}{0.21}
\newcommand{\RcPhTotEightHundredGuidedMinusRandom}{\ensuremath{-}0.003}
\newcommand{\RcPhTotEightHundredGuidedMinusRandomLo}{\ensuremath{-}0.004}
\newcommand{\RcPhTotEightHundredGuidedMinusRandomHi}{\ensuremath{-}0.002}
\newcommand{\RcPhAloneEightHundredMinusMeans}{\ensuremath{-}0.001}
\newcommand{\RcPhAloneEightHundredMinusMeansLo}{\ensuremath{-}0.033}
\newcommand{\RcPhAloneEightHundredMinusMeansHi}{0.020}
\newcommand{\RcPhSavingFifty}{31}
\newcommand{\RcPhSavingFiftyLo}{25}
\newcommand{\RcPhSavingFiftyHi}{34}
\newcommand{\RcPhRandomCellsFifty}{73}
\newcommand{\RcPhRandomCellsFiftyLo}{66}
\newcommand{\RcPhRandomCellsFiftyHi}{76}
\newcommand{\RcPhSavingHundred}{22}
\newcommand{\RcPhSavingHundredLo}{18}
\newcommand{\RcPhSavingHundredHi}{25}
\newcommand{\RcPhRandomCellsHundred}{128}
\newcommand{\RcPhRandomCellsHundredLo}{122}
\newcommand{\RcPhRandomCellsHundredHi}{133}
\newcommand{\RcPhSavingTwoHundred}{18}
\newcommand{\RcPhSavingTwoHundredLo}{14}
\newcommand{\RcPhSavingTwoHundredHi}{20}
\newcommand{\RcPhRandomCellsTwoHundred}{243}
\newcommand{\RcPhRandomCellsTwoHundredLo}{233}
\newcommand{\RcPhRandomCellsTwoHundredHi}{251}
\newcommand{\RcPhSavingFourHundred}{13}
\newcommand{\RcPhSavingFourHundredLo}{10}
\newcommand{\RcPhSavingFourHundredHi}{17}
\newcommand{\RcPhRandomCellsFourHundred}{462}
\newcommand{\RcPhRandomCellsFourHundredLo}{443}
\newcommand{\RcPhRandomCellsFourHundredHi}{480}
\newcommand{\RcPhIdAll}{0.84}
\newcommand{\RcPhIdAllLo}{0.78}
\newcommand{\RcPhIdAllHi}{0.90}
\newcommand{\RcPhIdInAll}{0.85}
\newcommand{\RcPhIdInAllLo}{0.80}
\newcommand{\RcPhIdInAllHi}{0.90}
\newcommand{\RcPhIdFamAll}{0.52}
\newcommand{\RcPhIdFamAllLo}{0.51}
\newcommand{\RcPhIdFamAllHi}{0.53}
\newcommand{\RcPhIdDiffAll}{0.030}
\newcommand{\RcPhIdDiffAllLo}{0.017}
\newcommand{\RcPhIdDiffAllHi}{0.045}
\newcommand{\RcPhPrimaryAll}{0.87}
\newcommand{\RcPhPrimaryAllLo}{0.81}
\newcommand{\RcPhPrimaryAllHi}{0.92}
\newcommand{\RcPhIdHi}{0.96}
\newcommand{\RcPhIdHiLo}{0.90}
\newcommand{\RcPhIdHiHi}{0.99}
\newcommand{\RcPhIdInHi}{0.96}
\newcommand{\RcPhIdInHiLo}{0.91}
\newcommand{\RcPhIdInHiHi}{0.99}
\newcommand{\RcPhIdFamHi}{0.94}
\newcommand{\RcPhIdFamHiLo}{0.88}
\newcommand{\RcPhIdFamHiHi}{0.96}
\newcommand{\RcPhIdDiffHi}{0.029}
\newcommand{\RcPhIdDiffHiLo}{0.002}
\newcommand{\RcPhIdDiffHiHi}{0.082}
\newcommand{\RcPhPrimaryHi}{0.99}
\newcommand{\RcPhPrimaryHiLo}{0.98}
\newcommand{\RcPhPrimaryHiHi}{1.00}
\newcommand{\RcPhLofoPOneHi}{0.99}
\newcommand{\RcPhLofoPOneHiLo}{0.97}
\newcommand{\RcPhLofoPOneHiHi}{0.99}
\newcommand{\RcPhLofoPOneAll}{0.87}
\newcommand{\RcPhLofoPOneAllLo}{0.80}
\newcommand{\RcPhLofoPOneAllHi}{0.92}
\newcommand{\RcPhLofoPOneFamily}{0.89}
\newcommand{\RcPhLofoPOneFamilyLo}{0.69}
\newcommand{\RcPhLofoPOneFamilyHi}{0.97}
\newcommand{\RcPhPOneWithinPairs}{24}
\newcommand{\RcPhPOneWithinFailures}{18}
\newcommand{\RcPhPOneWithinPrimary}{0.77}
\newcommand{\RcPhPOneWithinPrimaryLo}{0.51}
\newcommand{\RcPhPOneWithinPrimaryHi}{0.88}
\newcommand{\RcPhPOneWithinCoverage}{0.94}
\newcommand{\RcPhPOneWithinCoverageLo}{0.85}
\newcommand{\RcPhPOneWithinCoverageHi}{1.00}
\newcommand{\RcPhFamLcPrimary}{0.88}
\newcommand{\RcPhFamLcCoverage}{0.85}
\newcommand{\RcPhFamFrozenPrimary}{0.84}
\newcommand{\RcPhFamFrozenCoverage}{0.96}
\newcommand{\RcPhFamPOnePrimary}{0.92}
\newcommand{\RcPhFamPOneCoverage}{0.98}
\newcommand{\RcPhFamColonPrimary}{0.80}
\newcommand{\RcPhFamColonCoverage}{0.94}
\newcommand{\RcPhLabRna}{93}
\newcommand{\RcPhLabProtein}{94}
\newcommand{\RcPhLabAgree}{92}
\newcommand{\RcPhLabErr}{0.21}
\newcommand{\RcPhLabErrJoint}{0.21}
\newcommand{\RcPhLabDiff}{\ensuremath{-}0.002}
\newcommand{\RcPhLabDiffLo}{\ensuremath{-}0.009}
\newcommand{\RcPhLabDiffHi}{0.002}
\newcommand{\RcPhLabBetter}{12}
\newcommand{\RcPhLabAucAll}{0.87}
\newcommand{\RcPhLabAucHi}{0.99}
\newcommand{\RcPhLabAucDiff}{0.000}
\newcommand{\RcSimRows}{144}
\newcommand{\RcSimNuBiasMax}{0.02}
\newcommand{\RcSimNuMaeMax}{0.07}
\newcommand{\RcSimNuMaeMin}{0.03}
\newcommand{\RcSimCouplingMax}{6}
\newcommand{\RcSimErrMeasMin}{0.12}
\newcommand{\RcSimErrMeasMax}{1.01}
\newcommand{\RcSimErrCorrMin}{0.06}
\newcommand{\RcSimErrCorrMax}{0.38}
\newcommand{\RcSimCorrRatioMin}{0.22}
\newcommand{\RcSimCorrRatioMax}{0.93}
\newcommand{\RcSimMeasRatioMin}{1.33}
\newcommand{\RcSimMeasRatioMax}{3.03}
\newcommand{\RcSimFrozenRatioMin}{0.22}
\newcommand{\RcSimFrozenRatioMax}{0.63}
\newcommand{\RcSimBracket}{431}
\newcommand{\RcSimBracketTotal}{432}
\newcommand{\RcSimFrozenTrue}{1.51}
\newcommand{\RcSimFrozenMeas}{2.62}
\newcommand{\RcSimFrozenCorr}{0.43}
\newcommand{\RcSimFrozenTrueAbove}{24}
\newcommand{\RcSimFrozenCorrAbove}{2}
\newcommand{\RcSimDOne}{24}
\newcommand{\RcSimOracleMin}{0.96}
\newcommand{\RcSimOracleMax}{0.97}
\newcommand{\RcSimPOneCorrAbove}{141}
\newcommand{\RcSimPOneTrueAbove}{144}

% Generated by extension/perturbation/make_assets.py. Do not edit.
\newcommand{\PtCells}{218,331}
\newcommand{\PtTargets}{248}
\newcommand{\PtHeldTargets}{62}
\newcommand{\PtTrainTargets}{186}
\newcommand{\PtSingleCells}{111,453}
\newcommand{\PtNTCells}{16,319}
\newcommand{\PtGenes}{200}
\newcommand{\PtEncodingGenes}{19}
\newcommand{\PtProteins}{20}
\newcommand{\PtTrainPops}{505}
\newcommand{\PtGroups}{141}
\newcommand{\PtHeldEligible}{51}
\newcommand{\PtDimS}{1}
\newcommand{\PtDimSPooled}{2}
\newcommand{\PtRfPrimary}{6.8}
\newcommand{\PtRfPrimaryLo}{4.1}
\newcommand{\PtRfPrimaryHi}{9.4}
\newcommand{\PtRfPrimaryLatent}{7.0}
\newcommand{\PtRfPrimaryLatentLo}{4.4}
\newcommand{\PtRfPrimaryLatentHi}{9.5}
\newcommand{\PtRfRegression}{9.2}
\newcommand{\PtRfRegressionLo}{7.9}
\newcommand{\PtRfRegressionHi}{10.5}
\newcommand{\PtRfCondMarg}{15}
\newcommand{\PtRfCondMargLo}{13}
\newcommand{\PtRfCondMargHi}{17}
\newcommand{\PtRfPooled}{\ensuremath{-}36}
\newcommand{\PtRfPooledLo}{\ensuremath{-}42}
\newcommand{\PtRfPooledHi}{\ensuremath{-}32}
\newcommand{\PtRfPaired}{37}
\newcommand{\PtRfPairedLo}{32}
\newcommand{\PtRfPairedHi}{42}
\newcommand{\PtRfPairedLatent}{40}
\newcommand{\PtRfPairedLatentLo}{34}
\newcommand{\PtRfPairedLatentHi}{44}
\newcommand{\PtRfRefReg}{53}
\newcommand{\PtRfRefRegLo}{49}
\newcommand{\PtRfRefRegHi}{57}
\newcommand{\PtRfTransCorr}{97}
\newcommand{\PtRfTransCorrLo}{93}
\newcommand{\PtRfTransCorrHi}{101}
\newcommand{\PtRfBench}{24}
\newcommand{\PtRfBenchLo}{20}
\newcommand{\PtRfBenchHi}{28}
\newcommand{\PtRfPseudo}{\ensuremath{-}0.1}
\newcommand{\PtRfPseudoLo}{\ensuremath{-}0.1}
\newcommand{\PtRfPseudoHi}{\ensuremath{-}0.1}
\newcommand{\PtRfDerange}{\ensuremath{-}0.4}
\newcommand{\PtRfDerangeLo}{\ensuremath{-}0.6}
\newcommand{\PtRfDerangeHi}{\ensuremath{-}0.3}
\newcommand{\PtRfPooledPseudo}{\ensuremath{-}46}
\newcommand{\PtRfPooledPseudoLo}{\ensuremath{-}52}
\newcommand{\PtRfPooledPseudoHi}{\ensuremath{-}41}
\newcommand{\PtRfPooledDerange}{\ensuremath{-}75}
\newcommand{\PtRfPooledDerangeLo}{\ensuremath{-}82}
\newcommand{\PtRfPooledDerangeHi}{\ensuremath{-}70}
\newcommand{\PtRfDiffPseudo}{6.9}
\newcommand{\PtRfDiffPseudoLo}{4.2}
\newcommand{\PtRfDiffPseudoHi}{9.5}
\newcommand{\PtRfDiffDerange}{7.2}
\newcommand{\PtRfDiffDerangeLo}{4.6}
\newcommand{\PtRfDiffDerangeHi}{9.7}
\newcommand{\PtRfShare}{18}
\newcommand{\PtRfShareLo}{12}
\newcommand{\PtRfShareHi}{25}
\newcommand{\PtRfDiffPooled}{43}
\newcommand{\PtRfDiffPooledLo}{40}
\newcommand{\PtRfDiffPooledHi}{47}
\newcommand{\PtCondAlonePrimary}{11}
\newcommand{\PtNtAlonePrimary}{15}
\newcommand{\PtCondAlonePaired}{24}
\newcommand{\PtNtAlonePaired}{93}
\newcommand{\PtCondAloneTransCorr}{98}
\newcommand{\PtNtAloneTransCorr}{98}
\newcommand{\PtCondIfngPrimary}{16}
\newcommand{\PtNtIfngPrimary}{28}
\newcommand{\PtCondIfngPaired}{43}
\newcommand{\PtNtIfngPaired}{97}
\newcommand{\PtCondIfngTransCorr}{97}
\newcommand{\PtNtIfngTransCorr}{99}
\newcommand{\PtCondCocultPrimary}{\ensuremath{-}5.4}
\newcommand{\PtNtCocultPrimary}{13}
\newcommand{\PtCondCocultPaired}{43}
\newcommand{\PtNtCocultPaired}{94}
\newcommand{\PtCondCocultTransCorr}{96}
\newcommand{\PtNtCocultTransCorr}{98}
\newcommand{\PtNtPrimary}{19}
\newcommand{\PtNtPrimaryLo}{16}
\newcommand{\PtNtPrimaryHi}{21}
\newcommand{\PtNtPaired}{94}
\newcommand{\PtNtPairedLo}{93}
\newcommand{\PtNtPairedHi}{96}
\newcommand{\PtNtPooled}{\ensuremath{-}12}
\newcommand{\PtNtPooledLo}{\ensuremath{-}16}
\newcommand{\PtNtPooledHi}{\ensuremath{-}8}
\newcommand{\PtNtTransCorr}{98}
\newcommand{\PtNtTransCorrLo}{97}
\newcommand{\PtNtTransCorrHi}{100}
\newcommand{\PtNtRefReg}{91}
\newcommand{\PtNtRefRegLo}{90}
\newcommand{\PtNtRefRegHi}{93}
\newcommand{\PtNtBench}{90}
\newcommand{\PtNtBenchLo}{88}
\newcommand{\PtNtBenchHi}{91}
\newcommand{\PtNtShare}{20}
\newcommand{\PtNtShareLo}{17}
\newcommand{\PtNtShareHi}{22}
\newcommand{\PtCoverage}{2.4}
\newcommand{\PtCoverageNt}{3.6}
\newcommand{\PtAdaptMedian}{88}
\newcommand{\PtAdaptNtMin}{1,989}
\newcommand{\PtAdaptNtMax}{3,222}
\newcommand{\PtDevPrimary}{10}
\newcommand{\PtDevPooled}{\ensuremath{-}33}
\newcommand{\PtDevPaired}{38}
\newcommand{\PtHOne}{met}
\newcommand{\PtHThree}{met}
\newcommand{\PtHFour}{not met}
\newcommand{\PtHTwo}{met}
\newcommand{\PgIfnGenes}{34}
\newcommand{\PgIfnProteinCount}{3}
\newcommand{\PgIfnPf}{43}
\newcommand{\PgIfnPfLo}{32}
\newcommand{\PgIfnPfHi}{51}
\newcommand{\PgIfnPaired}{64}
\newcommand{\PgIfnPairedLo}{58}
\newcommand{\PgIfnPairedHi}{70}
\newcommand{\PgIfnTrans}{94}
\newcommand{\PgIfnTransLo}{88}
\newcommand{\PgIfnTransHi}{102}
\newcommand{\PgIfnRefReg}{69}
\newcommand{\PgIfnRefRegLo}{63}
\newcommand{\PgIfnRefRegHi}{74}
\newcommand{\PgIfnBench}{48}
\newcommand{\PgIfnBenchLo}{44}
\newcommand{\PgIfnBenchHi}{53}
\newcommand{\PgIfnRepl}{\ensuremath{-}24}
\newcommand{\PgIfnReplLo}{\ensuremath{-}40}
\newcommand{\PgIfnReplHi}{\ensuremath{-}10}
\newcommand{\PgIfnShare}{66}
\newcommand{\PgIfnShareLo}{54}
\newcommand{\PgIfnShareHi}{77}
\newcommand{\PgIfnShareTrans}{45}
\newcommand{\PgIfnShareTransLo}{35}
\newcommand{\PgIfnShareTransHi}{54}
\newcommand{\PgAlongPf}{27}
\newcommand{\PgAlongPfLo}{\ensuremath{-}1}
\newcommand{\PgAlongPfHi}{50}
\newcommand{\PgAlongPaired}{69}
\newcommand{\PgAlongPairedLo}{56}
\newcommand{\PgAlongPairedHi}{84}
\newcommand{\PgAlongShare}{39}
\newcommand{\PgAlongShareLo}{\ensuremath{-}1}
\newcommand{\PgAlongShareHi}{65}
\newcommand{\PgCompPf}{5.7}
\newcommand{\PgCompPfLo}{3.9}
\newcommand{\PgCompPfHi}{7.3}
\newcommand{\PgCompPaired}{36}
\newcommand{\PgCompPairedLo}{29}
\newcommand{\PgCompPairedHi}{43}
\newcommand{\PgCompTrans}{98}
\newcommand{\PgCompTransLo}{92}
\newcommand{\PgCompTransHi}{104}
\newcommand{\PgDiff}{37}
\newcommand{\PgDiffLo}{26}
\newcommand{\PgDiffHi}{46}
\newcommand{\PgIfnGenesPf}{16}
\newcommand{\PgIfnGenesPfLo}{6}
\newcommand{\PgIfnGenesPfHi}{26}
\newcommand{\PgIfnGenesPaired}{33}
\newcommand{\PgIfnGenesPairedLo}{26}
\newcommand{\PgIfnGenesPairedHi}{41}
\newcommand{\PgIfnRel}{0.29}
\newcommand{\PgCompRel}{0.11}
\newcommand{\PgIfnAlonePf}{19}
\newcommand{\PgIfnAlonePaired}{16}
\newcommand{\PgIfnIfngPf}{65}
\newcommand{\PgIfnIfngPaired}{79}
\newcommand{\PgIfnCocultPf}{19}
\newcommand{\PgIfnCocultPaired}{62}
\newcommand{\PgNtIfnPf}{61}
\newcommand{\PgNtCompPf}{12}
\newcommand{\PgNtIfnPaired}{93}
\newcommand{\PgNtCompPaired}{95}
\newcommand{\PgNtIfnTrans}{98}
\newcommand{\PgNtCompTrans}{99}
\newcommand{\PgNtIfnRepl}{94}
\newcommand{\PgNtCompRepl}{88}
\newcommand{\PgNtIfnRel}{0.95}
\newcommand{\PgOtherMax}{7.8}
\newcommand{\PgOtherCount}{5}
\newcommand{\PgSupportIfn}{0.094}
\newcommand{\PgSupportOtherMax}{0.021}
\newcommand{\PgPOne}{met}
\newcommand{\PgPTwo}{not met}
\newcommand{\PgPThree}{met}
\newcommand{\DsNoneAll}{6.8}
\newcommand{\DsNoneIfn}{43}
\newcommand{\DsFiveAllExp}{17}
\newcommand{\DsFiveAllStr}{3.7}
\newcommand{\DsFiveAllRand}{1.3}
\newcommand{\DsFiveAllRandLow}{\ensuremath{-}24}
\newcommand{\DsFiveAllDiff}{\ensuremath{-}15}
\newcommand{\DsFiveAllDiffLo}{\ensuremath{-}17}
\newcommand{\DsFiveAllDiffHi}{\ensuremath{-}13}
\newcommand{\DsFiveAllBelow}{50}
\newcommand{\DsFiveAllDiffStr}{\ensuremath{-}2.4}
\newcommand{\DsFiveAllDiffStrLo}{\ensuremath{-}5.3}
\newcommand{\DsFiveAllDiffStrHi}{0.1}
\newcommand{\DsFiveIfnExp}{41}
\newcommand{\DsFiveIfnStr}{7.9}
\newcommand{\DsFiveIfnRand}{36}
\newcommand{\DsFiveIfnRandLow}{6.0}
\newcommand{\DsFiveIfnDiff}{\ensuremath{-}4.1}
\newcommand{\DsFiveIfnDiffLo}{\ensuremath{-}12.8}
\newcommand{\DsFiveIfnDiffHi}{3.4}
\newcommand{\DsFiveIfnBelow}{9}
\newcommand{\DsFiveIfnDiffStr}{29}
\newcommand{\DsFiveIfnDiffStrLo}{18}
\newcommand{\DsFiveIfnDiffStrHi}{38}
\newcommand{\DsFiveDimExp}{0}
\newcommand{\DsFiveDimRand}{2.5}
\newcommand{\DsFiveZeroRand}{0}
\newcommand{\DsTenAllExp}{3.5}
\newcommand{\DsTenAllStr}{3.2}
\newcommand{\DsTenAllRand}{2.3}
\newcommand{\DsTenAllRandLow}{\ensuremath{-}24}
\newcommand{\DsTenAllDiff}{\ensuremath{-}1.2}
\newcommand{\DsTenAllDiffLo}{\ensuremath{-}4.2}
\newcommand{\DsTenAllDiffHi}{1.3}
\newcommand{\DsTenAllBelow}{7}
\newcommand{\DsTenAllDiffStr}{\ensuremath{-}0.9}
\newcommand{\DsTenAllDiffStrLo}{\ensuremath{-}4.0}
\newcommand{\DsTenAllDiffStrHi}{1.6}
\newcommand{\DsTenIfnExp}{8.0}
\newcommand{\DsTenIfnStr}{4.7}
\newcommand{\DsTenIfnRand}{38}
\newcommand{\DsTenIfnRandLow}{7.6}
\newcommand{\DsTenIfnDiff}{30}
\newcommand{\DsTenIfnDiffLo}{19}
\newcommand{\DsTenIfnDiffHi}{39}
\newcommand{\DsTenIfnBelow}{3}
\newcommand{\DsTenIfnDiffStr}{33}
\newcommand{\DsTenIfnDiffStrLo}{22}
\newcommand{\DsTenIfnDiffStrHi}{42}
\newcommand{\DsTenDimExp}{0}
\newcommand{\DsTenDimRand}{2.1}
\newcommand{\DsTenZeroRand}{0}
\newcommand{\DsTwentyAllExp}{3.5}
\newcommand{\DsTwentyAllStr}{3.4}
\newcommand{\DsTwentyAllRand}{2.5}
\newcommand{\DsTwentyAllRandLow}{\ensuremath{-}26}
\newcommand{\DsTwentyAllDiff}{\ensuremath{-}1.0}
\newcommand{\DsTwentyAllDiffLo}{\ensuremath{-}3.8}
\newcommand{\DsTwentyAllDiffHi}{1.5}
\newcommand{\DsTwentyAllBelow}{7}
\newcommand{\DsTwentyAllDiffStr}{\ensuremath{-}0.9}
\newcommand{\DsTwentyAllDiffStrLo}{\ensuremath{-}3.7}
\newcommand{\DsTwentyAllDiffStrHi}{1.7}
\newcommand{\DsTwentyIfnExp}{8.4}
\newcommand{\DsTwentyIfnStr}{5.1}
\newcommand{\DsTwentyIfnRand}{38}
\newcommand{\DsTwentyIfnRandLow}{5.2}
\newcommand{\DsTwentyIfnDiff}{29}
\newcommand{\DsTwentyIfnDiffLo}{19}
\newcommand{\DsTwentyIfnDiffHi}{39}
\newcommand{\DsTwentyIfnBelow}{5}
\newcommand{\DsTwentyIfnDiffStr}{33}
\newcommand{\DsTwentyIfnDiffStrLo}{22}
\newcommand{\DsTwentyIfnDiffStrHi}{42}
\newcommand{\DsTwentyDimExp}{0}
\newcommand{\DsTwentyDimRand}{2.0}
\newcommand{\DsTwentyZeroRand}{0}
\newcommand{\DsTopFour}{JAK2, IFNGR1, JAK1 and IFNGR2}
\newcommand{\DsTopFive}{JAK2, IFNGR1, JAK1, IFNGR2 and TMED10}
\newcommand{\DsTopStrengthFive}{JAK2, JAK1, IFNGR1, IFNGR2 and CTPS1}
\newcommand{\DsExposureGap}{6}
\newcommand{\DsLooSpearman}{\ensuremath{-}0.09}
\newcommand{\DsLooP}{0.87}
\newcommand{\DsLooTop}{JAK2 and JAK1}
\newcommand{\DsDraws}{50}
\newcommand{\DsDOne}{not met}
\newcommand{\PtTopGenes}{B2M, HLA-B, CD74, GBP1, HLA-DRA, HLA-A}

% Generated by extension/perturbation_replication/make_assets.py. Do not edit.
\newcommand{\RpCells}{20,652}
\newcommand{\RpTargets}{25}
\newcommand{\RpGroups}{55}
\newcommand{\RpTestTargets}{21}
\newcommand{\RpReplicates}{3}
\newcommand{\RpGenes}{200}
\newcommand{\RpProteins}{4}
\newcommand{\RpIfnGenes}{16}
\newcommand{\RpDimSMedian}{18}
\newcommand{\RpDimSMin}{18}
\newcommand{\RpDimSMax}{61}
\newcommand{\RpPrimary}{\ensuremath{-}212}
\newcommand{\RpPrimaryLo}{\ensuremath{-}287}
\newcommand{\RpPrimaryHi}{\ensuremath{-}157}
\newcommand{\RpDiffPseudo}{\ensuremath{-}202}
\newcommand{\RpDiffPseudoLo}{\ensuremath{-}275}
\newcommand{\RpDiffPseudoHi}{\ensuremath{-}151}
\newcommand{\RpDiffDerange}{\ensuremath{-}167}
\newcommand{\RpDiffDerangeLo}{\ensuremath{-}234}
\newcommand{\RpDiffDerangeHi}{\ensuremath{-}119}
\newcommand{\RpIfnShare}{\ensuremath{-}3.5}
\newcommand{\RpIfnShareLo}{\ensuremath{-}41.3}
\newcommand{\RpIfnShareHi}{24.3}
\newcommand{\RpRemDiff}{\ensuremath{-}44}
\newcommand{\RpRemDiffLo}{\ensuremath{-}91}
\newcommand{\RpRemDiffHi}{\ensuremath{-}3}
\newcommand{\RpPaired}{32}
\newcommand{\RpPairedLo}{21}
\newcommand{\RpPairedHi}{43}
\newcommand{\RpIfnPrimary}{\ensuremath{-}3.0}
\newcommand{\RpIfnPrimaryLo}{\ensuremath{-}35.6}
\newcommand{\RpIfnPrimaryHi}{19.5}
\newcommand{\RpIfnPaired}{84}
\newcommand{\RpIfnPairedLo}{73}
\newcommand{\RpIfnPairedHi}{100}
\newcommand{\RpShare}{\ensuremath{-}656}
\newcommand{\RpShareLo}{\ensuremath{-}1195}
\newcommand{\RpShareHi}{\ensuremath{-}429}
\newcommand{\RpIfnRemDiff}{1.9}
\newcommand{\RpIfnRemDiffLo}{\ensuremath{-}21.1}
\newcommand{\RpIfnRemDiffHi}{20.2}
\newcommand{\RpAlongPrimary}{\ensuremath{-}124}
\newcommand{\RpAlongPrimaryLo}{\ensuremath{-}188}
\newcommand{\RpAlongPrimaryHi}{\ensuremath{-}75}
\newcommand{\RpAlongShare}{\ensuremath{-}171}
\newcommand{\RpAlongShareLo}{\ensuremath{-}266}
\newcommand{\RpAlongShareHi}{\ensuremath{-}107}
\newcommand{\RpIfnTrans}{94}
\newcommand{\RpIfnTransLo}{80}
\newcommand{\RpIfnTransHi}{115}
\newcommand{\RpIfnReplicate}{19}
\newcommand{\RpIfnReplicateLo}{\ensuremath{-}2}
\newcommand{\RpIfnReplicateHi}{36}
\newcommand{\RpAllPf}{\ensuremath{-}212}
\newcommand{\RpAllPaired}{32}
\newcommand{\RpAllTrans}{81}
\newcommand{\RpAllRefReg}{37}
\newcommand{\RpAllBench}{22}
\newcommand{\RpAllRepl}{\ensuremath{-}205}
\newcommand{\RpAllPseudo}{\ensuremath{-}10}
\newcommand{\RpAllDerange}{\ensuremath{-}46}
\newcommand{\RpAllRemExp}{\ensuremath{-}172}
\newcommand{\RpAllRemStr}{\ensuremath{-}219}
\newcommand{\RpAllRemRand}{\ensuremath{-}216}
\newcommand{\RpAllRel}{0.14}
\newcommand{\RpIfnPf}{\ensuremath{-}3.0}
\newcommand{\RpIfnRefReg}{84}
\newcommand{\RpIfnBench}{61}
\newcommand{\RpIfnRepl}{19}
\newcommand{\RpIfnPseudo}{\ensuremath{-}3.2}
\newcommand{\RpIfnDerange}{\ensuremath{-}12}
\newcommand{\RpIfnRemExp}{\ensuremath{-}5.6}
\newcommand{\RpIfnRemStr}{\ensuremath{-}23}
\newcommand{\RpIfnRemRand}{\ensuremath{-}3.7}
\newcommand{\RpIfnRel}{0.34}
\newcommand{\RpCompPf}{\ensuremath{-}201}
\newcommand{\RpCompPaired}{25}
\newcommand{\RpCompTrans}{84}
\newcommand{\RpCompRefReg}{31}
\newcommand{\RpCompBench}{3.1}
\newcommand{\RpCompRepl}{\ensuremath{-}390}
\newcommand{\RpCompPseudo}{\ensuremath{-}16}
\newcommand{\RpCompDerange}{\ensuremath{-}54}
\newcommand{\RpCompRemExp}{\ensuremath{-}164}
\newcommand{\RpCompRemStr}{\ensuremath{-}173}
\newcommand{\RpCompRemRand}{\ensuremath{-}227}
\newcommand{\RpCompRel}{0.09}
\newcommand{\RpAlongPf}{\ensuremath{-}124}
\newcommand{\RpAlongPaired}{72}
\newcommand{\RpAlongTrans}{87}
\newcommand{\RpAlongRefReg}{73}
\newcommand{\RpAlongBench}{50}
\newcommand{\RpAlongRepl}{\ensuremath{-}5.0}
\newcommand{\RpAlongPseudo}{\ensuremath{-}8.1}
\newcommand{\RpAlongDerange}{\ensuremath{-}37}
\newcommand{\RpAlongRemExp}{\ensuremath{-}114}
\newcommand{\RpAlongRemStr}{\ensuremath{-}155}
\newcommand{\RpAlongRemRand}{\ensuremath{-}123}
\newcommand{\RpAlongRel}{0.30}
\newcommand{\RpRepOnePf}{\ensuremath{-}98}
\newcommand{\RpRepOneIfnPf}{25}
\newcommand{\RpRepOneIfnPaired}{87}
\newcommand{\RpRepThreePf}{\ensuremath{-}298}
\newcommand{\RpRepThreeIfnPf}{8.6}
\newcommand{\RpRepThreeIfnPaired}{74}
\newcommand{\RpRepFourPf}{\ensuremath{-}393}
\newcommand{\RpRepFourIfnPf}{\ensuremath{-}129}
\newcommand{\RpRepFourIfnPaired}{93}
\newcommand{\RpNtPrimary}{\ensuremath{-}90}
\newcommand{\RpNtPrimaryLo}{\ensuremath{-}140}
\newcommand{\RpNtPrimaryHi}{\ensuremath{-}52}
\newcommand{\RpNtPaired}{82}
\newcommand{\RpNtIfnPf}{21}
\newcommand{\RpNtIfnPaired}{89}
\newcommand{\RpNtIfnRepl}{88}
\newcommand{\RpTopRemoved}{CUL3, STAT1, SMAD4, IFNGR2, JAK2}
\newcommand{\RpFixedPointOneAll}{\ensuremath{-}76}
\newcommand{\RpFixedPointOneIfn}{32}
\newcommand{\RpFixedPointOneDim}{1}
\newcommand{\RpFixedOneAll}{4.6}
\newcommand{\RpFixedOneIfn}{45}
\newcommand{\RpFixedOneDim}{0}
\newcommand{\RpROne}{not met}
\newcommand{\RpRThree}{not met}
\newcommand{\RpRFour}{not met}
\newcommand{\RpRFive}{not met}
\newcommand{\RpRTwo}{not met}

% Generated by extension/semipaired/make_assets.py. Do not edit.
\newcommand{\HyFrBlockHundred}{38}
\newcommand{\HyFrCurrent}{3.6}
\newcommand{\HyFrHybridHundred}{35}
\newcommand{\HyFrHybridPcHundred}{36}
\newcommand{\HyFrNoiseAware}{\ensuremath{-}12}
\newcommand{\HyFrPairedJsHundred}{16}
\newcommand{\HyPaBlockHundred}{40}
\newcommand{\HyPaCurrent}{\ensuremath{-}154}
\newcommand{\HyPaHybridHundred}{25}
\newcommand{\HyPaHybridPcHundred}{38}
\newcommand{\HyPaNoiseAware}{\ensuremath{-}314}
\newcommand{\HyPaPairedJsHundred}{22}
\newcommand{\SpBoundBelowIdealized}{yes}
\newcommand{\SpBoundBudgetMaxFr}{800}
\newcommand{\SpBoundBudgetMaxOc}{400}
\newcommand{\SpBoundBudgetMaxPa}{800}
\newcommand{\SpBoundCellsFr}{21,171}
\newcommand{\SpBoundCellsOc}{457}
\newcommand{\SpBoundCellsPa}{2,256}
\newcommand{\SpBoundConditionFr}{yes}
\newcommand{\SpBoundConditionOc}{yes}
\newcommand{\SpBoundConditionPa}{yes}
\newcommand{\SpBoundCoveredFr}{yes}
\newcommand{\SpBoundCoveredOc}{yes}
\newcommand{\SpBoundCoveredPa}{yes}
\newcommand{\SpBoundDraws}{200}
\newcommand{\SpBoundEffDimMin}{4.4}
\newcommand{\SpBoundRatioMax}{3.3}
\newcommand{\SpBoundRatioMin}{1.3}
\newcommand{\SpDevChampText}{In Frangieh, Champollion recovered 40\% at 100 paired cells against 48\% for the proposed estimator, and the proposed estimator needed 1.4 to 1.7 times fewer paired cells to reach recovered fractions of 0.3 to 0.6 (lowest 95\% bound, 1.4). In Papalexi (7 held-out targets, 21 groups), Champollion recovered 41\% at 100 paired cells against 29\% for the proposed estimator, so it reached recovered fractions of 0.3 and 0.4 with fewer paired cells (ratios 0.55 and 0.69; 95\% CI, 0.35--0.79 and 0.35--1.00), and the two were similar at 0.5 and 0.6 (0.91 and 0.97). Champollion's lasso weight was chosen per budget on the evaluation groups, which favours it.}
\newcommand{\SpEffDimBudgetFr}{1,600}
\newcommand{\SpEffDimBudgetPa}{1,600}
\newcommand{\SpEffDimMin}{4.2}
\newcommand{\SpEffDimSmallMin}{3.3}
\newcommand{\SpExBestDiffHundred}{8.1}
\newcommand{\SpExBestDiffHundredHi}{13}
\newcommand{\SpExBestDiffHundredLo}{2.4}
\newcommand{\SpExBestHundred}{ridge\_p/mapped}
\newcommand{\SpExCells}{908}
\newcommand{\SpExCellsAll}{111}
\newcommand{\SpExCellsAllFour}{186}
\newcommand{\SpExCellsAllTwo}{58}
\newcommand{\SpExCellsChamp}{131}
\newcommand{\SpExCellsChampFour}{190}
\newcommand{\SpExCellsChampTwo}{87}
\newcommand{\SpExCellsOriginal}{115}
\newcommand{\SpExCellsOriginalFour}{187}
\newcommand{\SpExCellsOriginalTwo}{61}
\newcommand{\SpExCellsPairedOnly}{357}
\newcommand{\SpExCellsPairedOnlyFour}{400}
\newcommand{\SpExCellsPairedOnlyTwo}{196}
\newcommand{\SpExCellsProposed}{73}
\newcommand{\SpExCellsProposedRel}{}
\newcommand{\SpExCellsSemi}{122}
\newcommand{\SpExCellsSemiFour}{236}
\newcommand{\SpExCellsSemiTwo}{64}
\newcommand{\SpExChampFifty}{8.6}
\newcommand{\SpExChampFourHundred}{54}
\newcommand{\SpExChampHundred}{23}
\newcommand{\SpExChampTwentyFive}{0.9}
\newcommand{\SpExChampTwoHundred}{41}
\newcommand{\SpExCrossPairedOnlyHiBudget}{400}
\newcommand{\SpExCrossPairedOnlyHiRf}{32}
\newcommand{\SpExCrossPairedOnlyLoBudget}{200}
\newcommand{\SpExCrossPairedOnlyLoRf}{20}
\newcommand{\SpExCrossProposedHiBudget}{100}
\newcommand{\SpExCrossProposedHiRf}{36}
\newcommand{\SpExCrossProposedLoBudget}{50}
\newcommand{\SpExCrossProposedLoRf}{23}
\newcommand{\SpExFrozenArmAllMet}{yes}
\newcommand{\SpExFrozenArmCells}{101}
\newcommand{\SpExFrozenArmChamp}{1.3}
\newcommand{\SpExFrozenArmChampHi}{1.5}
\newcommand{\SpExFrozenArmChampLo}{0.96}
\newcommand{\SpExFrozenArmPairedOnly}{3.5}
\newcommand{\SpExFrozenArmPairedOnlyHi}{4.1}
\newcommand{\SpExFrozenArmPairedOnlyLo}{2.6}
\newcommand{\SpExFrozenArmSemi}{1.2}
\newcommand{\SpExFrozenArmSemiHi}{1.3}
\newcommand{\SpExFrozenArmSemiLo}{1.09}
\newcommand{\SpExFrozenAt}{14:49 EDT on 27 September 2026}
\newcommand{\SpExHOne}{met}
\newcommand{\SpExHThree}{met}
\newcommand{\SpExHTwo}{met}
\newcommand{\SpExLossBestOtherHundred}{0.85}
\newcommand{\SpExLossChampHundred}{0.87}
\newcommand{\SpExLossFifty}{0.87}
\newcommand{\SpExLossFourHundred}{0.68}
\newcommand{\SpExLossHundred}{0.80}
\newcommand{\SpExLossIndependence}{1.00}
\newcommand{\SpExLossPairedOnlyHundred}{0.93}
\newcommand{\SpExLossSemiHundred}{0.85}
\newcommand{\SpExLossTwentyFive}{0.97}
\newcommand{\SpExLossTwoHundred}{0.74}
\newcommand{\SpExOrfs}{11}
\newcommand{\SpExPairedAll}{64}
\newcommand{\SpExPairedOnlyFifty}{5.7}
\newcommand{\SpExPairedOnlyFourHundred}{32}
\newcommand{\SpExPairedOnlyHundred}{12}
\newcommand{\SpExPairedOnlyTwentyFive}{2.6}
\newcommand{\SpExPairedOnlyTwoHundred}{20}
\newcommand{\SpExPoolPops}{28}
\newcommand{\SpExReservoir}{457}
\newcommand{\SpExRfFifty}{23}
\newcommand{\SpExRfFiftyHi}{28}
\newcommand{\SpExRfFiftyLo}{15}
\newcommand{\SpExRfFourHundred}{57}
\newcommand{\SpExRfFourHundredHi}{68}
\newcommand{\SpExRfFourHundredLo}{47}
\newcommand{\SpExRfHundred}{36}
\newcommand{\SpExRfHundredHi}{44}
\newcommand{\SpExRfHundredLo}{27}
\newcommand{\SpExRfTwentyFive}{5.2}
\newcommand{\SpExRfTwentyFiveHi}{11}
\newcommand{\SpExRfTwentyFiveLo}{\ensuremath{-}6.6}
\newcommand{\SpExRfTwoHundred}{47}
\newcommand{\SpExRfTwoHundredHi}{56}
\newcommand{\SpExRfTwoHundredLo}{38}
\newcommand{\SpExSaveAll}{1.5}
\newcommand{\SpExSaveAllFour}{1.4}
\newcommand{\SpExSaveAllFourHi}{1.9}
\newcommand{\SpExSaveAllFourLo}{1.02}
\newcommand{\SpExSaveAllFourRel}{}
\newcommand{\SpExSaveAllHi}{1.8}
\newcommand{\SpExSaveAllLo}{1.11}
\newcommand{\SpExSaveAllRel}{}
\newcommand{\SpExSaveAllTwo}{1.3}
\newcommand{\SpExSaveAllTwoHi}{1.5}
\newcommand{\SpExSaveAllTwoLo}{1.02}
\newcommand{\SpExSaveAllTwoRel}{}
\newcommand{\SpExSaveChamp}{1.8}
\newcommand{\SpExSaveChampFour}{1.5}
\newcommand{\SpExSaveChampFourHi}{2.0}
\newcommand{\SpExSaveChampFourLo}{1.02}
\newcommand{\SpExSaveChampFourRel}{}
\newcommand{\SpExSaveChampHi}{2.2}
\newcommand{\SpExSaveChampLo}{1.2}
\newcommand{\SpExSaveChampRel}{}
\newcommand{\SpExSaveChampTwo}{1.9}
\newcommand{\SpExSaveChampTwoHi}{2.4}
\newcommand{\SpExSaveChampTwoLo}{1.4}
\newcommand{\SpExSaveChampTwoRel}{}
\newcommand{\SpExSaveOriginal}{1.6}
\newcommand{\SpExSaveOriginalFour}{1.5}
\newcommand{\SpExSaveOriginalFourHi}{1.9}
\newcommand{\SpExSaveOriginalFourLo}{1.02}
\newcommand{\SpExSaveOriginalFourRel}{}
\newcommand{\SpExSaveOriginalHi}{1.9}
\newcommand{\SpExSaveOriginalLo}{1.12}
\newcommand{\SpExSaveOriginalRel}{}
\newcommand{\SpExSaveOriginalTwo}{1.4}
\newcommand{\SpExSaveOriginalTwoHi}{1.6}
\newcommand{\SpExSaveOriginalTwoLo}{1.03}
\newcommand{\SpExSaveOriginalTwoRel}{}
\newcommand{\SpExSavePairedOnly}{4.9}
\newcommand{\SpExSavePairedOnlyFour}{3.1}
\newcommand{\SpExSavePairedOnlyFourHi}{4.7}
\newcommand{\SpExSavePairedOnlyFourLo}{1.7}
\newcommand{\SpExSavePairedOnlyFourRel}{more than }
\newcommand{\SpExSavePairedOnlyHi}{6.0}
\newcommand{\SpExSavePairedOnlyLo}{3.2}
\newcommand{\SpExSavePairedOnlyRel}{}
\newcommand{\SpExSavePairedOnlyTwo}{4.4}
\newcommand{\SpExSavePairedOnlyTwoHi}{5.2}
\newcommand{\SpExSavePairedOnlyTwoLo}{3.2}
\newcommand{\SpExSavePairedOnlyTwoRel}{}
\newcommand{\SpExSaveSemi}{1.7}
\newcommand{\SpExSaveSemiFour}{1.8}
\newcommand{\SpExSaveSemiFourHi}{2.5}
\newcommand{\SpExSaveSemiFourLo}{1.5}
\newcommand{\SpExSaveSemiFourRel}{}
\newcommand{\SpExSaveSemiHi}{2.0}
\newcommand{\SpExSaveSemiLo}{1.3}
\newcommand{\SpExSaveSemiRel}{}
\newcommand{\SpExSaveSemiTwo}{1.4}
\newcommand{\SpExSaveSemiTwoHi}{1.7}
\newcommand{\SpExSaveSemiTwoLo}{1.2}
\newcommand{\SpExSaveSemiTwoRel}{}
\newcommand{\SpExSecondsChamp}{8.7}
\newcommand{\SpExSecondsChampFit}{5.5}
\newcommand{\SpExSecondsProposed}{14}
\newcommand{\SpExSemiFifty}{16}
\newcommand{\SpExSemiFourHundred}{47}
\newcommand{\SpExSemiHundred}{27}
\newcommand{\SpExSemiTwentyFive}{12}
\newcommand{\SpExSemiTwoHundred}{38}
\newcommand{\SpExTarget}{0.3}
\newcommand{\SpFixedCriterion}{60.8}
\newcommand{\SpFixedFrFifty}{37}
\newcommand{\SpFixedPaFifty}{16}
\newcommand{\SpFrBestDiffHundred}{16}
\newcommand{\SpFrBestDiffHundredHi}{17}
\newcommand{\SpFrBestDiffHundredLo}{15}
\newcommand{\SpFrBestWins}{7}
\newcommand{\SpFrBudgets}{7}
\newcommand{\SpFrCellsFour}{59}
\newcommand{\SpFrChampBudgets}{4}
\newcommand{\SpFrChampFolds}{1}
\newcommand{\SpFrChampGroups}{141}
\newcommand{\SpFrChampHundred}{40}
\newcommand{\SpFrChampPropHundred}{48}
\newcommand{\SpFrChampSaveFour}{1.7}
\newcommand{\SpFrChampSaveFourHi}{1.9}
\newcommand{\SpFrChampSaveFourLo}{1.5}
\newcommand{\SpFrChampSaveMax}{1.7}
\newcommand{\SpFrChampSaveMin}{1.4}
\newcommand{\SpFrChampWins}{4}
\newcommand{\SpFrDoseEightThousandHundred}{47}
\newcommand{\SpFrDoseFiveHundredHundred}{42}
\newcommand{\SpFrDoseTwoThousandHundred}{45}
\newcommand{\SpFrGroups}{141}
\newcommand{\SpFrLossLowestEverywhere}{yes}
\newcommand{\SpFrMapGain}{3.5}
\newcommand{\SpFrMapMeasuredHundred}{43}
\newcommand{\SpFrMapPairedHundred}{\ensuremath{-}398}
\newcommand{\SpFrMarkerHundred}{15}
\newcommand{\SpFrOneSidedHundred}{39}
\newcommand{\SpFrPairBasisHundred}{11}
\newcommand{\SpFrPairedAll}{97}
\newcommand{\SpFrPairedOnlyCellsFour}{347}
\newcommand{\SpFrPairedOnlyFourHundred}{43}
\newcommand{\SpFrPairedOnlyHundred}{16}
\newcommand{\SpFrPercondHundred}{26}
\newcommand{\SpFrPooledHundred}{45}
\newcommand{\SpFrRfFifty}{37}
\newcommand{\SpFrRfFourHundred}{72}
\newcommand{\SpFrRfHundred}{48}
\newcommand{\SpFrRidgeHundred}{27}
\newcommand{\SpFrSaveBestCiMin}{1.7}
\newcommand{\SpFrSaveBestFour}{2.4}
\newcommand{\SpFrSaveBestFourHi}{2.6}
\newcommand{\SpFrSaveBestFourLo}{2.2}
\newcommand{\SpFrSaveBestMax}{2.4}
\newcommand{\SpFrSaveBestMin}{1.8}
\newcommand{\SpFrSavePairedOnlyCiMin}{4.1}
\newcommand{\SpFrSavePairedOnlyFour}{5.9}
\newcommand{\SpFrSavePairedOnlyFourHi}{6.5}
\newcommand{\SpFrSavePairedOnlyFourLo}{5.3}
\newcommand{\SpFrSavePairedOnlyMax}{5.9}
\newcommand{\SpFrSavePairedOnlyMin}{4.4}
\newcommand{\SpFrSavePublishedCiMin}{1.8}
\newcommand{\SpFrSavePublishedFour}{2.7}
\newcommand{\SpFrSavePublishedFourHi}{2.9}
\newcommand{\SpFrSavePublishedFourLo}{2.4}
\newcommand{\SpFrSavePublishedMax}{2.9}
\newcommand{\SpFrSavePublishedMin}{1.9}
\newcommand{\SpFrSemiccaHundred}{31}
\newcommand{\SpFrTargets}{51}
\newcommand{\SpPaBestDiffHundred}{4.6}
\newcommand{\SpPaBestDiffHundredHi}{11}
\newcommand{\SpPaBestDiffHundredLo}{\ensuremath{-}2.2}
\newcommand{\SpPaBestWins}{0}
\newcommand{\SpPaBudgets}{6}
\newcommand{\SpPaCellsFour}{112}
\newcommand{\SpPaChampBudgets}{4}
\newcommand{\SpPaChampFolds}{7}
\newcommand{\SpPaChampGroups}{21}
\newcommand{\SpPaChampHundred}{41}
\newcommand{\SpPaChampPropHundred}{29}
\newcommand{\SpPaChampSaveFour}{0.69}
\newcommand{\SpPaChampSaveFourHi}{1.00}
\newcommand{\SpPaChampSaveFourLo}{0.35}
\newcommand{\SpPaChampSaveMax}{0.97}
\newcommand{\SpPaChampSaveMin}{0.55}
\newcommand{\SpPaChampWins}{0}
\newcommand{\SpPaDoseEightThousandHundred}{37}
\newcommand{\SpPaDoseFiveHundredHundred}{35}
\newcommand{\SpPaDoseTwoThousandHundred}{37}
\newcommand{\SpPaGroups}{55}
\newcommand{\SpPaLossLowestEverywhere}{no}
\newcommand{\SpPaMapGain}{1.4}
\newcommand{\SpPaMapMeasuredHundred}{34}
\newcommand{\SpPaMapPairedHundred}{\ensuremath{-}306}
\newcommand{\SpPaMarkerHundred}{19}
\newcommand{\SpPaOneSidedHundred}{37}
\newcommand{\SpPaPairBasisHundred}{16}
\newcommand{\SpPaPairedAll}{81}
\newcommand{\SpPaPairedOnlyCellsFour}{236}
\newcommand{\SpPaPairedOnlyFourHundred}{51}
\newcommand{\SpPaPairedOnlyHundred}{21}
\newcommand{\SpPaPercondHundred}{14}
\newcommand{\SpPaPooledHundred}{36}
\newcommand{\SpPaRfFifty}{16}
\newcommand{\SpPaRfFourHundred}{67}
\newcommand{\SpPaRfHundred}{37}
\newcommand{\SpPaRidgeHundred}{31}
\newcommand{\SpPaSaveBestCiMin}{0.65}
\newcommand{\SpPaSaveBestFour}{1.12}
\newcommand{\SpPaSaveBestFourHi}{1.3}
\newcommand{\SpPaSaveBestFourLo}{0.88}
\newcommand{\SpPaSaveBestMax}{1.12}
\newcommand{\SpPaSaveBestMin}{0.99}
\newcommand{\SpPaSavePairedOnlyCiMin}{1.5}
\newcommand{\SpPaSavePairedOnlyFour}{2.1}
\newcommand{\SpPaSavePairedOnlyFourHi}{2.6}
\newcommand{\SpPaSavePairedOnlyFourLo}{1.7}
\newcommand{\SpPaSavePairedOnlyMax}{2.5}
\newcommand{\SpPaSavePairedOnlyMin}{1.9}
\newcommand{\SpPaSavePublishedCiMin}{0.82}
\newcommand{\SpPaSavePublishedFour}{1.2}
\newcommand{\SpPaSavePublishedFourHi}{1.5}
\newcommand{\SpPaSavePublishedFourLo}{0.94}
\newcommand{\SpPaSavePublishedMax}{1.3}
\newcommand{\SpPaSavePublishedMin}{1.14}
\newcommand{\SpPaSemiccaHundred}{28}
\newcommand{\SpPaTargets}{21}
\newcommand{\SpProvArrays}{200}
\newcommand{\SpProvManifestAt}{15:06:29 EDT}
\newcommand{\SpProvMaxDiff}{0}
\newcommand{\SpSelfCriterion}{61.9}
\newcommand{\SpSelfFrFifty}{41}
\newcommand{\SpSelfFullCriterion}{57.3}
\newcommand{\SpSelfPaFifty}{19}
\newcommand{\SpSelfPartitionCriterion}{57.4}

% Generated by extension/generality/make_assets.py. Do not edit.
\newcommand{\GenCellsBanc}{2,487}
\newcommand{\GenCellsBmCite}{36,000}
\newcommand{\GenCellsBmMulti}{36,450}
\newcommand{\GenCellsColon}{15,953}
\newcommand{\GenCellsGouwens}{3,395}
\newcommand{\GenCellsHao}{72,000}
\newcommand{\GenCellsMaleCns}{3,009}
\newcommand{\GenCellsPairedOnlyBanc}{90}
\newcommand{\GenCellsPairedOnlyBmCite}{$>$800}
\newcommand{\GenCellsPairedOnlyBmMulti}{$>$800}
\newcommand{\GenCellsPairedOnlyColon}{$>$800}
\newcommand{\GenCellsPairedOnlyGouwens}{219}
\newcommand{\GenCellsPairedOnlyHao}{$>$800}
\newcommand{\GenCellsPairedOnlyMaleCns}{73}
\newcommand{\GenCellsPairedOnlyScala}{$\leq$25}
\newcommand{\GenCellsPairedOnlyStephenson}{$>$800}
\newcommand{\GenCellsProposedBanc}{25}
\newcommand{\GenCellsProposedBmCite}{168}
\newcommand{\GenCellsProposedBmMulti}{285}
\newcommand{\GenCellsProposedColon}{126}
\newcommand{\GenCellsProposedGouwens}{157}
\newcommand{\GenCellsProposedHao}{408}
\newcommand{\GenCellsProposedMaleCns}{$\leq$25}
\newcommand{\GenCellsProposedScala}{$>$100}
\newcommand{\GenCellsProposedStephenson}{390}
\newcommand{\GenCellsScala}{1,203}
\newcommand{\GenCellsStephenson}{70,800}
\newcommand{\GenFrozenGOne}{1.5}
\newcommand{\GenFrozenGOneHi}{1.6}
\newcommand{\GenFrozenGOneLo}{1.4}
\newcommand{\GenFrozenGOneMet}{met}
\newcommand{\GenFrozenGTwo}{1.3}
\newcommand{\GenFrozenGTwoHi}{1.4}
\newcommand{\GenFrozenGTwoLo}{1.05}
\newcommand{\GenFrozenGTwoMet}{met}
\newcommand{\GenFrozenNTrivial}{five}
\newcommand{\GenFrozenSaveBanc}{3.7}
\newcommand{\GenFrozenSaveBmCite}{2.5}
\newcommand{\GenFrozenSaveBmMulti}{1.00}
\newcommand{\GenFrozenSaveColon}{1.00}
\newcommand{\GenFrozenSaveGouwens}{1.4}
\newcommand{\GenFrozenSaveHao}{1.00}
\newcommand{\GenFrozenSaveMaleCns}{3.1}
\newcommand{\GenFrozenSaveScala}{1.00}
\newcommand{\GenFrozenSaveStephenson}{1.00}
\newcommand{\GenGOne}{2.9}
\newcommand{\GenGOneAll}{2.2}
\newcommand{\GenGOneAllHi}{2.8}
\newcommand{\GenGOneAllLo}{1.8}
\newcommand{\GenGOneAllMet}{met}
\newcommand{\GenGOneHi}{3.0}
\newcommand{\GenGOneLo}{2.3}
\newcommand{\GenGOneMet}{met}
\newcommand{\GenGOneQuarter}{2.9}
\newcommand{\GenGOneQuarterAll}{2.2}
\newcommand{\GenGOneQuarterAllHi}{2.7}
\newcommand{\GenGOneQuarterAllLo}{1.7}
\newcommand{\GenGOneQuarterHi}{3.1}
\newcommand{\GenGOneQuarterLo}{2.2}
\newcommand{\GenGOneThreeQuarters}{1.9}
\newcommand{\GenGOneThreeQuartersAll}{1.5}
\newcommand{\GenGOneThreeQuartersAllHi}{1.9}
\newcommand{\GenGOneThreeQuartersAllLo}{1.4}
\newcommand{\GenGOneThreeQuartersHi}{2.0}
\newcommand{\GenGOneThreeQuartersLo}{1.7}
\newcommand{\GenGOriginal}{1.4}
\newcommand{\GenGOriginalAll}{1.13}
\newcommand{\GenGOriginalAllHi}{1.3}
\newcommand{\GenGOriginalAllLo}{0.94}
\newcommand{\GenGOriginalHi}{1.4}
\newcommand{\GenGOriginalLo}{1.10}
\newcommand{\GenGThree}{0.58}
\newcommand{\GenGThreeAll}{0.61}
\newcommand{\GenGThreeAllHi}{0.71}
\newcommand{\GenGThreeAllLo}{0.59}
\newcommand{\GenGThreeHi}{0.69}
\newcommand{\GenGThreeLo}{0.55}
\newcommand{\GenGTwo}{2.3}
\newcommand{\GenGTwoAll}{2.1}
\newcommand{\GenGTwoAllHi}{2.2}
\newcommand{\GenGTwoAllLo}{1.6}
\newcommand{\GenGTwoAllMet}{met}
\newcommand{\GenGTwoHi}{2.4}
\newcommand{\GenGTwoLo}{1.8}
\newcommand{\GenGTwoMet}{met}
\newcommand{\GenLawC}{0.92}
\newcommand{\GenLawCP}{0.008}
\newcommand{\GenLawCPairs}{24}
\newcommand{\GenLawFrozenMet}{not met}
\newcommand{\GenLawFrozenN}{12}
\newcommand{\GenLawFrozenP}{0.627}
\newcommand{\GenLawFrozenRho}{-0.09}
\newcommand{\GenLawMet}{not met}
\newcommand{\GenLawN}{11}
\newcommand{\GenLawNCensored}{six}
\newcommand{\GenLawNConsistent}{five}
\newcommand{\GenLawNUncensored}{five}
\newcommand{\GenLawP}{0.104}
\newcommand{\GenLawRatioMax}{1.3}
\newcommand{\GenLawRatioMin}{0.6}
\newcommand{\GenLawRho}{0.42}
\newcommand{\GenLawUncensoredRho}{0.90}
\newcommand{\GenLevelBanc}{68}
\newcommand{\GenLevelBancX}{49}
\newcommand{\GenLevelBmCite}{46}
\newcommand{\GenLevelBmMulti}{25}
\newcommand{\GenLevelColon}{25}
\newcommand{\GenLevelGouwens}{42}
\newcommand{\GenLevelHao}{35}
\newcommand{\GenLevelMaleCns}{69}
\newcommand{\GenLevelScala}{0.0}
\newcommand{\GenLevelStephenson}{13}
\newcommand{\GenMaxBudgetBanc}{400}
\newcommand{\GenMaxBudgetBmCite}{800}
\newcommand{\GenMaxBudgetBmMulti}{800}
\newcommand{\GenMaxBudgetColon}{800}
\newcommand{\GenMaxBudgetGouwens}{400}
\newcommand{\GenMaxBudgetHao}{800}
\newcommand{\GenMaxBudgetMaleCns}{400}
\newcommand{\GenMaxBudgetScala}{100}
\newcommand{\GenMaxBudgetStephenson}{800}
\newcommand{\GenNInformative}{eight}
\newcommand{\GenNLbAboveOne}{five}
\newcommand{\GenNLbAboveOneSemi}{five}
\newcommand{\GenNotInformative}{motor cortex}
\newcommand{\GenPairedOnlyHundredBanc}{36}
\newcommand{\GenPairedOnlyHundredBancX}{\ensuremath{-}52}
\newcommand{\GenPairedOnlyHundredBmCite}{1.1}
\newcommand{\GenPairedOnlyHundredBmMulti}{0.0}
\newcommand{\GenPairedOnlyHundredColon}{1.2}
\newcommand{\GenPairedOnlyHundredGouwens}{12}
\newcommand{\GenPairedOnlyHundredHao}{0.7}
\newcommand{\GenPairedOnlyHundredMaleCns}{40}
\newcommand{\GenPairedOnlyHundredScala}{0.0}
\newcommand{\GenPairedOnlyHundredStephenson}{0.4}
\newcommand{\GenPairedOnlyMaxBanc}{55}
\newcommand{\GenPairedOnlyMaxBmCite}{5.9}
\newcommand{\GenPairedOnlyMaxBmMulti}{0.6}
\newcommand{\GenPairedOnlyMaxColon}{7.7}
\newcommand{\GenPairedOnlyMaxGouwens}{29}
\newcommand{\GenPairedOnlyMaxHao}{2.1}
\newcommand{\GenPairedOnlyMaxMaleCns}{56}
\newcommand{\GenPairedOnlyMaxScala}{0.0}
\newcommand{\GenPairedOnlyMaxStephenson}{3.7}
\newcommand{\GenPropMaxBancX}{\ensuremath{-}30}
\newcommand{\GenPropPooledMaxBancX}{\ensuremath{-}9.1}
\newcommand{\GenRefBanc}{76}
\newcommand{\GenRefBmCite}{66}
\newcommand{\GenRefBmMulti}{\ensuremath{-}148}
\newcommand{\GenRefColon}{\ensuremath{-}237}
\newcommand{\GenRefGouwens}{43}
\newcommand{\GenRefHao}{70}
\newcommand{\GenRefMaleCns}{77}
\newcommand{\GenRefMax}{77}
\newcommand{\GenRefScala}{\ensuremath{-}274}
\newcommand{\GenRefStephenson}{\ensuremath{-}1042}
\newcommand{\GenRfHundredBanc}{56}
\newcommand{\GenRfHundredBancX}{\ensuremath{-}78}
\newcommand{\GenRfHundredBmCite}{16}
\newcommand{\GenRfHundredBmMulti}{\ensuremath{-}3.7}
\newcommand{\GenRfHundredColon}{11}
\newcommand{\GenRfHundredGouwens}{15}
\newcommand{\GenRfHundredHao}{4.3}
\newcommand{\GenRfHundredMaleCns}{54}
\newcommand{\GenRfHundredScala}{\ensuremath{-}11}
\newcommand{\GenRfHundredStephenson}{\ensuremath{-}0.7}
\newcommand{\GenRfMaxBanc}{68}
\newcommand{\GenRfMaxBmCite}{46}
\newcommand{\GenRfMaxBmMulti}{25}
\newcommand{\GenRfMaxColon}{25}
\newcommand{\GenRfMaxGouwens}{32}
\newcommand{\GenRfMaxHao}{30}
\newcommand{\GenRfMaxMaleCns}{67}
\newcommand{\GenRfMaxScala}{\ensuremath{-}11}
\newcommand{\GenRfMaxStephenson}{9.3}
\newcommand{\GenRidgeHundredBancX}{31}
\newcommand{\GenRidgeMinBancX}{11}
\newcommand{\GenSaveMax}{6.4}
\newcommand{\GenSaveMin}{1.4}
\newcommand{\GenSavePairedOnlyBanc}{3.6}
\newcommand{\GenSavePairedOnlyBancHi}{3.8}
\newcommand{\GenSavePairedOnlyBancLo}{3.1}
\newcommand{\GenSavePairedOnlyBancX}{$\geq$1.00}
\newcommand{\GenSavePairedOnlyBancXHi}{1.00}
\newcommand{\GenSavePairedOnlyBancXLo}{1.00}
\newcommand{\GenSavePairedOnlyBmCite}{$\geq$4.7}
\newcommand{\GenSavePairedOnlyBmCiteHi}{6.0}
\newcommand{\GenSavePairedOnlyBmCiteLo}{1.00}
\newcommand{\GenSavePairedOnlyBmMulti}{$\geq$2.8}
\newcommand{\GenSavePairedOnlyBmMultiHi}{3.0}
\newcommand{\GenSavePairedOnlyBmMultiLo}{1.7}
\newcommand{\GenSavePairedOnlyColon}{$\geq$6.4}
\newcommand{\GenSavePairedOnlyColonHi}{7.7}
\newcommand{\GenSavePairedOnlyColonLo}{3.3}
\newcommand{\GenSavePairedOnlyGouwens}{1.4}
\newcommand{\GenSavePairedOnlyGouwensHi}{1.6}
\newcommand{\GenSavePairedOnlyGouwensLo}{1.2}
\newcommand{\GenSavePairedOnlyHao}{$\geq$2.0}
\newcommand{\GenSavePairedOnlyHaoHi}{2.1}
\newcommand{\GenSavePairedOnlyHaoLo}{1.00}
\newcommand{\GenSavePairedOnlyMaleCns}{2.9}
\newcommand{\GenSavePairedOnlyMaleCnsHi}{3.1}
\newcommand{\GenSavePairedOnlyMaleCnsLo}{2.6}
\newcommand{\GenSavePairedOnlyScala}{$\leq$0.25}
\newcommand{\GenSavePairedOnlyScalaHi}{2.7}
\newcommand{\GenSavePairedOnlyScalaLo}{0.25}
\newcommand{\GenSavePairedOnlyStephenson}{$\geq$2.0}
\newcommand{\GenSavePairedOnlyStephensonHi}{2.6}
\newcommand{\GenSavePairedOnlyStephensonLo}{1.00}
\newcommand{\GenSaveSemiBanc}{2.3}
\newcommand{\GenSaveSemiBancHi}{2.4}
\newcommand{\GenSaveSemiBancLo}{2.0}
\newcommand{\GenSaveSemiBmCite}{3.7}
\newcommand{\GenSaveSemiBmCiteHi}{4.6}
\newcommand{\GenSaveSemiBmCiteLo}{1.00}
\newcommand{\GenSaveSemiBmMulti}{$\geq$2.8}
\newcommand{\GenSaveSemiBmMultiHi}{3.0}
\newcommand{\GenSaveSemiBmMultiLo}{1.7}
\newcommand{\GenSaveSemiColon}{$\geq$6.4}
\newcommand{\GenSaveSemiColonHi}{7.7}
\newcommand{\GenSaveSemiColonLo}{1.9}
\newcommand{\GenSaveSemiGouwens}{1.3}
\newcommand{\GenSaveSemiGouwensHi}{1.5}
\newcommand{\GenSaveSemiGouwensLo}{1.10}
\newcommand{\GenSaveSemiHao}{$\geq$2.0}
\newcommand{\GenSaveSemiHaoHi}{2.1}
\newcommand{\GenSaveSemiHaoLo}{1.00}
\newcommand{\GenSaveSemiMaleCns}{1.9}
\newcommand{\GenSaveSemiMaleCnsHi}{2.1}
\newcommand{\GenSaveSemiMaleCnsLo}{1.7}
\newcommand{\GenSaveSemiScala}{$\geq$1.00}
\newcommand{\GenSaveSemiScalaHi}{1.3}
\newcommand{\GenSaveSemiScalaLo}{0.25}
\newcommand{\GenSaveSemiStephenson}{1.02}
\newcommand{\GenSaveSemiStephensonHi}{1.3}
\newcommand{\GenSaveSemiStephensonLo}{0.53}
\newcommand{\GenSelfBanc}{0.62}
\newcommand{\GenSelfBetter}{none}
\newcommand{\GenSelfBmCite}{0.39}
\newcommand{\GenSelfBmMulti}{0.81}
\newcommand{\GenSelfColon}{0.59}
\newcommand{\GenSelfGouwens}{0.59}
\newcommand{\GenSelfHao}{0.54}
\newcommand{\GenSelfMaleCns}{0.69}
\newcommand{\GenSelfPooledHundredBancX}{7.0}
\newcommand{\GenSelfScala}{1.00}
\newcommand{\GenSelfStephenson}{0.49}
\newcommand{\GenSemiHundredBanc}{44}
\newcommand{\GenSemiHundredBancX}{28}
\newcommand{\GenSemiHundredBmCite}{7.7}
\newcommand{\GenSemiHundredBmMulti}{\ensuremath{-}1.9}
\newcommand{\GenSemiHundredColon}{1.6}
\newcommand{\GenSemiHundredGouwens}{11}
\newcommand{\GenSemiHundredHao}{6.5}
\newcommand{\GenSemiHundredMaleCns}{47}
\newcommand{\GenSemiHundredScala}{\ensuremath{-}3.6}
\newcommand{\GenSemiHundredStephenson}{1.6}
\newcommand{\GenUnitsBanc}{913}
\newcommand{\GenUnitsBmCite}{9}
\newcommand{\GenUnitsBmMulti}{10}
\newcommand{\GenUnitsColon}{12}
\newcommand{\GenUnitsGouwens}{362}
\newcommand{\GenUnitsHao}{8}
\newcommand{\GenUnitsMaleCns}{1,010}
\newcommand{\GenUnitsScala}{263}
\newcommand{\GenUnitsStephenson}{118}

% Generated by extension/predictability/make_assets.py. Do not edit.
\newcommand{\PredBootErr}{6}
\newcommand{\PredBootErrSeven}{8}
\newcommand{\PredBootErrThree}{33}
\newcommand{\PredBootRatio}{1.02}
\newcommand{\PredCellsBlocksRatio}{1.23}
\newcommand{\PredCellsJsRatio}{1.04}
\newcommand{\PredCensAgree}{34}
\newcommand{\PredCensTotal}{37}
\newcommand{\PredCheckCases}{72}
\newcommand{\PredCheckMaxRatio}{0.10}
\newcommand{\PredCurveErrBlocks}{4.2}
\newcommand{\PredCurveErrJs}{0.8}
\newcommand{\PredDataAt}{21:56 EDT}
\newcommand{\PredDevLawErr}{7}
\newcommand{\PredDevLawErrMax}{22}
\newcommand{\PredDevN}{14}
\newcommand{\PredDevSaveMax}{14.7}
\newcommand{\PredDevSaveMin}{1.4}
\newcommand{\PredFrozenAt}{21:39 EDT}
\newcommand{\PredHFive}{met}
\newcommand{\PredHFour}{met}
\newcommand{\PredHOne}{met}
\newcommand{\PredHSix}{met}
\newcommand{\PredHThree}{met}
\newcommand{\PredHTwo}{met}
\newcommand{\PredLawErr}{21}
\newcommand{\PredLawErrHi}{29}
\newcommand{\PredLawErrLo}{11}
\newcommand{\PredLawErrSeven}{4}
\newcommand{\PredLawErrThree}{74}
\newcommand{\PredLawRatio}{0.83}
\newcommand{\PredLawRatioSeven}{0.96}
\newcommand{\PredLawRatioThree}{0.57}
\newcommand{\PredLawWithin}{60}
\newcommand{\PredNCensored}{19}
\newcommand{\PredNData}{ten}
\newcommand{\PredNEasy}{three}
\newcommand{\PredNEstimable}{10}
\newcommand{\PredNEstimableData}{four}
\newcommand{\PredNProblems}{29}
\newcommand{\PredNSeven}{five}
\newcommand{\PredNThree}{eight}
\newcommand{\PredPilotErr}{12}
\newcommand{\PredPilotErrHi}{147}
\newcommand{\PredPilotErrLo}{4}
\newcommand{\PredPilotErrSeven}{2}
\newcommand{\PredPilotErrThree}{96}
\newcommand{\PredPilotWithin}{70}
\newcommand{\PredProfileN}{eight}
\newcommand{\PredProfileTop}{73}
\newcommand{\PredProfileTopTwenty}{94}
\newcommand{\PredRandomLaw}{1.01}
\newcommand{\PredRandomN}{11}
\newcommand{\PredRandomObs}{1.00}
\newcommand{\PredRho}{0.98}
\newcommand{\PredRhoRel}{<0.001}
\newcommand{\PredSaveMax}{11.6}
\newcommand{\PredSaveMin}{1.9}
\newcommand{\PredSnareLawMax}{408}
\newcommand{\PredSnareLawMin}{396}
\newcommand{\PredSnareObsMax}{1,357}
\newcommand{\PredSnareObsMin}{783}
\newcommand{\PredTransferErr}{40}
\newcommand{\PredTransferErrHi}{45}
\newcommand{\PredTransferErrLo}{31}
\newcommand{\PredTransferWithin}{10}

% Generated by extension/synthesis/make_assets.py. Do not edit.
\newcommand{\AlNDonors}{18}
\newcommand{\AlNPools}{seven}
\newcommand{\AlNSamples}{34}
\newcommand{\AlNTwoTimepoints}{16}
\newcommand{\AltCellsAtlas}{147}
\newcommand{\AltCellsAtlasHi}{185}
\newcommand{\AltCellsAtlasLo}{125}
\newcommand{\AltCellsChamp}{180}
\newcommand{\AltCellsChampAtlas}{147}
\newcommand{\AltCellsChampAtlasHi}{189}
\newcommand{\AltCellsChampAtlasLo}{123}
\newcommand{\AltCellsChampHi}{198}
\newcommand{\AltCellsChampLo}{165}
\newcommand{\AltCellsOther}{70}
\newcommand{\AltCellsOtherHi}{85}
\newcommand{\AltCellsOtherLo}{58}
\newcommand{\AltCellsOwn}{110}
\newcommand{\AltCellsOwnHi}{165}
\newcommand{\AltCellsOwnLo}{83}
\newcommand{\AltCellsPo}{$>$800}
\newcommand{\AltCellsPoHi}{800}
\newcommand{\AltCellsPoLo}{800}
\newcommand{\AltCellsRidge}{169}
\newcommand{\AltCellsRidgeHi}{236}
\newcommand{\AltCellsRidgeLo}{131}
\newcommand{\AltCellsRidgePMapped}{235}
\newcommand{\AltCellsRidgePMappedHi}{275}
\newcommand{\AltCellsRidgePMappedLo}{203}
\newcommand{\AltCellsRidgePPercond}{197}
\newcommand{\AltCellsRidgePPercondHi}{241}
\newcommand{\AltCellsRidgePPercondLo}{158}
\newcommand{\AltCellsRidgePPooled}{$>$800}
\newcommand{\AltCellsRidgePPooledHi}{800}
\newcommand{\AltCellsRidgePPooledLo}{800}
\newcommand{\AltCellsRidgeUMapped}{214}
\newcommand{\AltCellsRidgeUMappedHi}{248}
\newcommand{\AltCellsRidgeUMappedLo}{186}
\newcommand{\AltCellsRidgeUPercond}{175}
\newcommand{\AltCellsRidgeUPercondHi}{414}
\newcommand{\AltCellsRidgeUPercondLo}{136}
\newcommand{\AltCellsRidgeUPooled}{$>$800}
\newcommand{\AltCellsRidgeUPooledHi}{800}
\newcommand{\AltCellsRidgeUPooledLo}{800}
\newcommand{\AltCellsSemi}{496}
\newcommand{\AltCellsSemiHi}{800}
\newcommand{\AltCellsSemiLo}{308}
\newcommand{\AltCellsSemiMapped}{579}
\newcommand{\AltCellsSemiMappedHi}{800}
\newcommand{\AltCellsSemiMappedLo}{417}
\newcommand{\AltCellsSemiPercond}{$>$800}
\newcommand{\AltCellsSemiPercondHi}{800}
\newcommand{\AltCellsSemiPercondLo}{308}
\newcommand{\AltCellsSemiPooled}{$>$800}
\newcommand{\AltCellsSemiPooledHi}{800}
\newcommand{\AltCellsSemiPooledLo}{800}
\newcommand{\AltChampFit}{50}
\newcommand{\AltChampFolds}{34}
\newcommand{\AltCompared}{11,220}
\newcommand{\AltDiffOneChampEightHundred}{0.7}
\newcommand{\AltDiffOneChampEightHundredHi}{1.1}
\newcommand{\AltDiffOneChampEightHundredLo}{0.4}
\newcommand{\AltDiffOneChampFifty}{\ensuremath{-}1.3}
\newcommand{\AltDiffOneChampFiftyHi}{0.4}
\newcommand{\AltDiffOneChampFiftyLo}{\ensuremath{-}3.0}
\newcommand{\AltDiffOneChampFourHundred}{0.2}
\newcommand{\AltDiffOneChampFourHundredHi}{0.7}
\newcommand{\AltDiffOneChampFourHundredLo}{\ensuremath{-}0.3}
\newcommand{\AltDiffOneChampHundred}{\ensuremath{-}1.2}
\newcommand{\AltDiffOneChampHundredHi}{\ensuremath{-}0.0}
\newcommand{\AltDiffOneChampHundredLo}{\ensuremath{-}2.4}
\newcommand{\AltDiffOneChampTwentyFive}{4.0}
\newcommand{\AltDiffOneChampTwentyFiveHi}{6.9}
\newcommand{\AltDiffOneChampTwentyFiveLo}{1.0}
\newcommand{\AltDiffOneChampTwoHundred}{\ensuremath{-}0.3}
\newcommand{\AltDiffOneChampTwoHundredHi}{0.3}
\newcommand{\AltDiffOneChampTwoHundredLo}{\ensuremath{-}0.9}
\newcommand{\AltDiffOnePoEightHundred}{1.5}
\newcommand{\AltDiffOnePoEightHundredHi}{1.8}
\newcommand{\AltDiffOnePoEightHundredLo}{1.0}
\newcommand{\AltDiffOnePoFifty}{12.9}
\newcommand{\AltDiffOnePoFiftyHi}{13.8}
\newcommand{\AltDiffOnePoFiftyLo}{12.0}
\newcommand{\AltDiffOnePoFourHundred}{4.1}
\newcommand{\AltDiffOnePoFourHundredHi}{4.7}
\newcommand{\AltDiffOnePoFourHundredLo}{3.4}
\newcommand{\AltDiffOnePoHundred}{11.2}
\newcommand{\AltDiffOnePoHundredHi}{12.0}
\newcommand{\AltDiffOnePoHundredLo}{10.3}
\newcommand{\AltDiffOnePoOneFifty}{9.1}
\newcommand{\AltDiffOnePoOneFiftyHi}{9.9}
\newcommand{\AltDiffOnePoOneFiftyLo}{8.3}
\newcommand{\AltDiffOnePoTwentyFive}{9.2}
\newcommand{\AltDiffOnePoTwentyFiveHi}{9.9}
\newcommand{\AltDiffOnePoTwentyFiveLo}{8.5}
\newcommand{\AltDiffOnePoTwoHundred}{7.6}
\newcommand{\AltDiffOnePoTwoHundredHi}{8.2}
\newcommand{\AltDiffOnePoTwoHundredLo}{6.8}
\newcommand{\AltDiffOneRidgeEightHundred}{\ensuremath{-}1.5}
\newcommand{\AltDiffOneRidgeEightHundredHi}{\ensuremath{-}1.3}
\newcommand{\AltDiffOneRidgeEightHundredLo}{\ensuremath{-}1.7}
\newcommand{\AltDiffOneRidgeFifty}{0.0}
\newcommand{\AltDiffOneRidgeFiftyHi}{0.3}
\newcommand{\AltDiffOneRidgeFiftyLo}{\ensuremath{-}0.2}
\newcommand{\AltDiffOneRidgeFourHundred}{\ensuremath{-}1.3}
\newcommand{\AltDiffOneRidgeFourHundredHi}{\ensuremath{-}1.0}
\newcommand{\AltDiffOneRidgeFourHundredLo}{\ensuremath{-}1.6}
\newcommand{\AltDiffOneRidgeHundred}{\ensuremath{-}0.1}
\newcommand{\AltDiffOneRidgeHundredHi}{0.3}
\newcommand{\AltDiffOneRidgeHundredLo}{\ensuremath{-}0.5}
\newcommand{\AltDiffOneRidgeOneFifty}{\ensuremath{-}0.6}
\newcommand{\AltDiffOneRidgeOneFiftyHi}{\ensuremath{-}0.2}
\newcommand{\AltDiffOneRidgeOneFiftyLo}{\ensuremath{-}1.0}
\newcommand{\AltDiffOneRidgeTwentyFive}{\ensuremath{-}1.7}
\newcommand{\AltDiffOneRidgeTwentyFiveHi}{\ensuremath{-}1.1}
\newcommand{\AltDiffOneRidgeTwentyFiveLo}{\ensuremath{-}2.5}
\newcommand{\AltDiffOneRidgeTwoHundred}{\ensuremath{-}0.9}
\newcommand{\AltDiffOneRidgeTwoHundredHi}{\ensuremath{-}0.6}
\newcommand{\AltDiffOneRidgeTwoHundredLo}{\ensuremath{-}1.2}
\newcommand{\AltDiffOneSemiEightHundred}{2.4}
\newcommand{\AltDiffOneSemiEightHundredHi}{2.9}
\newcommand{\AltDiffOneSemiEightHundredLo}{1.9}
\newcommand{\AltDiffOneSemiFifty}{11.1}
\newcommand{\AltDiffOneSemiFiftyHi}{12.5}
\newcommand{\AltDiffOneSemiFiftyLo}{9.7}
\newcommand{\AltDiffOneSemiFourHundred}{5.1}
\newcommand{\AltDiffOneSemiFourHundredHi}{5.6}
\newcommand{\AltDiffOneSemiFourHundredLo}{4.6}
\newcommand{\AltDiffOneSemiHundred}{9.6}
\newcommand{\AltDiffOneSemiHundredHi}{11.1}
\newcommand{\AltDiffOneSemiHundredLo}{8.3}
\newcommand{\AltDiffOneSemiOneFifty}{9.0}
\newcommand{\AltDiffOneSemiOneFiftyHi}{9.7}
\newcommand{\AltDiffOneSemiOneFiftyLo}{7.9}
\newcommand{\AltDiffOneSemiTwentyFive}{11.2}
\newcommand{\AltDiffOneSemiTwentyFiveHi}{12.7}
\newcommand{\AltDiffOneSemiTwentyFiveLo}{9.8}
\newcommand{\AltDiffOneSemiTwoHundred}{7.2}
\newcommand{\AltDiffOneSemiTwoHundredHi}{8.2}
\newcommand{\AltDiffOneSemiTwoHundredLo}{6.2}
\newcommand{\AltDiffThreePoEightHundred}{8.8}
\newcommand{\AltDiffThreePoEightHundredHi}{10.5}
\newcommand{\AltDiffThreePoEightHundredLo}{7.0}
\newcommand{\AltDiffThreePoFifty}{13.0}
\newcommand{\AltDiffThreePoFiftyHi}{14.1}
\newcommand{\AltDiffThreePoFiftyLo}{11.5}
\newcommand{\AltDiffThreePoFourHundred}{10.2}
\newcommand{\AltDiffThreePoFourHundredHi}{11.6}
\newcommand{\AltDiffThreePoFourHundredLo}{8.5}
\newcommand{\AltDiffThreePoHundred}{14.2}
\newcommand{\AltDiffThreePoHundredHi}{15.7}
\newcommand{\AltDiffThreePoHundredLo}{12.5}
\newcommand{\AltDiffThreePoOneFifty}{14.4}
\newcommand{\AltDiffThreePoOneFiftyHi}{15.4}
\newcommand{\AltDiffThreePoOneFiftyLo}{12.9}
\newcommand{\AltDiffThreePoTwentyFive}{7.9}
\newcommand{\AltDiffThreePoTwentyFiveHi}{9.0}
\newcommand{\AltDiffThreePoTwentyFiveLo}{7.0}
\newcommand{\AltDiffThreePoTwoHundred}{14.1}
\newcommand{\AltDiffThreePoTwoHundredHi}{15.2}
\newcommand{\AltDiffThreePoTwoHundredLo}{12.7}
\newcommand{\AltDiffThreeRidgeEightHundred}{17.5}
\newcommand{\AltDiffThreeRidgeEightHundredHi}{18.3}
\newcommand{\AltDiffThreeRidgeEightHundredLo}{15.7}
\newcommand{\AltDiffThreeRidgeFifty}{6.9}
\newcommand{\AltDiffThreeRidgeFiftyHi}{7.7}
\newcommand{\AltDiffThreeRidgeFiftyLo}{6.1}
\newcommand{\AltDiffThreeRidgeFourHundred}{10.9}
\newcommand{\AltDiffThreeRidgeFourHundredHi}{12.0}
\newcommand{\AltDiffThreeRidgeFourHundredLo}{10.1}
\newcommand{\AltDiffThreeRidgeHundred}{11.3}
\newcommand{\AltDiffThreeRidgeHundredHi}{12.0}
\newcommand{\AltDiffThreeRidgeHundredLo}{10.5}
\newcommand{\AltDiffThreeRidgeOneFifty}{13.6}
\newcommand{\AltDiffThreeRidgeOneFiftyHi}{14.5}
\newcommand{\AltDiffThreeRidgeOneFiftyLo}{12.2}
\newcommand{\AltDiffThreeRidgeTwentyFive}{2.4}
\newcommand{\AltDiffThreeRidgeTwentyFiveHi}{3.1}
\newcommand{\AltDiffThreeRidgeTwentyFiveLo}{1.5}
\newcommand{\AltDiffThreeRidgeTwoHundred}{12.8}
\newcommand{\AltDiffThreeRidgeTwoHundredHi}{14.3}
\newcommand{\AltDiffThreeRidgeTwoHundredLo}{11.3}
\newcommand{\AltDiffThreeSemiEightHundred}{27.8}
\newcommand{\AltDiffThreeSemiEightHundredHi}{29.5}
\newcommand{\AltDiffThreeSemiEightHundredLo}{26.1}
\newcommand{\AltDiffThreeSemiFifty}{19.3}
\newcommand{\AltDiffThreeSemiFiftyHi}{21.1}
\newcommand{\AltDiffThreeSemiFiftyLo}{17.4}
\newcommand{\AltDiffThreeSemiFourHundred}{30.2}
\newcommand{\AltDiffThreeSemiFourHundredHi}{31.5}
\newcommand{\AltDiffThreeSemiFourHundredLo}{29.0}
\newcommand{\AltDiffThreeSemiHundred}{26.0}
\newcommand{\AltDiffThreeSemiHundredHi}{27.8}
\newcommand{\AltDiffThreeSemiHundredLo}{24.1}
\newcommand{\AltDiffThreeSemiOneFifty}{26.0}
\newcommand{\AltDiffThreeSemiOneFiftyHi}{27.7}
\newcommand{\AltDiffThreeSemiOneFiftyLo}{24.4}
\newcommand{\AltDiffThreeSemiTwentyFive}{11.4}
\newcommand{\AltDiffThreeSemiTwentyFiveHi}{13.1}
\newcommand{\AltDiffThreeSemiTwentyFiveLo}{9.7}
\newcommand{\AltDiffThreeSemiTwoHundred}{26.7}
\newcommand{\AltDiffThreeSemiTwoHundredHi}{28.4}
\newcommand{\AltDiffThreeSemiTwoHundredLo}{25.2}
\newcommand{\AltDiffTwoChampEightHundred}{2.3}
\newcommand{\AltDiffTwoChampEightHundredHi}{3.5}
\newcommand{\AltDiffTwoChampEightHundredLo}{1.3}
\newcommand{\AltDiffTwoChampFifty}{7.8}
\newcommand{\AltDiffTwoChampFiftyHi}{9.7}
\newcommand{\AltDiffTwoChampFiftyLo}{6.0}
\newcommand{\AltDiffTwoChampFourHundred}{3.7}
\newcommand{\AltDiffTwoChampFourHundredHi}{5.0}
\newcommand{\AltDiffTwoChampFourHundredLo}{2.5}
\newcommand{\AltDiffTwoChampHundred}{4.7}
\newcommand{\AltDiffTwoChampHundredHi}{6.8}
\newcommand{\AltDiffTwoChampHundredLo}{2.6}
\newcommand{\AltDiffTwoChampTwentyFive}{11.4}
\newcommand{\AltDiffTwoChampTwentyFiveHi}{13.8}
\newcommand{\AltDiffTwoChampTwentyFiveLo}{9.2}
\newcommand{\AltDiffTwoChampTwoHundred}{6.2}
\newcommand{\AltDiffTwoChampTwoHundredHi}{7.6}
\newcommand{\AltDiffTwoChampTwoHundredLo}{4.9}
\newcommand{\AltDiffTwoPoEightHundred}{5.6}
\newcommand{\AltDiffTwoPoEightHundredHi}{6.4}
\newcommand{\AltDiffTwoPoEightHundredLo}{4.8}
\newcommand{\AltDiffTwoPoFifty}{20.1}
\newcommand{\AltDiffTwoPoFiftyHi}{21.3}
\newcommand{\AltDiffTwoPoFiftyLo}{18.9}
\newcommand{\AltDiffTwoPoFourHundred}{9.9}
\newcommand{\AltDiffTwoPoFourHundredHi}{10.9}
\newcommand{\AltDiffTwoPoFourHundredLo}{9.0}
\newcommand{\AltDiffTwoPoHundred}{19.3}
\newcommand{\AltDiffTwoPoHundredHi}{20.4}
\newcommand{\AltDiffTwoPoHundredLo}{18.0}
\newcommand{\AltDiffTwoPoOneFifty}{16.8}
\newcommand{\AltDiffTwoPoOneFiftyHi}{17.9}
\newcommand{\AltDiffTwoPoOneFiftyLo}{15.6}
\newcommand{\AltDiffTwoPoTwentyFive}{15.7}
\newcommand{\AltDiffTwoPoTwentyFiveHi}{16.6}
\newcommand{\AltDiffTwoPoTwentyFiveLo}{14.8}
\newcommand{\AltDiffTwoPoTwoHundred}{15.2}
\newcommand{\AltDiffTwoPoTwoHundredHi}{16.0}
\newcommand{\AltDiffTwoPoTwoHundredLo}{14.2}
\newcommand{\AltDiffTwoRidgeEightHundred}{\ensuremath{-}4.1}
\newcommand{\AltDiffTwoRidgeEightHundredHi}{\ensuremath{-}3.1}
\newcommand{\AltDiffTwoRidgeEightHundredLo}{\ensuremath{-}5.0}
\newcommand{\AltDiffTwoRidgeFifty}{3.0}
\newcommand{\AltDiffTwoRidgeFiftyHi}{3.9}
\newcommand{\AltDiffTwoRidgeFiftyLo}{1.8}
\newcommand{\AltDiffTwoRidgeFourHundred}{\ensuremath{-}5.1}
\newcommand{\AltDiffTwoRidgeFourHundredHi}{\ensuremath{-}4.2}
\newcommand{\AltDiffTwoRidgeFourHundredLo}{\ensuremath{-}6.0}
\newcommand{\AltDiffTwoRidgeHundred}{\ensuremath{-}2.8}
\newcommand{\AltDiffTwoRidgeHundredHi}{\ensuremath{-}1.9}
\newcommand{\AltDiffTwoRidgeHundredLo}{\ensuremath{-}3.6}
\newcommand{\AltDiffTwoRidgeOneFifty}{\ensuremath{-}4.3}
\newcommand{\AltDiffTwoRidgeOneFiftyHi}{\ensuremath{-}3.4}
\newcommand{\AltDiffTwoRidgeOneFiftyLo}{\ensuremath{-}5.0}
\newcommand{\AltDiffTwoRidgeTwentyFive}{1.7}
\newcommand{\AltDiffTwoRidgeTwentyFiveHi}{2.8}
\newcommand{\AltDiffTwoRidgeTwentyFiveLo}{0.5}
\newcommand{\AltDiffTwoRidgeTwoHundred}{\ensuremath{-}4.5}
\newcommand{\AltDiffTwoRidgeTwoHundredHi}{\ensuremath{-}3.6}
\newcommand{\AltDiffTwoRidgeTwoHundredLo}{\ensuremath{-}5.5}
\newcommand{\AltDiffTwoSemiEightHundred}{3.1}
\newcommand{\AltDiffTwoSemiEightHundredHi}{4.2}
\newcommand{\AltDiffTwoSemiEightHundredLo}{2.2}
\newcommand{\AltDiffTwoSemiFifty}{12.5}
\newcommand{\AltDiffTwoSemiFiftyHi}{13.9}
\newcommand{\AltDiffTwoSemiFiftyLo}{11.2}
\newcommand{\AltDiffTwoSemiFourHundred}{10.0}
\newcommand{\AltDiffTwoSemiFourHundredHi}{11.5}
\newcommand{\AltDiffTwoSemiFourHundredLo}{8.6}
\newcommand{\AltDiffTwoSemiHundred}{15.4}
\newcommand{\AltDiffTwoSemiHundredHi}{17.1}
\newcommand{\AltDiffTwoSemiHundredLo}{13.8}
\newcommand{\AltDiffTwoSemiOneFifty}{14.8}
\newcommand{\AltDiffTwoSemiOneFiftyHi}{15.9}
\newcommand{\AltDiffTwoSemiOneFiftyLo}{13.8}
\newcommand{\AltDiffTwoSemiTwentyFive}{9.8}
\newcommand{\AltDiffTwoSemiTwentyFiveHi}{11.5}
\newcommand{\AltDiffTwoSemiTwentyFiveLo}{8.1}
\newcommand{\AltDiffTwoSemiTwoHundred}{14.5}
\newcommand{\AltDiffTwoSemiTwoHundredHi}{15.5}
\newcommand{\AltDiffTwoSemiTwoHundredLo}{13.7}
\newcommand{\AltFrozenAt}{16:31 EDT on 29 September 2026}
\newcommand{\AltIntegrity}{$1.2\times10^{-7}$}
\newcommand{\AltKOne}{met}
\newcommand{\AltKTwo}{met}
\newcommand{\AltOneAtlasEightHundred}{94}
\newcommand{\AltOneAtlasFifty}{80}
\newcommand{\AltOneAtlasFourHundred}{92}
\newcommand{\AltOneAtlasHundred}{86}
\newcommand{\AltOneAtlasOneFifty}{89}
\newcommand{\AltOneAtlasTwentyFive}{71}
\newcommand{\AltOneAtlasTwoHundred}{90}
\newcommand{\AltOneChampAtlasEightHundred}{94}
\newcommand{\AltOneChampAtlasFifty}{80}
\newcommand{\AltOneChampAtlasFourHundred}{92}
\newcommand{\AltOneChampAtlasHundred}{87}
\newcommand{\AltOneChampAtlasTwentyFive}{69}
\newcommand{\AltOneChampAtlasTwoHundred}{90}
\newcommand{\AltOneChampEightHundred}{93}
\newcommand{\AltOneChampFifty}{82}
\newcommand{\AltOneChampFourHundred}{92}
\newcommand{\AltOneChampHundred}{88}
\newcommand{\AltOneChampTwentyFive}{65}
\newcommand{\AltOneChampTwoHundred}{90}
\newcommand{\AltOneOtherEightHundred}{96}
\newcommand{\AltOneOtherFifty}{85}
\newcommand{\AltOneOtherFourHundred}{94}
\newcommand{\AltOneOtherHundred}{89}
\newcommand{\AltOneOtherOneFifty}{91}
\newcommand{\AltOneOtherTwentyFive}{77}
\newcommand{\AltOneOtherTwoHundred}{92}
\newcommand{\AltOneOwnEightHundred}{94}
\newcommand{\AltOneOwnFifty}{84}
\newcommand{\AltOneOwnFourHundred}{92}
\newcommand{\AltOneOwnHundred}{88}
\newcommand{\AltOneOwnOneFifty}{89}
\newcommand{\AltOneOwnTwentyFive}{76}
\newcommand{\AltOneOwnTwoHundred}{90}
\newcommand{\AltOnePoEightHundred}{92}
\newcommand{\AltOnePoFifty}{67}
\newcommand{\AltOnePoFourHundred}{88}
\newcommand{\AltOnePoHundred}{75}
\newcommand{\AltOnePoOneFifty}{80}
\newcommand{\AltOnePoTwentyFive}{61}
\newcommand{\AltOnePoTwoHundred}{82}
\newcommand{\AltOneRidgeEightHundred}{95}
\newcommand{\AltOneRidgeFifty}{80}
\newcommand{\AltOneRidgeFourHundred}{94}
\newcommand{\AltOneRidgeHundred}{86}
\newcommand{\AltOneRidgeOneFifty}{89}
\newcommand{\AltOneRidgePMappedEightHundred}{93}
\newcommand{\AltOneRidgePMappedFifty}{80}
\newcommand{\AltOneRidgePMappedFourHundred}{91}
\newcommand{\AltOneRidgePMappedHundred}{85}
\newcommand{\AltOneRidgePMappedOneFifty}{88}
\newcommand{\AltOneRidgePMappedTwentyFive}{72}
\newcommand{\AltOneRidgePMappedTwoHundred}{89}
\newcommand{\AltOneRidgePPercondEightHundred}{95}
\newcommand{\AltOneRidgePPercondFifty}{76}
\newcommand{\AltOneRidgePPercondFourHundred}{93}
\newcommand{\AltOneRidgePPercondHundred}{85}
\newcommand{\AltOneRidgePPercondOneFifty}{88}
\newcommand{\AltOneRidgePPercondTwentyFive}{57}
\newcommand{\AltOneRidgePPercondTwoHundred}{89}
\newcommand{\AltOneRidgePPooledEightHundred}{88}
\newcommand{\AltOneRidgePPooledFifty}{76}
\newcommand{\AltOneRidgePPooledFourHundred}{86}
\newcommand{\AltOneRidgePPooledHundred}{80}
\newcommand{\AltOneRidgePPooledOneFifty}{83}
\newcommand{\AltOneRidgePPooledTwentyFive}{69}
\newcommand{\AltOneRidgePPooledTwoHundred}{84}
\newcommand{\AltOneRidgeTwentyFive}{72}
\newcommand{\AltOneRidgeTwoHundred}{91}
\newcommand{\AltOneRidgeUMappedEightHundred}{93}
\newcommand{\AltOneRidgeUMappedFifty}{80}
\newcommand{\AltOneRidgeUMappedFourHundred}{91}
\newcommand{\AltOneRidgeUMappedHundred}{85}
\newcommand{\AltOneRidgeUMappedOneFifty}{88}
\newcommand{\AltOneRidgeUMappedTwentyFive}{72}
\newcommand{\AltOneRidgeUMappedTwoHundred}{89}
\newcommand{\AltOneRidgeUPercondEightHundred}{95}
\newcommand{\AltOneRidgeUPercondFifty}{76}
\newcommand{\AltOneRidgeUPercondFourHundred}{94}
\newcommand{\AltOneRidgeUPercondHundred}{86}
\newcommand{\AltOneRidgeUPercondOneFifty}{89}
\newcommand{\AltOneRidgeUPercondTwentyFive}{57}
\newcommand{\AltOneRidgeUPercondTwoHundred}{91}
\newcommand{\AltOneRidgeUPooledEightHundred}{88}
\newcommand{\AltOneRidgeUPooledFifty}{76}
\newcommand{\AltOneRidgeUPooledFourHundred}{87}
\newcommand{\AltOneRidgeUPooledHundred}{81}
\newcommand{\AltOneRidgeUPooledOneFifty}{83}
\newcommand{\AltOneRidgeUPooledTwentyFive}{70}
\newcommand{\AltOneRidgeUPooledTwoHundred}{84}
\newcommand{\AltOneSemiEightHundred}{92}
\newcommand{\AltOneSemiFifty}{69}
\newcommand{\AltOneSemiFourHundred}{87}
\newcommand{\AltOneSemiHundred}{76}
\newcommand{\AltOneSemiMappedEightHundred}{89}
\newcommand{\AltOneSemiMappedFifty}{67}
\newcommand{\AltOneSemiMappedFourHundred}{87}
\newcommand{\AltOneSemiMappedHundred}{73}
\newcommand{\AltOneSemiMappedOneFifty}{79}
\newcommand{\AltOneSemiMappedTwentyFive}{60}
\newcommand{\AltOneSemiMappedTwoHundred}{83}
\newcommand{\AltOneSemiOneFifty}{80}
\newcommand{\AltOneSemiPercondEightHundred}{92}
\newcommand{\AltOneSemiPercondFifty}{69}
\newcommand{\AltOneSemiPercondFourHundred}{87}
\newcommand{\AltOneSemiPercondHundred}{76}
\newcommand{\AltOneSemiPercondOneFifty}{80}
\newcommand{\AltOneSemiPercondTwentyFive}{56}
\newcommand{\AltOneSemiPercondTwoHundred}{81}
\newcommand{\AltOneSemiPooledEightHundred}{85}
\newcommand{\AltOneSemiPooledFifty}{65}
\newcommand{\AltOneSemiPooledFourHundred}{83}
\newcommand{\AltOneSemiPooledHundred}{72}
\newcommand{\AltOneSemiPooledOneFifty}{76}
\newcommand{\AltOneSemiPooledTwentyFive}{59}
\newcommand{\AltOneSemiPooledTwoHundred}{79}
\newcommand{\AltOneSemiTwentyFive}{60}
\newcommand{\AltOneSemiTwoHundred}{83}
\newcommand{\AltRatioChamp}{1.23}
\newcommand{\AltRatioChampHi}{1.44}
\newcommand{\AltRatioChampLo}{0.98}
\newcommand{\AltRatioPoOverAtlas}{at least 5.4}
\newcommand{\AltRatioPoOverAtlasHi}{6.4}
\newcommand{\AltRatioPoOverAtlasLo}{4.3}
\newcommand{\AltRatioPoOverRidge}{at least 4.7}
\newcommand{\AltRatioPoOverRidgeHi}{6.1}
\newcommand{\AltRatioPoOverRidgeLo}{3.4}
\newcommand{\AltRatioPoOverSemi}{at least 1.61}
\newcommand{\AltRatioPoOverSemiHi}{2.6}
\newcommand{\AltRatioPoOverSemiLo}{1.00}
\newcommand{\AltRatioRidge}{1.15}
\newcommand{\AltRatioRidgeHi}{1.36}
\newcommand{\AltRatioRidgeLo}{1.003}
\newcommand{\AltRatioRidgeOverOwn}{1.54}
\newcommand{\AltRatioRidgeOverOwnHi}{1.87}
\newcommand{\AltRatioRidgeOverOwnLo}{1.27}
\newcommand{\AltRatioSemi}{3.4}
\newcommand{\AltRatioSemiHi}{4.6}
\newcommand{\AltRatioSemiLo}{2.4}
\newcommand{\AltRatioSemiOverOwn}{4.5}
\newcommand{\AltRatioSemiOverOwnHi}{5.8}
\newcommand{\AltRatioSemiOverOwnLo}{3.4}
\newcommand{\AltReproduces}{$3.3\times10^{-6}$}
\newcommand{\AltRfAtlasEightHundred}{48}
\newcommand{\AltRfAtlasFifty}{11}
\newcommand{\AltRfAtlasFourHundred}{42}
\newcommand{\AltRfAtlasHundred}{23}
\newcommand{\AltRfAtlasOneFifty}{29}
\newcommand{\AltRfAtlasTwentyFive}{0.5}
\newcommand{\AltRfAtlasTwoHundred}{33}
\newcommand{\AltRfChampAtlasEightHundred}{48}
\newcommand{\AltRfChampAtlasFifty}{11}
\newcommand{\AltRfChampAtlasFourHundred}{42}
\newcommand{\AltRfChampAtlasHundred}{22}
\newcommand{\AltRfChampAtlasTwentyFive}{\ensuremath{-}1.3}
\newcommand{\AltRfChampAtlasTwoHundred}{33}
\newcommand{\AltRfChampEightHundred}{40}
\newcommand{\AltRfChampFifty}{11}
\newcommand{\AltRfChampFourHundred}{35}
\newcommand{\AltRfChampHundred}{20}
\newcommand{\AltRfChampTwentyFive}{0.0}
\newcommand{\AltRfChampTwoHundred}{30}
\newcommand{\AltRfMaxAtlas}{48}
\newcommand{\AltRfMaxChamp}{40}
\newcommand{\AltRfMaxChampAtlas}{48}
\newcommand{\AltRfMaxOther}{57}
\newcommand{\AltRfMaxOwn}{45}
\newcommand{\AltRfMaxPo}{14}
\newcommand{\AltRfMaxRidge}{45}
\newcommand{\AltRfMaxRidgePMapped}{41}
\newcommand{\AltRfMaxRidgePPercond}{43}
\newcommand{\AltRfMaxRidgePPooled}{18}
\newcommand{\AltRfMaxRidgeUMapped}{45}
\newcommand{\AltRfMaxRidgeUPercond}{42}
\newcommand{\AltRfMaxRidgeUPooled}{18}
\newcommand{\AltRfMaxSemi}{32}
\newcommand{\AltRfMaxSemiMapped}{32}
\newcommand{\AltRfMaxSemiPercond}{28}
\newcommand{\AltRfMaxSemiPooled}{11}
\newcommand{\AltRfOtherEightHundred}{57}
\newcommand{\AltRfOtherFifty}{23}
\newcommand{\AltRfOtherFourHundred}{50}
\newcommand{\AltRfOtherHundred}{34}
\newcommand{\AltRfOtherOneFifty}{39}
\newcommand{\AltRfOtherTwentyFive}{8.8}
\newcommand{\AltRfOtherTwoHundred}{43}
\newcommand{\AltRfOwnEightHundred}{45}
\newcommand{\AltRfOwnFifty}{19}
\newcommand{\AltRfOwnFourHundred}{40}
\newcommand{\AltRfOwnHundred}{27}
\newcommand{\AltRfOwnOneFifty}{31}
\newcommand{\AltRfOwnTwentyFive}{5.6}
\newcommand{\AltRfOwnTwoHundred}{34}
\newcommand{\AltRfPoEightHundred}{13}
\newcommand{\AltRfPoFifty}{2.4}
\newcommand{\AltRfPoFourHundred}{14}
\newcommand{\AltRfPoHundred}{3.5}
\newcommand{\AltRfPoOneFifty}{5.1}
\newcommand{\AltRfPoTwentyFive}{0.5}
\newcommand{\AltRfPoTwoHundred}{7.7}
\newcommand{\AltRfRidgeEightHundred}{45}
\newcommand{\AltRfRidgeFifty}{14}
\newcommand{\AltRfRidgeFourHundred}{38}
\newcommand{\AltRfRidgeHundred}{22}
\newcommand{\AltRfRidgeOneFifty}{28}
\newcommand{\AltRfRidgePMappedEightHundred}{41}
\newcommand{\AltRfRidgePMappedFifty}{11}
\newcommand{\AltRfRidgePMappedFourHundred}{36}
\newcommand{\AltRfRidgePMappedHundred}{19}
\newcommand{\AltRfRidgePMappedOneFifty}{23}
\newcommand{\AltRfRidgePMappedTwentyFive}{2.1}
\newcommand{\AltRfRidgePMappedTwoHundred}{26}
\newcommand{\AltRfRidgePPercondEightHundred}{43}
\newcommand{\AltRfRidgePPercondFifty}{14}
\newcommand{\AltRfRidgePPercondFourHundred}{37}
\newcommand{\AltRfRidgePPercondHundred}{22}
\newcommand{\AltRfRidgePPercondOneFifty}{26}
\newcommand{\AltRfRidgePPercondTwentyFive}{3.2}
\newcommand{\AltRfRidgePPercondTwoHundred}{28}
\newcommand{\AltRfRidgePPooledEightHundred}{18}
\newcommand{\AltRfRidgePPooledFifty}{5.2}
\newcommand{\AltRfRidgePPooledFourHundred}{16}
\newcommand{\AltRfRidgePPooledHundred}{8.6}
\newcommand{\AltRfRidgePPooledOneFifty}{10}
\newcommand{\AltRfRidgePPooledTwentyFive}{0.7}
\newcommand{\AltRfRidgePPooledTwoHundred}{12}
\newcommand{\AltRfRidgeTwentyFive}{5.6}
\newcommand{\AltRfRidgeTwoHundred}{30}
\newcommand{\AltRfRidgeUMappedEightHundred}{45}
\newcommand{\AltRfRidgeUMappedFifty}{11}
\newcommand{\AltRfRidgeUMappedFourHundred}{38}
\newcommand{\AltRfRidgeUMappedHundred}{20}
\newcommand{\AltRfRidgeUMappedOneFifty}{25}
\newcommand{\AltRfRidgeUMappedTwentyFive}{5.6}
\newcommand{\AltRfRidgeUMappedTwoHundred}{27}
\newcommand{\AltRfRidgeUPercondEightHundred}{42}
\newcommand{\AltRfRidgeUPercondFifty}{11}
\newcommand{\AltRfRidgeUPercondFourHundred}{35}
\newcommand{\AltRfRidgeUPercondHundred}{20}
\newcommand{\AltRfRidgeUPercondOneFifty}{28}
\newcommand{\AltRfRidgeUPercondTwentyFive}{2.6}
\newcommand{\AltRfRidgeUPercondTwoHundred}{30}
\newcommand{\AltRfRidgeUPooledEightHundred}{18}
\newcommand{\AltRfRidgeUPooledFifty}{5.2}
\newcommand{\AltRfRidgeUPooledFourHundred}{15}
\newcommand{\AltRfRidgeUPooledHundred}{9.2}
\newcommand{\AltRfRidgeUPooledOneFifty}{11}
\newcommand{\AltRfRidgeUPooledTwentyFive}{2.6}
\newcommand{\AltRfRidgeUPooledTwoHundred}{12}
\newcommand{\AltRfSemiEightHundred}{32}
\newcommand{\AltRfSemiFifty}{11}
\newcommand{\AltRfSemiFourHundred}{27}
\newcommand{\AltRfSemiHundred}{17}
\newcommand{\AltRfSemiMappedEightHundred}{32}
\newcommand{\AltRfSemiMappedFifty}{\ensuremath{-}0.2}
\newcommand{\AltRfSemiMappedFourHundred}{24}
\newcommand{\AltRfSemiMappedHundred}{3.6}
\newcommand{\AltRfSemiMappedOneFifty}{7.4}
\newcommand{\AltRfSemiMappedTwentyFive}{1.1}
\newcommand{\AltRfSemiMappedTwoHundred}{13}
\newcommand{\AltRfSemiOneFifty}{20}
\newcommand{\AltRfSemiPercondEightHundred}{28}
\newcommand{\AltRfSemiPercondFifty}{11}
\newcommand{\AltRfSemiPercondFourHundred}{27}
\newcommand{\AltRfSemiPercondHundred}{17}
\newcommand{\AltRfSemiPercondOneFifty}{20}
\newcommand{\AltRfSemiPercondTwentyFive}{3.1}
\newcommand{\AltRfSemiPercondTwoHundred}{20}
\newcommand{\AltRfSemiPooledEightHundred}{11}
\newcommand{\AltRfSemiPooledFifty}{0.3}
\newcommand{\AltRfSemiPooledFourHundred}{8.4}
\newcommand{\AltRfSemiPooledHundred}{1.8}
\newcommand{\AltRfSemiPooledOneFifty}{2.9}
\newcommand{\AltRfSemiPooledTwentyFive}{0.6}
\newcommand{\AltRfSemiPooledTwoHundred}{4.5}
\newcommand{\AltRfSemiTwentyFive}{3.1}
\newcommand{\AltRfSemiTwoHundred}{20}
\newcommand{\AltTarget}{28}
\newcommand{\AltThreeAtlasEightHundred}{73}
\newcommand{\AltThreeAtlasFifty}{24}
\newcommand{\AltThreeAtlasFourHundred}{64}
\newcommand{\AltThreeAtlasHundred}{39}
\newcommand{\AltThreeAtlasOneFifty}{47}
\newcommand{\AltThreeAtlasTwentyFive}{13}
\newcommand{\AltThreeAtlasTwoHundred}{53}
\newcommand{\AltThreeOtherEightHundred}{74}
\newcommand{\AltThreeOtherFifty}{24}
\newcommand{\AltThreeOtherFourHundred}{63}
\newcommand{\AltThreeOtherHundred}{36}
\newcommand{\AltThreeOtherOneFifty}{44}
\newcommand{\AltThreeOtherTwentyFive}{14}
\newcommand{\AltThreeOtherTwoHundred}{50}
\newcommand{\AltThreeOwnEightHundred}{67}
\newcommand{\AltThreeOwnFifty}{23}
\newcommand{\AltThreeOwnFourHundred}{57}
\newcommand{\AltThreeOwnHundred}{34}
\newcommand{\AltThreeOwnOneFifty}{40}
\newcommand{\AltThreeOwnTwentyFive}{14}
\newcommand{\AltThreeOwnTwoHundred}{45}
\newcommand{\AltThreePoEightHundred}{65}
\newcommand{\AltThreePoFifty}{11}
\newcommand{\AltThreePoFourHundred}{54}
\newcommand{\AltThreePoHundred}{25}
\newcommand{\AltThreePoOneFifty}{33}
\newcommand{\AltThreePoTwentyFive}{4.7}
\newcommand{\AltThreePoTwoHundred}{39}
\newcommand{\AltThreeRidgeEightHundred}{56}
\newcommand{\AltThreeRidgeFifty}{17}
\newcommand{\AltThreeRidgeFourHundred}{53}
\newcommand{\AltThreeRidgeHundred}{28}
\newcommand{\AltThreeRidgeOneFifty}{33}
\newcommand{\AltThreeRidgePMappedEightHundred}{52}
\newcommand{\AltThreeRidgePMappedFifty}{17}
\newcommand{\AltThreeRidgePMappedFourHundred}{49}
\newcommand{\AltThreeRidgePMappedHundred}{28}
\newcommand{\AltThreeRidgePMappedOneFifty}{33}
\newcommand{\AltThreeRidgePMappedTwentyFive}{10}
\newcommand{\AltThreeRidgePMappedTwoHundred}{40}
\newcommand{\AltThreeRidgePPercondEightHundred}{51}
\newcommand{\AltThreeRidgePPercondFifty}{14}
\newcommand{\AltThreeRidgePPercondFourHundred}{45}
\newcommand{\AltThreeRidgePPercondHundred}{26}
\newcommand{\AltThreeRidgePPercondOneFifty}{32}
\newcommand{\AltThreeRidgePPercondTwentyFive}{3.3}
\newcommand{\AltThreeRidgePPercondTwoHundred}{36}
\newcommand{\AltThreeRidgePPooledEightHundred}{46}
\newcommand{\AltThreeRidgePPooledFifty}{15}
\newcommand{\AltThreeRidgePPooledFourHundred}{43}
\newcommand{\AltThreeRidgePPooledHundred}{24}
\newcommand{\AltThreeRidgePPooledOneFifty}{29}
\newcommand{\AltThreeRidgePPooledTwentyFive}{9.0}
\newcommand{\AltThreeRidgePPooledTwoHundred}{35}
\newcommand{\AltThreeRidgeTwentyFive}{10}
\newcommand{\AltThreeRidgeTwoHundred}{40}
\newcommand{\AltThreeRidgeUMappedEightHundred}{56}
\newcommand{\AltThreeRidgeUMappedFifty}{13}
\newcommand{\AltThreeRidgeUMappedFourHundred}{53}
\newcommand{\AltThreeRidgeUMappedHundred}{24}
\newcommand{\AltThreeRidgeUMappedOneFifty}{30}
\newcommand{\AltThreeRidgeUMappedTwentyFive}{7.0}
\newcommand{\AltThreeRidgeUMappedTwoHundred}{34}
\newcommand{\AltThreeRidgeUPercondEightHundred}{56}
\newcommand{\AltThreeRidgeUPercondFifty}{13}
\newcommand{\AltThreeRidgeUPercondFourHundred}{49}
\newcommand{\AltThreeRidgeUPercondHundred}{26}
\newcommand{\AltThreeRidgeUPercondOneFifty}{32}
\newcommand{\AltThreeRidgeUPercondTwentyFive}{2.7}
\newcommand{\AltThreeRidgeUPercondTwoHundred}{35}
\newcommand{\AltThreeRidgeUPooledEightHundred}{48}
\newcommand{\AltThreeRidgeUPooledFifty}{11}
\newcommand{\AltThreeRidgeUPooledFourHundred}{47}
\newcommand{\AltThreeRidgeUPooledHundred}{21}
\newcommand{\AltThreeRidgeUPooledOneFifty}{26}
\newcommand{\AltThreeRidgeUPooledTwentyFive}{6.0}
\newcommand{\AltThreeRidgeUPooledTwoHundred}{29}
\newcommand{\AltThreeSemiEightHundred}{46}
\newcommand{\AltThreeSemiFifty}{5.0}
\newcommand{\AltThreeSemiFourHundred}{34}
\newcommand{\AltThreeSemiHundred}{13}
\newcommand{\AltThreeSemiMappedEightHundred}{38}
\newcommand{\AltThreeSemiMappedFifty}{4.6}
\newcommand{\AltThreeSemiMappedFourHundred}{34}
\newcommand{\AltThreeSemiMappedHundred}{13}
\newcommand{\AltThreeSemiMappedOneFifty}{21}
\newcommand{\AltThreeSemiMappedTwentyFive}{1.1}
\newcommand{\AltThreeSemiMappedTwoHundred}{26}
\newcommand{\AltThreeSemiOneFifty}{21}
\newcommand{\AltThreeSemiPercondEightHundred}{46}
\newcommand{\AltThreeSemiPercondFifty}{5.0}
\newcommand{\AltThreeSemiPercondFourHundred}{32}
\newcommand{\AltThreeSemiPercondHundred}{9.8}
\newcommand{\AltThreeSemiPercondOneFifty}{14}
\newcommand{\AltThreeSemiPercondTwentyFive}{1.2}
\newcommand{\AltThreeSemiPercondTwoHundred}{19}
\newcommand{\AltThreeSemiPooledEightHundred}{33}
\newcommand{\AltThreeSemiPooledFifty}{3.7}
\newcommand{\AltThreeSemiPooledFourHundred}{31}
\newcommand{\AltThreeSemiPooledHundred}{11}
\newcommand{\AltThreeSemiPooledOneFifty}{18}
\newcommand{\AltThreeSemiPooledTwentyFive}{0.8}
\newcommand{\AltThreeSemiPooledTwoHundred}{23}
\newcommand{\AltThreeSemiTwentyFive}{1.2}
\newcommand{\AltThreeSemiTwoHundred}{26}
\newcommand{\AltTwoAtlasEightHundred}{67}
\newcommand{\AltTwoAtlasFifty}{35}
\newcommand{\AltTwoAtlasFourHundred}{62}
\newcommand{\AltTwoAtlasHundred}{47}
\newcommand{\AltTwoAtlasOneFifty}{52}
\newcommand{\AltTwoAtlasTwentyFive}{23}
\newcommand{\AltTwoAtlasTwoHundred}{55}
\newcommand{\AltTwoChampAtlasEightHundred}{67}
\newcommand{\AltTwoChampAtlasFifty}{36}
\newcommand{\AltTwoChampAtlasFourHundred}{62}
\newcommand{\AltTwoChampAtlasHundred}{47}
\newcommand{\AltTwoChampAtlasTwentyFive}{23}
\newcommand{\AltTwoChampAtlasTwoHundred}{55}
\newcommand{\AltTwoChampEightHundred}{64}
\newcommand{\AltTwoChampFifty}{28}
\newcommand{\AltTwoChampFourHundred}{58}
\newcommand{\AltTwoChampHundred}{42}
\newcommand{\AltTwoChampTwentyFive}{11}
\newcommand{\AltTwoChampTwoHundred}{49}
\newcommand{\AltTwoOtherEightHundred}{67}
\newcommand{\AltTwoOtherFifty}{38}
\newcommand{\AltTwoOtherFourHundred}{61}
\newcommand{\AltTwoOtherHundred}{46}
\newcommand{\AltTwoOtherOneFifty}{50}
\newcommand{\AltTwoOtherTwentyFive}{28}
\newcommand{\AltTwoOtherTwoHundred}{54}
\newcommand{\AltTwoOwnEightHundred}{57}
\newcommand{\AltTwoOwnFifty}{36}
\newcommand{\AltTwoOwnFourHundred}{52}
\newcommand{\AltTwoOwnHundred}{41}
\newcommand{\AltTwoOwnOneFifty}{44}
\newcommand{\AltTwoOwnTwentyFive}{27}
\newcommand{\AltTwoOwnTwoHundred}{46}
\newcommand{\AltTwoPoEightHundred}{61}
\newcommand{\AltTwoPoFifty}{15}
\newcommand{\AltTwoPoFourHundred}{52}
\newcommand{\AltTwoPoHundred}{28}
\newcommand{\AltTwoPoOneFifty}{35}
\newcommand{\AltTwoPoTwentyFive}{6.9}
\newcommand{\AltTwoPoTwoHundred}{40}
\newcommand{\AltTwoRidgeEightHundred}{71}
\newcommand{\AltTwoRidgeFifty}{32}
\newcommand{\AltTwoRidgeFourHundred}{67}
\newcommand{\AltTwoRidgeHundred}{50}
\newcommand{\AltTwoRidgeOneFifty}{57}
\newcommand{\AltTwoRidgePMappedEightHundred}{63}
\newcommand{\AltTwoRidgePMappedFifty}{32}
\newcommand{\AltTwoRidgePMappedFourHundred}{58}
\newcommand{\AltTwoRidgePMappedHundred}{43}
\newcommand{\AltTwoRidgePMappedOneFifty}{48}
\newcommand{\AltTwoRidgePMappedTwentyFive}{21}
\newcommand{\AltTwoRidgePMappedTwoHundred}{51}
\newcommand{\AltTwoRidgePPercondEightHundred}{69}
\newcommand{\AltTwoRidgePPercondFifty}{32}
\newcommand{\AltTwoRidgePPercondFourHundred}{65}
\newcommand{\AltTwoRidgePPercondHundred}{48}
\newcommand{\AltTwoRidgePPercondOneFifty}{55}
\newcommand{\AltTwoRidgePPercondTwentyFive}{9.3}
\newcommand{\AltTwoRidgePPercondTwoHundred}{58}
\newcommand{\AltTwoRidgePPooledEightHundred}{33}
\newcommand{\AltTwoRidgePPooledFifty}{23}
\newcommand{\AltTwoRidgePPooledFourHundred}{32}
\newcommand{\AltTwoRidgePPooledHundred}{28}
\newcommand{\AltTwoRidgePPooledOneFifty}{30}
\newcommand{\AltTwoRidgePPooledTwentyFive}{17}
\newcommand{\AltTwoRidgePPooledTwoHundred}{31}
\newcommand{\AltTwoRidgeTwentyFive}{21}
\newcommand{\AltTwoRidgeTwoHundred}{60}
\newcommand{\AltTwoRidgeUMappedEightHundred}{64}
\newcommand{\AltTwoRidgeUMappedFifty}{32}
\newcommand{\AltTwoRidgeUMappedFourHundred}{60}
\newcommand{\AltTwoRidgeUMappedHundred}{44}
\newcommand{\AltTwoRidgeUMappedOneFifty}{49}
\newcommand{\AltTwoRidgeUMappedTwentyFive}{21}
\newcommand{\AltTwoRidgeUMappedTwoHundred}{52}
\newcommand{\AltTwoRidgeUPercondEightHundred}{71}
\newcommand{\AltTwoRidgeUPercondFifty}{32}
\newcommand{\AltTwoRidgeUPercondFourHundred}{67}
\newcommand{\AltTwoRidgeUPercondHundred}{50}
\newcommand{\AltTwoRidgeUPercondOneFifty}{57}
\newcommand{\AltTwoRidgeUPercondTwentyFive}{9.5}
\newcommand{\AltTwoRidgeUPercondTwoHundred}{60}
\newcommand{\AltTwoRidgeUPooledEightHundred}{32}
\newcommand{\AltTwoRidgeUPooledFifty}{23}
\newcommand{\AltTwoRidgeUPooledFourHundred}{32}
\newcommand{\AltTwoRidgeUPooledHundred}{28}
\newcommand{\AltTwoRidgeUPooledOneFifty}{29}
\newcommand{\AltTwoRidgeUPooledTwentyFive}{17}
\newcommand{\AltTwoRidgeUPooledTwoHundred}{30}
\newcommand{\AltTwoSemiEightHundred}{63}
\newcommand{\AltTwoSemiFifty}{23}
\newcommand{\AltTwoSemiFourHundred}{52}
\newcommand{\AltTwoSemiHundred}{32}
\newcommand{\AltTwoSemiMappedEightHundred}{49}
\newcommand{\AltTwoSemiMappedFifty}{16}
\newcommand{\AltTwoSemiMappedFourHundred}{45}
\newcommand{\AltTwoSemiMappedHundred}{22}
\newcommand{\AltTwoSemiMappedOneFifty}{29}
\newcommand{\AltTwoSemiMappedTwentyFive}{13}
\newcommand{\AltTwoSemiMappedTwoHundred}{36}
\newcommand{\AltTwoSemiOneFifty}{37}
\newcommand{\AltTwoSemiPercondEightHundred}{63}
\newcommand{\AltTwoSemiPercondFifty}{23}
\newcommand{\AltTwoSemiPercondFourHundred}{52}
\newcommand{\AltTwoSemiPercondHundred}{32}
\newcommand{\AltTwoSemiPercondOneFifty}{37}
\newcommand{\AltTwoSemiPercondTwentyFive}{8.2}
\newcommand{\AltTwoSemiPercondTwoHundred}{41}
\newcommand{\AltTwoSemiPooledEightHundred}{33}
\newcommand{\AltTwoSemiPooledFifty}{14}
\newcommand{\AltTwoSemiPooledFourHundred}{32}
\newcommand{\AltTwoSemiPooledHundred}{20}
\newcommand{\AltTwoSemiPooledOneFifty}{26}
\newcommand{\AltTwoSemiPooledTwentyFive}{10}
\newcommand{\AltTwoSemiPooledTwoHundred}{29}
\newcommand{\AltTwoSemiTwentyFive}{13}
\newcommand{\AltTwoSemiTwoHundred}{41}
\newcommand{\BoBOne}{met}
\newcommand{\BoBThree}{met}
\newcommand{\BoBTwo}{met}
\newcommand{\BoCellsOneAtlas}{29}
\newcommand{\BoCellsOneOther}{$\leq$25}
\newcommand{\BoCellsOneOwn}{$\leq$25}
\newcommand{\BoCellsOnePo}{84}
\newcommand{\BoCellsTwoAtlas}{45}
\newcommand{\BoCellsTwoOther}{37}
\newcommand{\BoCellsTwoOwn}{42}
\newcommand{\BoCellsTwoPo}{136}
\newcommand{\BoCogAtlasEightHundred}{99}
\newcommand{\BoCogAtlasEightHundredHi}{100}
\newcommand{\BoCogAtlasEightHundredLo}{99}
\newcommand{\BoCogAtlasFifty}{85}
\newcommand{\BoCogAtlasFiftyHi}{87}
\newcommand{\BoCogAtlasFiftyLo}{83}
\newcommand{\BoCogAtlasFourHundred}{99}
\newcommand{\BoCogAtlasFourHundredHi}{99}
\newcommand{\BoCogAtlasFourHundredLo}{98}
\newcommand{\BoCogAtlasHundred}{93}
\newcommand{\BoCogAtlasHundredHi}{94}
\newcommand{\BoCogAtlasHundredLo}{92}
\newcommand{\BoCogAtlasOneFifty}{96}
\newcommand{\BoCogAtlasOneFiftyHi}{96}
\newcommand{\BoCogAtlasOneFiftyLo}{95}
\newcommand{\BoCogAtlasTwentyFive}{76}
\newcommand{\BoCogAtlasTwentyFiveHi}{78}
\newcommand{\BoCogAtlasTwentyFiveLo}{75}
\newcommand{\BoCogAtlasTwoHundred}{97}
\newcommand{\BoCogAtlasTwoHundredHi}{97}
\newcommand{\BoCogAtlasTwoHundredLo}{96}
\newcommand{\BoCogOtherEightHundred}{99}
\newcommand{\BoCogOtherEightHundredHi}{100}
\newcommand{\BoCogOtherEightHundredLo}{99}
\newcommand{\BoCogOtherFifty}{87}
\newcommand{\BoCogOtherFiftyHi}{88}
\newcommand{\BoCogOtherFiftyLo}{86}
\newcommand{\BoCogOtherFourHundred}{99}
\newcommand{\BoCogOtherFourHundredHi}{99}
\newcommand{\BoCogOtherFourHundredLo}{98}
\newcommand{\BoCogOtherHundred}{93}
\newcommand{\BoCogOtherHundredHi}{94}
\newcommand{\BoCogOtherHundredLo}{92}
\newcommand{\BoCogOtherOneFifty}{96}
\newcommand{\BoCogOtherOneFiftyHi}{96}
\newcommand{\BoCogOtherOneFiftyLo}{95}
\newcommand{\BoCogOtherTwentyFive}{78}
\newcommand{\BoCogOtherTwentyFiveHi}{80}
\newcommand{\BoCogOtherTwentyFiveLo}{77}
\newcommand{\BoCogOtherTwoHundred}{97}
\newcommand{\BoCogOtherTwoHundredHi}{98}
\newcommand{\BoCogOtherTwoHundredLo}{96}
\newcommand{\BoCogOwnEightHundred}{99}
\newcommand{\BoCogOwnEightHundredHi}{100}
\newcommand{\BoCogOwnEightHundredLo}{99}
\newcommand{\BoCogOwnFifty}{85}
\newcommand{\BoCogOwnFiftyHi}{87}
\newcommand{\BoCogOwnFiftyLo}{84}
\newcommand{\BoCogOwnFourHundred}{98}
\newcommand{\BoCogOwnFourHundredHi}{99}
\newcommand{\BoCogOwnFourHundredLo}{98}
\newcommand{\BoCogOwnHundred}{92}
\newcommand{\BoCogOwnHundredHi}{93}
\newcommand{\BoCogOwnHundredLo}{90}
\newcommand{\BoCogOwnOneFifty}{94}
\newcommand{\BoCogOwnOneFiftyHi}{95}
\newcommand{\BoCogOwnOneFiftyLo}{93}
\newcommand{\BoCogOwnTwentyFive}{78}
\newcommand{\BoCogOwnTwentyFiveHi}{80}
\newcommand{\BoCogOwnTwentyFiveLo}{76}
\newcommand{\BoCogOwnTwoHundred}{96}
\newcommand{\BoCogOwnTwoHundredHi}{97}
\newcommand{\BoCogOwnTwoHundredLo}{95}
\newcommand{\BoCogPoEightHundred}{99}
\newcommand{\BoCogPoEightHundredHi}{100}
\newcommand{\BoCogPoEightHundredLo}{99}
\newcommand{\BoCogPoFifty}{90}
\newcommand{\BoCogPoFiftyHi}{92}
\newcommand{\BoCogPoFiftyLo}{87}
\newcommand{\BoCogPoFourHundred}{99}
\newcommand{\BoCogPoFourHundredHi}{99}
\newcommand{\BoCogPoFourHundredLo}{99}
\newcommand{\BoCogPoHundred}{95}
\newcommand{\BoCogPoHundredHi}{96}
\newcommand{\BoCogPoHundredLo}{94}
\newcommand{\BoCogPoOneFifty}{97}
\newcommand{\BoCogPoOneFiftyHi}{98}
\newcommand{\BoCogPoOneFiftyLo}{96}
\newcommand{\BoCogPoTwentyFive}{82}
\newcommand{\BoCogPoTwentyFiveHi}{84}
\newcommand{\BoCogPoTwentyFiveLo}{80}
\newcommand{\BoCogPoTwoHundred}{98}
\newcommand{\BoCogPoTwoHundredHi}{98}
\newcommand{\BoCogPoTwoHundredLo}{97}
\newcommand{\BoCompared}{11,220}
\newcommand{\BoDiffOneFifty}{12.9}
\newcommand{\BoDiffOneFiftyHi}{13.7}
\newcommand{\BoDiffOneFiftyLo}{12.1}
\newcommand{\BoDiffOneHundred}{11.2}
\newcommand{\BoDiffOneHundredHi}{12.0}
\newcommand{\BoDiffOneHundredLo}{10.3}
\newcommand{\BoDiffOneOneFifty}{9.1}
\newcommand{\BoDiffOneOneFiftyHi}{9.8}
\newcommand{\BoDiffOneOneFiftyLo}{8.4}
\newcommand{\BoDiffThreeHundred}{14.2}
\newcommand{\BoDiffThreeHundredHi}{15.8}
\newcommand{\BoDiffThreeHundredLo}{12.5}
\newcommand{\BoDiffThreeOneFifty}{14.4}
\newcommand{\BoDiffThreeOneFiftyHi}{15.3}
\newcommand{\BoDiffThreeOneFiftyLo}{12.9}
\newcommand{\BoDiffTwoFifty}{20.1}
\newcommand{\BoDiffTwoFiftyHi}{21.3}
\newcommand{\BoDiffTwoFiftyLo}{18.9}
\newcommand{\BoDiffTwoHundred}{19.3}
\newcommand{\BoDiffTwoHundredHi}{20.5}
\newcommand{\BoDiffTwoHundredLo}{18.0}
\newcommand{\BoDiffTwoOneFifty}{16.8}
\newcommand{\BoDiffTwoOneFiftyHi}{18.0}
\newcommand{\BoDiffTwoOneFiftyLo}{15.6}
\newcommand{\BoFrozenAt}{07:06 EDT on 29 September 2026}
\newcommand{\BoIntegrity}{$1.2\times10^{-7}$}
\newcommand{\BoNCognate}{16}
\newcommand{\BoNPairs}{22,200}
\newcommand{\BoNRep}{120,649}
\newcommand{\BoNUnits}{170}
\newcommand{\BoOneAtlasEightHundred}{94}
\newcommand{\BoOneAtlasEightHundredHi}{95}
\newcommand{\BoOneAtlasEightHundredLo}{93}
\newcommand{\BoOneAtlasFifty}{80}
\newcommand{\BoOneAtlasFiftyHi}{82}
\newcommand{\BoOneAtlasFiftyLo}{79}
\newcommand{\BoOneAtlasFourHundred}{92}
\newcommand{\BoOneAtlasFourHundredHi}{93}
\newcommand{\BoOneAtlasFourHundredLo}{91}
\newcommand{\BoOneAtlasHundred}{86}
\newcommand{\BoOneAtlasHundredHi}{87}
\newcommand{\BoOneAtlasHundredLo}{85}
\newcommand{\BoOneAtlasOneFifty}{89}
\newcommand{\BoOneAtlasOneFiftyHi}{90}
\newcommand{\BoOneAtlasOneFiftyLo}{87}
\newcommand{\BoOneAtlasTwentyFive}{71}
\newcommand{\BoOneAtlasTwentyFiveHi}{72}
\newcommand{\BoOneAtlasTwentyFiveLo}{70}
\newcommand{\BoOneAtlasTwoHundred}{90}
\newcommand{\BoOneAtlasTwoHundredHi}{91}
\newcommand{\BoOneAtlasTwoHundredLo}{89}
\newcommand{\BoOneOtherEightHundred}{96}
\newcommand{\BoOneOtherEightHundredHi}{96}
\newcommand{\BoOneOtherEightHundredLo}{95}
\newcommand{\BoOneOtherFifty}{85}
\newcommand{\BoOneOtherFiftyHi}{86}
\newcommand{\BoOneOtherFiftyLo}{84}
\newcommand{\BoOneOtherFourHundred}{94}
\newcommand{\BoOneOtherFourHundredHi}{95}
\newcommand{\BoOneOtherFourHundredLo}{93}
\newcommand{\BoOneOtherHundred}{89}
\newcommand{\BoOneOtherHundredHi}{90}
\newcommand{\BoOneOtherHundredLo}{88}
\newcommand{\BoOneOtherOneFifty}{91}
\newcommand{\BoOneOtherOneFiftyHi}{92}
\newcommand{\BoOneOtherOneFiftyLo}{90}
\newcommand{\BoOneOtherTwentyFive}{77}
\newcommand{\BoOneOtherTwentyFiveHi}{78}
\newcommand{\BoOneOtherTwentyFiveLo}{75}
\newcommand{\BoOneOtherTwoHundred}{92}
\newcommand{\BoOneOtherTwoHundredHi}{93}
\newcommand{\BoOneOtherTwoHundredLo}{91}
\newcommand{\BoOneOwnEightHundred}{94}
\newcommand{\BoOneOwnEightHundredHi}{95}
\newcommand{\BoOneOwnEightHundredLo}{93}
\newcommand{\BoOneOwnFifty}{84}
\newcommand{\BoOneOwnFiftyHi}{85}
\newcommand{\BoOneOwnFiftyLo}{83}
\newcommand{\BoOneOwnFourHundred}{92}
\newcommand{\BoOneOwnFourHundredHi}{93}
\newcommand{\BoOneOwnFourHundredLo}{92}
\newcommand{\BoOneOwnHundred}{88}
\newcommand{\BoOneOwnHundredHi}{89}
\newcommand{\BoOneOwnHundredLo}{87}
\newcommand{\BoOneOwnOneFifty}{89}
\newcommand{\BoOneOwnOneFiftyHi}{90}
\newcommand{\BoOneOwnOneFiftyLo}{88}
\newcommand{\BoOneOwnTwentyFive}{76}
\newcommand{\BoOneOwnTwentyFiveHi}{77}
\newcommand{\BoOneOwnTwentyFiveLo}{74}
\newcommand{\BoOneOwnTwoHundred}{90}
\newcommand{\BoOneOwnTwoHundredHi}{91}
\newcommand{\BoOneOwnTwoHundredLo}{89}
\newcommand{\BoOnePoEightHundred}{92}
\newcommand{\BoOnePoEightHundredHi}{93}
\newcommand{\BoOnePoEightHundredLo}{92}
\newcommand{\BoOnePoFifty}{67}
\newcommand{\BoOnePoFiftyHi}{69}
\newcommand{\BoOnePoFiftyLo}{66}
\newcommand{\BoOnePoFourHundred}{88}
\newcommand{\BoOnePoFourHundredHi}{89}
\newcommand{\BoOnePoFourHundredLo}{87}
\newcommand{\BoOnePoHundred}{75}
\newcommand{\BoOnePoHundredHi}{76}
\newcommand{\BoOnePoHundredLo}{74}
\newcommand{\BoOnePoOneFifty}{80}
\newcommand{\BoOnePoOneFiftyHi}{81}
\newcommand{\BoOnePoOneFiftyLo}{78}
\newcommand{\BoOnePoTwentyFive}{61}
\newcommand{\BoOnePoTwentyFiveHi}{63}
\newcommand{\BoOnePoTwentyFiveLo}{61}
\newcommand{\BoOnePoTwoHundred}{82}
\newcommand{\BoOnePoTwoHundredHi}{83}
\newcommand{\BoOnePoTwoHundredLo}{81}
\newcommand{\BoPerfectTwo}{97}
\newcommand{\BoRandomTwo}{1.6}
\newcommand{\BoRepB}{5,077}
\newcommand{\BoRepCDEight}{56,637}
\newcommand{\BoRepCDFour}{27,388}
\newcommand{\BoRepMono}{21,021}
\newcommand{\BoRepNK}{10,526}
\newcommand{\BoRepShare}{3.2}
\newcommand{\BoSaveOne}{2.9}
\newcommand{\BoSaveOneHi}{3.1}
\newcommand{\BoSaveOneLo}{2.6}
\newcommand{\BoSaveTwo}{3.0}
\newcommand{\BoSaveTwoHi}{3.2}
\newcommand{\BoSaveTwoLo}{2.8}
\newcommand{\BoTargetOne}{73}
\newcommand{\BoTargetTwo}{34}
\newcommand{\BoThreeAtlasEightHundred}{73}
\newcommand{\BoThreeAtlasEightHundredHi}{76}
\newcommand{\BoThreeAtlasEightHundredLo}{71}
\newcommand{\BoThreeAtlasFifty}{24}
\newcommand{\BoThreeAtlasFiftyHi}{27}
\newcommand{\BoThreeAtlasFiftyLo}{22}
\newcommand{\BoThreeAtlasFourHundred}{64}
\newcommand{\BoThreeAtlasFourHundredHi}{67}
\newcommand{\BoThreeAtlasFourHundredLo}{62}
\newcommand{\BoThreeAtlasHundred}{39}
\newcommand{\BoThreeAtlasHundredHi}{41}
\newcommand{\BoThreeAtlasHundredLo}{36}
\newcommand{\BoThreeAtlasOneFifty}{47}
\newcommand{\BoThreeAtlasOneFiftyHi}{50}
\newcommand{\BoThreeAtlasOneFiftyLo}{44}
\newcommand{\BoThreeAtlasTwentyFive}{13}
\newcommand{\BoThreeAtlasTwentyFiveHi}{14}
\newcommand{\BoThreeAtlasTwentyFiveLo}{11}
\newcommand{\BoThreeAtlasTwoHundred}{53}
\newcommand{\BoThreeAtlasTwoHundredHi}{56}
\newcommand{\BoThreeAtlasTwoHundredLo}{50}
\newcommand{\BoThreeOtherEightHundred}{74}
\newcommand{\BoThreeOtherEightHundredHi}{77}
\newcommand{\BoThreeOtherEightHundredLo}{70}
\newcommand{\BoThreeOtherFifty}{24}
\newcommand{\BoThreeOtherFiftyHi}{26}
\newcommand{\BoThreeOtherFiftyLo}{22}
\newcommand{\BoThreeOtherFourHundred}{63}
\newcommand{\BoThreeOtherFourHundredHi}{66}
\newcommand{\BoThreeOtherFourHundredLo}{60}
\newcommand{\BoThreeOtherHundred}{36}
\newcommand{\BoThreeOtherHundredHi}{39}
\newcommand{\BoThreeOtherHundredLo}{33}
\newcommand{\BoThreeOtherOneFifty}{44}
\newcommand{\BoThreeOtherOneFiftyHi}{47}
\newcommand{\BoThreeOtherOneFiftyLo}{40}
\newcommand{\BoThreeOtherTwentyFive}{14}
\newcommand{\BoThreeOtherTwentyFiveHi}{16}
\newcommand{\BoThreeOtherTwentyFiveLo}{13}
\newcommand{\BoThreeOtherTwoHundred}{50}
\newcommand{\BoThreeOtherTwoHundredHi}{53}
\newcommand{\BoThreeOtherTwoHundredLo}{47}
\newcommand{\BoThreeOwnEightHundred}{67}
\newcommand{\BoThreeOwnEightHundredHi}{70}
\newcommand{\BoThreeOwnEightHundredLo}{63}
\newcommand{\BoThreeOwnFifty}{23}
\newcommand{\BoThreeOwnFiftyHi}{25}
\newcommand{\BoThreeOwnFiftyLo}{21}
\newcommand{\BoThreeOwnFourHundred}{57}
\newcommand{\BoThreeOwnFourHundredHi}{60}
\newcommand{\BoThreeOwnFourHundredLo}{53}
\newcommand{\BoThreeOwnHundred}{34}
\newcommand{\BoThreeOwnHundredHi}{36}
\newcommand{\BoThreeOwnHundredLo}{31}
\newcommand{\BoThreeOwnOneFifty}{40}
\newcommand{\BoThreeOwnOneFiftyHi}{43}
\newcommand{\BoThreeOwnOneFiftyLo}{37}
\newcommand{\BoThreeOwnTwentyFive}{14}
\newcommand{\BoThreeOwnTwentyFiveHi}{16}
\newcommand{\BoThreeOwnTwentyFiveLo}{13}
\newcommand{\BoThreeOwnTwoHundred}{45}
\newcommand{\BoThreeOwnTwoHundredHi}{48}
\newcommand{\BoThreeOwnTwoHundredLo}{42}
\newcommand{\BoThreePoEightHundred}{65}
\newcommand{\BoThreePoEightHundredHi}{67}
\newcommand{\BoThreePoEightHundredLo}{61}
\newcommand{\BoThreePoFifty}{11}
\newcommand{\BoThreePoFiftyHi}{13}
\newcommand{\BoThreePoFiftyLo}{9.9}
\newcommand{\BoThreePoFourHundred}{54}
\newcommand{\BoThreePoFourHundredHi}{58}
\newcommand{\BoThreePoFourHundredLo}{50}
\newcommand{\BoThreePoHundred}{25}
\newcommand{\BoThreePoHundredHi}{27}
\newcommand{\BoThreePoHundredLo}{23}
\newcommand{\BoThreePoOneFifty}{33}
\newcommand{\BoThreePoOneFiftyHi}{36}
\newcommand{\BoThreePoOneFiftyLo}{30}
\newcommand{\BoThreePoTwentyFive}{4.7}
\newcommand{\BoThreePoTwentyFiveHi}{5.6}
\newcommand{\BoThreePoTwentyFiveLo}{3.8}
\newcommand{\BoThreePoTwoHundred}{39}
\newcommand{\BoThreePoTwoHundredHi}{42}
\newcommand{\BoThreePoTwoHundredLo}{36}
\newcommand{\BoTwoAtlasEightHundred}{67}
\newcommand{\BoTwoAtlasEightHundredHi}{68}
\newcommand{\BoTwoAtlasEightHundredLo}{65}
\newcommand{\BoTwoAtlasFifty}{35}
\newcommand{\BoTwoAtlasFiftyHi}{37}
\newcommand{\BoTwoAtlasFiftyLo}{34}
\newcommand{\BoTwoAtlasFourHundred}{62}
\newcommand{\BoTwoAtlasFourHundredHi}{63}
\newcommand{\BoTwoAtlasFourHundredLo}{60}
\newcommand{\BoTwoAtlasHundred}{47}
\newcommand{\BoTwoAtlasHundredHi}{49}
\newcommand{\BoTwoAtlasHundredLo}{46}
\newcommand{\BoTwoAtlasOneFifty}{52}
\newcommand{\BoTwoAtlasOneFiftyHi}{54}
\newcommand{\BoTwoAtlasOneFiftyLo}{51}
\newcommand{\BoTwoAtlasTwentyFive}{23}
\newcommand{\BoTwoAtlasTwentyFiveHi}{24}
\newcommand{\BoTwoAtlasTwentyFiveLo}{21}
\newcommand{\BoTwoAtlasTwoHundred}{55}
\newcommand{\BoTwoAtlasTwoHundredHi}{56}
\newcommand{\BoTwoAtlasTwoHundredLo}{54}
\newcommand{\BoTwoOtherEightHundred}{67}
\newcommand{\BoTwoOtherEightHundredHi}{69}
\newcommand{\BoTwoOtherEightHundredLo}{66}
\newcommand{\BoTwoOtherFifty}{38}
\newcommand{\BoTwoOtherFiftyHi}{40}
\newcommand{\BoTwoOtherFiftyLo}{36}
\newcommand{\BoTwoOtherFourHundred}{61}
\newcommand{\BoTwoOtherFourHundredHi}{63}
\newcommand{\BoTwoOtherFourHundredLo}{60}
\newcommand{\BoTwoOtherHundred}{46}
\newcommand{\BoTwoOtherHundredHi}{48}
\newcommand{\BoTwoOtherHundredLo}{44}
\newcommand{\BoTwoOtherOneFifty}{50}
\newcommand{\BoTwoOtherOneFiftyHi}{52}
\newcommand{\BoTwoOtherOneFiftyLo}{49}
\newcommand{\BoTwoOtherTwentyFive}{28}
\newcommand{\BoTwoOtherTwentyFiveHi}{30}
\newcommand{\BoTwoOtherTwentyFiveLo}{27}
\newcommand{\BoTwoOtherTwoHundred}{54}
\newcommand{\BoTwoOtherTwoHundredHi}{55}
\newcommand{\BoTwoOtherTwoHundredLo}{52}
\newcommand{\BoTwoOwnEightHundred}{57}
\newcommand{\BoTwoOwnEightHundredHi}{59}
\newcommand{\BoTwoOwnEightHundredLo}{54}
\newcommand{\BoTwoOwnFifty}{36}
\newcommand{\BoTwoOwnFiftyHi}{37}
\newcommand{\BoTwoOwnFiftyLo}{35}
\newcommand{\BoTwoOwnFourHundred}{52}
\newcommand{\BoTwoOwnFourHundredHi}{54}
\newcommand{\BoTwoOwnFourHundredLo}{50}
\newcommand{\BoTwoOwnHundred}{41}
\newcommand{\BoTwoOwnHundredHi}{43}
\newcommand{\BoTwoOwnHundredLo}{40}
\newcommand{\BoTwoOwnOneFifty}{44}
\newcommand{\BoTwoOwnOneFiftyHi}{46}
\newcommand{\BoTwoOwnOneFiftyLo}{43}
\newcommand{\BoTwoOwnTwentyFive}{27}
\newcommand{\BoTwoOwnTwentyFiveHi}{28}
\newcommand{\BoTwoOwnTwentyFiveLo}{25}
\newcommand{\BoTwoOwnTwoHundred}{46}
\newcommand{\BoTwoOwnTwoHundredHi}{48}
\newcommand{\BoTwoOwnTwoHundredLo}{45}
\newcommand{\BoTwoPoEightHundred}{61}
\newcommand{\BoTwoPoEightHundredHi}{62}
\newcommand{\BoTwoPoEightHundredLo}{60}
\newcommand{\BoTwoPoFifty}{15}
\newcommand{\BoTwoPoFiftyHi}{16}
\newcommand{\BoTwoPoFiftyLo}{14}
\newcommand{\BoTwoPoFourHundred}{52}
\newcommand{\BoTwoPoFourHundredHi}{53}
\newcommand{\BoTwoPoFourHundredLo}{51}
\newcommand{\BoTwoPoHundred}{28}
\newcommand{\BoTwoPoHundredHi}{29}
\newcommand{\BoTwoPoHundredLo}{27}
\newcommand{\BoTwoPoOneFifty}{35}
\newcommand{\BoTwoPoOneFiftyHi}{36}
\newcommand{\BoTwoPoOneFiftyLo}{34}
\newcommand{\BoTwoPoTwentyFive}{6.9}
\newcommand{\BoTwoPoTwentyFiveHi}{7.8}
\newcommand{\BoTwoPoTwentyFiveLo}{6.0}
\newcommand{\BoTwoPoTwoHundred}{40}
\newcommand{\BoTwoPoTwoHundredHi}{41}
\newcommand{\BoTwoPoTwoHundredLo}{39}
\newcommand{\BoTypeBAtlasFifty}{63}
\newcommand{\BoTypeBAtlasHundred}{67}
\newcommand{\BoTypeBAtlasOneFifty}{70}
\newcommand{\BoTypeBPoFifty}{57}
\newcommand{\BoTypeBPoHundred}{61}
\newcommand{\BoTypeBPoOneFifty}{63}
\newcommand{\BoTypeCDEightAtlasFifty}{85}
\newcommand{\BoTypeCDEightAtlasHundred}{91}
\newcommand{\BoTypeCDEightAtlasOneFifty}{94}
\newcommand{\BoTypeCDEightPoFifty}{70}
\newcommand{\BoTypeCDEightPoHundred}{79}
\newcommand{\BoTypeCDEightPoOneFifty}{83}
\newcommand{\BoTypeCDFourAtlasFifty}{87}
\newcommand{\BoTypeCDFourAtlasHundred}{92}
\newcommand{\BoTypeCDFourAtlasOneFifty}{93}
\newcommand{\BoTypeCDFourPoFifty}{71}
\newcommand{\BoTypeCDFourPoHundred}{77}
\newcommand{\BoTypeCDFourPoOneFifty}{81}
\newcommand{\BoTypeMonoAtlasFifty}{66}
\newcommand{\BoTypeMonoAtlasHundred}{73}
\newcommand{\BoTypeMonoAtlasOneFifty}{76}
\newcommand{\BoTypeMonoPoFifty}{64}
\newcommand{\BoTypeMonoPoHundred}{72}
\newcommand{\BoTypeMonoPoOneFifty}{77}
\newcommand{\BoTypeNKAtlasFifty}{71}
\newcommand{\BoTypeNKAtlasHundred}{78}
\newcommand{\BoTypeNKAtlasOneFifty}{82}
\newcommand{\BoTypeNKPoFifty}{64}
\newcommand{\BoTypeNKPoHundred}{69}
\newcommand{\BoTypeNKPoOneFifty}{73}
\newcommand{\BoUnitsB}{34}
\newcommand{\BoUnitsCDEight}{34}
\newcommand{\BoUnitsCDFour}{34}
\newcommand{\BoUnitsMono}{34}
\newcommand{\BoUnitsNK}{34}
\newcommand{\CalCon}{1.00}
\newcommand{\CalDfMax}{0.95}
\newcommand{\CalDfMin}{0.64}
\newcommand{\CbBoot}{20,000}
\newcommand{\CbDiffOneLoMin}{7.1}
\newcommand{\CbDiffOneMax}{11.8}
\newcommand{\CbDiffOneMin}{8.3}
\newcommand{\CbDiffTwoLoMin}{15.2}
\newcommand{\CbDiffTwoMax}{21.0}
\newcommand{\CbDiffTwoMin}{16.5}
\newcommand{\CbGlobalThreshold}{$1.1\times10^{-3}$}
\newcommand{\CbHolmMax}{$\leq2\times10^{-4}$}
\newcommand{\CbHolmOne}{$\leq2\times10^{-4}$}
\newcommand{\CbHolmThree}{$\leq2\times10^{-4}$}
\newcommand{\CbHolmTwo}{$\leq2\times10^{-4}$}
\newcommand{\CbMajFdrAllPooled}{58}
\newcommand{\CbMajFdrAllSample}{58}
\newcommand{\CbMajFdrAllType}{60}
\newcommand{\CbMajFdrAllUnit}{61}
\newcommand{\CbMajFdrUnitPooled}{57}
\newcommand{\CbMajFdrUnitSample}{58}
\newcommand{\CbMajFdrUnitType}{60}
\newcommand{\CbMajFdrUnitUnit}{62}
\newcommand{\CbMajPrimaryPooled}{57}
\newcommand{\CbMajPrimarySample}{57}
\newcommand{\CbMajPrimaryType}{58}
\newcommand{\CbMajPrimaryUnit}{58}
\newcommand{\CbMajThousandthPooled}{58}
\newcommand{\CbMajThousandthSample}{58}
\newcommand{\CbMajThousandthType}{60}
\newcommand{\CbMajThousandthUnit}{61}
\newcommand{\CbMajTwentiethPooled}{55}
\newcommand{\CbMajTwentiethSample}{56}
\newcommand{\CbMajTwentiethType}{56}
\newcommand{\CbMajTwentiethUnit}{56}
\newcommand{\CbNFdrAll}{80,156}
\newcommand{\CbNFdrUnit}{83,863}
\newcommand{\CbNPrimary}{120,649}
\newcommand{\CbNThousandth}{79,414}
\newcommand{\CbNTwentieth}{189,472}
\newcommand{\CbOneFdrAllPooledAtlasFifty}{83}
\newcommand{\CbOneFdrAllPooledAtlasHundred}{89}
\newcommand{\CbOneFdrAllPooledAtlasOneFifty}{91}
\newcommand{\CbOneFdrAllPooledDiffFifty}{12.7}
\newcommand{\CbOneFdrAllPooledDiffFiftyHi}{13.6}
\newcommand{\CbOneFdrAllPooledDiffFiftyLo}{11.7}
\newcommand{\CbOneFdrAllPooledDiffHundred}{10.3}
\newcommand{\CbOneFdrAllPooledDiffHundredHi}{11.1}
\newcommand{\CbOneFdrAllPooledDiffHundredLo}{9.3}
\newcommand{\CbOneFdrAllPooledDiffOneFifty}{7.8}
\newcommand{\CbOneFdrAllPooledDiffOneFiftyHi}{8.5}
\newcommand{\CbOneFdrAllPooledDiffOneFiftyLo}{7.0}
\newcommand{\CbOneFdrAllPooledPoFifty}{70}
\newcommand{\CbOneFdrAllPooledPoHundred}{79}
\newcommand{\CbOneFdrAllPooledPoOneFifty}{84}
\newcommand{\CbOneFdrAllSampleAtlasFifty}{83}
\newcommand{\CbOneFdrAllSampleAtlasHundred}{89}
\newcommand{\CbOneFdrAllSampleAtlasOneFifty}{91}
\newcommand{\CbOneFdrAllSampleDiffFifty}{12.5}
\newcommand{\CbOneFdrAllSampleDiffFiftyHi}{13.5}
\newcommand{\CbOneFdrAllSampleDiffFiftyLo}{11.5}
\newcommand{\CbOneFdrAllSampleDiffHundred}{10.0}
\newcommand{\CbOneFdrAllSampleDiffHundredHi}{10.9}
\newcommand{\CbOneFdrAllSampleDiffHundredLo}{9.0}
\newcommand{\CbOneFdrAllSampleDiffOneFifty}{7.6}
\newcommand{\CbOneFdrAllSampleDiffOneFiftyHi}{8.4}
\newcommand{\CbOneFdrAllSampleDiffOneFiftyLo}{6.8}
\newcommand{\CbOneFdrAllSamplePoFifty}{71}
\newcommand{\CbOneFdrAllSamplePoHundred}{79}
\newcommand{\CbOneFdrAllSamplePoOneFifty}{84}
\newcommand{\CbOneFdrAllTypeAtlasFifty}{76}
\newcommand{\CbOneFdrAllTypeAtlasHundred}{82}
\newcommand{\CbOneFdrAllTypeAtlasOneFifty}{85}
\newcommand{\CbOneFdrAllTypeDiffFifty}{9.8}
\newcommand{\CbOneFdrAllTypeDiffFiftyHi}{10.4}
\newcommand{\CbOneFdrAllTypeDiffFiftyLo}{9.2}
\newcommand{\CbOneFdrAllTypeDiffHundred}{8.9}
\newcommand{\CbOneFdrAllTypeDiffHundredHi}{9.6}
\newcommand{\CbOneFdrAllTypeDiffHundredLo}{8.1}
\newcommand{\CbOneFdrAllTypeDiffOneFifty}{6.9}
\newcommand{\CbOneFdrAllTypeDiffOneFiftyHi}{7.5}
\newcommand{\CbOneFdrAllTypeDiffOneFiftyLo}{6.2}
\newcommand{\CbOneFdrAllTypePoFifty}{67}
\newcommand{\CbOneFdrAllTypePoHundred}{73}
\newcommand{\CbOneFdrAllTypePoOneFifty}{78}
\newcommand{\CbOneFdrAllUnitAtlasFifty}{77}
\newcommand{\CbOneFdrAllUnitAtlasHundred}{83}
\newcommand{\CbOneFdrAllUnitAtlasOneFifty}{86}
\newcommand{\CbOneFdrAllUnitDiffFifty}{9.9}
\newcommand{\CbOneFdrAllUnitDiffFiftyHi}{10.7}
\newcommand{\CbOneFdrAllUnitDiffFiftyLo}{9.2}
\newcommand{\CbOneFdrAllUnitDiffHundred}{8.7}
\newcommand{\CbOneFdrAllUnitDiffHundredHi}{9.5}
\newcommand{\CbOneFdrAllUnitDiffHundredLo}{8.0}
\newcommand{\CbOneFdrAllUnitDiffOneFifty}{6.5}
\newcommand{\CbOneFdrAllUnitDiffOneFiftyHi}{7.1}
\newcommand{\CbOneFdrAllUnitDiffOneFiftyLo}{6.0}
\newcommand{\CbOneFdrAllUnitPoFifty}{67}
\newcommand{\CbOneFdrAllUnitPoHundred}{74}
\newcommand{\CbOneFdrAllUnitPoOneFifty}{79}
\newcommand{\CbOneFdrUnitPooledAtlasFifty}{83}
\newcommand{\CbOneFdrUnitPooledAtlasHundred}{89}
\newcommand{\CbOneFdrUnitPooledAtlasOneFifty}{92}
\newcommand{\CbOneFdrUnitPooledDiffFifty}{13.4}
\newcommand{\CbOneFdrUnitPooledDiffFiftyHi}{14.3}
\newcommand{\CbOneFdrUnitPooledDiffFiftyLo}{12.4}
\newcommand{\CbOneFdrUnitPooledDiffHundred}{10.9}
\newcommand{\CbOneFdrUnitPooledDiffHundredHi}{11.7}
\newcommand{\CbOneFdrUnitPooledDiffHundredLo}{9.9}
\newcommand{\CbOneFdrUnitPooledDiffOneFifty}{8.5}
\newcommand{\CbOneFdrUnitPooledDiffOneFiftyHi}{9.3}
\newcommand{\CbOneFdrUnitPooledDiffOneFiftyLo}{7.6}
\newcommand{\CbOneFdrUnitPooledPoFifty}{70}
\newcommand{\CbOneFdrUnitPooledPoHundred}{78}
\newcommand{\CbOneFdrUnitPooledPoOneFifty}{83}
\newcommand{\CbOneFdrUnitSampleAtlasFifty}{84}
\newcommand{\CbOneFdrUnitSampleAtlasHundred}{89}
\newcommand{\CbOneFdrUnitSampleAtlasOneFifty}{92}
\newcommand{\CbOneFdrUnitSampleDiffFifty}{13.1}
\newcommand{\CbOneFdrUnitSampleDiffFiftyHi}{14.1}
\newcommand{\CbOneFdrUnitSampleDiffFiftyLo}{12.1}
\newcommand{\CbOneFdrUnitSampleDiffHundred}{10.6}
\newcommand{\CbOneFdrUnitSampleDiffHundredHi}{11.6}
\newcommand{\CbOneFdrUnitSampleDiffHundredLo}{9.5}
\newcommand{\CbOneFdrUnitSampleDiffOneFifty}{8.2}
\newcommand{\CbOneFdrUnitSampleDiffOneFiftyHi}{9.1}
\newcommand{\CbOneFdrUnitSampleDiffOneFiftyLo}{7.2}
\newcommand{\CbOneFdrUnitSamplePoFifty}{71}
\newcommand{\CbOneFdrUnitSamplePoHundred}{79}
\newcommand{\CbOneFdrUnitSamplePoOneFifty}{84}
\newcommand{\CbOneFdrUnitTypeAtlasFifty}{76}
\newcommand{\CbOneFdrUnitTypeAtlasHundred}{82}
\newcommand{\CbOneFdrUnitTypeAtlasOneFifty}{85}
\newcommand{\CbOneFdrUnitTypeDiffFifty}{9.8}
\newcommand{\CbOneFdrUnitTypeDiffFiftyHi}{10.6}
\newcommand{\CbOneFdrUnitTypeDiffFiftyLo}{9.1}
\newcommand{\CbOneFdrUnitTypeDiffHundred}{8.7}
\newcommand{\CbOneFdrUnitTypeDiffHundredHi}{9.6}
\newcommand{\CbOneFdrUnitTypeDiffHundredLo}{7.7}
\newcommand{\CbOneFdrUnitTypeDiffOneFifty}{6.8}
\newcommand{\CbOneFdrUnitTypeDiffOneFiftyHi}{7.6}
\newcommand{\CbOneFdrUnitTypeDiffOneFiftyLo}{5.9}
\newcommand{\CbOneFdrUnitTypePoFifty}{67}
\newcommand{\CbOneFdrUnitTypePoHundred}{74}
\newcommand{\CbOneFdrUnitTypePoOneFifty}{78}
\newcommand{\CbOneFdrUnitUnitAtlasFifty}{78}
\newcommand{\CbOneFdrUnitUnitAtlasHundred}{84}
\newcommand{\CbOneFdrUnitUnitAtlasOneFifty}{86}
\newcommand{\CbOneFdrUnitUnitDiffFifty}{10.4}
\newcommand{\CbOneFdrUnitUnitDiffFiftyHi}{11.9}
\newcommand{\CbOneFdrUnitUnitDiffFiftyLo}{9.0}
\newcommand{\CbOneFdrUnitUnitDiffHundred}{8.3}
\newcommand{\CbOneFdrUnitUnitDiffHundredHi}{9.6}
\newcommand{\CbOneFdrUnitUnitDiffHundredLo}{7.1}
\newcommand{\CbOneFdrUnitUnitDiffOneFifty}{6.0}
\newcommand{\CbOneFdrUnitUnitDiffOneFiftyHi}{6.9}
\newcommand{\CbOneFdrUnitUnitDiffOneFiftyLo}{5.0}
\newcommand{\CbOneFdrUnitUnitPoFifty}{68}
\newcommand{\CbOneFdrUnitUnitPoHundred}{75}
\newcommand{\CbOneFdrUnitUnitPoOneFifty}{80}
\newcommand{\CbOnePrimaryPooledAtlasFifty}{80}
\newcommand{\CbOnePrimaryPooledAtlasHundred}{86}
\newcommand{\CbOnePrimaryPooledAtlasOneFifty}{89}
\newcommand{\CbOnePrimaryPooledDiffFifty}{12.9}
\newcommand{\CbOnePrimaryPooledDiffFiftyHi}{13.7}
\newcommand{\CbOnePrimaryPooledDiffFiftyLo}{12.1}
\newcommand{\CbOnePrimaryPooledDiffHundred}{11.2}
\newcommand{\CbOnePrimaryPooledDiffHundredHi}{12.0}
\newcommand{\CbOnePrimaryPooledDiffHundredLo}{10.3}
\newcommand{\CbOnePrimaryPooledDiffOneFifty}{9.1}
\newcommand{\CbOnePrimaryPooledDiffOneFiftyHi}{9.8}
\newcommand{\CbOnePrimaryPooledDiffOneFiftyLo}{8.4}
\newcommand{\CbOnePrimaryPooledPoFifty}{67}
\newcommand{\CbOnePrimaryPooledPoHundred}{75}
\newcommand{\CbOnePrimaryPooledPoOneFifty}{80}
\newcommand{\CbOnePrimarySampleAtlasFifty}{80}
\newcommand{\CbOnePrimarySampleAtlasHundred}{86}
\newcommand{\CbOnePrimarySampleAtlasOneFifty}{89}
\newcommand{\CbOnePrimarySampleDiffFifty}{12.8}
\newcommand{\CbOnePrimarySampleDiffFiftyHi}{13.7}
\newcommand{\CbOnePrimarySampleDiffFiftyLo}{11.9}
\newcommand{\CbOnePrimarySampleDiffHundred}{11.1}
\newcommand{\CbOnePrimarySampleDiffHundredHi}{11.9}
\newcommand{\CbOnePrimarySampleDiffHundredLo}{10.1}
\newcommand{\CbOnePrimarySampleDiffOneFifty}{9.0}
\newcommand{\CbOnePrimarySampleDiffOneFiftyHi}{9.7}
\newcommand{\CbOnePrimarySampleDiffOneFiftyLo}{8.1}
\newcommand{\CbOnePrimarySamplePoFifty}{68}
\newcommand{\CbOnePrimarySamplePoHundred}{75}
\newcommand{\CbOnePrimarySamplePoOneFifty}{80}
\newcommand{\CbOnePrimaryTypeAtlasFifty}{75}
\newcommand{\CbOnePrimaryTypeAtlasHundred}{80}
\newcommand{\CbOnePrimaryTypeAtlasOneFifty}{83}
\newcommand{\CbOnePrimaryTypeDiffFifty}{10.3}
\newcommand{\CbOnePrimaryTypeDiffFiftyHi}{10.8}
\newcommand{\CbOnePrimaryTypeDiffFiftyLo}{9.7}
\newcommand{\CbOnePrimaryTypeDiffHundred}{9.9}
\newcommand{\CbOnePrimaryTypeDiffHundredHi}{10.4}
\newcommand{\CbOnePrimaryTypeDiffHundredLo}{9.3}
\newcommand{\CbOnePrimaryTypeDiffOneFifty}{8.2}
\newcommand{\CbOnePrimaryTypeDiffOneFiftyHi}{8.7}
\newcommand{\CbOnePrimaryTypeDiffOneFiftyLo}{7.6}
\newcommand{\CbOnePrimaryTypePoFifty}{64}
\newcommand{\CbOnePrimaryTypePoHundred}{70}
\newcommand{\CbOnePrimaryTypePoOneFifty}{75}
\newcommand{\CbOnePrimaryUnitAtlasFifty}{75}
\newcommand{\CbOnePrimaryUnitAtlasHundred}{81}
\newcommand{\CbOnePrimaryUnitAtlasOneFifty}{84}
\newcommand{\CbOnePrimaryUnitDiffFifty}{10.3}
\newcommand{\CbOnePrimaryUnitDiffFiftyHi}{10.9}
\newcommand{\CbOnePrimaryUnitDiffFiftyLo}{9.7}
\newcommand{\CbOnePrimaryUnitDiffHundred}{9.9}
\newcommand{\CbOnePrimaryUnitDiffHundredHi}{10.4}
\newcommand{\CbOnePrimaryUnitDiffHundredLo}{9.3}
\newcommand{\CbOnePrimaryUnitDiffOneFifty}{7.9}
\newcommand{\CbOnePrimaryUnitDiffOneFiftyHi}{8.3}
\newcommand{\CbOnePrimaryUnitDiffOneFiftyLo}{7.4}
\newcommand{\CbOnePrimaryUnitPoFifty}{65}
\newcommand{\CbOnePrimaryUnitPoHundred}{71}
\newcommand{\CbOnePrimaryUnitPoOneFifty}{76}
\newcommand{\CbOneThousandthPooledAtlasFifty}{83}
\newcommand{\CbOneThousandthPooledAtlasHundred}{89}
\newcommand{\CbOneThousandthPooledAtlasOneFifty}{91}
\newcommand{\CbOneThousandthPooledDiffFifty}{12.7}
\newcommand{\CbOneThousandthPooledDiffFiftyHi}{13.6}
\newcommand{\CbOneThousandthPooledDiffFiftyLo}{11.7}
\newcommand{\CbOneThousandthPooledDiffHundred}{10.2}
\newcommand{\CbOneThousandthPooledDiffHundredHi}{11.0}
\newcommand{\CbOneThousandthPooledDiffHundredLo}{9.3}
\newcommand{\CbOneThousandthPooledDiffOneFifty}{7.8}
\newcommand{\CbOneThousandthPooledDiffOneFiftyHi}{8.5}
\newcommand{\CbOneThousandthPooledDiffOneFiftyLo}{7.0}
\newcommand{\CbOneThousandthPooledPoFifty}{70}
\newcommand{\CbOneThousandthPooledPoHundred}{79}
\newcommand{\CbOneThousandthPooledPoOneFifty}{84}
\newcommand{\CbOneThousandthSampleAtlasFifty}{83}
\newcommand{\CbOneThousandthSampleAtlasHundred}{89}
\newcommand{\CbOneThousandthSampleAtlasOneFifty}{92}
\newcommand{\CbOneThousandthSampleDiffFifty}{12.5}
\newcommand{\CbOneThousandthSampleDiffFiftyHi}{13.5}
\newcommand{\CbOneThousandthSampleDiffFiftyLo}{11.5}
\newcommand{\CbOneThousandthSampleDiffHundred}{10.0}
\newcommand{\CbOneThousandthSampleDiffHundredHi}{10.9}
\newcommand{\CbOneThousandthSampleDiffHundredLo}{9.0}
\newcommand{\CbOneThousandthSampleDiffOneFifty}{7.6}
\newcommand{\CbOneThousandthSampleDiffOneFiftyHi}{8.4}
\newcommand{\CbOneThousandthSampleDiffOneFiftyLo}{6.8}
\newcommand{\CbOneThousandthSamplePoFifty}{71}
\newcommand{\CbOneThousandthSamplePoHundred}{79}
\newcommand{\CbOneThousandthSamplePoOneFifty}{84}
\newcommand{\CbOneThousandthTypeAtlasFifty}{76}
\newcommand{\CbOneThousandthTypeAtlasHundred}{82}
\newcommand{\CbOneThousandthTypeAtlasOneFifty}{85}
\newcommand{\CbOneThousandthTypeDiffFifty}{9.7}
\newcommand{\CbOneThousandthTypeDiffFiftyHi}{10.3}
\newcommand{\CbOneThousandthTypeDiffFiftyLo}{9.1}
\newcommand{\CbOneThousandthTypeDiffHundred}{8.8}
\newcommand{\CbOneThousandthTypeDiffHundredHi}{9.5}
\newcommand{\CbOneThousandthTypeDiffHundredLo}{8.0}
\newcommand{\CbOneThousandthTypeDiffOneFifty}{6.8}
\newcommand{\CbOneThousandthTypeDiffOneFiftyHi}{7.4}
\newcommand{\CbOneThousandthTypeDiffOneFiftyLo}{6.1}
\newcommand{\CbOneThousandthTypePoFifty}{67}
\newcommand{\CbOneThousandthTypePoHundred}{73}
\newcommand{\CbOneThousandthTypePoOneFifty}{78}
\newcommand{\CbOneThousandthUnitAtlasFifty}{77}
\newcommand{\CbOneThousandthUnitAtlasHundred}{83}
\newcommand{\CbOneThousandthUnitAtlasOneFifty}{86}
\newcommand{\CbOneThousandthUnitDiffFifty}{9.9}
\newcommand{\CbOneThousandthUnitDiffFiftyHi}{10.6}
\newcommand{\CbOneThousandthUnitDiffFiftyLo}{9.2}
\newcommand{\CbOneThousandthUnitDiffHundred}{8.6}
\newcommand{\CbOneThousandthUnitDiffHundredHi}{9.4}
\newcommand{\CbOneThousandthUnitDiffHundredLo}{7.9}
\newcommand{\CbOneThousandthUnitDiffOneFifty}{6.5}
\newcommand{\CbOneThousandthUnitDiffOneFiftyHi}{7.1}
\newcommand{\CbOneThousandthUnitDiffOneFiftyLo}{5.9}
\newcommand{\CbOneThousandthUnitPoFifty}{67}
\newcommand{\CbOneThousandthUnitPoHundred}{75}
\newcommand{\CbOneThousandthUnitPoOneFifty}{79}
\newcommand{\CbOneTwentiethPooledAtlasFifty}{76}
\newcommand{\CbOneTwentiethPooledAtlasHundred}{82}
\newcommand{\CbOneTwentiethPooledAtlasOneFifty}{85}
\newcommand{\CbOneTwentiethPooledDiffFifty}{12.5}
\newcommand{\CbOneTwentiethPooledDiffFiftyHi}{13.2}
\newcommand{\CbOneTwentiethPooledDiffFiftyLo}{11.7}
\newcommand{\CbOneTwentiethPooledDiffHundred}{11.8}
\newcommand{\CbOneTwentiethPooledDiffHundredHi}{12.4}
\newcommand{\CbOneTwentiethPooledDiffHundredLo}{11.0}
\newcommand{\CbOneTwentiethPooledDiffOneFifty}{10.2}
\newcommand{\CbOneTwentiethPooledDiffOneFiftyHi}{10.9}
\newcommand{\CbOneTwentiethPooledDiffOneFiftyLo}{9.5}
\newcommand{\CbOneTwentiethPooledPoFifty}{64}
\newcommand{\CbOneTwentiethPooledPoHundred}{70}
\newcommand{\CbOneTwentiethPooledPoOneFifty}{74}
\newcommand{\CbOneTwentiethSampleAtlasFifty}{77}
\newcommand{\CbOneTwentiethSampleAtlasHundred}{82}
\newcommand{\CbOneTwentiethSampleAtlasOneFifty}{85}
\newcommand{\CbOneTwentiethSampleDiffFifty}{12.4}
\newcommand{\CbOneTwentiethSampleDiffFiftyHi}{13.2}
\newcommand{\CbOneTwentiethSampleDiffFiftyLo}{11.6}
\newcommand{\CbOneTwentiethSampleDiffHundred}{11.7}
\newcommand{\CbOneTwentiethSampleDiffHundredHi}{12.4}
\newcommand{\CbOneTwentiethSampleDiffHundredLo}{10.9}
\newcommand{\CbOneTwentiethSampleDiffOneFifty}{10.1}
\newcommand{\CbOneTwentiethSampleDiffOneFiftyHi}{10.8}
\newcommand{\CbOneTwentiethSampleDiffOneFiftyLo}{9.4}
\newcommand{\CbOneTwentiethSamplePoFifty}{64}
\newcommand{\CbOneTwentiethSamplePoHundred}{70}
\newcommand{\CbOneTwentiethSamplePoOneFifty}{75}
\newcommand{\CbOneTwentiethTypeAtlasFifty}{72}
\newcommand{\CbOneTwentiethTypeAtlasHundred}{77}
\newcommand{\CbOneTwentiethTypeAtlasOneFifty}{80}
\newcommand{\CbOneTwentiethTypeDiffFifty}{10.2}
\newcommand{\CbOneTwentiethTypeDiffFiftyHi}{10.7}
\newcommand{\CbOneTwentiethTypeDiffFiftyLo}{9.7}
\newcommand{\CbOneTwentiethTypeDiffHundred}{10.6}
\newcommand{\CbOneTwentiethTypeDiffHundredHi}{11.1}
\newcommand{\CbOneTwentiethTypeDiffHundredLo}{10.0}
\newcommand{\CbOneTwentiethTypeDiffOneFifty}{9.3}
\newcommand{\CbOneTwentiethTypeDiffOneFiftyHi}{9.9}
\newcommand{\CbOneTwentiethTypeDiffOneFiftyLo}{8.8}
\newcommand{\CbOneTwentiethTypePoFifty}{62}
\newcommand{\CbOneTwentiethTypePoHundred}{67}
\newcommand{\CbOneTwentiethTypePoOneFifty}{71}
\newcommand{\CbOneTwentiethUnitAtlasFifty}{72}
\newcommand{\CbOneTwentiethUnitAtlasHundred}{78}
\newcommand{\CbOneTwentiethUnitAtlasOneFifty}{80}
\newcommand{\CbOneTwentiethUnitDiffFifty}{10.2}
\newcommand{\CbOneTwentiethUnitDiffFiftyHi}{10.7}
\newcommand{\CbOneTwentiethUnitDiffFiftyLo}{9.7}
\newcommand{\CbOneTwentiethUnitDiffHundred}{10.5}
\newcommand{\CbOneTwentiethUnitDiffHundredHi}{11.0}
\newcommand{\CbOneTwentiethUnitDiffHundredLo}{10.0}
\newcommand{\CbOneTwentiethUnitDiffOneFifty}{9.1}
\newcommand{\CbOneTwentiethUnitDiffOneFiftyHi}{9.6}
\newcommand{\CbOneTwentiethUnitDiffOneFiftyLo}{8.6}
\newcommand{\CbOneTwentiethUnitPoFifty}{62}
\newcommand{\CbOneTwentiethUnitPoHundred}{67}
\newcommand{\CbOneTwentiethUnitPoOneFifty}{71}
\newcommand{\CbOverallFdrAll}{58}
\newcommand{\CbOverallFdrUnit}{57}
\newcommand{\CbOverallPrimary}{57}
\newcommand{\CbOverallThousandth}{58}
\newcommand{\CbOverallTwentieth}{55}
\newcommand{\CbPResolution}{$5.0\times10^{-5}$}
\newcommand{\CbPlanMaxP}{$\leq5\times10^{-5}$}
\newcommand{\CbPlanOneFiftyP}{$\leq5\times10^{-5}$}
\newcommand{\CbPlanOneFiftySimHi}{14.1}
\newcommand{\CbPlanOneFiftySimLo}{11.7}
\newcommand{\CbPlanOneHundredP}{$\leq5\times10^{-5}$}
\newcommand{\CbPlanOneHundredSimHi}{12.2}
\newcommand{\CbPlanOneHundredSimLo}{10.1}
\newcommand{\CbPlanOneOneFiftyP}{$\leq5\times10^{-5}$}
\newcommand{\CbPlanOneOneFiftySimHi}{10.2}
\newcommand{\CbPlanOneOneFiftySimLo}{8.0}
\newcommand{\CbPlanThreeHundredP}{$\leq5\times10^{-5}$}
\newcommand{\CbPlanThreeHundredSimHi}{16.2}
\newcommand{\CbPlanThreeHundredSimLo}{11.9}
\newcommand{\CbPlanThreeOneFiftyP}{$\leq5\times10^{-5}$}
\newcommand{\CbPlanThreeOneFiftySimHi}{15.8}
\newcommand{\CbPlanThreeOneFiftySimLo}{12.4}
\newcommand{\CbPlanTwoFiftyP}{$\leq5\times10^{-5}$}
\newcommand{\CbPlanTwoFiftySimHi}{21.8}
\newcommand{\CbPlanTwoFiftySimLo}{18.5}
\newcommand{\CbPlanTwoHundredP}{$\leq5\times10^{-5}$}
\newcommand{\CbPlanTwoHundredSimHi}{20.9}
\newcommand{\CbPlanTwoHundredSimLo}{17.4}
\newcommand{\CbPlanTwoOneFiftyP}{$\leq5\times10^{-5}$}
\newcommand{\CbPlanTwoOneFiftySimHi}{18.3}
\newcommand{\CbPlanTwoOneFiftySimLo}{15.1}
\newcommand{\CbPosFdrAll}{58}
\newcommand{\CbPosFdrUnit}{57}
\newcommand{\CbPosPrimary}{57}
\newcommand{\CbPosThousandth}{58}
\newcommand{\CbPosTwentieth}{55}
\newcommand{\CbRandomTwoFdrAll}{1.1}
\newcommand{\CbRandomTwoFdrUnit}{1.1}
\newcommand{\CbRandomTwoPrimary}{1.6}
\newcommand{\CbRandomTwoThousandth}{1.1}
\newcommand{\CbRandomTwoTwentieth}{2.5}
\newcommand{\CbShareCDEightFdrAll}{52}
\newcommand{\CbShareCDEightFdrUnit}{57}
\newcommand{\CbShareCDEightPrimary}{47}
\newcommand{\CbShareCDEightThousandth}{53}
\newcommand{\CbShareCDEightTwentieth}{41}
\newcommand{\CbTwoFdrAllPooledAtlasFifty}{31}
\newcommand{\CbTwoFdrAllPooledAtlasHundred}{42}
\newcommand{\CbTwoFdrAllPooledAtlasOneFifty}{47}
\newcommand{\CbTwoFdrAllPooledDiffFifty}{17.7}
\newcommand{\CbTwoFdrAllPooledDiffFiftyHi}{18.9}
\newcommand{\CbTwoFdrAllPooledDiffFiftyLo}{16.4}
\newcommand{\CbTwoFdrAllPooledDiffHundred}{17.2}
\newcommand{\CbTwoFdrAllPooledDiffHundredHi}{18.5}
\newcommand{\CbTwoFdrAllPooledDiffHundredLo}{15.9}
\newcommand{\CbTwoFdrAllPooledDiffOneFifty}{14.8}
\newcommand{\CbTwoFdrAllPooledDiffOneFiftyHi}{16.0}
\newcommand{\CbTwoFdrAllPooledDiffOneFiftyLo}{13.5}
\newcommand{\CbTwoFdrAllPooledPoFifty}{14}
\newcommand{\CbTwoFdrAllPooledPoHundred}{25}
\newcommand{\CbTwoFdrAllPooledPoOneFifty}{32}
\newcommand{\CbTwoFdrAllSampleAtlasFifty}{31}
\newcommand{\CbTwoFdrAllSampleAtlasHundred}{42}
\newcommand{\CbTwoFdrAllSampleAtlasOneFifty}{47}
\newcommand{\CbTwoFdrAllSampleDiffFifty}{17.7}
\newcommand{\CbTwoFdrAllSampleDiffFiftyHi}{18.9}
\newcommand{\CbTwoFdrAllSampleDiffFiftyLo}{16.4}
\newcommand{\CbTwoFdrAllSampleDiffHundred}{17.2}
\newcommand{\CbTwoFdrAllSampleDiffHundredHi}{18.5}
\newcommand{\CbTwoFdrAllSampleDiffHundredLo}{15.9}
\newcommand{\CbTwoFdrAllSampleDiffOneFifty}{14.8}
\newcommand{\CbTwoFdrAllSampleDiffOneFiftyHi}{16.0}
\newcommand{\CbTwoFdrAllSampleDiffOneFiftyLo}{13.5}
\newcommand{\CbTwoFdrAllSamplePoFifty}{14}
\newcommand{\CbTwoFdrAllSamplePoHundred}{25}
\newcommand{\CbTwoFdrAllSamplePoOneFifty}{32}
\newcommand{\CbTwoFdrAllTypeAtlasFifty}{31}
\newcommand{\CbTwoFdrAllTypeAtlasHundred}{42}
\newcommand{\CbTwoFdrAllTypeAtlasOneFifty}{47}
\newcommand{\CbTwoFdrAllTypeDiffFifty}{17.7}
\newcommand{\CbTwoFdrAllTypeDiffFiftyHi}{18.9}
\newcommand{\CbTwoFdrAllTypeDiffFiftyLo}{16.4}
\newcommand{\CbTwoFdrAllTypeDiffHundred}{17.2}
\newcommand{\CbTwoFdrAllTypeDiffHundredHi}{18.5}
\newcommand{\CbTwoFdrAllTypeDiffHundredLo}{15.9}
\newcommand{\CbTwoFdrAllTypeDiffOneFifty}{14.8}
\newcommand{\CbTwoFdrAllTypeDiffOneFiftyHi}{16.0}
\newcommand{\CbTwoFdrAllTypeDiffOneFiftyLo}{13.5}
\newcommand{\CbTwoFdrAllTypePoFifty}{14}
\newcommand{\CbTwoFdrAllTypePoHundred}{25}
\newcommand{\CbTwoFdrAllTypePoOneFifty}{32}
\newcommand{\CbTwoFdrAllUnitAtlasFifty}{31}
\newcommand{\CbTwoFdrAllUnitAtlasHundred}{42}
\newcommand{\CbTwoFdrAllUnitAtlasOneFifty}{47}
\newcommand{\CbTwoFdrAllUnitDiffFifty}{17.7}
\newcommand{\CbTwoFdrAllUnitDiffFiftyHi}{18.9}
\newcommand{\CbTwoFdrAllUnitDiffFiftyLo}{16.4}
\newcommand{\CbTwoFdrAllUnitDiffHundred}{17.2}
\newcommand{\CbTwoFdrAllUnitDiffHundredHi}{18.5}
\newcommand{\CbTwoFdrAllUnitDiffHundredLo}{15.9}
\newcommand{\CbTwoFdrAllUnitDiffOneFifty}{14.8}
\newcommand{\CbTwoFdrAllUnitDiffOneFiftyHi}{16.0}
\newcommand{\CbTwoFdrAllUnitDiffOneFiftyLo}{13.5}
\newcommand{\CbTwoFdrAllUnitPoFifty}{14}
\newcommand{\CbTwoFdrAllUnitPoHundred}{25}
\newcommand{\CbTwoFdrAllUnitPoOneFifty}{32}
\newcommand{\CbTwoFdrUnitPooledAtlasFifty}{31}
\newcommand{\CbTwoFdrUnitPooledAtlasHundred}{41}
\newcommand{\CbTwoFdrUnitPooledAtlasOneFifty}{46}
\newcommand{\CbTwoFdrUnitPooledDiffFifty}{17.2}
\newcommand{\CbTwoFdrUnitPooledDiffFiftyHi}{18.6}
\newcommand{\CbTwoFdrUnitPooledDiffFiftyLo}{15.9}
\newcommand{\CbTwoFdrUnitPooledDiffHundred}{16.5}
\newcommand{\CbTwoFdrUnitPooledDiffHundredHi}{17.8}
\newcommand{\CbTwoFdrUnitPooledDiffHundredLo}{15.2}
\newcommand{\CbTwoFdrUnitPooledDiffOneFifty}{14.2}
\newcommand{\CbTwoFdrUnitPooledDiffOneFiftyHi}{15.5}
\newcommand{\CbTwoFdrUnitPooledDiffOneFiftyLo}{12.8}
\newcommand{\CbTwoFdrUnitPooledPoFifty}{14}
\newcommand{\CbTwoFdrUnitPooledPoHundred}{25}
\newcommand{\CbTwoFdrUnitPooledPoOneFifty}{32}
\newcommand{\CbTwoFdrUnitSampleAtlasFifty}{31}
\newcommand{\CbTwoFdrUnitSampleAtlasHundred}{41}
\newcommand{\CbTwoFdrUnitSampleAtlasOneFifty}{46}
\newcommand{\CbTwoFdrUnitSampleDiffFifty}{17.2}
\newcommand{\CbTwoFdrUnitSampleDiffFiftyHi}{18.6}
\newcommand{\CbTwoFdrUnitSampleDiffFiftyLo}{15.9}
\newcommand{\CbTwoFdrUnitSampleDiffHundred}{16.5}
\newcommand{\CbTwoFdrUnitSampleDiffHundredHi}{17.8}
\newcommand{\CbTwoFdrUnitSampleDiffHundredLo}{15.2}
\newcommand{\CbTwoFdrUnitSampleDiffOneFifty}{14.2}
\newcommand{\CbTwoFdrUnitSampleDiffOneFiftyHi}{15.5}
\newcommand{\CbTwoFdrUnitSampleDiffOneFiftyLo}{12.8}
\newcommand{\CbTwoFdrUnitSamplePoFifty}{14}
\newcommand{\CbTwoFdrUnitSamplePoHundred}{25}
\newcommand{\CbTwoFdrUnitSamplePoOneFifty}{32}
\newcommand{\CbTwoFdrUnitTypeAtlasFifty}{31}
\newcommand{\CbTwoFdrUnitTypeAtlasHundred}{41}
\newcommand{\CbTwoFdrUnitTypeAtlasOneFifty}{46}
\newcommand{\CbTwoFdrUnitTypeDiffFifty}{17.2}
\newcommand{\CbTwoFdrUnitTypeDiffFiftyHi}{18.6}
\newcommand{\CbTwoFdrUnitTypeDiffFiftyLo}{15.9}
\newcommand{\CbTwoFdrUnitTypeDiffHundred}{16.5}
\newcommand{\CbTwoFdrUnitTypeDiffHundredHi}{17.8}
\newcommand{\CbTwoFdrUnitTypeDiffHundredLo}{15.2}
\newcommand{\CbTwoFdrUnitTypeDiffOneFifty}{14.2}
\newcommand{\CbTwoFdrUnitTypeDiffOneFiftyHi}{15.5}
\newcommand{\CbTwoFdrUnitTypeDiffOneFiftyLo}{12.8}
\newcommand{\CbTwoFdrUnitTypePoFifty}{14}
\newcommand{\CbTwoFdrUnitTypePoHundred}{25}
\newcommand{\CbTwoFdrUnitTypePoOneFifty}{32}
\newcommand{\CbTwoFdrUnitUnitAtlasFifty}{31}
\newcommand{\CbTwoFdrUnitUnitAtlasHundred}{41}
\newcommand{\CbTwoFdrUnitUnitAtlasOneFifty}{46}
\newcommand{\CbTwoFdrUnitUnitDiffFifty}{17.2}
\newcommand{\CbTwoFdrUnitUnitDiffFiftyHi}{18.6}
\newcommand{\CbTwoFdrUnitUnitDiffFiftyLo}{15.9}
\newcommand{\CbTwoFdrUnitUnitDiffHundred}{16.5}
\newcommand{\CbTwoFdrUnitUnitDiffHundredHi}{17.8}
\newcommand{\CbTwoFdrUnitUnitDiffHundredLo}{15.2}
\newcommand{\CbTwoFdrUnitUnitDiffOneFifty}{14.2}
\newcommand{\CbTwoFdrUnitUnitDiffOneFiftyHi}{15.5}
\newcommand{\CbTwoFdrUnitUnitDiffOneFiftyLo}{12.8}
\newcommand{\CbTwoFdrUnitUnitPoFifty}{14}
\newcommand{\CbTwoFdrUnitUnitPoHundred}{25}
\newcommand{\CbTwoFdrUnitUnitPoOneFifty}{32}
\newcommand{\CbTwoPrimaryPooledAtlasFifty}{35}
\newcommand{\CbTwoPrimaryPooledAtlasHundred}{47}
\newcommand{\CbTwoPrimaryPooledAtlasOneFifty}{52}
\newcommand{\CbTwoPrimaryPooledDiffFifty}{20.1}
\newcommand{\CbTwoPrimaryPooledDiffFiftyHi}{21.3}
\newcommand{\CbTwoPrimaryPooledDiffFiftyLo}{18.9}
\newcommand{\CbTwoPrimaryPooledDiffHundred}{19.3}
\newcommand{\CbTwoPrimaryPooledDiffHundredHi}{20.5}
\newcommand{\CbTwoPrimaryPooledDiffHundredLo}{18.0}
\newcommand{\CbTwoPrimaryPooledDiffOneFifty}{16.8}
\newcommand{\CbTwoPrimaryPooledDiffOneFiftyHi}{18.0}
\newcommand{\CbTwoPrimaryPooledDiffOneFiftyLo}{15.6}
\newcommand{\CbTwoPrimaryPooledPoFifty}{15}
\newcommand{\CbTwoPrimaryPooledPoHundred}{28}
\newcommand{\CbTwoPrimaryPooledPoOneFifty}{35}
\newcommand{\CbTwoPrimarySampleAtlasFifty}{35}
\newcommand{\CbTwoPrimarySampleAtlasHundred}{47}
\newcommand{\CbTwoPrimarySampleAtlasOneFifty}{52}
\newcommand{\CbTwoPrimarySampleDiffFifty}{20.1}
\newcommand{\CbTwoPrimarySampleDiffFiftyHi}{21.3}
\newcommand{\CbTwoPrimarySampleDiffFiftyLo}{18.9}
\newcommand{\CbTwoPrimarySampleDiffHundred}{19.3}
\newcommand{\CbTwoPrimarySampleDiffHundredHi}{20.5}
\newcommand{\CbTwoPrimarySampleDiffHundredLo}{18.0}
\newcommand{\CbTwoPrimarySampleDiffOneFifty}{16.8}
\newcommand{\CbTwoPrimarySampleDiffOneFiftyHi}{18.0}
\newcommand{\CbTwoPrimarySampleDiffOneFiftyLo}{15.6}
\newcommand{\CbTwoPrimarySamplePoFifty}{15}
\newcommand{\CbTwoPrimarySamplePoHundred}{28}
\newcommand{\CbTwoPrimarySamplePoOneFifty}{35}
\newcommand{\CbTwoPrimaryTypeAtlasFifty}{35}
\newcommand{\CbTwoPrimaryTypeAtlasHundred}{47}
\newcommand{\CbTwoPrimaryTypeAtlasOneFifty}{52}
\newcommand{\CbTwoPrimaryTypeDiffFifty}{20.1}
\newcommand{\CbTwoPrimaryTypeDiffFiftyHi}{21.3}
\newcommand{\CbTwoPrimaryTypeDiffFiftyLo}{18.9}
\newcommand{\CbTwoPrimaryTypeDiffHundred}{19.3}
\newcommand{\CbTwoPrimaryTypeDiffHundredHi}{20.5}
\newcommand{\CbTwoPrimaryTypeDiffHundredLo}{18.0}
\newcommand{\CbTwoPrimaryTypeDiffOneFifty}{16.8}
\newcommand{\CbTwoPrimaryTypeDiffOneFiftyHi}{18.0}
\newcommand{\CbTwoPrimaryTypeDiffOneFiftyLo}{15.6}
\newcommand{\CbTwoPrimaryTypePoFifty}{15}
\newcommand{\CbTwoPrimaryTypePoHundred}{28}
\newcommand{\CbTwoPrimaryTypePoOneFifty}{35}
\newcommand{\CbTwoPrimaryUnitAtlasFifty}{35}
\newcommand{\CbTwoPrimaryUnitAtlasHundred}{47}
\newcommand{\CbTwoPrimaryUnitAtlasOneFifty}{52}
\newcommand{\CbTwoPrimaryUnitDiffFifty}{20.1}
\newcommand{\CbTwoPrimaryUnitDiffFiftyHi}{21.3}
\newcommand{\CbTwoPrimaryUnitDiffFiftyLo}{18.9}
\newcommand{\CbTwoPrimaryUnitDiffHundred}{19.3}
\newcommand{\CbTwoPrimaryUnitDiffHundredHi}{20.5}
\newcommand{\CbTwoPrimaryUnitDiffHundredLo}{18.0}
\newcommand{\CbTwoPrimaryUnitDiffOneFifty}{16.8}
\newcommand{\CbTwoPrimaryUnitDiffOneFiftyHi}{18.0}
\newcommand{\CbTwoPrimaryUnitDiffOneFiftyLo}{15.6}
\newcommand{\CbTwoPrimaryUnitPoFifty}{15}
\newcommand{\CbTwoPrimaryUnitPoHundred}{28}
\newcommand{\CbTwoPrimaryUnitPoOneFifty}{35}
\newcommand{\CbTwoThousandthPooledAtlasFifty}{31}
\newcommand{\CbTwoThousandthPooledAtlasHundred}{42}
\newcommand{\CbTwoThousandthPooledAtlasOneFifty}{47}
\newcommand{\CbTwoThousandthPooledDiffFifty}{17.6}
\newcommand{\CbTwoThousandthPooledDiffFiftyHi}{18.9}
\newcommand{\CbTwoThousandthPooledDiffFiftyLo}{16.4}
\newcommand{\CbTwoThousandthPooledDiffHundred}{17.2}
\newcommand{\CbTwoThousandthPooledDiffHundredHi}{18.4}
\newcommand{\CbTwoThousandthPooledDiffHundredLo}{15.8}
\newcommand{\CbTwoThousandthPooledDiffOneFifty}{14.8}
\newcommand{\CbTwoThousandthPooledDiffOneFiftyHi}{16.0}
\newcommand{\CbTwoThousandthPooledDiffOneFiftyLo}{13.4}
\newcommand{\CbTwoThousandthPooledPoFifty}{13}
\newcommand{\CbTwoThousandthPooledPoHundred}{25}
\newcommand{\CbTwoThousandthPooledPoOneFifty}{32}
\newcommand{\CbTwoThousandthSampleAtlasFifty}{31}
\newcommand{\CbTwoThousandthSampleAtlasHundred}{42}
\newcommand{\CbTwoThousandthSampleAtlasOneFifty}{47}
\newcommand{\CbTwoThousandthSampleDiffFifty}{17.6}
\newcommand{\CbTwoThousandthSampleDiffFiftyHi}{18.9}
\newcommand{\CbTwoThousandthSampleDiffFiftyLo}{16.4}
\newcommand{\CbTwoThousandthSampleDiffHundred}{17.2}
\newcommand{\CbTwoThousandthSampleDiffHundredHi}{18.4}
\newcommand{\CbTwoThousandthSampleDiffHundredLo}{15.8}
\newcommand{\CbTwoThousandthSampleDiffOneFifty}{14.8}
\newcommand{\CbTwoThousandthSampleDiffOneFiftyHi}{16.0}
\newcommand{\CbTwoThousandthSampleDiffOneFiftyLo}{13.4}
\newcommand{\CbTwoThousandthSamplePoFifty}{13}
\newcommand{\CbTwoThousandthSamplePoHundred}{25}
\newcommand{\CbTwoThousandthSamplePoOneFifty}{32}
\newcommand{\CbTwoThousandthTypeAtlasFifty}{31}
\newcommand{\CbTwoThousandthTypeAtlasHundred}{42}
\newcommand{\CbTwoThousandthTypeAtlasOneFifty}{47}
\newcommand{\CbTwoThousandthTypeDiffFifty}{17.6}
\newcommand{\CbTwoThousandthTypeDiffFiftyHi}{18.9}
\newcommand{\CbTwoThousandthTypeDiffFiftyLo}{16.4}
\newcommand{\CbTwoThousandthTypeDiffHundred}{17.2}
\newcommand{\CbTwoThousandthTypeDiffHundredHi}{18.4}
\newcommand{\CbTwoThousandthTypeDiffHundredLo}{15.8}
\newcommand{\CbTwoThousandthTypeDiffOneFifty}{14.8}
\newcommand{\CbTwoThousandthTypeDiffOneFiftyHi}{16.0}
\newcommand{\CbTwoThousandthTypeDiffOneFiftyLo}{13.4}
\newcommand{\CbTwoThousandthTypePoFifty}{13}
\newcommand{\CbTwoThousandthTypePoHundred}{25}
\newcommand{\CbTwoThousandthTypePoOneFifty}{32}
\newcommand{\CbTwoThousandthUnitAtlasFifty}{31}
\newcommand{\CbTwoThousandthUnitAtlasHundred}{42}
\newcommand{\CbTwoThousandthUnitAtlasOneFifty}{47}
\newcommand{\CbTwoThousandthUnitDiffFifty}{17.6}
\newcommand{\CbTwoThousandthUnitDiffFiftyHi}{18.9}
\newcommand{\CbTwoThousandthUnitDiffFiftyLo}{16.4}
\newcommand{\CbTwoThousandthUnitDiffHundred}{17.2}
\newcommand{\CbTwoThousandthUnitDiffHundredHi}{18.4}
\newcommand{\CbTwoThousandthUnitDiffHundredLo}{15.8}
\newcommand{\CbTwoThousandthUnitDiffOneFifty}{14.8}
\newcommand{\CbTwoThousandthUnitDiffOneFiftyHi}{16.0}
\newcommand{\CbTwoThousandthUnitDiffOneFiftyLo}{13.4}
\newcommand{\CbTwoThousandthUnitPoFifty}{13}
\newcommand{\CbTwoThousandthUnitPoHundred}{25}
\newcommand{\CbTwoThousandthUnitPoOneFifty}{32}
\newcommand{\CbTwoTwentiethPooledAtlasFifty}{40}
\newcommand{\CbTwoTwentiethPooledAtlasHundred}{52}
\newcommand{\CbTwoTwentiethPooledAtlasOneFifty}{57}
\newcommand{\CbTwoTwentiethPooledDiffFifty}{22.2}
\newcommand{\CbTwoTwentiethPooledDiffFiftyHi}{23.3}
\newcommand{\CbTwoTwentiethPooledDiffFiftyLo}{21.1}
\newcommand{\CbTwoTwentiethPooledDiffHundred}{21.0}
\newcommand{\CbTwoTwentiethPooledDiffHundredHi}{22.0}
\newcommand{\CbTwoTwentiethPooledDiffHundredLo}{19.9}
\newcommand{\CbTwoTwentiethPooledDiffOneFifty}{18.1}
\newcommand{\CbTwoTwentiethPooledDiffOneFiftyHi}{19.1}
\newcommand{\CbTwoTwentiethPooledDiffOneFiftyLo}{17.1}
\newcommand{\CbTwoTwentiethPooledPoFifty}{18}
\newcommand{\CbTwoTwentiethPooledPoHundred}{31}
\newcommand{\CbTwoTwentiethPooledPoOneFifty}{39}
\newcommand{\CbTwoTwentiethSampleAtlasFifty}{40}
\newcommand{\CbTwoTwentiethSampleAtlasHundred}{52}
\newcommand{\CbTwoTwentiethSampleAtlasOneFifty}{57}
\newcommand{\CbTwoTwentiethSampleDiffFifty}{22.2}
\newcommand{\CbTwoTwentiethSampleDiffFiftyHi}{23.3}
\newcommand{\CbTwoTwentiethSampleDiffFiftyLo}{21.1}
\newcommand{\CbTwoTwentiethSampleDiffHundred}{21.0}
\newcommand{\CbTwoTwentiethSampleDiffHundredHi}{22.0}
\newcommand{\CbTwoTwentiethSampleDiffHundredLo}{19.9}
\newcommand{\CbTwoTwentiethSampleDiffOneFifty}{18.1}
\newcommand{\CbTwoTwentiethSampleDiffOneFiftyHi}{19.1}
\newcommand{\CbTwoTwentiethSampleDiffOneFiftyLo}{17.1}
\newcommand{\CbTwoTwentiethSamplePoFifty}{18}
\newcommand{\CbTwoTwentiethSamplePoHundred}{31}
\newcommand{\CbTwoTwentiethSamplePoOneFifty}{39}
\newcommand{\CbTwoTwentiethTypeAtlasFifty}{40}
\newcommand{\CbTwoTwentiethTypeAtlasHundred}{52}
\newcommand{\CbTwoTwentiethTypeAtlasOneFifty}{57}
\newcommand{\CbTwoTwentiethTypeDiffFifty}{22.2}
\newcommand{\CbTwoTwentiethTypeDiffFiftyHi}{23.3}
\newcommand{\CbTwoTwentiethTypeDiffFiftyLo}{21.1}
\newcommand{\CbTwoTwentiethTypeDiffHundred}{21.0}
\newcommand{\CbTwoTwentiethTypeDiffHundredHi}{22.0}
\newcommand{\CbTwoTwentiethTypeDiffHundredLo}{19.9}
\newcommand{\CbTwoTwentiethTypeDiffOneFifty}{18.1}
\newcommand{\CbTwoTwentiethTypeDiffOneFiftyHi}{19.1}
\newcommand{\CbTwoTwentiethTypeDiffOneFiftyLo}{17.1}
\newcommand{\CbTwoTwentiethTypePoFifty}{18}
\newcommand{\CbTwoTwentiethTypePoHundred}{31}
\newcommand{\CbTwoTwentiethTypePoOneFifty}{39}
\newcommand{\CbTwoTwentiethUnitAtlasFifty}{40}
\newcommand{\CbTwoTwentiethUnitAtlasHundred}{52}
\newcommand{\CbTwoTwentiethUnitAtlasOneFifty}{57}
\newcommand{\CbTwoTwentiethUnitDiffFifty}{22.2}
\newcommand{\CbTwoTwentiethUnitDiffFiftyHi}{23.3}
\newcommand{\CbTwoTwentiethUnitDiffFiftyLo}{21.1}
\newcommand{\CbTwoTwentiethUnitDiffHundred}{21.0}
\newcommand{\CbTwoTwentiethUnitDiffHundredHi}{22.0}
\newcommand{\CbTwoTwentiethUnitDiffHundredLo}{19.9}
\newcommand{\CbTwoTwentiethUnitDiffOneFifty}{18.1}
\newcommand{\CbTwoTwentiethUnitDiffOneFiftyHi}{19.1}
\newcommand{\CbTwoTwentiethUnitDiffOneFiftyLo}{17.1}
\newcommand{\CbTwoTwentiethUnitPoFifty}{18}
\newcommand{\CbTwoTwentiethUnitPoHundred}{31}
\newcommand{\CbTwoTwentiethUnitPoOneFifty}{39}
\newcommand{\CbTypeOneBDiffFifty}{6.1}
\newcommand{\CbTypeOneBDiffFiftyHi}{7.3}
\newcommand{\CbTypeOneBDiffFiftyLo}{5.1}
\newcommand{\CbTypeOneBDiffHundred}{6.6}
\newcommand{\CbTypeOneBDiffHundredHi}{7.7}
\newcommand{\CbTypeOneBDiffHundredLo}{5.5}
\newcommand{\CbTypeOneBDiffOneFifty}{7.1}
\newcommand{\CbTypeOneBDiffOneFiftyHi}{8.3}
\newcommand{\CbTypeOneBDiffOneFiftyLo}{5.2}
\newcommand{\CbTypeOneCDEightDiffFifty}{15.2}
\newcommand{\CbTypeOneCDEightDiffFiftyHi}{16.1}
\newcommand{\CbTypeOneCDEightDiffFiftyLo}{14.1}
\newcommand{\CbTypeOneCDEightDiffHundred}{12.8}
\newcommand{\CbTypeOneCDEightDiffHundredHi}{13.9}
\newcommand{\CbTypeOneCDEightDiffHundredLo}{11.6}
\newcommand{\CbTypeOneCDEightDiffOneFifty}{10.4}
\newcommand{\CbTypeOneCDEightDiffOneFiftyHi}{11.4}
\newcommand{\CbTypeOneCDEightDiffOneFiftyLo}{9.4}
\newcommand{\CbTypeOneCDFourDiffFifty}{16.5}
\newcommand{\CbTypeOneCDFourDiffFiftyHi}{18.0}
\newcommand{\CbTypeOneCDFourDiffFiftyLo}{15.0}
\newcommand{\CbTypeOneCDFourDiffHundred}{15.2}
\newcommand{\CbTypeOneCDFourDiffHundredHi}{16.1}
\newcommand{\CbTypeOneCDFourDiffHundredLo}{13.5}
\newcommand{\CbTypeOneCDFourDiffOneFifty}{12.8}
\newcommand{\CbTypeOneCDFourDiffOneFiftyHi}{13.9}
\newcommand{\CbTypeOneCDFourDiffOneFiftyLo}{11.6}
\newcommand{\CbTypeOneMonoDiffFifty}{1.6}
\newcommand{\CbTypeOneMonoDiffFiftyHi}{2.6}
\newcommand{\CbTypeOneMonoDiffFiftyLo}{0.6}
\newcommand{\CbTypeOneMonoDiffHundred}{1.3}
\newcommand{\CbTypeOneMonoDiffHundredHi}{1.9}
\newcommand{\CbTypeOneMonoDiffHundredLo}{0.5}
\newcommand{\CbTypeOneMonoDiffOneFifty}{\ensuremath{-}0.4}
\newcommand{\CbTypeOneMonoDiffOneFiftyHi}{0.5}
\newcommand{\CbTypeOneMonoDiffOneFiftyLo}{\ensuremath{-}1.2}
\newcommand{\CbTypeOneNKDiffFifty}{7.3}
\newcommand{\CbTypeOneNKDiffFiftyHi}{8.6}
\newcommand{\CbTypeOneNKDiffFiftyLo}{6.2}
\newcommand{\CbTypeOneNKDiffHundred}{9.5}
\newcommand{\CbTypeOneNKDiffHundredHi}{10.8}
\newcommand{\CbTypeOneNKDiffHundredLo}{8.4}
\newcommand{\CbTypeOneNKDiffOneFifty}{9.5}
\newcommand{\CbTypeOneNKDiffOneFiftyHi}{10.8}
\newcommand{\CbTypeOneNKDiffOneFiftyLo}{7.9}
\newcommand{\CbTypeTwoBDiffFifty}{0.4}
\newcommand{\CbTypeTwoBDiffFiftyHi}{1.0}
\newcommand{\CbTypeTwoBDiffFiftyLo}{\ensuremath{-}0.2}
\newcommand{\CbTypeTwoBDiffHundred}{1.1}
\newcommand{\CbTypeTwoBDiffHundredHi}{2.1}
\newcommand{\CbTypeTwoBDiffHundredLo}{0.2}
\newcommand{\CbTypeTwoBDiffOneFifty}{0.7}
\newcommand{\CbTypeTwoBDiffOneFiftyHi}{1.8}
\newcommand{\CbTypeTwoBDiffOneFiftyLo}{\ensuremath{-}0.4}
\newcommand{\CbTypeTwoCDEightDiffFifty}{24.0}
\newcommand{\CbTypeTwoCDEightDiffFiftyHi}{27.3}
\newcommand{\CbTypeTwoCDEightDiffFiftyLo}{20.8}
\newcommand{\CbTypeTwoCDEightDiffHundred}{10.7}
\newcommand{\CbTypeTwoCDEightDiffHundredHi}{13.4}
\newcommand{\CbTypeTwoCDEightDiffHundredLo}{8.2}
\newcommand{\CbTypeTwoCDEightDiffOneFifty}{5.4}
\newcommand{\CbTypeTwoCDEightDiffOneFiftyHi}{7.0}
\newcommand{\CbTypeTwoCDEightDiffOneFiftyLo}{4.0}
\newcommand{\CbTypeTwoCDFourDiffFifty}{43.8}
\newcommand{\CbTypeTwoCDFourDiffFiftyHi}{47.1}
\newcommand{\CbTypeTwoCDFourDiffFiftyLo}{40.0}
\newcommand{\CbTypeTwoCDFourDiffHundred}{42.3}
\newcommand{\CbTypeTwoCDFourDiffHundredHi}{45.3}
\newcommand{\CbTypeTwoCDFourDiffHundredLo}{39.0}
\newcommand{\CbTypeTwoCDFourDiffOneFifty}{33.5}
\newcommand{\CbTypeTwoCDFourDiffOneFiftyHi}{36.2}
\newcommand{\CbTypeTwoCDFourDiffOneFiftyLo}{30.7}
\newcommand{\CbTypeTwoMonoDiffFifty}{16.2}
\newcommand{\CbTypeTwoMonoDiffFiftyHi}{18.8}
\newcommand{\CbTypeTwoMonoDiffFiftyLo}{13.3}
\newcommand{\CbTypeTwoMonoDiffHundred}{23.8}
\newcommand{\CbTypeTwoMonoDiffHundredHi}{26.0}
\newcommand{\CbTypeTwoMonoDiffHundredLo}{21.7}
\newcommand{\CbTypeTwoMonoDiffOneFifty}{25.7}
\newcommand{\CbTypeTwoMonoDiffOneFiftyHi}{27.7}
\newcommand{\CbTypeTwoMonoDiffOneFiftyLo}{23.4}
\newcommand{\CbTypeTwoNKDiffFifty}{9.2}
\newcommand{\CbTypeTwoNKDiffFiftyHi}{10.3}
\newcommand{\CbTypeTwoNKDiffFiftyLo}{8.1}
\newcommand{\CbTypeTwoNKDiffHundred}{12.2}
\newcommand{\CbTypeTwoNKDiffHundredHi}{14.6}
\newcommand{\CbTypeTwoNKDiffHundredLo}{9.8}
\newcommand{\CbTypeTwoNKDiffOneFifty}{14.0}
\newcommand{\CbTypeTwoNKDiffOneFiftyHi}{16.2}
\newcommand{\CbTypeTwoNKDiffOneFiftyLo}{11.8}
\newcommand{\CtAssayMax}{7.3}
\newcommand{\CtCellsOther}{125}
\newcommand{\CtCellsSame}{185}
\newcommand{\CtFrozenOtherFirst}{\ensuremath{-}535}
\newcommand{\CtFrozenOtherHundred}{\ensuremath{-}3.9}
\newcommand{\CtOtherFifty}{7.4}
\newcommand{\CtOtherFirst}{0.6}
\newcommand{\CtOtherHundred}{17}
\newcommand{\CtOtherMax}{35}
\newcommand{\CtOverSame}{0.68}
\newcommand{\CtOverSameHi}{0.93}
\newcommand{\CtOverSameLo}{0.61}
\newcommand{\CtPairedFirst}{0.3}
\newcommand{\CtPairedMax}{5.9}
\newcommand{\CtSaveOther}{6.4}
\newcommand{\CtSaveOtherHi}{8.3}
\newcommand{\CtSaveOtherLo}{3.7}
\newcommand{\CtVsFrozen}{1.6}
\newcommand{\CtVsFrozenHi}{1.9}
\newcommand{\CtVsFrozenLo}{1.2}
\newcommand{\DpCalConMax}{1.00}
\newcommand{\DpCalConMin}{1.00}
\newcommand{\DpCalDfMax}{0.91}
\newcommand{\DpCalDfMin}{0.64}
\newcommand{\DpCalDraws}{20}
\newcommand{\DpCellsOther}{205}
\newcommand{\DpCellsSame}{235}
\newcommand{\DpCellsSamePool}{169}
\newcommand{\DpDOne}{met}
\newcommand{\DpDThree}{met}
\newcommand{\DpDTwo}{not met}
\newcommand{\DpFirstBudget}{25}
\newcommand{\DpFrozenAt}{19:35 EDT on 28 September 2026}
\newcommand{\DpLevel}{36}
\newcommand{\DpMaxBudget}{800}
\newcommand{\DpNFolds}{twelve}
\newcommand{\DpNSites}{four}
\newcommand{\DpNonInf}{0.87}
\newcommand{\DpNonInfHi}{0.96}
\newcommand{\DpNonInfLo}{0.81}
\newcommand{\DpNonInfQuarter}{0.93}
\newcommand{\DpNonInfQuarterHi}{0.99}
\newcommand{\DpNonInfQuarterLo}{0.89}
\newcommand{\DpNonInfThreeQ}{0.83}
\newcommand{\DpNonInfThreeQHi}{0.88}
\newcommand{\DpNonInfThreeQLo}{0.75}
\newcommand{\DpOverlapAssay}{0.53}
\newcommand{\DpOverlapAssayTwenty}{0.56}
\newcommand{\DpOverlapOther}{0.74}
\newcommand{\DpOverlapOtherTwenty}{0.71}
\newcommand{\DpOwnNegativeBelow}{200}
\newcommand{\DpOwnNegativeMax}{100}
\newcommand{\DpPoConFifty}{0.9}
\newcommand{\DpPoConHundred}{2.1}
\newcommand{\DpPoConTwentyFive}{0.3}
\newcommand{\DpPoDfFifty}{0.0}
\newcommand{\DpPoDfHundred}{0.9}
\newcommand{\DpPoDfTwentyFive}{\ensuremath{-}16}
\newcommand{\DpPoolOtherMax}{17,621}
\newcommand{\DpPoolOtherMin}{17,369}
\newcommand{\DpPoolSameMax}{4,514}
\newcommand{\DpPoolSameMin}{4,207}
\newcommand{\DpRAssay}{0.59}
\newcommand{\DpROther}{0.80}
\newcommand{\DpReservoir}{1,471}
\newcommand{\DpRfAssayMax}{7.3}
\newcommand{\DpRfMeansBest}{\ensuremath{-}443}
\newcommand{\DpRfOtherMax}{36}
\newcommand{\DpRfPairedMax}{5.9}
\newcommand{\DpRfSameMax}{34}
\newcommand{\DpRfSamePoolFirst}{1.9}
\newcommand{\DpSaveOther}{3.9}
\newcommand{\DpSaveOtherHi}{4.5}
\newcommand{\DpSaveOtherLo}{3.3}
\newcommand{\DpSaveOtherQuarter}{5.3}
\newcommand{\DpSaveOtherQuarterHi}{5.9}
\newcommand{\DpSaveOtherQuarterLo}{4.7}
\newcommand{\DpSaveOtherThreeQ}{2.3}
\newcommand{\DpSaveOtherThreeQHi}{2.6}
\newcommand{\DpSaveOtherThreeQLo}{2.1}
\newcommand{\DpSaveSame}{3.4}
\newcommand{\DpSaveSamePool}{4.7}
\newcommand{\DpSiteMax}{4.4}
\newcommand{\DpSiteMin}{2.0}
\newcommand{\DpTarget}{18}
\newcommand{\FxDiffPaired}{1.4}
\newcommand{\FxDiffProposed}{0.5}
\newcommand{\FxOwnHundred}{60}
\newcommand{\FxOwnMax}{74}
\newcommand{\FxPairedOwnHundred}{36}
\newcommand{\FxPairedPoolHundred}{\ensuremath{-}52}
\newcommand{\FxPoolHundred}{\ensuremath{-}78}
\newcommand{\FxSameHundred}{56}
\newcommand{\LwAssay}{3.0}
\newcommand{\LwAssayHi}{4.3}
\newcommand{\LwAssayLo}{2.0}
\newcommand{\LwOther}{4.3}
\newcommand{\LwOtherHi}{6.2}
\newcommand{\LwOtherLo}{3.3}
\newcommand{\LwPiAssay}{11}
\newcommand{\LwPiOther}{14}
\newcommand{\LwPiSame}{15}
\newcommand{\LwSame}{6.9}
\newcommand{\LwSameHi}{10.5}
\newcommand{\LwSameLo}{4.0}
\newcommand{\LwWAssay}{45}
\newcommand{\LwWOther}{68}
\newcommand{\LwWSame}{77}
\newcommand{\MpAssayMapped}{7.3}
\newcommand{\MpAssayOtherMap}{26}
\newcommand{\MpAssayUnmapped}{6.9}
\newcommand{\MpOtherMapped}{35}
\newcommand{\MpOtherUnmapped}{7.0}
\newcommand{\MpSameUnmapped}{6.0}
\newcommand{\PrNAbove}{10}
\newcommand{\PrNBelow}{9}
\newcommand{\PrNNeither}{6}
\newcommand{\PsSmallFifty}{29}
\newcommand{\PsSmallTwentyFive}{60}
\newcommand{\RbCellCountFewBMono}{47,683}
\newcommand{\RbCellCountFewT}{34,942}
\newcommand{\RbCellCountFull}{63,796}
\newcommand{\RbCellCountPatientsFifteen}{8,231}
\newcommand{\RbCellCountPatientsFiftyNine}{32,350}
\newcommand{\RbCellCountPatientsFour}{2,132}
\newcommand{\RbCellCountPatientsSeven}{3,823}
\newcommand{\RbCellCountPatientsThirty}{16,510}
\newcommand{\RbCellCountWithoutB}{55,371}
\newcommand{\RbCellCountWithoutCDEightT}{52,250}
\newcommand{\RbCellCountWithoutCDFourT}{42,339}
\newcommand{\RbCellCountWithoutMono}{53,494}
\newcommand{\RbCellCountWithoutNK}{51,730}
\newcommand{\RbCellsFewBMono}{$>$800}
\newcommand{\RbCellsFewBMonoHi}{800}
\newcommand{\RbCellsFewBMonoLo}{800}
\newcommand{\RbCellsFewT}{$>$800}
\newcommand{\RbCellsFewTHi}{800}
\newcommand{\RbCellsFewTLo}{800}
\newcommand{\RbCellsFull}{147}
\newcommand{\RbCellsFullHi}{182}
\newcommand{\RbCellsFullLo}{126}
\newcommand{\RbCellsPatientsFifteen}{185}
\newcommand{\RbCellsPatientsFifteenHi}{246}
\newcommand{\RbCellsPatientsFifteenLo}{155}
\newcommand{\RbCellsPatientsFiftyNine}{174}
\newcommand{\RbCellsPatientsFiftyNineHi}{217}
\newcommand{\RbCellsPatientsFiftyNineLo}{147}
\newcommand{\RbCellsPatientsFour}{$>$800}
\newcommand{\RbCellsPatientsFourHi}{800}
\newcommand{\RbCellsPatientsFourLo}{800}
\newcommand{\RbCellsPatientsSeven}{385}
\newcommand{\RbCellsPatientsSevenHi}{800}
\newcommand{\RbCellsPatientsSevenLo}{281}
\newcommand{\RbCellsPatientsThirty}{177}
\newcommand{\RbCellsPatientsThirtyHi}{219}
\newcommand{\RbCellsPatientsThirtyLo}{151}
\newcommand{\RbCellsWithoutB}{174}
\newcommand{\RbCellsWithoutBHi}{222}
\newcommand{\RbCellsWithoutBLo}{147}
\newcommand{\RbCellsWithoutCDEightT}{$>$800}
\newcommand{\RbCellsWithoutCDEightTHi}{800}
\newcommand{\RbCellsWithoutCDEightTLo}{800}
\newcommand{\RbCellsWithoutCDFourT}{159}
\newcommand{\RbCellsWithoutCDFourTHi}{201}
\newcommand{\RbCellsWithoutCDFourTLo}{132}
\newcommand{\RbCellsWithoutMono}{209}
\newcommand{\RbCellsWithoutMonoHi}{316}
\newcommand{\RbCellsWithoutMonoLo}{162}
\newcommand{\RbCellsWithoutNK}{184}
\newcommand{\RbCellsWithoutNKHi}{241}
\newcommand{\RbCellsWithoutNKLo}{152}
\newcommand{\RbDonorAtlasMax}{741}
\newcommand{\RbDonorMax}{8.1}
\newcommand{\RbDonorMin}{1.08}
\newcommand{\RbDonorsCensored}{15}
\newcommand{\RbDonorsPoReached}{three}
\newcommand{\RbLopoMax}{5.6}
\newcommand{\RbLopoMin}{5.2}
\newcommand{\RbLopoOwnMax}{1.42}
\newcommand{\RbLopoOwnMin}{1.24}
\newcommand{\RbMappedFewBMono}{20}
\newcommand{\RbMappedFewT}{\ensuremath{-}2.7}
\newcommand{\RbMappedFull}{48}
\newcommand{\RbMappedMaxFewBMono}{20}
\newcommand{\RbMappedMaxFewT}{14}
\newcommand{\RbMappedMaxFull}{48}
\newcommand{\RbNDonors}{18}
\newcommand{\RbPoolMax}{8.4}
\newcommand{\RbPoolMin}{3.4}
\newcommand{\RbPopsFewBMono}{395}
\newcommand{\RbPopsFewT}{351}
\newcommand{\RbPopsFull}{566}
\newcommand{\RbPopsPatientsFifteen}{73}
\newcommand{\RbPopsPatientsFiftyNine}{281}
\newcommand{\RbPopsPatientsFour}{17}
\newcommand{\RbPopsPatientsSeven}{34}
\newcommand{\RbPopsPatientsThirty}{145}
\newcommand{\RbPopsWithoutB}{457}
\newcommand{\RbPopsWithoutCDEightT}{444}
\newcommand{\RbPopsWithoutCDFourT}{440}
\newcommand{\RbPopsWithoutMono}{480}
\newcommand{\RbPopsWithoutNK}{443}
\newcommand{\RbReproduces}{0}
\newcommand{\RbRfHundredFewBMono}{2.8}
\newcommand{\RbRfHundredFewT}{14}
\newcommand{\RbRfHundredFull}{23}
\newcommand{\RbRfHundredPatientsFifteen}{20}
\newcommand{\RbRfHundredPatientsFiftyNine}{20}
\newcommand{\RbRfHundredPatientsFour}{9.5}
\newcommand{\RbRfHundredPatientsSeven}{16}
\newcommand{\RbRfHundredPatientsThirty}{20}
\newcommand{\RbRfHundredWithoutB}{21}
\newcommand{\RbRfHundredWithoutCDEightT}{9.8}
\newcommand{\RbRfHundredWithoutCDFourT}{23}
\newcommand{\RbRfHundredWithoutMono}{20}
\newcommand{\RbRfHundredWithoutNK}{21}
\newcommand{\RbRfMaxFewBMono}{20}
\newcommand{\RbRfMaxFewT}{14}
\newcommand{\RbRfMaxFull}{48}
\newcommand{\RbRfMaxPatientsFifteen}{41}
\newcommand{\RbRfMaxPatientsFiftyNine}{46}
\newcommand{\RbRfMaxPatientsFour}{15}
\newcommand{\RbRfMaxPatientsSeven}{31}
\newcommand{\RbRfMaxPatientsThirty}{44}
\newcommand{\RbRfMaxWithoutB}{43}
\newcommand{\RbRfMaxWithoutCDEightT}{25}
\newcommand{\RbRfMaxWithoutCDFourT}{43}
\newcommand{\RbRfMaxWithoutMono}{38}
\newcommand{\RbRfMaxWithoutNK}{41}
\newcommand{\RbSaveFewBMono}{none}
\newcommand{\RbSaveFewBMonoHi}{1.00}
\newcommand{\RbSaveFewBMonoLo}{1.00}
\newcommand{\RbSaveFewT}{none}
\newcommand{\RbSaveFewTHi}{1.00}
\newcommand{\RbSaveFewTLo}{1.00}
\newcommand{\RbSaveFull}{at least 5.4}
\newcommand{\RbSaveFullHi}{6.3}
\newcommand{\RbSaveFullLo}{4.4}
\newcommand{\RbSavePatientsFifteen}{at least 4.3}
\newcommand{\RbSavePatientsFifteenHi}{5.2}
\newcommand{\RbSavePatientsFifteenLo}{3.3}
\newcommand{\RbSavePatientsFiftyNine}{at least 4.6}
\newcommand{\RbSavePatientsFiftyNineHi}{5.4}
\newcommand{\RbSavePatientsFiftyNineLo}{3.7}
\newcommand{\RbSavePatientsFour}{none}
\newcommand{\RbSavePatientsFourHi}{1.00}
\newcommand{\RbSavePatientsFourLo}{1.00}
\newcommand{\RbSavePatientsSeven}{at least 2.1}
\newcommand{\RbSavePatientsSevenHi}{2.8}
\newcommand{\RbSavePatientsSevenLo}{1.00}
\newcommand{\RbSavePatientsThirty}{at least 4.5}
\newcommand{\RbSavePatientsThirtyHi}{5.3}
\newcommand{\RbSavePatientsThirtyLo}{3.6}
\newcommand{\RbSaveWithoutB}{at least 4.6}
\newcommand{\RbSaveWithoutBHi}{5.5}
\newcommand{\RbSaveWithoutBLo}{3.6}
\newcommand{\RbSaveWithoutCDEightT}{none}
\newcommand{\RbSaveWithoutCDEightTHi}{1.00}
\newcommand{\RbSaveWithoutCDEightTLo}{1.00}
\newcommand{\RbSaveWithoutCDFourT}{at least 5.0}
\newcommand{\RbSaveWithoutCDFourTHi}{6.1}
\newcommand{\RbSaveWithoutCDFourTLo}{4.0}
\newcommand{\RbSaveWithoutMono}{at least 3.8}
\newcommand{\RbSaveWithoutMonoHi}{4.9}
\newcommand{\RbSaveWithoutMonoLo}{2.5}
\newcommand{\RbSaveWithoutNK}{at least 4.4}
\newcommand{\RbSaveWithoutNKHi}{5.3}
\newcommand{\RbSaveWithoutNKLo}{3.3}
\newcommand{\RbUnmapFewBMono}{14}
\newcommand{\RbUnmapFewT}{13}
\newcommand{\RbUnmapFull}{14}
\newcommand{\RvAtlasBaseCells}{147}
\newcommand{\RvAtlasBaseKTwo}{1.13}
\newcommand{\RvAtlasCellsMax}{170}
\newcommand{\RvAtlasCellsMin}{141}
\newcommand{\RvAtlasKTwoAbove}{8}
\newcommand{\RvAtlasKTwoHi}{1.40}
\newcommand{\RvAtlasKTwoLo}{0.82}
\newcommand{\RvAtlasKTwoMax}{1.27}
\newcommand{\RvAtlasKTwoMed}{1.07}
\newcommand{\RvAtlasKTwoMin}{0.99}
\newcommand{\RvAtlasKTwoShare}{61}
\newcommand{\RvAtlasReps}{10}
\newcommand{\RvAtlasVOneHi}{6.3}
\newcommand{\RvAtlasVOneLo}{3.9}
\newcommand{\RvAtlasVOneMax}{5.7}
\newcommand{\RvAtlasVOneMin}{4.7}
\newcommand{\RvAxesCells}{more than 800}
\newcommand{\RvAxesCellsHi}{800}
\newcommand{\RvAxesCellsLo}{800}
\newcommand{\RvAxesEight}{14}
\newcommand{\RvAxesEightHi}{16}
\newcommand{\RvAxesEightLo}{12}
\newcommand{\RvAxesFifty}{4.0}
\newcommand{\RvAxesFiftyHi}{4.8}
\newcommand{\RvAxesFiftyLo}{3.2}
\newcommand{\RvAxesFourHundred}{13}
\newcommand{\RvAxesFourHundredHi}{14}
\newcommand{\RvAxesFourHundredLo}{11}
\newcommand{\RvAxesHundred}{7.8}
\newcommand{\RvAxesHundredHi}{8.9}
\newcommand{\RvAxesHundredLo}{6.6}
\newcommand{\RvAxesOneFifty}{9.3}
\newcommand{\RvAxesOneFiftyHi}{11}
\newcommand{\RvAxesOneFiftyLo}{7.6}
\newcommand{\RvAxesSave}{at least 1.0}
\newcommand{\RvAxesSaveHi}{1.0}
\newcommand{\RvAxesSaveLo}{1.0}
\newcommand{\RvAxesTwentyFive}{\ensuremath{-}0.7}
\newcommand{\RvAxesTwentyFiveHi}{0.5}
\newcommand{\RvAxesTwentyFiveLo}{\ensuremath{-}2.5}
\newcommand{\RvAxesTwoHundred}{10}
\newcommand{\RvAxesTwoHundredHi}{12}
\newcommand{\RvAxesTwoHundredLo}{8.8}
\newcommand{\RvBothCells}{145}
\newcommand{\RvBothCellsHi}{183}
\newcommand{\RvBothCellsLo}{124}
\newcommand{\RvBothEight}{48}
\newcommand{\RvBothEightHi}{52}
\newcommand{\RvBothEightLo}{44}
\newcommand{\RvBothFifty}{11}
\newcommand{\RvBothFiftyHi}{14}
\newcommand{\RvBothFiftyLo}{8.8}
\newcommand{\RvBothFourHundred}{42}
\newcommand{\RvBothFourHundredHi}{45}
\newcommand{\RvBothFourHundredLo}{38}
\newcommand{\RvBothHundred}{23}
\newcommand{\RvBothHundredHi}{25}
\newcommand{\RvBothHundredLo}{20}
\newcommand{\RvBothOneFifty}{29}
\newcommand{\RvBothOneFiftyHi}{32}
\newcommand{\RvBothOneFiftyLo}{25}
\newcommand{\RvBothSave}{at least 5.5}
\newcommand{\RvBothSaveHi}{6.5}
\newcommand{\RvBothSaveLo}{4.4}
\newcommand{\RvBothTwentyFive}{0.5}
\newcommand{\RvBothTwentyFiveHi}{2.3}
\newcommand{\RvBothTwentyFiveLo}{\ensuremath{-}2.0}
\newcommand{\RvBothTwoHundred}{33}
\newcommand{\RvBothTwoHundredHi}{36}
\newcommand{\RvBothTwoHundredLo}{30}
\newcommand{\RvCalBmBlockMax}{1.31}
\newcommand{\RvCalBmBlockMin}{1.00}
\newcommand{\RvCalBmJackGainMax}{6.0}
\newcommand{\RvCalBmJackGainMin}{0.9}
\newcommand{\RvCalBmJackMax}{1.54}
\newcommand{\RvCalBmJackMin}{1.05}
\newcommand{\RvCalBmMax}{1.029}
\newcommand{\RvCalBmMin}{1.005}
\newcommand{\RvCalBmOrderCvEight}{4.4}
\newcommand{\RvCalBmOrderCvHundred}{9.7}
\newcommand{\RvCalBmOrderCvMax}{12}
\newcommand{\RvCalBmOrderFactorEight}{0.015}
\newcommand{\RvCalBmOrderFactorHundred}{0.11}
\newcommand{\RvCalBmOrderFactorMax}{0.22}
\newcommand{\RvCalBmOrthoMax}{1.034}
\newcommand{\RvCalBmOrthoMin}{1.005}
\newcommand{\RvCalBmOwnMax}{1.01}
\newcommand{\RvCalBmOwnMin}{0.65}
\newcommand{\RvCalValBlockMax}{1.07}
\newcommand{\RvCalValBlockMin}{1.00}
\newcommand{\RvCalValJackGainMax}{2.9}
\newcommand{\RvCalValJackGainMin}{0.6}
\newcommand{\RvCalValJackMax}{1.48}
\newcommand{\RvCalValJackMin}{1.03}
\newcommand{\RvCalValMax}{1.016}
\newcommand{\RvCalValMin}{1.005}
\newcommand{\RvCalValOrderCvEight}{2.8}
\newcommand{\RvCalValOrderCvHundred}{8.5}
\newcommand{\RvCalValOrderCvMax}{12}
\newcommand{\RvCalValOrderFactorEight}{0.016}
\newcommand{\RvCalValOrderFactorHundred}{0.10}
\newcommand{\RvCalValOrderFactorMax}{0.22}
\newcommand{\RvCalValOrthoMax}{1.010}
\newcommand{\RvCalValOrthoMin}{1.003}
\newcommand{\RvCalValOwnMax}{1.00}
\newcommand{\RvCalValOwnMin}{0.67}
\newcommand{\RvDonorBetterPo}{18}
\newcommand{\RvDonorBetterReg}{10}
\newcommand{\RvDonorMinusPoMax}{16}
\newcommand{\RvDonorMinusPoMin}{6.8}
\newcommand{\RvDonorMinusRegMax}{1.8}
\newcommand{\RvDonorMinusRegMin}{\ensuremath{-}2.0}
\newcommand{\RvDonorN}{18}
\newcommand{\RvDrawCellsMax}{154}
\newcommand{\RvDrawCellsMin}{133}
\newcommand{\RvDrawFileIdentical}{identical}
\newcommand{\RvDrawInflMax}{537}
\newcommand{\RvDrawInflMedianMax}{29}
\newcommand{\RvDrawInflMedianMin}{2}
\newcommand{\RvDrawKTwoHi}{1.45}
\newcommand{\RvDrawKTwoLo}{0.85}
\newcommand{\RvDrawKTwoMax}{1.17}
\newcommand{\RvDrawKTwoMed}{1.04}
\newcommand{\RvDrawKTwoMin}{0.95}
\newcommand{\RvDrawKTwoShare}{63}
\newcommand{\RvDrawMaxDiff}{$4\times10^{-9}$}
\newcommand{\RvDrawTests}{23}
\newcommand{\RvDrawVOneHi}{6.7}
\newcommand{\RvDrawVOneLo}{4.2}
\newcommand{\RvDrawVOneMax}{6.0}
\newcommand{\RvDrawVOneMed}{5.5}
\newcommand{\RvDrawVOneMin}{5.2}
\newcommand{\RvDrawVOneShare}{100}
\newcommand{\RvFewBMonoAxesCells}{more than 800}
\newcommand{\RvFewBMonoAxesCellsHi}{800}
\newcommand{\RvFewBMonoAxesCellsLo}{800}
\newcommand{\RvFewBMonoAxesEight}{14}
\newcommand{\RvFewBMonoAxesEightHi}{16}
\newcommand{\RvFewBMonoAxesEightLo}{12}
\newcommand{\RvFewBMonoAxesFifty}{4.2}
\newcommand{\RvFewBMonoAxesFiftyHi}{5.1}
\newcommand{\RvFewBMonoAxesFiftyLo}{3.4}
\newcommand{\RvFewBMonoAxesFourHundred}{12}
\newcommand{\RvFewBMonoAxesFourHundredHi}{14}
\newcommand{\RvFewBMonoAxesFourHundredLo}{11}
\newcommand{\RvFewBMonoAxesHundred}{8.0}
\newcommand{\RvFewBMonoAxesHundredHi}{9.1}
\newcommand{\RvFewBMonoAxesHundredLo}{6.7}
\newcommand{\RvFewBMonoAxesOneFifty}{9.3}
\newcommand{\RvFewBMonoAxesOneFiftyHi}{11}
\newcommand{\RvFewBMonoAxesOneFiftyLo}{7.5}
\newcommand{\RvFewBMonoAxesSave}{at least 1.0}
\newcommand{\RvFewBMonoAxesSaveHi}{1.0}
\newcommand{\RvFewBMonoAxesSaveLo}{1.0}
\newcommand{\RvFewBMonoAxesTwentyFive}{\ensuremath{-}0.5}
\newcommand{\RvFewBMonoAxesTwentyFiveHi}{0.8}
\newcommand{\RvFewBMonoAxesTwentyFiveLo}{\ensuremath{-}2.3}
\newcommand{\RvFewBMonoAxesTwoHundred}{10}
\newcommand{\RvFewBMonoAxesTwoHundredHi}{12}
\newcommand{\RvFewBMonoAxesTwoHundredLo}{8.6}
\newcommand{\RvFewBMonoMappedCells}{more than 800}
\newcommand{\RvFewBMonoMappedCellsHi}{800}
\newcommand{\RvFewBMonoMappedCellsLo}{800}
\newcommand{\RvFewBMonoMappedEight}{20}
\newcommand{\RvFewBMonoMappedEightHi}{26}
\newcommand{\RvFewBMonoMappedEightLo}{13}
\newcommand{\RvFewBMonoMappedFifty}{\ensuremath{-}6.1}
\newcommand{\RvFewBMonoMappedFiftyHi}{\ensuremath{-}2.4}
\newcommand{\RvFewBMonoMappedFiftyLo}{\ensuremath{-}11}
\newcommand{\RvFewBMonoMappedFourHundred}{17}
\newcommand{\RvFewBMonoMappedFourHundredHi}{22}
\newcommand{\RvFewBMonoMappedFourHundredLo}{12}
\newcommand{\RvFewBMonoMappedHundred}{2.8}
\newcommand{\RvFewBMonoMappedHundredHi}{5.9}
\newcommand{\RvFewBMonoMappedHundredLo}{\ensuremath{-}0.6}
\newcommand{\RvFewBMonoMappedOneFifty}{5.7}
\newcommand{\RvFewBMonoMappedOneFiftyHi}{11}
\newcommand{\RvFewBMonoMappedOneFiftyLo}{\ensuremath{-}0.5}
\newcommand{\RvFewBMonoMappedSave}{at least 1.0}
\newcommand{\RvFewBMonoMappedSaveHi}{1.0}
\newcommand{\RvFewBMonoMappedSaveLo}{1.0}
\newcommand{\RvFewBMonoMappedTwentyFive}{\ensuremath{-}13}
\newcommand{\RvFewBMonoMappedTwentyFiveHi}{\ensuremath{-}9.9}
\newcommand{\RvFewBMonoMappedTwentyFiveLo}{\ensuremath{-}17}
\newcommand{\RvFewBMonoMappedTwoHundred}{10.0}
\newcommand{\RvFewBMonoMappedTwoHundredHi}{15}
\newcommand{\RvFewBMonoMappedTwoHundredLo}{5.1}
\newcommand{\RvFewBMonoRwCells}{148}
\newcommand{\RvFewBMonoRwCellsHi}{192}
\newcommand{\RvFewBMonoRwCellsLo}{122}
\newcommand{\RvFewBMonoRwEight}{46}
\newcommand{\RvFewBMonoRwEightHi}{49}
\newcommand{\RvFewBMonoRwEightLo}{42}
\newcommand{\RvFewBMonoRwFifty}{10.0}
\newcommand{\RvFewBMonoRwFiftyHi}{13}
\newcommand{\RvFewBMonoRwFiftyLo}{6.9}
\newcommand{\RvFewBMonoRwFourHundred}{40}
\newcommand{\RvFewBMonoRwFourHundredHi}{43}
\newcommand{\RvFewBMonoRwFourHundredLo}{37}
\newcommand{\RvFewBMonoRwHundred}{23}
\newcommand{\RvFewBMonoRwHundredHi}{25}
\newcommand{\RvFewBMonoRwHundredLo}{20}
\newcommand{\RvFewBMonoRwOneFifty}{29}
\newcommand{\RvFewBMonoRwOneFiftyHi}{32}
\newcommand{\RvFewBMonoRwOneFiftyLo}{25}
\newcommand{\RvFewBMonoRwSave}{at least 5.4}
\newcommand{\RvFewBMonoRwSaveHi}{6.5}
\newcommand{\RvFewBMonoRwSaveLo}{4.2}
\newcommand{\RvFewBMonoRwTwentyFive}{\ensuremath{-}25}
\newcommand{\RvFewBMonoRwTwentyFiveHi}{\ensuremath{-}4.3}
\newcommand{\RvFewBMonoRwTwentyFiveLo}{\ensuremath{-}64}
\newcommand{\RvFewBMonoRwTwoHundred}{32}
\newcommand{\RvFewBMonoRwTwoHundredHi}{35}
\newcommand{\RvFewBMonoRwTwoHundredLo}{29}
\newcommand{\RvFewTAxesCells}{more than 800}
\newcommand{\RvFewTAxesCellsHi}{800}
\newcommand{\RvFewTAxesCellsLo}{800}
\newcommand{\RvFewTAxesEight}{13}
\newcommand{\RvFewTAxesEightHi}{15}
\newcommand{\RvFewTAxesEightLo}{11}
\newcommand{\RvFewTAxesFifty}{2.5}
\newcommand{\RvFewTAxesFiftyHi}{3.2}
\newcommand{\RvFewTAxesFiftyLo}{1.8}
\newcommand{\RvFewTAxesFourHundred}{12}
\newcommand{\RvFewTAxesFourHundredHi}{13}
\newcommand{\RvFewTAxesFourHundredLo}{10.0}
\newcommand{\RvFewTAxesHundred}{5.9}
\newcommand{\RvFewTAxesHundredHi}{6.9}
\newcommand{\RvFewTAxesHundredLo}{4.9}
\newcommand{\RvFewTAxesOneFifty}{7.6}
\newcommand{\RvFewTAxesOneFiftyHi}{8.8}
\newcommand{\RvFewTAxesOneFiftyLo}{6.1}
\newcommand{\RvFewTAxesSave}{at least 1.0}
\newcommand{\RvFewTAxesSaveHi}{1.0}
\newcommand{\RvFewTAxesSaveLo}{1.0}
\newcommand{\RvFewTAxesTwentyFive}{\ensuremath{-}1.3}
\newcommand{\RvFewTAxesTwentyFiveHi}{\ensuremath{-}0.3}
\newcommand{\RvFewTAxesTwentyFiveLo}{\ensuremath{-}2.7}
\newcommand{\RvFewTAxesTwoHundred}{8.9}
\newcommand{\RvFewTAxesTwoHundredHi}{10}
\newcommand{\RvFewTAxesTwoHundredLo}{7.5}
\newcommand{\RvFewTMappedCells}{more than 800}
\newcommand{\RvFewTMappedCellsHi}{800}
\newcommand{\RvFewTMappedCellsLo}{800}
\newcommand{\RvFewTMappedEight}{\ensuremath{-}2.7}
\newcommand{\RvFewTMappedEightHi}{12}
\newcommand{\RvFewTMappedEightLo}{\ensuremath{-}19}
\newcommand{\RvFewTMappedFifty}{7.2}
\newcommand{\RvFewTMappedFiftyHi}{10.0}
\newcommand{\RvFewTMappedFiftyLo}{4.0}
\newcommand{\RvFewTMappedFourHundred}{8.8}
\newcommand{\RvFewTMappedFourHundredHi}{20}
\newcommand{\RvFewTMappedFourHundredLo}{\ensuremath{-}3.9}
\newcommand{\RvFewTMappedHundred}{14}
\newcommand{\RvFewTMappedHundredHi}{18}
\newcommand{\RvFewTMappedHundredLo}{9.3}
\newcommand{\RvFewTMappedOneFifty}{15}
\newcommand{\RvFewTMappedOneFiftyHi}{21}
\newcommand{\RvFewTMappedOneFiftyLo}{7.0}
\newcommand{\RvFewTMappedSave}{at least 1.0}
\newcommand{\RvFewTMappedSaveHi}{1.0}
\newcommand{\RvFewTMappedSaveLo}{1.0}
\newcommand{\RvFewTMappedTwentyFive}{\ensuremath{-}3.3}
\newcommand{\RvFewTMappedTwentyFiveHi}{0.3}
\newcommand{\RvFewTMappedTwentyFiveLo}{\ensuremath{-}8.1}
\newcommand{\RvFewTMappedTwoHundred}{14}
\newcommand{\RvFewTMappedTwoHundredHi}{22}
\newcommand{\RvFewTMappedTwoHundredLo}{4.8}
\newcommand{\RvFewTRwCells}{311}
\newcommand{\RvFewTRwCellsHi}{378}
\newcommand{\RvFewTRwCellsLo}{266}
\newcommand{\RvFewTRwEight}{41}
\newcommand{\RvFewTRwEightHi}{43}
\newcommand{\RvFewTRwEightLo}{37}
\newcommand{\RvFewTRwFifty}{0.6}
\newcommand{\RvFewTRwFiftyHi}{3.6}
\newcommand{\RvFewTRwFiftyLo}{\ensuremath{-}3.2}
\newcommand{\RvFewTRwFourHundred}{32}
\newcommand{\RvFewTRwFourHundredHi}{35}
\newcommand{\RvFewTRwFourHundredLo}{29}
\newcommand{\RvFewTRwHundred}{12}
\newcommand{\RvFewTRwHundredHi}{14}
\newcommand{\RvFewTRwHundredLo}{9.6}
\newcommand{\RvFewTRwOneFifty}{18}
\newcommand{\RvFewTRwOneFiftyHi}{20}
\newcommand{\RvFewTRwOneFiftyLo}{14}
\newcommand{\RvFewTRwSave}{at least 2.6}
\newcommand{\RvFewTRwSaveHi}{3.0}
\newcommand{\RvFewTRwSaveLo}{2.1}
\newcommand{\RvFewTRwTwentyFive}{\ensuremath{-}35}
\newcommand{\RvFewTRwTwentyFiveHi}{\ensuremath{-}11}
\newcommand{\RvFewTRwTwentyFiveLo}{\ensuremath{-}72}
\newcommand{\RvFewTRwTwoHundred}{22}
\newcommand{\RvFewTRwTwoHundredHi}{24}
\newcommand{\RvFewTRwTwoHundredLo}{19}
\newcommand{\RvFineOneAtlas}{74}
\newcommand{\RvFineOneMinusPo}{2.8}
\newcommand{\RvFineOneMinusPoHi}{3.4}
\newcommand{\RvFineOneMinusPoLo}{2.1}
\newcommand{\RvFineOneMinusReg}{\ensuremath{-}2.0}
\newcommand{\RvFineOneMinusRegHi}{\ensuremath{-}1.3}
\newcommand{\RvFineOneMinusRegLo}{\ensuremath{-}2.7}
\newcommand{\RvFineOnePo}{71}
\newcommand{\RvFineOneReg}{76}
\newcommand{\RvFineReps}{31,509}
\newcommand{\RvFineRfAtlas}{\ensuremath{-}78}
\newcommand{\RvFineRfMinusPo}{\ensuremath{-}71}
\newcommand{\RvFineRfMinusPoHi}{\ensuremath{-}58}
\newcommand{\RvFineRfMinusPoLo}{\ensuremath{-}87}
\newcommand{\RvFineRfMinusReg}{\ensuremath{-}70}
\newcommand{\RvFineRfMinusRegHi}{\ensuremath{-}57}
\newcommand{\RvFineRfMinusRegLo}{\ensuremath{-}85}
\newcommand{\RvFineRfPo}{\ensuremath{-}6.6}
\newcommand{\RvFineRfReg}{\ensuremath{-}7.6}
\newcommand{\RvFineTwoAtlas}{17}
\newcommand{\RvFineTwoMinusPo}{6.5}
\newcommand{\RvFineTwoMinusPoHi}{7.0}
\newcommand{\RvFineTwoMinusPoLo}{6.1}
\newcommand{\RvFineTwoMinusReg}{\ensuremath{-}1.9}
\newcommand{\RvFineTwoMinusRegHi}{\ensuremath{-}1.6}
\newcommand{\RvFineTwoMinusRegLo}{\ensuremath{-}2.2}
\newcommand{\RvFineTwoPo}{10}
\newcommand{\RvFineTwoReg}{19}
\newcommand{\RvFineUnits}{170}
\newcommand{\RvFisherFive}{4.90}
\newcommand{\RvFisherOne}{1.07}
\newcommand{\RvFisherOneMax}{1.10}
\newcommand{\RvFisherOneMin}{1.06}
\newcommand{\RvFisherPerm}{20}
\newcommand{\RvFisherSmallFive}{4.54}
\newcommand{\RvFisherSmallLo}{40}
\newcommand{\RvFisherSmallOne}{1.07}
\newcommand{\RvFisherSmallRange}{40 to 79}
\newcommand{\RvFisherSmallTenth}{0.18}
\newcommand{\RvFisherTenth}{0.149}
\newcommand{\RvFisherUnits}{170}
\newcommand{\RvFixedSameRunningMax}{none}
\newcommand{\RvFixedSets}{10}
\newcommand{\RvFixedTenAbove}{6}
\newcommand{\RvFixedTenBelow}{2}
\newcommand{\RvFixedTenBothAbove}{two}
\newcommand{\RvFixedTenEstBelow}{no}
\newcommand{\RvFixedTenGm}{1.7}
\newcommand{\RvFixedTenMax}{2.6}
\newcommand{\RvFixedTenMin}{1.2}
\newcommand{\RvFixedTenN}{2}
\newcommand{\RvFixedTenPoAbove}{four}
\newcommand{\RvFixedThirtyAbove}{7}
\newcommand{\RvFixedThirtyBelow}{2}
\newcommand{\RvFixedThirtyBothAbove}{four}
\newcommand{\RvFixedThirtyEstBelow}{two}
\newcommand{\RvFixedThirtyGm}{4.9}
\newcommand{\RvFixedThirtyMax}{4.9}
\newcommand{\RvFixedThirtyMin}{4.9}
\newcommand{\RvFixedThirtyN}{1}
\newcommand{\RvFixedThirtyPoAbove}{three}
\newcommand{\RvFixedTwentyAbove}{6}
\newcommand{\RvFixedTwentyBelow}{2}
\newcommand{\RvFixedTwentyBothAbove}{two}
\newcommand{\RvFixedTwentyEstBelow}{two}
\newcommand{\RvFixedTwentyGm}{2.5}
\newcommand{\RvFixedTwentyMax}{4.4}
\newcommand{\RvFixedTwentyMin}{1.4}
\newcommand{\RvFixedTwentyN}{2}
\newcommand{\RvFixedTwentyPoAbove}{four}
\newcommand{\RvFrozenAt}{04:30 EDT on 30 September 2026}
\newcommand{\RvFullRwCells}{131}
\newcommand{\RvFullRwCellsHi}{162}
\newcommand{\RvFullRwCellsLo}{114}
\newcommand{\RvFullRwEight}{50}
\newcommand{\RvFullRwEightHi}{53}
\newcommand{\RvFullRwEightLo}{47}
\newcommand{\RvFullRwFifty}{8.3}
\newcommand{\RvFullRwFiftyHi}{12}
\newcommand{\RvFullRwFiftyLo}{2.9}
\newcommand{\RvFullRwFourHundred}{44}
\newcommand{\RvFullRwFourHundredHi}{47}
\newcommand{\RvFullRwFourHundredLo}{40}
\newcommand{\RvFullRwHundred}{24}
\newcommand{\RvFullRwHundredHi}{26}
\newcommand{\RvFullRwHundredLo}{21}
\newcommand{\RvFullRwOneFifty}{31}
\newcommand{\RvFullRwOneFiftyHi}{34}
\newcommand{\RvFullRwOneFiftyLo}{27}
\newcommand{\RvFullRwSave}{at least 6.1}
\newcommand{\RvFullRwSaveHi}{7.0}
\newcommand{\RvFullRwSaveLo}{4.9}
\newcommand{\RvFullRwTwentyFive}{\ensuremath{-}37}
\newcommand{\RvFullRwTwentyFiveHi}{\ensuremath{-}8.2}
\newcommand{\RvFullRwTwentyFiveLo}{\ensuremath{-}84}
\newcommand{\RvFullRwTwoHundred}{35}
\newcommand{\RvFullRwTwoHundredHi}{38}
\newcommand{\RvFullRwTwoHundredLo}{32}
\newcommand{\RvLawHalfBelow}{8}
\newcommand{\RvLawHalfErr}{21}
\newcommand{\RvLawHalfErrHi}{29}
\newcommand{\RvLawHalfErrLo}{11}
\newcommand{\RvLawHalfErrStHi}{26}
\newcommand{\RvLawHalfErrStLo}{13}
\newcommand{\RvLawHalfLooMax}{26}
\newcommand{\RvLawHalfLooMin}{17}
\newcommand{\RvLawHalfMeansP}{0.042}
\newcommand{\RvLawHalfMeansRho}{1.00}
\newcommand{\RvLawHalfN}{10}
\newcommand{\RvLawHalfNonN}{19}
\newcommand{\RvLawHalfNotReached}{10}
\newcommand{\RvLawHalfOneAbove}{1}
\newcommand{\RvLawHalfOver}{4}
\newcommand{\RvLawHalfSets}{4}
\newcommand{\RvLawHalfSlope}{1.08}
\newcommand{\RvLawHalfSlopeHi}{1.26}
\newcommand{\RvLawHalfSlopeLo}{0.99}
\newcommand{\RvLawHalfWithinP}{0.007}
\newcommand{\RvLawHalfWithinPMin}{0.007}
\newcommand{\RvLawHalfWithinPerm}{144}
\newcommand{\RvLawSevenBelow}{0}
\newcommand{\RvLawSevenErr}{4}
\newcommand{\RvLawSevenErrHi}{5}
\newcommand{\RvLawSevenErrLo}{2}
\newcommand{\RvLawSevenErrStHi}{4}
\newcommand{\RvLawSevenErrStLo}{4}
\newcommand{\RvLawSevenLooMax}{5}
\newcommand{\RvLawSevenLooMin}{2}
\newcommand{\RvLawSevenN}{5}
\newcommand{\RvLawSevenNonN}{24}
\newcommand{\RvLawSevenNotReached}{15}
\newcommand{\RvLawSevenOneAbove}{9}
\newcommand{\RvLawSevenOver}{1}
\newcommand{\RvLawSevenSets}{2}
\newcommand{\RvLawSevenSlope}{1.02}
\newcommand{\RvLawSevenSlopeHi}{1.15}
\newcommand{\RvLawSevenSlopeLo}{0.99}
\newcommand{\RvLawSevenWithinP}{0.083}
\newcommand{\RvLawSevenWithinPMin}{0.083}
\newcommand{\RvLawSevenWithinPerm}{12}
\newcommand{\RvLawThreeBelow}{9}
\newcommand{\RvLawThreeErr}{74}
\newcommand{\RvLawThreeErrHi}{144}
\newcommand{\RvLawThreeErrLo}{54}
\newcommand{\RvLawThreeErrStHi}{144}
\newcommand{\RvLawThreeErrStLo}{61}
\newcommand{\RvLawThreeLooMax}{99}
\newcommand{\RvLawThreeLooMin}{61}
\newcommand{\RvLawThreeMeansP}{0.167}
\newcommand{\RvLawThreeMeansRho}{1.00}
\newcommand{\RvLawThreeN}{8}
\newcommand{\RvLawThreeNonN}{21}
\newcommand{\RvLawThreeNotReached}{7}
\newcommand{\RvLawThreeOneAbove}{5}
\newcommand{\RvLawThreeOver}{3}
\newcommand{\RvLawThreeSets}{3}
\newcommand{\RvLawThreeSlope}{0.98}
\newcommand{\RvLawThreeSlopeHi}{1.22}
\newcommand{\RvLawThreeSlopeLo}{-2.39}
\newcommand{\RvLawThreeWithinP}{0.028}
\newcommand{\RvLawThreeWithinPMin}{0.014}
\newcommand{\RvLawThreeWithinPerm}{72}
\newcommand{\RvMapOnlyCells}{more than 800}
\newcommand{\RvMapOnlyCellsHi}{800}
\newcommand{\RvMapOnlyCellsLo}{683}
\newcommand{\RvMapOnlyEight}{28}
\newcommand{\RvMapOnlyEightHi}{31}
\newcommand{\RvMapOnlyEightLo}{25}
\newcommand{\RvMapOnlyFifty}{1.8}
\newcommand{\RvMapOnlyFiftyHi}{2.3}
\newcommand{\RvMapOnlyFiftyLo}{1.3}
\newcommand{\RvMapOnlyFourHundred}{18}
\newcommand{\RvMapOnlyFourHundredHi}{20}
\newcommand{\RvMapOnlyFourHundredLo}{15}
\newcommand{\RvMapOnlyHundred}{4.7}
\newcommand{\RvMapOnlyHundredHi}{5.5}
\newcommand{\RvMapOnlyHundredLo}{3.9}
\newcommand{\RvMapOnlyOneFifty}{7.1}
\newcommand{\RvMapOnlyOneFiftyHi}{8.5}
\newcommand{\RvMapOnlyOneFiftyLo}{5.5}
\newcommand{\RvMapOnlySave}{at least 1.0}
\newcommand{\RvMapOnlySaveHi}{1.2}
\newcommand{\RvMapOnlySaveLo}{1.0}
\newcommand{\RvMapOnlyTwentyFive}{\ensuremath{-}0.2}
\newcommand{\RvMapOnlyTwentyFiveHi}{0.3}
\newcommand{\RvMapOnlyTwentyFiveLo}{\ensuremath{-}0.8}
\newcommand{\RvMapOnlyTwoHundred}{9.5}
\newcommand{\RvMapOnlyTwoHundredHi}{11}
\newcommand{\RvMapOnlyTwoHundredLo}{8.0}
\newcommand{\RvNeitherCells}{more than 800}
\newcommand{\RvNeitherCellsHi}{800}
\newcommand{\RvNeitherCellsLo}{800}
\newcommand{\RvNeitherEight}{9.9}
\newcommand{\RvNeitherEightHi}{11}
\newcommand{\RvNeitherEightLo}{8.5}
\newcommand{\RvNeitherFifty}{1.1}
\newcommand{\RvNeitherFiftyHi}{1.2}
\newcommand{\RvNeitherFiftyLo}{0.9}
\newcommand{\RvNeitherFourHundred}{7.1}
\newcommand{\RvNeitherFourHundredHi}{7.9}
\newcommand{\RvNeitherFourHundredLo}{6.2}
\newcommand{\RvNeitherHundred}{2.4}
\newcommand{\RvNeitherHundredHi}{2.7}
\newcommand{\RvNeitherHundredLo}{2.0}
\newcommand{\RvNeitherOneFifty}{3.4}
\newcommand{\RvNeitherOneFiftyHi}{3.9}
\newcommand{\RvNeitherOneFiftyLo}{2.8}
\newcommand{\RvNeitherTwentyFive}{0.0}
\newcommand{\RvNeitherTwentyFiveHi}{0.3}
\newcommand{\RvNeitherTwentyFiveLo}{\ensuremath{-}0.3}
\newcommand{\RvNeitherTwoHundred}{4.4}
\newcommand{\RvNeitherTwoHundredHi}{5.0}
\newcommand{\RvNeitherTwoHundredLo}{3.8}
\newcommand{\RvOtherAxesCells}{more than 800}
\newcommand{\RvOtherAxesCellsHi}{800}
\newcommand{\RvOtherAxesCellsLo}{800}
\newcommand{\RvOtherAxesEight}{14}
\newcommand{\RvOtherAxesEightHi}{17}
\newcommand{\RvOtherAxesEightLo}{12}
\newcommand{\RvOtherAxesFifty}{7.5}
\newcommand{\RvOtherAxesFiftyHi}{9.1}
\newcommand{\RvOtherAxesFiftyLo}{5.8}
\newcommand{\RvOtherAxesFourHundred}{13}
\newcommand{\RvOtherAxesFourHundredHi}{15}
\newcommand{\RvOtherAxesFourHundredLo}{11}
\newcommand{\RvOtherAxesHundred}{10}
\newcommand{\RvOtherAxesHundredHi}{12}
\newcommand{\RvOtherAxesHundredLo}{8.7}
\newcommand{\RvOtherAxesOneFifty}{11}
\newcommand{\RvOtherAxesOneFiftyHi}{13}
\newcommand{\RvOtherAxesOneFiftyLo}{8.9}
\newcommand{\RvOtherAxesSave}{at least 1.0}
\newcommand{\RvOtherAxesSaveHi}{1.0}
\newcommand{\RvOtherAxesSaveLo}{1.0}
\newcommand{\RvOtherAxesTwentyFive}{1.6}
\newcommand{\RvOtherAxesTwentyFiveHi}{3.5}
\newcommand{\RvOtherAxesTwentyFiveLo}{\ensuremath{-}1.0}
\newcommand{\RvOtherAxesTwoHundred}{12}
\newcommand{\RvOtherAxesTwoHundredHi}{14}
\newcommand{\RvOtherAxesTwoHundredLo}{9.8}
\newcommand{\RvOwnAxesCells}{more than 800}
\newcommand{\RvOwnAxesCellsHi}{800}
\newcommand{\RvOwnAxesCellsLo}{800}
\newcommand{\RvOwnAxesEight}{13}
\newcommand{\RvOwnAxesEightHi}{15}
\newcommand{\RvOwnAxesEightLo}{10}
\newcommand{\RvOwnAxesFifty}{6.5}
\newcommand{\RvOwnAxesFiftyHi}{8.3}
\newcommand{\RvOwnAxesFiftyLo}{4.9}
\newcommand{\RvOwnAxesFourHundred}{12}
\newcommand{\RvOwnAxesFourHundredHi}{14}
\newcommand{\RvOwnAxesFourHundredLo}{9.6}
\newcommand{\RvOwnAxesHundred}{9.2}
\newcommand{\RvOwnAxesHundredHi}{11}
\newcommand{\RvOwnAxesHundredLo}{7.3}
\newcommand{\RvOwnAxesOneFifty}{9.7}
\newcommand{\RvOwnAxesOneFiftyHi}{12}
\newcommand{\RvOwnAxesOneFiftyLo}{7.2}
\newcommand{\RvOwnAxesSave}{at least 1.0}
\newcommand{\RvOwnAxesSaveHi}{1.0}
\newcommand{\RvOwnAxesSaveLo}{1.0}
\newcommand{\RvOwnAxesTwentyFive}{0.5}
\newcommand{\RvOwnAxesTwentyFiveHi}{3.0}
\newcommand{\RvOwnAxesTwentyFiveLo}{\ensuremath{-}3.2}
\newcommand{\RvOwnAxesTwoHundred}{10}
\newcommand{\RvOwnAxesTwoHundredHi}{13}
\newcommand{\RvOwnAxesTwoHundredLo}{8.1}
\newcommand{\RvPilotHalfFiftyAchieved}{54}
\newcommand{\RvPilotHalfFiftyBudget}{84}
\newcommand{\RvPilotHalfFiftyFlagged}{5}
\newcommand{\RvPilotHalfFiftyProblems}{29}
\newcommand{\RvPilotHalfFiftyRatio}{1.16}
\newcommand{\RvPilotHalfFiftyReached}{64}
\newcommand{\RvPilotHalfFiftyWithin}{72}
\newcommand{\RvPilotHalfHundredAchieved}{53}
\newcommand{\RvPilotHalfHundredBudget}{100}
\newcommand{\RvPilotHalfHundredFlagged}{8}
\newcommand{\RvPilotHalfHundredProblems}{29}
\newcommand{\RvPilotHalfHundredRatio}{1.12}
\newcommand{\RvPilotHalfHundredReached}{65}
\newcommand{\RvPilotHalfHundredWithin}{75}
\newcommand{\RvPilotHalfResAchieved}{46}
\newcommand{\RvPilotHalfResBudget}{86}
\newcommand{\RvPilotHalfResFlagged}{0}
\newcommand{\RvPilotHalfResProblems}{29}
\newcommand{\RvPilotHalfResRatio}{0.77}
\newcommand{\RvPilotHalfResReached}{26}
\newcommand{\RvPilotHalfResWithin}{50}
\newcommand{\RvPilotHalfTwoHundredAchieved}{53}
\newcommand{\RvPilotHalfTwoHundredBudget}{200}
\newcommand{\RvPilotHalfTwoHundredFlagged}{10}
\newcommand{\RvPilotHalfTwoHundredProblems}{29}
\newcommand{\RvPilotHalfTwoHundredRatio}{0.87}
\newcommand{\RvPilotHalfTwoHundredReached}{65}
\newcommand{\RvPilotHalfTwoHundredWithin}{77}
\newcommand{\RvPilotSevenFiftyAchieved}{68}
\newcommand{\RvPilotSevenFiftyBudget}{181}
\newcommand{\RvPilotSevenFiftyFlagged}{22}
\newcommand{\RvPilotSevenFiftyProblems}{29}
\newcommand{\RvPilotSevenFiftyRatio}{0.82}
\newcommand{\RvPilotSevenFiftyReached}{46}
\newcommand{\RvPilotSevenFiftyWithin}{56}
\newcommand{\RvPilotSevenHundredAchieved}{69}
\newcommand{\RvPilotSevenHundredBudget}{181}
\newcommand{\RvPilotSevenHundredFlagged}{24}
\newcommand{\RvPilotSevenHundredProblems}{29}
\newcommand{\RvPilotSevenHundredRatio}{0.85}
\newcommand{\RvPilotSevenHundredReached}{48}
\newcommand{\RvPilotSevenHundredWithin}{59}
\newcommand{\RvPilotSevenResAchieved}{67}
\newcommand{\RvPilotSevenResBudget}{181}
\newcommand{\RvPilotSevenResFlagged}{0}
\newcommand{\RvPilotSevenResProblems}{29}
\newcommand{\RvPilotSevenResRatio}{0.82}
\newcommand{\RvPilotSevenResReached}{23}
\newcommand{\RvPilotSevenResWithin}{49}
\newcommand{\RvPilotSevenTwoHundredAchieved}{70}
\newcommand{\RvPilotSevenTwoHundredBudget}{200}
\newcommand{\RvPilotSevenTwoHundredFlagged}{25}
\newcommand{\RvPilotSevenTwoHundredProblems}{29}
\newcommand{\RvPilotSevenTwoHundredRatio}{0.99}
\newcommand{\RvPilotSevenTwoHundredReached}{56}
\newcommand{\RvPilotSevenTwoHundredWithin}{61}
\newcommand{\RvPilotThreeFiftyAchieved}{41}
\newcommand{\RvPilotThreeFiftyBudget}{50}
\newcommand{\RvPilotThreeFiftyFlagged}{0}
\newcommand{\RvPilotThreeFiftyProblems}{29}
\newcommand{\RvPilotThreeFiftyRatio}{0.99}
\newcommand{\RvPilotThreeFiftyReached}{73}
\newcommand{\RvPilotThreeFiftyWithin}{78}
\newcommand{\RvPilotThreeHundredAchieved}{41}
\newcommand{\RvPilotThreeHundredBudget}{100}
\newcommand{\RvPilotThreeHundredFlagged}{1}
\newcommand{\RvPilotThreeHundredProblems}{29}
\newcommand{\RvPilotThreeHundredRatio}{1.02}
\newcommand{\RvPilotThreeHundredReached}{77}
\newcommand{\RvPilotThreeHundredWithin}{82}
\newcommand{\RvPilotThreeResAchieved}{23}
\newcommand{\RvPilotThreeResBudget}{32}
\newcommand{\RvPilotThreeResFlagged}{0}
\newcommand{\RvPilotThreeResProblems}{29}
\newcommand{\RvPilotThreeResRatio}{0.57}
\newcommand{\RvPilotThreeResReached}{34}
\newcommand{\RvPilotThreeResWithin}{45}
\newcommand{\RvPilotThreeTwoHundredAchieved}{48}
\newcommand{\RvPilotThreeTwoHundredBudget}{200}
\newcommand{\RvPilotThreeTwoHundredFlagged}{1}
\newcommand{\RvPilotThreeTwoHundredProblems}{29}
\newcommand{\RvPilotThreeTwoHundredRatio}{1.16}
\newcommand{\RvPilotThreeTwoHundredReached}{83}
\newcommand{\RvPilotThreeTwoHundredWithin}{87}
\newcommand{\RvPoCells}{more than 800}
\newcommand{\RvPoCellsHi}{800}
\newcommand{\RvPoCellsLo}{800}
\newcommand{\RvPoEight}{13}
\newcommand{\RvPoEightHi}{20}
\newcommand{\RvPoEightLo}{11}
\newcommand{\RvPoFifty}{2.4}
\newcommand{\RvPoFiftyHi}{3.8}
\newcommand{\RvPoFiftyLo}{1.2}
\newcommand{\RvPoFourHundred}{14}
\newcommand{\RvPoFourHundredHi}{19}
\newcommand{\RvPoFourHundredLo}{7.9}
\newcommand{\RvPoHundred}{3.5}
\newcommand{\RvPoHundredHi}{6.1}
\newcommand{\RvPoHundredLo}{2.9}
\newcommand{\RvPoOneFifty}{5.1}
\newcommand{\RvPoOneFiftyHi}{9.3}
\newcommand{\RvPoOneFiftyLo}{3.5}
\newcommand{\RvPoTwentyFive}{0.5}
\newcommand{\RvPoTwentyFiveHi}{0.9}
\newcommand{\RvPoTwentyFiveLo}{0.4}
\newcommand{\RvPoTwoHundred}{7.7}
\newcommand{\RvPoTwoHundredHi}{12}
\newcommand{\RvPoTwoHundredLo}{4.5}
\newcommand{\RvPrimaryOneAtlas}{86}
\newcommand{\RvPrimaryOneMinusPo}{11}
\newcommand{\RvPrimaryOneMinusPoHi}{12}
\newcommand{\RvPrimaryOneMinusPoLo}{10}
\newcommand{\RvPrimaryOneMinusReg}{\ensuremath{-}0.1}
\newcommand{\RvPrimaryOneMinusRegHi}{0.2}
\newcommand{\RvPrimaryOneMinusRegLo}{\ensuremath{-}0.5}
\newcommand{\RvPrimaryOnePo}{75}
\newcommand{\RvPrimaryOneReg}{86}
\newcommand{\RvPrimaryReps}{120,649}
\newcommand{\RvPrimaryRfAtlas}{23}
\newcommand{\RvPrimaryRfMinusPo}{19}
\newcommand{\RvPrimaryRfMinusPoHi}{20}
\newcommand{\RvPrimaryRfMinusPoLo}{17}
\newcommand{\RvPrimaryRfMinusReg}{0.1}
\newcommand{\RvPrimaryRfMinusRegHi}{1.6}
\newcommand{\RvPrimaryRfMinusRegLo}{\ensuremath{-}1.3}
\newcommand{\RvPrimaryRfPo}{3.5}
\newcommand{\RvPrimaryRfReg}{22}
\newcommand{\RvPrimaryTwoAtlas}{47}
\newcommand{\RvPrimaryTwoMinusPo}{19}
\newcommand{\RvPrimaryTwoMinusPoHi}{20}
\newcommand{\RvPrimaryTwoMinusPoLo}{18}
\newcommand{\RvPrimaryTwoMinusReg}{\ensuremath{-}2.8}
\newcommand{\RvPrimaryTwoMinusRegHi}{\ensuremath{-}1.9}
\newcommand{\RvPrimaryTwoMinusRegLo}{\ensuremath{-}3.6}
\newcommand{\RvPrimaryTwoPo}{28}
\newcommand{\RvPrimaryTwoReg}{50}
\newcommand{\RvPrimaryUnits}{170}
\newcommand{\RvRegCells}{152}
\newcommand{\RvRegCellsHi}{231}
\newcommand{\RvRegCellsLo}{123}
\newcommand{\RvRegEight}{45}
\newcommand{\RvRegEightHi}{48}
\newcommand{\RvRegEightLo}{42}
\newcommand{\RvRegFifty}{14}
\newcommand{\RvRegFiftyHi}{16}
\newcommand{\RvRegFiftyLo}{11}
\newcommand{\RvRegFourHundred}{38}
\newcommand{\RvRegFourHundredHi}{41}
\newcommand{\RvRegFourHundredLo}{35}
\newcommand{\RvRegHundred}{22}
\newcommand{\RvRegHundredHi}{24}
\newcommand{\RvRegHundredLo}{20}
\newcommand{\RvRegOneFifty}{28}
\newcommand{\RvRegOneFiftyHi}{32}
\newcommand{\RvRegOneFiftyLo}{24}
\newcommand{\RvRegSave}{at least 5.3}
\newcommand{\RvRegSaveHi}{6.5}
\newcommand{\RvRegSaveLo}{3.5}
\newcommand{\RvRegTwentyFive}{5.6}
\newcommand{\RvRegTwentyFiveHi}{6.5}
\newcommand{\RvRegTwentyFiveLo}{4.7}
\newcommand{\RvRegTwoHundred}{30}
\newcommand{\RvRegTwoHundredHi}{35}
\newcommand{\RvRegTwoHundredLo}{26}
\newcommand{\RvRelBetterPo}{7}
\newcommand{\RvRelBetterReg}{7}
\newcommand{\RvRelDonorWinPo}{81}
\newcommand{\RvRelDonorWinReg}{77}
\newcommand{\RvRelEightAtlasPred}{0.11}
\newcommand{\RvRelEightAtlasRf}{55}
\newcommand{\RvRelEightAtlasSign}{100}
\newcommand{\RvRelEightDonorsPo}{18 of 18}
\newcommand{\RvRelEightDonorsReg}{18 of 18}
\newcommand{\RvRelEightPoPred}{0.06}
\newcommand{\RvRelEightPoRf}{29}
\newcommand{\RvRelEightPoSign}{96}
\newcommand{\RvRelEightRegPred}{0.08}
\newcommand{\RvRelEightRegRf}{38}
\newcommand{\RvRelEightRegSign}{100}
\newcommand{\RvRelEightTruth}{0.33}
\newcommand{\RvRelEightTruthShare}{100}
\newcommand{\RvRelEightUnits}{34}
\newcommand{\RvRelFiveAtlasPred}{\ensuremath{-}0.28}
\newcommand{\RvRelFiveAtlasRf}{69}
\newcommand{\RvRelFiveAtlasSign}{99}
\newcommand{\RvRelFiveDonorsPo}{16 of 18}
\newcommand{\RvRelFiveDonorsReg}{10 of 18}
\newcommand{\RvRelFivePoPred}{\ensuremath{-}0.22}
\newcommand{\RvRelFivePoRf}{56}
\newcommand{\RvRelFivePoSign}{99}
\newcommand{\RvRelFiveRegPred}{\ensuremath{-}0.32}
\newcommand{\RvRelFiveRegRf}{66}
\newcommand{\RvRelFiveRegSign}{100}
\newcommand{\RvRelFiveTruth}{\ensuremath{-}0.49}
\newcommand{\RvRelFiveTruthShare}{100}
\newcommand{\RvRelFiveUnits}{34}
\newcommand{\RvRelFourAtlasPred}{0.28}
\newcommand{\RvRelFourAtlasRf}{65}
\newcommand{\RvRelFourAtlasSign}{99}
\newcommand{\RvRelFourDonorsPo}{12 of 18}
\newcommand{\RvRelFourDonorsReg}{12 of 18}
\newcommand{\RvRelFourPoPred}{0.28}
\newcommand{\RvRelFourPoRf}{57}
\newcommand{\RvRelFourPoSign}{94}
\newcommand{\RvRelFourRegPred}{0.25}
\newcommand{\RvRelFourRegRf}{63}
\newcommand{\RvRelFourRegSign}{93}
\newcommand{\RvRelFourTruth}{0.50}
\newcommand{\RvRelFourTruthShare}{100}
\newcommand{\RvRelFourUnits}{34}
\newcommand{\RvRelN}{8}
\newcommand{\RvRelOneAtlasPred}{0.02}
\newcommand{\RvRelOneAtlasRf}{27}
\newcommand{\RvRelOneAtlasSign}{72}
\newcommand{\RvRelOneDonorsPo}{15 of 18}
\newcommand{\RvRelOneDonorsReg}{14 of 18}
\newcommand{\RvRelOnePoPred}{0.01}
\newcommand{\RvRelOnePoRf}{15}
\newcommand{\RvRelOnePoSign}{76}
\newcommand{\RvRelOneRegPred}{0.01}
\newcommand{\RvRelOneRegRf}{21}
\newcommand{\RvRelOneRegSign}{79}
\newcommand{\RvRelOneTruth}{0.07}
\newcommand{\RvRelOneTruthShare}{100}
\newcommand{\RvRelOneUnits}{34}
\newcommand{\RvRelSevenAtlasPred}{0.01}
\newcommand{\RvRelSevenAtlasRf}{\ensuremath{-}81}
\newcommand{\RvRelSevenAtlasSign}{36}
\newcommand{\RvRelSevenDonorsPo}{3 of 18}
\newcommand{\RvRelSevenDonorsReg}{4 of 18}
\newcommand{\RvRelSevenPoPred}{0.00}
\newcommand{\RvRelSevenPoRf}{0.0}
\newcommand{\RvRelSevenPoSign}{0}
\newcommand{\RvRelSevenRegPred}{\ensuremath{-}0.03}
\newcommand{\RvRelSevenRegRf}{\ensuremath{-}20}
\newcommand{\RvRelSevenRegSign}{46}
\newcommand{\RvRelSevenTruth}{\ensuremath{-}0.05}
\newcommand{\RvRelSevenTruthShare}{68}
\newcommand{\RvRelSevenUnits}{34}
\newcommand{\RvRelSixAtlasPred}{0.13}
\newcommand{\RvRelSixAtlasRf}{53}
\newcommand{\RvRelSixAtlasSign}{98}
\newcommand{\RvRelSixDonorsPo}{18 of 18}
\newcommand{\RvRelSixDonorsReg}{18 of 18}
\newcommand{\RvRelSixPoPred}{0.04}
\newcommand{\RvRelSixPoRf}{18}
\newcommand{\RvRelSixPoSign}{81}
\newcommand{\RvRelSixRegPred}{0.09}
\newcommand{\RvRelSixRegRf}{37}
\newcommand{\RvRelSixRegSign}{95}
\newcommand{\RvRelSixTruth}{0.38}
\newcommand{\RvRelSixTruthShare}{100}
\newcommand{\RvRelSixUnits}{34}
\newcommand{\RvRelThreeAtlasPred}{0.10}
\newcommand{\RvRelThreeAtlasRf}{81}
\newcommand{\RvRelThreeAtlasSign}{100}
\newcommand{\RvRelThreeDonorsPo}{17 of 18}
\newcommand{\RvRelThreeDonorsReg}{17 of 18}
\newcommand{\RvRelThreePoPred}{0.04}
\newcommand{\RvRelThreePoRf}{39}
\newcommand{\RvRelThreePoSign}{100}
\newcommand{\RvRelThreeRegPred}{0.05}
\newcommand{\RvRelThreeRegRf}{54}
\newcommand{\RvRelThreeRegSign}{100}
\newcommand{\RvRelThreeTruth}{0.15}
\newcommand{\RvRelThreeTruthShare}{100}
\newcommand{\RvRelThreeUnits}{34}
\newcommand{\RvRelTwoAtlasPred}{0.31}
\newcommand{\RvRelTwoAtlasRf}{80}
\newcommand{\RvRelTwoAtlasSign}{100}
\newcommand{\RvRelTwoDonorsPo}{18 of 18}
\newcommand{\RvRelTwoDonorsReg}{18 of 18}
\newcommand{\RvRelTwoPoPred}{0.05}
\newcommand{\RvRelTwoPoRf}{19}
\newcommand{\RvRelTwoPoSign}{93}
\newcommand{\RvRelTwoRegPred}{0.13}
\newcommand{\RvRelTwoRegRf}{46}
\newcommand{\RvRelTwoRegSign}{100}
\newcommand{\RvRelTwoTruth}{0.47}
\newcommand{\RvRelTwoTruthShare}{100}
\newcommand{\RvRelTwoUnits}{34}
\newcommand{\RvRelUnits}{272}
\newcommand{\RvReplFalseShare}{0.21}
\newcommand{\RvReplNominal}{5.0}
\newcommand{\RvReplNull}{6.8}
\newcommand{\RvReplObserved}{3.2}
\newcommand{\RvTechFineOneAtlas}{73}
\newcommand{\RvTechFineOneMinusPo}{3.2}
\newcommand{\RvTechFineOneMinusPoHi}{3.6}
\newcommand{\RvTechFineOneMinusPoLo}{2.6}
\newcommand{\RvTechFineOneMinusReg}{\ensuremath{-}0.6}
\newcommand{\RvTechFineOneMinusRegHi}{\ensuremath{-}0.0}
\newcommand{\RvTechFineOneMinusRegLo}{\ensuremath{-}1.2}
\newcommand{\RvTechFineOnePo}{70}
\newcommand{\RvTechFineOneReg}{74}
\newcommand{\RvTechFineReps}{26,750}
\newcommand{\RvTechFineRfAtlas}{\ensuremath{-}97}
\newcommand{\RvTechFineRfMinusPo}{\ensuremath{-}89}
\newcommand{\RvTechFineRfMinusPoHi}{\ensuremath{-}73}
\newcommand{\RvTechFineRfMinusPoLo}{\ensuremath{-}107}
\newcommand{\RvTechFineRfMinusReg}{\ensuremath{-}87}
\newcommand{\RvTechFineRfMinusRegHi}{\ensuremath{-}71}
\newcommand{\RvTechFineRfMinusRegLo}{\ensuremath{-}104}
\newcommand{\RvTechFineRfPo}{\ensuremath{-}8.2}
\newcommand{\RvTechFineRfReg}{\ensuremath{-}10}
\newcommand{\RvTechFineTwoAtlas}{15}
\newcommand{\RvTechFineTwoMinusPo}{5.6}
\newcommand{\RvTechFineTwoMinusPoHi}{6.1}
\newcommand{\RvTechFineTwoMinusPoLo}{5.1}
\newcommand{\RvTechFineTwoMinusReg}{\ensuremath{-}1.0}
\newcommand{\RvTechFineTwoMinusRegHi}{\ensuremath{-}0.7}
\newcommand{\RvTechFineTwoMinusRegLo}{\ensuremath{-}1.2}
\newcommand{\RvTechFineTwoPo}{9.6}
\newcommand{\RvTechFineTwoReg}{16}
\newcommand{\RvTechFineUnits}{170}
\newcommand{\RvTechOneAtlas}{86}
\newcommand{\RvTechOneMinusPo}{10}
\newcommand{\RvTechOneMinusPoHi}{11}
\newcommand{\RvTechOneMinusPoLo}{9.6}
\newcommand{\RvTechOneMinusReg}{0.1}
\newcommand{\RvTechOneMinusRegHi}{0.5}
\newcommand{\RvTechOneMinusRegLo}{\ensuremath{-}0.3}
\newcommand{\RvTechOnePo}{75}
\newcommand{\RvTechOneReg}{86}
\newcommand{\RvTechReps}{98,739}
\newcommand{\RvTechRfAtlas}{17}
\newcommand{\RvTechRfMinusPo}{13}
\newcommand{\RvTechRfMinusPoHi}{16}
\newcommand{\RvTechRfMinusPoLo}{11}
\newcommand{\RvTechRfMinusReg}{\ensuremath{-}3.9}
\newcommand{\RvTechRfMinusRegHi}{\ensuremath{-}2.0}
\newcommand{\RvTechRfMinusRegLo}{\ensuremath{-}5.8}
\newcommand{\RvTechRfPo}{3.2}
\newcommand{\RvTechRfReg}{20}
\newcommand{\RvTechTwoAtlas}{44}
\newcommand{\RvTechTwoMinusPo}{18}
\newcommand{\RvTechTwoMinusPoHi}{19}
\newcommand{\RvTechTwoMinusPoLo}{17}
\newcommand{\RvTechTwoMinusReg}{\ensuremath{-}1.7}
\newcommand{\RvTechTwoMinusRegHi}{\ensuremath{-}1.1}
\newcommand{\RvTechTwoMinusRegLo}{\ensuremath{-}2.3}
\newcommand{\RvTechTwoPo}{26}
\newcommand{\RvTechTwoReg}{46}
\newcommand{\RvTechUnits}{170}
\newcommand{\RvTruthMinusPoMax}{18}
\newcommand{\RvTruthMinusPoMin}{2.8}
\newcommand{\RvValMeanDiff}{$10^{-7}$}
\newcommand{\RvValMsqDiff}{machine precision}
\newcommand{\SyExt}{4.9}
\newcommand{\SyGeo}{3.1}
\newcommand{\SyGeoHi}{3.2}
\newcommand{\SyGeoLo}{2.4}
\newcommand{\SyMax}{6.4}
\newcommand{\SyMin}{1.4}
\newcommand{\SyNInf}{nine}
\newcommand{\SyNLb}{six}
\newcommand{\SyNTest}{ten}
\newcommand{\VlAtlasOverOther}{2.10}
\newcommand{\VlAtlasOverOtherHi}{2.26}
\newcommand{\VlAtlasOverOtherLo}{1.97}
\newcommand{\VlAtlasOverOwn}{1.34}
\newcommand{\VlAtlasOverOwnHi}{1.56}
\newcommand{\VlAtlasOverOwnLo}{1.15}
\newcommand{\VlCalConMax}{1.00}
\newcommand{\VlCalConMin}{1.00}
\newcommand{\VlCalDfMax}{0.95}
\newcommand{\VlCalDfMin}{0.67}
\newcommand{\VlCellsAtlas}{147}
\newcommand{\VlCellsDf}{178}
\newcommand{\VlCellsOther}{70}
\newcommand{\VlCellsOwn}{110}
\newcommand{\VlCellsOwnPool}{114}
\newcommand{\VlCellsPaired}{800}
\newcommand{\VlFrozenAt}{21:54 EDT on 28 September 2026}
\newcommand{\VlHashedAt}{22:16 EDT}
\newcommand{\VlLevel}{57}
\newcommand{\VlNAntibodies}{111}
\newcommand{\VlNDonors}{18}
\newcommand{\VlNFolds}{34}
\newcommand{\VlNPools}{seven}
\newcommand{\VlPoolMax}{8.2}
\newcommand{\VlPoolMin}{4.5}
\newcommand{\VlPopsFifty}{10}
\newcommand{\VlPopsHundred}{10}
\newcommand{\VlPopsTwentyFive}{9}
\newcommand{\VlRfAtlasEightHundred}{48}
\newcommand{\VlRfAtlasEightHundredHi}{52}
\newcommand{\VlRfAtlasEightHundredLo}{44}
\newcommand{\VlRfAtlasFifty}{11}
\newcommand{\VlRfAtlasFiftyHi}{14}
\newcommand{\VlRfAtlasFiftyLo}{8.8}
\newcommand{\VlRfAtlasHundred}{23}
\newcommand{\VlRfAtlasHundredHi}{25}
\newcommand{\VlRfAtlasHundredLo}{20}
\newcommand{\VlRfAtlasTwentyFive}{0.5}
\newcommand{\VlRfAtlasTwentyFiveHi}{2.3}
\newcommand{\VlRfAtlasTwentyFiveLo}{\ensuremath{-}2.0}
\newcommand{\VlRfDfEightHundred}{48}
\newcommand{\VlRfDfEightHundredHi}{52}
\newcommand{\VlRfDfEightHundredLo}{44}
\newcommand{\VlRfDfFifty}{\ensuremath{-}60}
\newcommand{\VlRfDfFiftyHi}{\ensuremath{-}52}
\newcommand{\VlRfDfFiftyLo}{\ensuremath{-}69}
\newcommand{\VlRfDfHundred}{13}
\newcommand{\VlRfDfHundredHi}{16}
\newcommand{\VlRfDfHundredLo}{9.3}
\newcommand{\VlRfDfTwentyFive}{\ensuremath{-}408}
\newcommand{\VlRfDfTwentyFiveHi}{\ensuremath{-}376}
\newcommand{\VlRfDfTwentyFiveLo}{\ensuremath{-}447}
\newcommand{\VlRfOtherEightHundred}{57}
\newcommand{\VlRfOtherEightHundredHi}{60}
\newcommand{\VlRfOtherEightHundredLo}{53}
\newcommand{\VlRfOtherFifty}{23}
\newcommand{\VlRfOtherFiftyHi}{26}
\newcommand{\VlRfOtherFiftyLo}{20}
\newcommand{\VlRfOtherHundred}{34}
\newcommand{\VlRfOtherHundredHi}{37}
\newcommand{\VlRfOtherHundredLo}{31}
\newcommand{\VlRfOtherTwentyFive}{8.8}
\newcommand{\VlRfOtherTwentyFiveHi}{12}
\newcommand{\VlRfOtherTwentyFiveLo}{5.6}
\newcommand{\VlRfOwnEightHundred}{45}
\newcommand{\VlRfOwnEightHundredHi}{49}
\newcommand{\VlRfOwnEightHundredLo}{41}
\newcommand{\VlRfOwnFifty}{19}
\newcommand{\VlRfOwnFiftyHi}{23}
\newcommand{\VlRfOwnFiftyLo}{15}
\newcommand{\VlRfOwnHundred}{27}
\newcommand{\VlRfOwnHundredHi}{31}
\newcommand{\VlRfOwnHundredLo}{24}
\newcommand{\VlRfOwnPoolEightHundred}{44}
\newcommand{\VlRfOwnPoolEightHundredHi}{48}
\newcommand{\VlRfOwnPoolEightHundredLo}{40}
\newcommand{\VlRfOwnPoolFifty}{19}
\newcommand{\VlRfOwnPoolFiftyHi}{23}
\newcommand{\VlRfOwnPoolFiftyLo}{14}
\newcommand{\VlRfOwnPoolHundred}{27}
\newcommand{\VlRfOwnPoolHundredHi}{31}
\newcommand{\VlRfOwnPoolHundredLo}{23}
\newcommand{\VlRfOwnPoolTwentyFive}{8.2}
\newcommand{\VlRfOwnPoolTwentyFiveHi}{13}
\newcommand{\VlRfOwnPoolTwentyFiveLo}{2.4}
\newcommand{\VlRfOwnTwentyFive}{5.6}
\newcommand{\VlRfOwnTwentyFiveHi}{9.1}
\newcommand{\VlRfOwnTwentyFiveLo}{1.5}
\newcommand{\VlRfPairedEightHundred}{13}
\newcommand{\VlRfPairedEightHundredHi}{20}
\newcommand{\VlRfPairedEightHundredLo}{11}
\newcommand{\VlRfPairedFifty}{2.4}
\newcommand{\VlRfPairedFiftyHi}{3.8}
\newcommand{\VlRfPairedFiftyLo}{1.1}
\newcommand{\VlRfPairedHundred}{3.5}
\newcommand{\VlRfPairedHundredHi}{6.2}
\newcommand{\VlRfPairedHundredLo}{2.9}
\newcommand{\VlRfPairedTwentyFive}{0.5}
\newcommand{\VlRfPairedTwentyFiveHi}{0.9}
\newcommand{\VlRfPairedTwentyFiveLo}{0.4}
\newcommand{\VlRowsAtlas}{137}
\newcommand{\VlRowsFifty}{40}
\newcommand{\VlRowsHundred}{90}
\newcommand{\VlRowsMax}{790}
\newcommand{\VlRowsTwentyFive}{16}
\newcommand{\VlSaveAtlas}{5.4}
\newcommand{\VlSaveAtlasHi}{6.0}
\newcommand{\VlSaveAtlasLo}{4.8}
\newcommand{\VlSaveAtlasQuarter}{13.4}
\newcommand{\VlSaveAtlasQuarterHi}{13.7}
\newcommand{\VlSaveAtlasQuarterLo}{4.5}
\newcommand{\VlSaveAtlasRows}{5.8}
\newcommand{\VlSaveAtlasThreeQ}{1.9}
\newcommand{\VlSaveAtlasThreeQHi}{2.1}
\newcommand{\VlSaveAtlasThreeQLo}{1.6}
\newcommand{\VlSaveDf}{4.5}
\newcommand{\VlSaveDfHi}{4.8}
\newcommand{\VlSaveDfLo}{4.1}
\newcommand{\VlSaveOther}{11.4}
\newcommand{\VlSaveOtherHi}{12.8}
\newcommand{\VlSaveOtherLo}{10.3}
\newcommand{\VlSaveOwn}{7.3}
\newcommand{\VlSaveOwnHi}{8.8}
\newcommand{\VlSaveOwnLo}{5.8}
\newcommand{\VlSaveOwnPool}{7.0}
\newcommand{\VlSaveOwnPoolHi}{8.9}
\newcommand{\VlSaveOwnPoolLo}{5.4}
\newcommand{\VlSinglesFifty}{0.9}
\newcommand{\VlSinglesHundred}{0.2}
\newcommand{\VlSinglesTwentyFive}{2.2}
\newcommand{\VlTarget}{28}
\newcommand{\VlVFour}{not met}
\newcommand{\VlVOne}{met}
\newcommand{\VlVThree}{met}
\newcommand{\VlVTwo}{met}

% Generated by extension/synthesis/si_refs.py. Do not edit.
\makeatletter
\expandafter\def\csname siref@prop:bjs\endcsname{S6}
\expandafter\def\csname siref@prop:contrast\endcsname{S12}
\expandafter\def\csname siref@prop:gaussian\endcsname{S4}
\expandafter\def\csname siref@prop:law\endcsname{S10}
\expandafter\def\csname siref@prop:map\endcsname{S7}
\expandafter\def\csname siref@prop:means\endcsname{S2}
\expandafter\def\csname siref@prop:minimax\endcsname{S8}
\expandafter\def\csname siref@prop:mixture\endcsname{S3}
\expandafter\def\csname siref@prop:noise\endcsname{S5}
\expandafter\def\csname siref@prop:oracle\endcsname{S9}
\expandafter\def\csname siref@prop:pairing\endcsname{S1}
\expandafter\def\csname siref@prop:reveal\endcsname{S11}
\expandafter\def\csname siref@sec:cross\endcsname{8}
\expandafter\def\csname siref@sec:deploy\endcsname{3}
\expandafter\def\csname siref@sec:gen\endcsname{6}
\expandafter\def\csname siref@sec:ident\endcsname{1}
\expandafter\def\csname siref@sec:noise\endcsname{2}
\expandafter\def\csname siref@sec:pert\endcsname{9}
\expandafter\def\csname siref@sec:pred\endcsname{7}
\expandafter\def\csname siref@sec:rep\endcsname{10}
\expandafter\def\csname siref@sec:review\endcsname{11}
\expandafter\def\csname siref@sec:semi\endcsname{5}
\expandafter\def\csname siref@sec:val\endcsname{4}
\expandafter\def\csname siref@sfig:entries\endcsname{3}
\expandafter\def\csname siref@sfig:ident\endcsname{1}
\expandafter\def\csname siref@sfig:transfer\endcsname{2}
\expandafter\def\csname siref@tab:allocation\endcsname{7}
\expandafter\def\csname siref@tab:alternatives\endcsname{13}
\expandafter\def\csname siref@tab:audit\endcsname{9}
\expandafter\def\csname siref@tab:axesmaps\endcsname{32}
\expandafter\def\csname siref@tab:calibration\endcsname{11}
\expandafter\def\csname siref@tab:cross\endcsname{24}
\expandafter\def\csname siref@tab:deploy\endcsname{3}
\expandafter\def\csname siref@tab:deploycon\endcsname{5}
\expandafter\def\csname siref@tab:deploysite\endcsname{4}
\expandafter\def\csname siref@tab:donors\endcsname{16}
\expandafter\def\csname siref@tab:draws\endcsname{29}
\expandafter\def\csname siref@tab:extarms\endcsname{19}
\expandafter\def\csname siref@tab:fisher\endcsname{35}
\expandafter\def\csname siref@tab:fixed\endcsname{39}
\expandafter\def\csname siref@tab:fixeddetail\endcsname{40}
\expandafter\def\csname siref@tab:fly\endcsname{6}
\expandafter\def\csname siref@tab:gen\endcsname{20}
\expandafter\def\csname siref@tab:lawunits\endcsname{37}
\expandafter\def\csname siref@tab:mixture\endcsname{25}
\expandafter\def\csname siref@tab:multiplicity\endcsname{12}
\expandafter\def\csname siref@tab:noiseblocks\endcsname{31}
\expandafter\def\csname siref@tab:noisecal\endcsname{30}
\expandafter\def\csname siref@tab:onedraw\endcsname{33}
\expandafter\def\csname siref@tab:pert\endcsname{26}
\expandafter\def\csname siref@tab:pilotdesign\endcsname{38}
\expandafter\def\csname siref@tab:pools\endcsname{15}
\expandafter\def\csname siref@tab:pred\endcsname{22}
\expandafter\def\csname siref@tab:programs\endcsname{27}
\expandafter\def\csname siref@tab:readout\endcsname{10}
\expandafter\def\csname siref@tab:registry\endcsname{1}
\expandafter\def\csname siref@tab:relationships\endcsname{34}
\expandafter\def\csname siref@tab:rep\endcsname{28}
\expandafter\def\csname siref@tab:rigor\endcsname{23}
\expandafter\def\csname siref@tab:robust\endcsname{14}
\expandafter\def\csname siref@tab:savings\endcsname{21}
\expandafter\def\csname siref@tab:semidev\endcsname{17}
\expandafter\def\csname siref@tab:semiext\endcsname{18}
\expandafter\def\csname siref@tab:sim\endcsname{2}
\expandafter\def\csname siref@tab:truths\endcsname{36}
\expandafter\def\csname siref@tab:val\endcsname{8}
\makeatother
\newcommand{\Sref}[1]{\ifcsname siref@#1\endcsname\csname siref@#1\endcsname\else\textbf{??}\typeout{Undefined supplementary label #1}\fi}

\setcitestyle{super,open={},close={}}
\newcolumntype{Y}{>{\raggedright\arraybackslash}X}
\unnumbered

\begin{document}
\raggedbottom
\ifdefined\ReaderVersion
  \newgeometry{left=25mm,right=25mm,top=23mm,bottom=25mm,twocolumn}
\else
  \newgeometry{left=25mm,right=25mm,top=23mm,bottom=25mm}
  \doublespacing
  \linenumbers
\fi
\title[Reducing the number of paired cells with unpaired measurements]{Using unpaired measurements to reduce the number of paired cells needed to estimate dependence between modalities}


\abstract{Measuring two modalities, such as RNA and protein, in the same cells shows how they vary together, but paired assays are costly. Unpaired measurements, which make up most single-cell data, cannot reveal this but provide each modality's principal axes and covariance within cell types. In a prespecified test on an independent blood study, our estimator with \VlCellsAtlas{} paired cells and another study's atlas estimated RNA--protein correlations within broad cell types better than paired-only methods with 800. With 100 paired cells, it predicted the direction of \BoOneAtlasHundred\% of reproducible gene--protein associations, against \BoOnePoHundred\% for paired-only methods (follow-up). Across \SyNInf{} further data sets, paired-only methods needed at least \SyGeo{} times as many paired cells as our estimator to reach data-set-specific targets (descriptive geometric mean). In a prespecified test, the distribution of correlations and noise along the axes predicted their saving in the \PredNEstimableData{} data sets from two studies where it could be measured. Without paired cells, population averages identify correlations only where populations differ.}
\maketitle

\section{Introduction}
A protein's abundance depends only partly on that of its messenger RNA, because translation and protein turnover are regulated separately and differ between cell types and states\cite{liu2016,schwanhausser2011,buccitelli2020}. Assays that measure RNA together with a second modality in the same cell show how the two vary together. Such paired assays exist for surface proteins\cite{stoeckius2017,peterson2017}, chromatin accessibility\cite{chen2019snare} and, in neurons, electrophysiology\cite{cadwell2016}, and they have resolved cell states that neither measurement shows alone\cite{hao2021,mimitou2021}. They are costly, however, and most single-cell data, including the atlases that combine many studies\cite{regev2017,tabulasapiens2022,dominguezconde2022}, measure one modality per cell. When two modalities are measured in separate cells, their individual patterns of variation remain observable, but the correspondence between RNA and protein profiles is lost\cite{vandereyken2023,baysoy2023}.

Computational methods restore this correspondence in two ways. Some learn a mapping from cells measured in both modalities, which we call paired cells, and apply it to cells measured in one\cite{stuart2019,zhou2020,totalvi2021,lotfollahi2022,hao2024}. Others align cells measured separately through shared factors, corresponding features or distances between cells\cite{welch2019,cao2022,demetci2022}. Knowing how each modality varies does not, however, determine how the two vary together. The same two distributions are compatible with many joint distributions\cite{dorazio2006,manski2003}, so any prediction from unpaired data alone rests on an assumption that picks one of them\cite{argelaguet2021}. Methods that combine a small paired data set with large unpaired ones instead learn how the modalities vary together from the paired cells\cite{hao2024,kimura2013,champollion}. Bridge integration has examined how the size and composition of a paired reference affect integration\cite{hao2024}. We address a different endpoint: how many paired cells are needed to estimate correlations within cell populations, and how unpaired covariance from the same or another study changes that requirement.

By dependence we mean the matrix of correlations between the features of two modalities within a population of cells, such as one cell type in one donor. Here we show that unpaired data make a small paired sample far more informative about the dependence, in two ways (Fig.~\ref{fig:overview}). First, the dependence concentrates along each modality's leading principal axes, the directions in which cells vary most, and unpaired cells reveal these axes. A small paired sample can then focus on the few leading axes, which carry most of the dependence. Second, each cell type's RNA covariance in unpaired cells turns an estimate pooled over cell types, which a few paired cells can support, into one for that type. We test this in a blood study that we had not analysed before, with an atlas from another study, and in data sets spanning tissues, species and pairs of modalities. We then show that how many paired cells the axes save can be predicted from where the dependence and the noise lie along them. Finally, we show that without paired cells, population averages identify the dependence only along the directions in which those averages differ.

\section{Results}

\subsection{Unpaired cells show where the dependence lies}
Unpaired cells cannot reveal the dependence (Proposition~\Sref{prop:pairing}), but they show where to look for it: along each modality's principal axes, ordered from the leading axes, which carry the most variation, to many trailing ones (Fig.~\ref{fig:overview}a). In the \PredProfileN{} benchmark data sets with recoverable dependence (described below; exploratory), the dependence concentrated along the leading axes. The five leading RNA axes carried \PredProfileTop\% of its sum of squares, and the trailing axes carried mostly noise (Supplementary Note~\Sref{sec:pred}).

Our estimator uses these axes in three steps (Fig.~\ref{fig:overview}b; Methods). First, it removes each population's average from its paired cells, so that only variation within populations remains. It subtracts the unpaired cells' averages when these come from the same populations and otherwise replaces the paired cells of each population by contrasts, differences between them, so that shifts in the average between sites, batches and studies cancel. Second, it computes the correlations between each gene's RNA and each protein in the paired cells and rotates them onto the principal axes of the unpaired cells. It groups the rotated correlations into blocks, from the leading axes to the trailing ones, and multiplies each block by a factor between zero and one, the James--Stein factor\cite{james1961,baranchik1970}, estimated from the paired cells (block shrinkage). A block whose entries are large compared with their noise is left nearly unchanged, and a block dominated by noise is shrunk to or near zero. Third, because a small paired sample cannot reliably estimate dependence separately within every cell type, a cell-type map turns this pooled estimate into one for each cell type, using that type's RNA covariance in the unpaired cells. The map assumes that protein depends on RNA through the same linear relation, which we call a channel, in every cell type; under that assumption it follows from an exact relation between covariances (Proposition~\Sref{prop:map}). For estimation, the map is regularized (Methods). The block partition, shrinkage rule and map construction were fixed during development; their inputs are estimated from the paired and unpaired cells without using held-out cells.

Block shrinkage along the unpaired axes concentrates estimation on blocks that carry dependence. With the best fixed multiplier for each block (oracle-linear shrinkage), a block without dependence adds no error however many genes and proteins it spans. The oracle budget is therefore set by the noise in the signal-bearing blocks. Proposition~\Sref{prop:oracle} bounds the difference between the oracle and implemented estimators' root mean-squared errors.

We scored each estimate in held-out cells, which were never used for estimation, by the fraction of the dependence it recovered. Like the fraction of variance that a regression explains, this is 100\% for exact recovery, 0\% for predicting zero correlation and negative for estimates worse than that (Methods). An estimate's budget is the number of paired cells it uses. The primary comparator in the empirical estimator tests is paired-only estimation: at each budget, the best of four established methods that use only paired cells. The saving is the number of paired cells that paired-only estimation needs to reach a given accuracy divided by the number the estimator needs. We label analyses prespecified (protocol registered before any held-out result was computed), follow-up (registered after a test was scored, before any of its own quantities was computed) or exploratory (Methods).

\begin{figure*}[tp]
\centering
\resizebox{\textwidth}{!}{% Conceptual overview (Fig. 1). Drawn by hand; contains no data. Values in b and c are illustrative.
\definecolor{rnacol}{HTML}{217D91}%
\definecolor{protcol}{HTML}{A64442}%
\definecolor{propcol}{HTML}{217D91}%
\definecolor{inkgray}{HTML}{5D6670}%
\begin{tikzpicture}[font=\sffamily\scriptsize,>={Stealth[length=4pt]},
  lab/.style={font=\sffamily\small\bfseries,anchor=base west},
  head/.style={font=\sffamily\small,anchor=base west},
  sub/.style={font=\sffamily\scriptsize,anchor=base},
  note/.style={font=\sffamily\scriptsize,text=inkgray,align=center},
  axl/.style={font=\sffamily\scriptsize,text=black},
  ax/.style={->,gray!70,line width=.5pt}]
%% ---------------- a: paired and unpaired cells
\begin{scope}[shift={(0,0)}]
\node[lab] at (-.2,4.6) {a};
\node[head] at (.2,4.6) {Few paired cells, many unpaired cells};
\begin{scope}[shift={(.45,1.05)}]
\node[sub] at (1.1,2.75) {paired: same cells};
\draw[ax] (0,0) -- (2.3,0) node[axl,below left] {RNA \textit{x}};
\draw[ax] (0,0) -- (0,2.4) node[axl,right,xshift=1pt,yshift=-3pt] {protein \textit{y}};
\foreach \u/\v in {.45/.52,.62/.78,.92/.95,1.12/1.20,1.35/1.38,1.55/1.60,1.78/1.85,.70/.55,1.30/1.05,1.62/1.48}
  \fill[inkgray!85] (\u,\v) circle (1.4pt);
\node[note] at (1.15,-.62) {dependence observed};
\end{scope}
\begin{scope}[shift={(4.2,1.05)}]
\node[sub] at (1.1,2.75) {unpaired: different cells};
\fill[rnacol!45] plot[smooth,domain=0.1:2.2,samples=40] (\x,{.02+.5*exp(-((\x-1.1)^2)/0.28)}) -- (2.2,.02) -- (0.1,.02) -- cycle;
\fill[protcol!45] plot[smooth,domain=0.1:2.2,samples=40] ({.02+.5*exp(-((\x-1.15)^2)/0.3)},\x) -- (.02,2.2) -- (.02,0.1) -- cycle;
\draw[ax] (0,0) -- (2.3,0) node[axl,below left] {RNA \textit{x}};
\draw[ax] (0,0) -- (0,2.4) node[axl,right,xshift=1pt,yshift=-3pt] {protein \textit{y}};
\node[font=\sffamily\large,text=inkgray] at (1.3,1.3) {?};
\node[note] at (1.15,-.62) {how each modality varies};
\end{scope}
\end{scope}
%% ---------------- b: the three steps of the estimator
\begin{scope}[shift={(8.55,0)}]
\node[lab] at (-.2,4.6) {b};
\node[head] at (.2,4.6) {Estimating the dependence from a few paired cells};
\node[font=\sffamily\scriptsize\bfseries,anchor=base west] at (.1,4.12) {1\enspace contrasts};
\node[font=\sffamily\scriptsize\bfseries,anchor=base west] at (2.1,4.12) {2\enspace block shrinkage};
\node[font=\sffamily\scriptsize\bfseries,anchor=base west] at (6.25,4.12) {3\enspace cell-type maps};
% 1: two populations, each with its own average removed
\begin{scope}[shift={(.15,1.0)}]
\foreach \u/\v in {.30/2.35,.55/2.55,.42/2.15,.68/2.30,.20/2.60}
  \fill[rnacol!80] (\u,\v) circle (1.3pt);
\foreach \u/\v in {1.35/2.75,1.60/2.95,1.48/2.55,1.75/2.70,1.25/3.00}
  \fill[protcol!80] (\u,\v) circle (1.3pt);
\draw[->,black!70,line width=.6pt] (.95,2.05) -- (.95,1.35);
\draw[ax] (.25,.75) -- (1.65,.75);
\draw[ax] (.95,.1) -- (.95,1.3);
\foreach \u/\v in {.72/.95,.98/.62,.83/.42,1.10/.88,.66/.70}
  \fill[rnacol!80] (\u,\v) circle (1.3pt);
\foreach \u/\v in {1.05/1.05,.80/.55,1.22/.66,.92/.98,1.15/.40}
  \fill[protcol!80] (\u,\v) circle (1.3pt);
\node[note,anchor=north] at (.95,.05) {each population's\\average removed};
\end{scope}
% 2: block shrinkage of the paired average in the reference's axes
\begin{scope}[shift={(2.1,-.12)}]
\pgfmathsetseed{11}
\foreach \j in {0,...,31}{
  \foreach \i in {0,...,17}{
    \pgfmathtruncatemacro{\rb}{\j<5 ? 0 : (\j<12 ? 1 : (\j<21 ? 2 : 3))}
    \pgfmathtruncatemacro{\cb}{\i<6 ? 0 : 1}
    \pgfmathtruncatemacro{\k}{2*\rb+\cb}
    \pgfmathsetmacro{\s}{{72,18,24,6,0,0,0,0}[\k]}
    \pgfmathsetmacro{\f}{{.92,.55,.5,.12,0,0,0,0}[\k]}
    \pgfmathtruncatemacro{\vl}{min(100,\s+30*rnd)}
    \pgfmathtruncatemacro{\vr}{\f*\vl}
    \fill[propcol!\vl] ({.35+.065*\i},{3.72-.09*\j}) rectangle ++(.065,-.09);
    \fill[propcol!\vr] ({2.75+.065*\i},{3.72-.09*\j}) rectangle ++(.065,-.09);
  }
}
\foreach \x in {.35,2.75}{
  \draw[white,line width=.8pt] (\x+.39,3.72) -- (\x+.39,.84);
  \foreach \y in {3.27,2.64,1.83}
    \draw[white,line width=.8pt] (\x,\y) -- (\x+1.17,\y);
  \draw[gray!60,line width=.4pt] (\x,.84) rectangle (\x+1.17,3.72);
}
\node[axl,rotate=90,anchor=south] at (.3,2.28) {RNA axes};
\node[axl,anchor=south] at (.94,3.74) {protein axes};
\draw[->,black!70,line width=.7pt] (1.62,2.28) -- (2.68,2.28);
\node[font=\sffamily\scriptsize,anchor=south,align=center] at (2.13,2.36) {shrink\\each\\block};
\node[note,anchor=north] at (.94,.78) {paired\\average $m$};
\node[note,anchor=north] at (3.34,.78) {pooled\\estimate};
\end{scope}
% 3: cell-type maps from the reference's covariance of each cell type
\begin{scope}[shift={(6.25,-.12)}]
\foreach \t/\y/\a in {B/3.05/.8,T/2.05/1.2,NK/1.05/.6}{
  \pgfmathsetseed{11}
  \foreach \j in {0,...,9}{
    \foreach \i in {0,...,8}{
      \pgfmathtruncatemacro{\v}{min(100,max(0,(\j<3 ? 70 : 12)*\a+25*rnd))}
      \fill[propcol!\v] ({.45+.075*\i},{\y+.6-.075*\j}) rectangle ++(.075,-.075);
    }
  }
  \draw[gray!60,line width=.4pt] (.45,\y-.15) rectangle (1.125,\y+.6);
  \node[font=\sffamily\scriptsize,anchor=west] at (1.17,\y+.22) {\t};
  \draw[->,black!55,line width=.5pt] (-.25,2.28) -- (.4,\y+.22);
}
\node[note,anchor=north,align=center] at (1.0,.78) {per type, via its\\covariance};
\end{scope}
\end{scope}
%% ---------------- c: where dependence lies relative to noise
\begin{scope}[shift={(0,-5.35)}]
\node[lab] at (-.2,4.6) {c};
\node[head] at (.2,4.6) {Where dependence lies relative to noise sets the saving};
\foreach \x/\ttl/\ws/\ps/\res in {.45/unpaired axes/{.80,.12,.03,.02,.01,.01,.005,.005}/{.33,.23,.08,.05,.10,.06,.09,.06}/1.8-fold saving,
    4.45/random axes/{.011,.039,.032,.118,.086,.314,.086,.314}/{.011,.039,.032,.118,.086,.314,.086,.314}/no saving}{
  \begin{scope}[shift={(\x,1.05)}]
  \node[sub] at (1.7,2.3) {\ttl};
  \draw[gray!60,line width=.5pt] (0,0) -- (3.4,0);
  \foreach \g in {0,...,7}{
    \pgfmathsetmacro{\hw}{2.6*{\ws}[\g]}
    \pgfmathsetmacro{\hp}{2.6*{\ps}[\g]}
    \fill[propcol] ({.12+.41*\g},0) rectangle ++(.16,\hw);
    \fill[inkgray!55] ({.29+.41*\g},0) rectangle ++(.16,\hp);
  }
  \node[note] at (1.7,-.28) {blocks};
  \node[font=\sffamily\scriptsize,anchor=base] at (1.7,-.72) {\res};
  \end{scope}
}
\fill[propcol] (1.0,3.98) rectangle ++(.16,.16);
\node[font=\sffamily\scriptsize,anchor=west] at (1.2,4.06) {dependence share $w_b$};
\fill[inkgray!55] (5.0,3.98) rectangle ++(.16,.16);
\node[font=\sffamily\scriptsize,anchor=west] at (5.2,4.06) {noise share $\pi_b$};
\node[note,anchor=base] at (4.35,-.15) {saving $=1$ if $w=\pi$;\enspace$\to1+\chi^2(w\,\|\,\pi)$ at low accuracy;\enspace$\le1/\min_b\pi_b$};
\end{scope}
%% ---------------- d: population differences determine a shared channel
\begin{scope}[shift={(8.55,-5.35)}]
\node[lab] at (-.2,4.6) {d};
\node[head] at (.2,4.6) {Without paired cells: population averages};
\begin{scope}[shift={(.55,.95)}]
\draw[ax] (0,0) -- (3.0,0) node[axl,below left] {RNA direction 1};
\draw[ax] (0,0) -- (0,3.0) node[axl,right,xshift=1pt,yshift=-3pt] {RNA direction 2};
\draw[rnacol,line width=1.1pt,->] (.35,.55) -- (2.75,2.35) node[right,font=\sffamily\scriptsize] {\textit{M}};
\draw[gray!55,dashed,line width=.8pt,->] (1.55,1.45) -- (.85,2.4) node[left,font=\sffamily\scriptsize,text=inkgray] {\textit{M}$^{\perp}$};
\foreach \u in {.55,.85,1.1,1.4,1.65,1.95,2.25,2.5}
  \fill[inkgray] (\u,{.55+(\u-.35)*0.75+0.06*sin(900*\u)}) circle (1.6pt);
\node[note,anchor=north] at (1.5,-.41) {RNA averages of populations};
\end{scope}
\begin{scope}[shift={(4.55,.95)}]
\draw[ax] (0,0) -- (3.0,0) node[axl,below left] {RNA along \textit{M}};
\draw[ax] (0,0) -- (0,3.0) node[axl,right,xshift=1pt,yshift=-3pt] {protein average};
\draw[protcol,line width=1.1pt] (.2,.35) -- (2.8,2.55);
\foreach \u in {.45,.75,1.0,1.3,1.55,1.85,2.15,2.4}
  \fill[inkgray] (\u,{.35+(\u-.2)*0.846+0.07*cos(700*\u)}) circle (1.6pt);
\node[font=\sffamily\scriptsize,text=protcol,anchor=east] at (2.95,1.05) {channel \textit{W}};
\node[note,anchor=north] at (1.5,-.41) {determined on \textit{M},\\open on \textit{M}$^{\perp}$};
\end{scope}
\end{scope}
\end{tikzpicture}%
}
\caption{\textbf{Estimating the dependence between modalities from a few paired and many unpaired cells.}
\textbf{a}, Paired assays measure RNA ($x$) and protein ($y$) in the same cells and so observe how they vary together; separate assays observe how each modality varies, but not how the two vary together.
\textbf{b}, The estimator. 1, The paired cells of each population are replaced by contrasts, differences between them that remove the population's average. 2, Their average cross-product $m$ is rotated onto the principal axes of unpaired cells, and each block of entries is multiplied by the James--Stein factor $(1-\hat t_b/\|m_b\|^2)_+$, where $\|m_b\|^2$ is the sum of the block's squared entries, $\hat t_b$ is its noise, estimated from the paired cells, and $(\cdot)_+$ replaces a negative value by zero: blocks along the leading axes, where the dependence concentrates, are kept, and blocks dominated by noise shrink towards zero. 3, Cell-type maps turn the pooled estimate into one for each cell type (here B, T and natural killer (NK) cells), using that type's covariance in the unpaired cells.
\textbf{c}, With oracle-linear shrinkage in each block, the saving from the unpaired axes is set by where the dependence lies (dependence shares $w_b$) relative to where the per-cell noise lies (noise shares $\pi_b$) across the blocks (Proposition~\Sref{prop:law}): it is one when $w=\pi$, approximately so for random axes, at most $1/\min_b\pi_b$, and tends to $1+\chi^2(w\,\|\,\pi)$, with $\chi^2$ the chi-squared divergence, at low accuracy.
\textbf{d}, If protein depends on RNA through the same linear relation $W$, a channel, in several populations, their average RNA and protein determine $W$ along the directions $M$ spanned by their average RNA profiles and leave it undetermined (open) along the perpendicular directions $M^\perp$.
Values in \textbf{b} and \textbf{c} are illustrative.}
\label{fig:overview}
\end{figure*}

\subsection{An atlas from another study cuts the budget of paired cells}
A laboratory planning a paired study usually has unpaired profiles at hand, from its own runs, from other studies or from atlases, large collections of cells from many donors. In a prespecified validation test, we assessed the benefit of an external atlas for a study of healthy blood that we had not analysed before (Fig.~\ref{fig:atlas}a,b). The study measured RNA and surface proteins in the same cells by CITE-seq (cellular indexing of transcriptomes and epitopes by sequencing)\cite{babdor2026}. In the primary arm, the unpaired cells came from the CITE-seq atlas of a COVID-19 study of 118 patients, profiled at other sites with another antibody panel\cite{stephenson2021}. RNA and protein covariances were estimated from disjoint sets of atlas cells, with cell types defined by the published annotations. We held out each of \VlNFolds{} samples from \VlNDonors{} donors in turn and scored the dependence within five broad cell types. Paired cells came from two other samples of the held-out sample's processing pool, a set of samples processed together. The held-out donor contributed none of the paired or unpaired cells used for estimation.

With the atlas, the estimator reached the test's accuracy target, half of the best recovery that any method reached at any budget, with \VlCellsAtlas{} paired cells. Paired-only estimation did not reach it within 800 (saving at least \VlSaveAtlas-fold; 95\% CI, \VlSaveAtlasLo--\VlSaveAtlasHi). With 100 paired cells, the estimator recovered \VlRfAtlasHundred\% of the dependence, against \VlRfPairedHundred\% for paired-only estimation (Fig.~\ref{fig:atlas}a). The saving held for every donor and processing pool (exploratory) and when each curve came from a single draw of paired cells rather than the average of five (follow-up; Supplementary Tables~\Sref{tab:pools}, \Sref{tab:donors} and~\Sref{tab:onedraw}). With unpaired cells from the study itself, measured with the same assay, the estimator needed fewer paired cells still: \VlCellsOther{} with cells of the study's other pools and donors and \VlCellsOwn{} with the remaining cells of the paired samples themselves (Fig.~\ref{fig:atlas}b).

Other methods given the same paired cells and atlas also gained from it (follow-up comparison; Fig.~\ref{fig:atlas}b; Supplementary Table~\Sref{tab:alternatives}). Reference regression, a ridge regression that predicts each protein from RNA in the paired cells and then converts the coefficients into RNA--protein correlations using the atlas's RNA correlations, reached the target with \AltCellsRidge{} paired cells. Champollion\cite{champollion}, which learns from paired cells how to match unpaired ones, and semi-supervised canonical correlation analysis (SemiCCA)\cite{kimura2013} needed more. Most of the saving thus came from the atlas's correlations, whichever method used them.

The gain extended to individual gene--protein associations. In a follow-up readout, we split the held-out cells of each sample and cell type into halves. We called a gene--protein pair a reproducible association if its correlation had the same sign in both halves, with $P<0.01$ in each, and \BoRepShare\% of pairs qualified. With 100 paired cells, the estimator and reference regression, each given the atlas, predicted the direction of \BoOneAtlasHundred\% of these associations correctly, against \BoOnePoHundred\% for paired-only estimation and 50\% for guessing (Fig.~\ref{fig:atlas}c). Of the estimator's 100 strongest pairs in each sample and cell type, \BoTwoAtlasHundred\% were reproducible associations with the predicted sign, against \BoTwoPoHundred\% for paired-only estimation (Fig.~\ref{fig:atlas}d). All eight planned comparisons with paired-only estimation favoured the estimator with the atlas at budgets of 50 to 150 paired cells after correction for multiple testing, and it predicted more directions correctly than paired-only estimation in \RvDonorBetterPo{} of \RvDonorN{} donors (Fig.~\ref{fig:atlas}e). The gains held at other significance thresholds and after technical covariates were removed, and the association test gave the expected rate of false positives when proteins were permuted among cells (Supplementary Tables~\Sref{tab:readout}, \Sref{tab:calibration}, \Sref{tab:multiplicity}, \Sref{tab:fisher} and~\Sref{tab:truths}).

To show which biological relationships become recoverable, we chose eight relationships between a gene program (a set of genes that act together) and a protein within a cell type from prior knowledge, before scoring them (follow-up; Fig.~\ref{fig:atlas}f; Supplementary Table~\Sref{tab:relationships}). With 100 paired cells, the estimator's recovery exceeded that of reference regression for \RvRelBetterReg{} of the 8 relationships and that of paired-only estimation for \RvRelBetterPo{} of the 8. For the association between MHC class II genes and HLA-DR protein in CD14 monocytes, it recovered \RvRelTwoAtlasRf\%, against \RvRelTwoRegRf\% for reference regression and \RvRelTwoPoRf\% for paired-only estimation.

\begin{figure*}[tp]
\centering
% Generated by extension/synthesis/make_assets.py. Do not edit.
\definecolor{propcol}{HTML}{217D91}
\definecolor{pairedcol}{HTML}{5D6670}
\definecolor{semicol}{HTML}{C9803A}
\definecolor{champcol}{HTML}{8E4B8C}
\definecolor{ridgecol}{HTML}{9AA3AB}
\definecolor{meancol}{HTML}{A64442}
\definecolor{sitecol}{HTML}{6A8FB3}
\begin{tikzpicture}[font=\sffamily]
\begin{axis}[name=a,scale only axis,width=.3\textwidth,height=3.2cm,xmode=log,log basis x=2,
 xmin=20,xmax=1000,ymin=-20,ymax=60,xtick={25,50,100,200,400,800},xticklabels={25,50,100,200,400,800},ytick={-20,0,20,40,60},
 xlabel={Paired cells},ylabel={Dependence recovered (\%)},
 title={\textbf{a}\enspace Blood study, atlas of another study},
 axis x line*=bottom,axis y line*=left,tick align=outside,ymajorgrids=true,grid style={gray!18},axis line style={gray!55},tick label style={font=\sffamily\scriptsize},label style={font=\sffamily\small},clip=false,title style={font=\sffamily\small,at={(0,1)},anchor=base west,yshift=5pt}]
\draw[dashed,black!45] (axis cs:20,28.34) -- (axis cs:1000,28.34);
\addplot[color=black,line width=.9pt,dashed,mark=*,mark size=1.3pt] coordinates { (25,8.22) (50,18.67) (100,27.05) (200,33.65) (400,39.65) (800,43.83) };
\addplot[color=sitecol,line width=.9pt,mark=diamond*,mark size=1.7pt] coordinates { (25,5.56) (50,19.06) (100,27.46) (200,34.07) (400,40.31) (800,45.08) };
\addplot[color=champcol,line width=.9pt,mark=triangle*,mark size=1.7pt] coordinates { (25,8.83) (50,23.04) (100,33.86) (200,42.57) (400,50.44) (800,56.67) };
\addplot[color=propcol,line width=1.2pt,mark=*,mark size=1.6pt] coordinates { (25,0.46) (50,11.22) (100,22.52) (200,32.92) (400,41.95) (800,48.15) };
\addplot[color=propcol,line width=.9pt,densely dotted,mark=o,mark size=1.5pt,mark options={solid}] coordinates { (25,-20.00) (50,-20.00) (100,12.78) (200,31.44) (400,41.69) (800,48.09) };
\addplot[color=pairedcol,line width=.8pt,mark=square*,mark size=1.3pt] coordinates { (25,0.52) (50,2.41) (100,3.52) (200,7.72) (400,13.93) (800,12.55) };
\end{axis}
\begin{axis}[at={(a.outer north east)},anchor=outer north west,xshift=.4cm,name=b,scale only axis,width=.2\textwidth,height=3.2cm,xmode=log,log basis x=2,
 xmin=20,xmax=1600,ymin=.4,ymax=9.6,xtick={25,100,400,1600},xticklabels={25,100,400,$>$800},
 ytick={1,2,3,4,5,6,7,8,9},
 yticklabels={{Paired cells only},{Champollion, atlas (1 draw)},{Reference regression, atlas},{SemiCCA, atlas},{Atlas, centred as in bone marrow},{Atlas of another study},{Other pools and donors},{Own samples},{Own samples, their averages}},
 xlabel={Paired cells to the target},
 title={\textbf{b}\enspace Paired cells needed},
 axis x line*=bottom,axis y line*=left,tick align=outside,xmajorgrids=true,grid style={gray!18},axis line style={gray!55},tick label style={font=\sffamily\scriptsize},label style={font=\sffamily\small},clip=false,title style={font=\sffamily\small,at={(0,1)},anchor=base west,yshift=5pt},y tick style={draw=none}]
\draw[black!25] (axis cs:800,.4) -- (axis cs:800,9.6);
\draw[->,color=pairedcol,line width=.8pt] (axis cs:800,1) -- (axis cs:1400,1);
\addplot[only marks,mark=o,mark size=1.8pt,color=pairedcol] coordinates { (800,1) };
\draw[color=meancol,line width=2.4pt] (axis cs:20,2) -- (axis cs:180.0,2);
\node[font=\sffamily\scriptsize,anchor=west] at (axis cs:194.4,2) {180};
\draw[color=ridgecol,line width=2.4pt] (axis cs:20,3) -- (axis cs:168.9,3);
\node[font=\sffamily\scriptsize,anchor=west] at (axis cs:182.4,3) {169};
\draw[color=semicol,line width=2.4pt] (axis cs:20,4) -- (axis cs:495.7,4);
\node[font=\sffamily\scriptsize,anchor=west] at (axis cs:535.4,4) {496};
\draw[color=propcol,line width=1.4pt,densely dotted] (axis cs:20,5) -- (axis cs:178.2,5);
\node[font=\sffamily\scriptsize,anchor=west] at (axis cs:192.5,5) {178};
\draw[color=propcol,line width=2.4pt] (axis cs:20,6) -- (axis cs:147.3,6);
\node[font=\sffamily\scriptsize,anchor=west] at (axis cs:159.1,6) {147};
\draw[color=champcol,line width=2.4pt] (axis cs:20,7) -- (axis cs:70.2,7);
\node[font=\sffamily\scriptsize,anchor=west] at (axis cs:75.8,7) {70};
\draw[color=sitecol,line width=2.4pt] (axis cs:20,8) -- (axis cs:109.6,8);
\node[font=\sffamily\scriptsize,anchor=west] at (axis cs:118.4,8) {110};
\draw[color=black,line width=1.6pt,densely dashed] (axis cs:20,9) -- (axis cs:114.5,9);
\node[font=\sffamily\scriptsize,anchor=west] at (axis cs:123.7,9) {114};
\end{axis}
\path (a.outer south west) -- (b.outer south east) coordinate[pos=.5] (legab);
\begin{axis}[hide axis,scale only axis,width=1pt,height=1pt,xmin=0,xmax=1,ymin=0,ymax=1,at={(legab)},anchor=north,yshift=-2pt,legend style={draw=none,font=\sffamily\scriptsize,at={(0.5,0.5)},anchor=north,legend columns=3,column sep=4pt,cells={anchor=west}}]
\addlegendimage{color=propcol,line width=1.2pt,mark=*,mark size=1.6pt}
\addlegendentry{Atlas of another study}
\addlegendimage{color=propcol,line width=.9pt,densely dotted,mark=o,mark size=1.5pt,mark options={solid}}
\addlegendentry{Atlas, centred as in bone marrow (to \ensuremath{-}408\% at 25)}
\addlegendimage{color=pairedcol,line width=.8pt,mark=square*,mark size=1.3pt}
\addlegendentry{Paired cells only}
\addlegendimage{color=champcol,line width=.9pt,mark=triangle*,mark size=1.7pt}
\addlegendentry{Other pools and donors}
\addlegendimage{color=black,line width=.9pt,dashed,mark=*,mark size=1.3pt}
\addlegendentry{Own samples, their averages}
\addlegendimage{color=sitecol,line width=.9pt,mark=diamond*,mark size=1.7pt}
\addlegendentry{Own samples}
\end{axis}
\begin{axis}[at={(a.outer south west)},anchor=outer north west,yshift=-1.3cm,name=c,scale only axis,width=.3\textwidth,height=3.2cm,xmode=log,log basis x=2,
 xmin=20,xmax=1000,ymin=45,ymax=100,xtick={25,50,100,200,400,800},xticklabels={25,50,100,200,400,800},ytick={50,60,70,80,90,100},
 xlabel={Paired cells},ylabel={Correct direction (\%)},
 title={\textbf{c}\enspace Direction of reproducible associations},
 axis x line*=bottom,axis y line*=left,tick align=outside,ymajorgrids=true,grid style={gray!18},axis line style={gray!55},tick label style={font=\sffamily\scriptsize},label style={font=\sffamily\small},clip=false,title style={font=\sffamily\small,at={(0,1)},anchor=base west,yshift=5pt}]
\draw[dashed,black!45] (axis cs:20,50.000) -- (axis cs:1000,50.000);
\addplot[color=sitecol,line width=.9pt,mark=diamond*,mark size=1.7pt] coordinates { (25,75.83) (50,83.74) (100,87.65) (150,89.36) (200,90.33) (400,92.47) (800,94.07) };
\addplot[color=champcol,line width=.9pt,mark=triangle*,mark size=1.7pt] coordinates { (25,76.63) (50,84.82) (100,89.01) (150,90.99) (200,91.99) (400,94.29) (800,95.83) };
\addplot[color=propcol,line width=1.2pt,mark=*,mark size=1.6pt] coordinates { (25,70.69) (50,80.24) (100,86.07) (150,88.64) (200,89.82) (400,92.22) (800,93.95) };
\addplot[color=pairedcol,line width=.8pt,mark=square*,mark size=1.3pt] coordinates { (25,61.50) (50,67.29) (100,74.85) (150,79.54) (200,82.25) (400,88.12) (800,92.48) };
\addplot[color=ridgecol!70!black,line width=1pt,mark=pentagon*,mark size=1.5pt] coordinates { (25,72.43) (50,80.20) (100,86.22) (150,89.20) (200,90.72) (400,93.54) (800,95.42) };
\end{axis}
\begin{axis}[at={(c.outer north east)},anchor=outer north west,xshift=.4cm,name=d,scale only axis,width=.3\textwidth,height=3.2cm,xmode=log,log basis x=2,
 xmin=20,xmax=1000,ymin=0,ymax=80,xtick={25,50,100,200,400,800},xticklabels={25,50,100,200,400,800},ytick={0,20,40,60,80},
 xlabel={Paired cells},ylabel={Reproducible (\%)},
 title={\textbf{d}\enspace Top 100 pairs per cell type},
 axis x line*=bottom,axis y line*=left,tick align=outside,ymajorgrids=true,grid style={gray!18},axis line style={gray!55},tick label style={font=\sffamily\scriptsize},label style={font=\sffamily\small},clip=false,title style={font=\sffamily\small,at={(0,1)},anchor=base west,yshift=5pt}]
\draw[dashed,black!45] (axis cs:20,1.598) -- (axis cs:1000,1.598);
\addplot[color=sitecol,line width=.9pt,mark=diamond*,mark size=1.7pt] coordinates { (25,26.63) (50,36.13) (100,41.30) (150,44.19) (200,46.12) (400,51.86) (800,56.55) };
\addplot[color=champcol,line width=.9pt,mark=triangle*,mark size=1.7pt] coordinates { (25,28.07) (50,38.18) (100,46.00) (150,50.47) (200,53.63) (400,61.23) (800,67.37) };
\addplot[color=propcol,line width=1.2pt,mark=*,mark size=1.6pt] coordinates { (25,22.66) (50,35.46) (100,47.19) (150,52.26) (200,55.16) (400,61.53) (800,66.52) };
\addplot[color=pairedcol,line width=.8pt,mark=square*,mark size=1.3pt] coordinates { (25,6.91) (50,15.33) (100,27.90) (150,35.47) (200,39.98) (400,51.59) (800,60.91) };
\addplot[color=ridgecol!70!black,line width=1pt,mark=pentagon*,mark size=1.5pt] coordinates { (25,20.98) (50,32.47) (100,49.99) (150,56.52) (200,59.71) (400,66.66) (800,70.61) };
\end{axis}
\path (c.outer south west) -- (d.outer south east) coordinate[pos=.5] (legcd);
\begin{axis}[hide axis,scale only axis,width=1pt,height=1pt,xmin=0,xmax=1,ymin=0,ymax=1,at={(legcd)},anchor=north,yshift=-2pt,legend style={draw=none,font=\sffamily\scriptsize,at={(0.5,0.5)},anchor=north,legend columns=3,column sep=6pt,cells={anchor=west}}]
\addlegendimage{color=sitecol,line width=.9pt,mark=diamond*,mark size=1.7pt}
\addlegendentry{Own samples}
\addlegendimage{color=champcol,line width=.9pt,mark=triangle*,mark size=1.7pt}
\addlegendentry{Other pools and donors}
\addlegendimage{color=propcol,line width=1.2pt,mark=*,mark size=1.6pt}
\addlegendentry{Atlas of another study}
\addlegendimage{color=ridgecol!70!black,line width=1pt,mark=pentagon*,mark size=1.5pt}
\addlegendentry{Reference regression, atlas}
\addlegendimage{color=pairedcol,line width=.8pt,mark=square*,mark size=1.3pt}
\addlegendentry{Paired cells only}
\end{axis}
\begin{axis}[at={(c.outer south west)},anchor=outer north west,yshift=-1.05cm,name=e,scale only axis,width=.16\textwidth,height=3.2cm,
 xmin=-10,xmax=20,ymin=.4,ymax=2.6,ytick={1,2},
 yticklabels={{Regression, atlas},{Paired only}},
 xlabel={Difference (points)},
 title={\textbf{e}\enspace Direction, per donor},
 axis x line*=bottom,axis y line*=left,tick align=outside,xmajorgrids=true,grid style={gray!18},axis line style={gray!55},tick label style={font=\sffamily\scriptsize},label style={font=\sffamily\small},clip=false,title style={font=\sffamily\small,at={(0,1)},anchor=base west,yshift=5pt},y tick style={draw=none}]
\draw[dashed,black!45] (axis cs:0,.4) -- (axis cs:0,2.6);
\addplot[only marks,mark=*,mark size=1.1pt,color=propcol!60] coordinates { (-1.97,0.800) (-1.75,0.824) (-1.29,0.847) (-0.74,0.871) (-0.61,0.894) (-0.59,0.918) (-0.17,0.941) (-0.13,0.965) (0.02,0.988) (0.06,1.012) (0.13,1.035) (0.19,1.059) (0.19,1.082) (0.32,1.106) (0.81,1.129) (1.30,1.153) (1.34,1.176) (1.78,1.200) };
\draw[color=black,line width=1.2pt] (axis cs:-0.06,0.7) -- (axis cs:-0.06,1.3);
\addplot[only marks,mark=*,mark size=1.1pt,color=propcol!60] coordinates { (6.79,1.800) (8.35,1.824) (9.32,1.847) (10.02,1.871) (10.21,1.894) (10.41,1.918) (10.42,1.941) (11.03,1.965) (11.28,1.988) (11.42,2.012) (11.69,2.035) (11.92,2.059) (12.07,2.082) (12.12,2.106) (12.26,2.129) (12.45,2.153) (13.12,2.176) (15.52,2.200) };
\draw[color=black,line width=1.2pt] (axis cs:11.13,1.7) -- (axis cs:11.13,2.3);
\end{axis}
\begin{axis}[at={(e.outer north east)},anchor=outer north west,xshift=.4cm,name=f,scale only axis,width=.2\textwidth,height=3.5cm,
 xmin=-.8,xmax=.8,ymin=.4,ymax=8.6,xtick={-.8,-.4,0,.4,.8},xticklabels={$-$0.8,$-$0.4,0,0.4,0.8},
 ytick={1,2,3,4,5,6,7,8},
 yticklabels={{CD4 T: LEF1, CD45RA},{NK: cytotoxic, CD56},{NK: cytotoxic, CD16},{CD8 T: cytotoxic, CD27},{CD8 T: cytotoxic, CD57},{B: MHC II, HLA-DR},{CD14 Mono: MHC II, HLA-DR},{CD14 Mono: interferon, CD169}},
 xlabel={Within-type correlation},
 title={\textbf{f}\enspace Named relationships},
 axis x line*=bottom,axis y line*=left,tick align=outside,xmajorgrids=true,grid style={gray!18},axis line style={gray!55},tick label style={font=\sffamily\scriptsize},label style={font=\sffamily\small},clip=false,title style={font=\sffamily\small,at={(0,1)},anchor=base west,yshift=5pt},y tick style={draw=none}]
\draw[dashed,black!45] (axis cs:0,.4) -- (axis cs:0,8.6);
\draw[black!55,line width=.7pt] (axis cs:0.229,1.24) -- (axis cs:0.448,1.24);
\addplot[only marks,mark=diamond*,mark size=2.2pt,color=black] coordinates { (0.332,1.24) };
\addplot[only marks,mark=*,mark size=1.8pt,color=propcol] coordinates { (0.114,1.08) };
\addplot[only marks,mark=pentagon*,mark size=1.8pt,color=ridgecol!70!black] coordinates { (0.075,0.92) };
\addplot[only marks,mark=square*,mark size=1.5pt,color=pairedcol] coordinates { (0.057,0.76) };
\draw[black!55,line width=.7pt] (axis cs:-0.257,2.24) -- (axis cs:0.140,2.24);
\addplot[only marks,mark=diamond*,mark size=2.2pt,color=black] coordinates { (-0.051,2.24) };
\addplot[only marks,mark=*,mark size=1.8pt,color=propcol] coordinates { (0.006,2.08) };
\addplot[only marks,mark=pentagon*,mark size=1.8pt,color=ridgecol!70!black] coordinates { (-0.028,1.92) };
\addplot[only marks,mark=square*,mark size=1.5pt,color=pairedcol] coordinates { (0.000,1.76) };
\draw[black!55,line width=.7pt] (axis cs:0.236,3.24) -- (axis cs:0.552,3.24);
\addplot[only marks,mark=diamond*,mark size=2.2pt,color=black] coordinates { (0.383,3.24) };
\addplot[only marks,mark=*,mark size=1.8pt,color=propcol] coordinates { (0.132,3.08) };
\addplot[only marks,mark=pentagon*,mark size=1.8pt,color=ridgecol!70!black] coordinates { (0.092,2.92) };
\addplot[only marks,mark=square*,mark size=1.5pt,color=pairedcol] coordinates { (0.045,2.76) };
\draw[black!55,line width=.7pt] (axis cs:-0.701,4.24) -- (axis cs:-0.248,4.24);
\addplot[only marks,mark=diamond*,mark size=2.2pt,color=black] coordinates { (-0.493,4.24) };
\addplot[only marks,mark=*,mark size=1.8pt,color=propcol] coordinates { (-0.276,4.08) };
\addplot[only marks,mark=pentagon*,mark size=1.8pt,color=ridgecol!70!black] coordinates { (-0.324,3.92) };
\addplot[only marks,mark=square*,mark size=1.5pt,color=pairedcol] coordinates { (-0.220,3.76) };
\draw[black!55,line width=.7pt] (axis cs:0.119,5.24) -- (axis cs:0.697,5.24);
\addplot[only marks,mark=diamond*,mark size=2.2pt,color=black] coordinates { (0.501,5.24) };
\addplot[only marks,mark=*,mark size=1.8pt,color=propcol] coordinates { (0.279,5.08) };
\addplot[only marks,mark=pentagon*,mark size=1.8pt,color=ridgecol!70!black] coordinates { (0.255,4.92) };
\addplot[only marks,mark=square*,mark size=1.5pt,color=pairedcol] coordinates { (0.276,4.76) };
\draw[black!55,line width=.7pt] (axis cs:0.039,6.24) -- (axis cs:0.230,6.24);
\addplot[only marks,mark=diamond*,mark size=2.2pt,color=black] coordinates { (0.149,6.24) };
\addplot[only marks,mark=*,mark size=1.8pt,color=propcol] coordinates { (0.102,6.08) };
\addplot[only marks,mark=pentagon*,mark size=1.8pt,color=ridgecol!70!black] coordinates { (0.052,5.92) };
\addplot[only marks,mark=square*,mark size=1.5pt,color=pairedcol] coordinates { (0.035,5.76) };
\draw[black!55,line width=.7pt] (axis cs:0.386,7.24) -- (axis cs:0.543,7.24);
\addplot[only marks,mark=diamond*,mark size=2.2pt,color=black] coordinates { (0.472,7.24) };
\addplot[only marks,mark=*,mark size=1.8pt,color=propcol] coordinates { (0.308,7.08) };
\addplot[only marks,mark=pentagon*,mark size=1.8pt,color=ridgecol!70!black] coordinates { (0.130,6.92) };
\addplot[only marks,mark=square*,mark size=1.5pt,color=pairedcol] coordinates { (0.051,6.76) };
\draw[black!55,line width=.7pt] (axis cs:0.016,8.24) -- (axis cs:0.125,8.24);
\addplot[only marks,mark=diamond*,mark size=2.2pt,color=black] coordinates { (0.071,8.24) };
\addplot[only marks,mark=*,mark size=1.8pt,color=propcol] coordinates { (0.019,8.08) };
\addplot[only marks,mark=pentagon*,mark size=1.8pt,color=ridgecol!70!black] coordinates { (0.011,7.92) };
\addplot[only marks,mark=square*,mark size=1.5pt,color=pairedcol] coordinates { (0.007,7.76) };
\end{axis}
\begin{axis}[hide axis,scale only axis,width=1pt,height=1pt,xmin=0,xmax=1,ymin=0,ymax=1,at={(f.outer south east)},anchor=north east,yshift=-2pt,legend style={draw=none,font=\sffamily\scriptsize,at={(1,0.5)},anchor=north east,legend columns=4,column sep=5pt,cells={anchor=west}}]
\addlegendimage{only marks,mark=diamond*,mark size=2.2pt,color=black}
\addlegendentry{Held out}
\addlegendimage{only marks,mark=*,mark size=1.8pt,color=propcol}
\addlegendentry{Estimator, atlas}
\addlegendimage{only marks,mark=pentagon*,mark size=1.8pt,color=ridgecol!70!black}
\addlegendentry{Reference regression, atlas}
\addlegendimage{only marks,mark=square*,mark size=1.5pt,color=pairedcol}
\addlegendentry{Paired cells only}
\end{axis}
\end{tikzpicture}

\caption{\textbf{An atlas from another study reduces the number of paired cells a new blood study needs.}
\textbf{a}, Validation test (prespecified): fraction of the dependence within broad cell types recovered in a held-out sample against the number of paired cells, drawn from two other samples of its processing pool and centred by contrasts, pooled over \VlNFolds{} held-out samples. Unpaired cells came from an atlas of another study, the study's other pools and donors, or the paired samples themselves; two further lines centre the paired cells on averages from the unpaired cells of the same samples or, with the atlas, on the drawn paired cells' own averages, as in the bone-marrow test. Dashed line, the accuracy target (half the best recovery); values below $-20\%$ are drawn at $-20\%$.
\textbf{b}, Paired cells needed to reach the target, including three other methods given the same paired cells and atlas (follow-up; Champollion on one draw of five); arrow, not reached within 800.
\textbf{c}, Follow-up readout: fraction of the held-out sample's reproducible associations (same sign in both halves of its cells, with $P<0.01$ in each; \BoRepShare\% of gene--protein pairs) whose direction the estimate predicted correctly; dashed line, guessing.
\textbf{d}, Fraction of the 100 pairs with the largest absolute estimated correlations in each held-out sample and cell type that were reproducible associations with the predicted sign; dashed line, random pairs. In \textbf{c} and \textbf{d}, paired-only estimation and reference regression are each the best of their methods or forms at each budget.
\textbf{e}, Per donor (points), the atlas estimate's share of correct directions with 100 paired cells minus that of paired-only estimation or of reference regression; bars, means.
\textbf{f}, Eight relationships between a gene program and a protein within a cell type (follow-up): the held-out value (diamond, mean over held-out samples; bar, range) and each method's mean prediction with 100 paired cells.}
\label{fig:atlas}
\end{figure*}

\subsection{A reference transfers through its axes and cell-type maps, not its averages}
The atlas contributes principal axes, along which the paired estimate is shrunk, and each cell type's covariance, which the maps use. In a follow-up analysis, we rescored the validation test with the atlas's axes alone, its maps alone or neither (Fig.~\ref{fig:reference}a). With 100 paired cells, the axes alone recovered \RvAxesHundred\% of the dependence within cell types and the maps alone \RvMapOnlyHundred\%, against \RvNeitherHundred\% with neither, whereas together they recovered \RvBothHundred\%. The axes make the pooled estimate accurate from few cells, and the maps make it specific to each cell type.

The maps also carry the atlas's mix of cell types: they use the atlas's covariance pooled over cell types, which weights each cell type by its proportion in the atlas rather than in the paired cells. We rebuilt the atlas with its B cells and CD14 monocytes, or its CD4 and CD8 T cells, taken from only 15 of its 118 patients (exploratory). With the pooled covariance reweighted to the cell-type proportions of the paired cells (follow-up), the estimator reached the target with \RvFewBMonoRwCells{} and \RvFewTRwCells{} paired cells, against \RvBothCells{} for the unchanged atlas on the same draws, or \RvFullRwCells{} with its maps reweighted; the axes of the skewed atlases worked as before (Fig.~\ref{fig:reference}b). An atlas containing every cell type but drawn from only 15 of the 118 patients let the estimator reach the target with \RbCellsPatientsFifteen{} paired cells (exploratory; Fig.~\ref{fig:reference}c).

Correlation structure carried over between sites and between animals, whereas averages did not (Supplementary Fig.~\Sref{sfig:transfer}). The bone-marrow test, a prespecified test on cells profiled at \DpNSites{} sites\cite{luecken2021sandbox}, centred each population's paired cells on their own average. With unpaired cells from other sites' donors, the estimator reached half of the best recovery with \DpCellsOther{} paired cells. Paired-only estimation did not reach it within \DpMaxBudget{} (saving at least \DpSaveOther-fold; 95\% CI, \DpSaveOtherLo--\DpSaveOtherHi), and with contrasts the estimator needed \CtCellsOther{} (exploratory). In the fly, whose connectomes map every neuron and synapse, neurons that connect the brain and the nerve cord have synaptic partners on both sides, which we treated as two modalities. A prespecified secondary analysis of the benchmark kept the paired neurons of one connectome\cite{bates2026} and took the unpaired neurons from two other animals, one reconstructed only in the brain\cite{dorkenwald2024,schlegel2024} and one only in the nerve cord\cite{takemura2024}. With 100 paired neurons centred by contrasts (exploratory), the estimator recovered \FxOwnHundred\% of the dependence, against \FxSameHundred\% with unpaired neurons from the same animal and \FxPairedOwnHundred\% for paired-only estimation.

An exploratory analysis of the bone-marrow test compared unpaired RNA from single nuclei with RNA from whole cells (Fig.~\ref{fig:reference}d). Without maps, the two recovered about the same. With whole-cell maps and 800 paired cells, nuclear axes recovered \MpAssayOtherMap\% of the dependence, against \MpOtherMapped\% with whole-cell axes and maps, so nuclear RNA supplied useful axes.

\begin{figure*}[tp]
\centering
% Generated by extension/synthesis/make_assets.py. Do not edit.
\definecolor{propcol}{HTML}{217D91}
\definecolor{pairedcol}{HTML}{5D6670}
\definecolor{semicol}{HTML}{C9803A}
\definecolor{champcol}{HTML}{8E4B8C}
\definecolor{ridgecol}{HTML}{9AA3AB}
\definecolor{meancol}{HTML}{A64442}
\definecolor{sitecol}{HTML}{6A8FB3}
\begin{tikzpicture}[font=\sffamily]
\begin{axis}[name=a,scale only axis,width=.3\textwidth,height=3.2cm,xmode=log,log basis x=2,
 xmin=20,xmax=1000,ymin=-5,ymax=55,xtick={25,50,100,200,400,800},xticklabels={25,50,100,200,400,800},ytick={0,10,20,30,40,50},
 xlabel={Paired cells},ylabel={Dependence recovered (\%)},
 title={\textbf{a}\enspace Axes and maps, validation test},
 axis x line*=bottom,axis y line*=left,tick align=outside,ymajorgrids=true,grid style={gray!18},axis line style={gray!55},tick label style={font=\sffamily\scriptsize},label style={font=\sffamily\small},clip=false,title style={font=\sffamily\small,at={(0,1)},anchor=base west,yshift=5pt}]
\draw[dashed,black!45] (axis cs:20,28.34) -- (axis cs:1000,28.34);
\addplot[color=propcol,line width=1.2pt,mark=*,mark size=1.6pt] coordinates { (25,0.46) (50,11.22) (100,22.52) (150,28.81) (200,32.92) (400,41.95) (800,48.15) };
\addplot[color=champcol,line width=.9pt,mark=triangle*,mark size=1.7pt] coordinates { (25,-0.16) (50,1.84) (100,4.75) (150,7.12) (200,9.52) (400,17.70) (800,27.95) };
\addplot[color=propcol,line width=.9pt,dashed,mark=o,mark size=1.5pt,mark options={solid}] coordinates { (25,-0.75) (50,4.00) (100,7.78) (150,9.28) (200,10.42) (400,12.66) (800,13.97) };
\addplot[color=pairedcol,line width=.8pt,mark=square*,mark size=1.3pt] coordinates { (25,0.04) (50,1.09) (100,2.39) (150,3.39) (200,4.39) (400,7.05) (800,9.90) };
\end{axis}
\begin{axis}[at={(a.outer north east)},anchor=outer north west,xshift=.4cm,name=b,scale only axis,width=.2\textwidth,height=3.2cm,xmode=log,log basis x=2,
 xmin=20,xmax=1600,ymin=.4,ymax=3.6,xtick={25,100,400,1600},xticklabels={25,100,400,$>$800},
 ytick={1,2,3},
 yticklabels={{CD4, CD8 T from 15},{B, CD14 Mono from 15},{All 118 patients}},
 xlabel={Paired cells to the target},
 title={\textbf{b}\enspace Mix of cell types in the atlas},
 axis x line*=bottom,axis y line*=left,tick align=outside,xmajorgrids=true,grid style={gray!18},axis line style={gray!55},tick label style={font=\sffamily\scriptsize},label style={font=\sffamily\small},clip=false,title style={font=\sffamily\small,at={(0,1)},anchor=base west,yshift=5pt},y tick style={draw=none}]
\draw[black!25] (axis cs:800,.4) -- (axis cs:800,3.6);
\draw[->,color=propcol,line width=.7pt] (axis cs:800,1.22) -- (axis cs:1400,1.22);
\addplot[only marks,mark=o,mark size=1.8pt,color=propcol] coordinates { (800,1.22) };
\draw[color=meancol,line width=.7pt] (axis cs:266.4,1.00) -- (axis cs:377.6,1.00);
\addplot[only marks,mark=triangle*,mark size=2pt,color=meancol] coordinates { (310.5,1.00) };
\draw[->,color=propcol,line width=.7pt] (axis cs:800,0.78) -- (axis cs:1400,0.78);
\addplot[only marks,mark=o,mark size=1.8pt,color=propcol] coordinates { (800,0.78) };
\draw[->,color=propcol,line width=.7pt] (axis cs:800,2.22) -- (axis cs:1400,2.22);
\addplot[only marks,mark=o,mark size=1.8pt,color=propcol] coordinates { (800,2.22) };
\draw[color=meancol,line width=.7pt] (axis cs:122.4,2.00) -- (axis cs:191.6,2.00);
\addplot[only marks,mark=triangle*,mark size=2pt,color=meancol] coordinates { (148.2,2.00) };
\draw[->,color=propcol,line width=.7pt] (axis cs:800,1.78) -- (axis cs:1400,1.78);
\addplot[only marks,mark=o,mark size=1.8pt,color=propcol] coordinates { (800,1.78) };
\draw[color=propcol,line width=.7pt] (axis cs:123.5,3.22) -- (axis cs:183.5,3.22);
\addplot[only marks,mark=*,mark size=1.8pt,color=propcol] coordinates { (145.5,3.22) };
\draw[color=meancol,line width=.7pt] (axis cs:113.5,3.00) -- (axis cs:162.0,3.00);
\addplot[only marks,mark=triangle*,mark size=2pt,color=meancol] coordinates { (131.4,3.00) };
\draw[->,color=propcol,line width=.7pt] (axis cs:800,2.78) -- (axis cs:1400,2.78);
\addplot[only marks,mark=o,mark size=1.8pt,color=propcol] coordinates { (800,2.78) };
\end{axis}
\begin{axis}[hide axis,scale only axis,width=1pt,height=1pt,xmin=0,xmax=1,ymin=0,ymax=1,at={(a.outer south west)},anchor=north west,yshift=-2pt,legend style={draw=none,font=\sffamily\scriptsize,at={(0,0.5)},anchor=north west,legend columns=2,column sep=5pt,cells={anchor=west}}]
\addlegendimage{color=propcol,line width=1.2pt,mark=*,mark size=1.6pt}
\addlegendentry{Axes and maps}
\addlegendimage{color=champcol,line width=.9pt,mark=triangle*,mark size=1.7pt}
\addlegendentry{Maps only}
\addlegendimage{color=propcol,line width=.9pt,dashed,mark=o,mark size=1.5pt,mark options={solid}}
\addlegendentry{Axes only}
\addlegendimage{color=pairedcol,line width=.8pt,mark=square*,mark size=1.3pt}
\addlegendentry{Neither}
\end{axis}
\begin{axis}[hide axis,scale only axis,width=1pt,height=1pt,xmin=0,xmax=1,ymin=0,ymax=1,at={(b.outer south east)},anchor=north east,yshift=-2pt,legend style={draw=none,font=\sffamily\scriptsize,at={(1,0.5)},anchor=north east,legend columns=2,column sep=5pt,cells={anchor=west}}]
\addlegendimage{only marks,mark=*,mark size=1.8pt,color=propcol}
\addlegendentry{Unweighted maps}
\addlegendimage{only marks,mark=triangle*,mark size=2pt,color=meancol}
\addlegendentry{Reweighted maps}
\addlegendimage{only marks,mark=o,mark size=1.8pt,color=propcol}
\addlegendentry{No maps}
\end{axis}
\begin{axis}[at={(a.outer south west)},anchor=outer north west,yshift=-1.1cm,name=c,scale only axis,width=.3\textwidth,height=3.2cm,xmode=log,log basis x=2,
 ymode=log,log basis y=2,xmin=3,xmax=160,ymin=50,ymax=1600,xtick={4,7,15,30,59,118},xticklabels={4,7,15,30,59,118},
 ytick={50,100,200,400,800,1600},yticklabels={50,100,200,400,800,$>$800},
 xlabel={Atlas patients in the reference},ylabel={Paired cells to the target},
 title={\textbf{c}\enspace Size of the atlas},
 axis x line*=bottom,axis y line*=left,tick align=outside,ymajorgrids=true,grid style={gray!18},axis line style={gray!55},tick label style={font=\sffamily\scriptsize},label style={font=\sffamily\small},clip=false,title style={font=\sffamily\small,at={(0,1)},anchor=base west,yshift=5pt}]
\draw[color=pairedcol,line width=.8pt,->] (axis cs:3,800) -- (axis cs:3,1400);
\draw[color=pairedcol,dashed,line width=.7pt] (axis cs:3,800) -- (axis cs:160,800);
\draw[color=sitecol,dashed,line width=.7pt] (axis cs:3,109.6) -- (axis cs:160,109.6);
\node[font=\sffamily\scriptsize,color=pairedcol,anchor=south east] at (axis cs:155,800) {paired cells only};
\node[font=\sffamily\scriptsize,color=sitecol,anchor=north west] at (axis cs:3.2,109.6) {own samples};
\draw[->,color=propcol,line width=.8pt] (axis cs:4,800) -- (axis cs:4,1400);
\draw[color=propcol,line width=.8pt] (axis cs:7,281.0) -- (axis cs:7,800.0);
\draw[color=propcol,line width=.8pt] (axis cs:15,154.8) -- (axis cs:15,245.7);
\draw[color=propcol,line width=.8pt] (axis cs:30,150.7) -- (axis cs:30,219.4);
\draw[color=propcol,line width=.8pt] (axis cs:59,147.3) -- (axis cs:59,217.3);
\draw[color=propcol,line width=.8pt] (axis cs:118,126.0) -- (axis cs:118,182.1);
\addplot[color=propcol,line width=1.2pt,mark=*,mark size=1.6pt] coordinates {(4,800.0) (7,384.6) (15,185.5) (30,176.5) (59,173.8) (118,147.3)};
\addplot[only marks,mark=o,mark size=1.8pt,color=propcol,fill=white] coordinates { (4,800) };
\end{axis}
\begin{axis}[at={(c.outer north east)},anchor=outer north west,xshift=.4cm,name=d,scale only axis,width=.3\textwidth,height=3.2cm,xmode=log,log basis x=2,
 xmin=20,xmax=1000,ymin=-5,ymax=40,xtick={25,50,100,200,400,800},xticklabels={25,50,100,200,400,800},ytick={0,10,20,30,40},
 xlabel={Paired cells},ylabel={Dependence recovered (\%)},
 title={\textbf{d}\enspace Bone marrow, RNA from nuclei},
 axis x line*=bottom,axis y line*=left,tick align=outside,ymajorgrids=true,grid style={gray!18},axis line style={gray!55},tick label style={font=\sffamily\scriptsize},label style={font=\sffamily\small},clip=false,title style={font=\sffamily\small,at={(0,1)},anchor=base west,yshift=5pt}]
\addplot[color=propcol,line width=1.2pt,mark=*,mark size=1.6pt] coordinates { (25,0.57) (50,7.41) (100,16.60) (200,20.75) (400,29.10) (800,35.40) };
\addplot[color=semicol,line width=.9pt,dashed,mark=triangle*,mark size=1.6pt,mark options={solid}] coordinates { (25,-4.10) (50,1.45) (100,8.78) (200,11.17) (400,18.77) (800,25.58) };
\addplot[color=semicol,line width=.9pt,mark=triangle*,mark size=1.6pt] coordinates { (25,-2.59) (50,-0.19) (100,4.61) (200,3.97) (400,6.32) (800,7.30) };
\addplot[color=propcol,line width=.9pt,dashed,mark=o,mark size=1.5pt,mark options={solid}] coordinates { (25,1.06) (50,3.02) (100,5.58) (200,5.56) (400,6.81) (800,7.03) };
\addplot[color=semicol,line width=.9pt,densely dotted,mark=triangle,mark size=1.6pt,mark options={solid}] coordinates { (25,0.49) (50,2.50) (100,4.89) (200,5.17) (400,6.53) (800,6.86) };
\end{axis}
\begin{axis}[hide axis,scale only axis,width=1pt,height=1pt,xmin=0,xmax=1,ymin=0,ymax=1,at={(d.outer south east)},anchor=north east,yshift=-2pt,legend style={draw=none,font=\sffamily\scriptsize,at={(1,0.5)},anchor=north east,legend columns=2,column sep=5pt,cells={anchor=west}}]
\addlegendimage{color=propcol,line width=1.2pt,mark=*,mark size=1.6pt}
\addlegendentry{Whole-cell axes and maps}
\addlegendimage{color=semicol,line width=.9pt,dashed,mark=triangle*,mark size=1.6pt,mark options={solid}}
\addlegendentry{Nuclear axes, whole-cell maps}
\addlegendimage{color=semicol,line width=.9pt,mark=triangle*,mark size=1.6pt}
\addlegendentry{Nuclear axes and maps}
\addlegendimage{color=propcol,line width=.9pt,dashed,mark=o,mark size=1.5pt,mark options={solid}}
\addlegendentry{Whole-cell axes only}
\addlegendimage{color=semicol,line width=.9pt,densely dotted,mark=triangle,mark size=1.6pt,mark options={solid}}
\addlegendentry{Nuclear axes only}
\end{axis}
\end{tikzpicture}

\caption{\textbf{What a reference contributes: axes and cell-type maps.} Follow-up and exploratory analyses.
\textbf{a}, Validation test: fraction of the dependence within cell types recovered with the atlas's axes and cell-type maps (the estimator), its axes only, its maps only (all entries shrunk as one block, then passed through the maps) or neither; dashed line, the accuracy target.
\textbf{b}, Paired cells needed to reach the target with the unchanged atlas and with atlases whose B cells and CD14 monocytes, or CD4 and CD8 T cells, came from 15 of the atlas's 118 patients, with the maps unweighted, reweighted to the paired cells' cell-type proportions or left out; bars, 95\% intervals from resampling donors; arrows, not reached within 800. Paired cells needed were read from regenerated draws that add a budget of 150 paired cells to the test's budgets, so the unchanged atlas needs 145 here rather than 147.
\textbf{c}, Paired cells needed with the atlas cells of fewer patients, with 95\% intervals; dashed lines, paired cells only and the study's own unpaired cells.
\textbf{d}, Bone-marrow test with contrasts: unpaired RNA from whole cells or from single nuclei, with and without their own cell-type maps, and nuclear axes with whole-cell maps.}
\label{fig:reference}
\end{figure*}

\subsection{A threefold saving across tissues, species and pairs of modalities}
We developed the estimator on two perturbation screens that measured RNA and surface proteins in cells carrying single-gene knockouts: Perturb-CITE-seq of melanoma cells\cite{frangieh2021} and ECCITE-seq of monocytic THP-1 cells\cite{papalexi2021}. With 100 paired cells, it recovered \SpFrRfHundred\% and \SpPaRfHundred\% of the dependence in cells whose knocked-out gene had been held out, against \SpFrPairedOnlyHundred\% and \SpPaPairedOnlyHundred\% for paired-only estimation (Fig.~\ref{fig:savings}a,d). The gain came from the unpaired axes. The same shrinkage along axes estimated from the paired cells, or in the original coordinates, recovered far less, whereas 500 unpaired cells per modality kept most of the gain (Fig.~\ref{fig:savings}d).

We then tested the fixed estimator on \SyNTest{} data sets not used to develop it. In a prespecified external test on primary T cells that each overexpress one gene\cite{legut2022}, it reached its target, 30\% recovery, with \SpExCellsProposed{} paired cells, whereas paired-only estimation needed \SpExCellsPairedOnly. It met all three prespecified comparisons, with paired-only estimation, SemiCCA and Champollion (Fig.~\ref{fig:savings}b). A prespecified benchmark then held out donors, patients, mice, recording days or neuron types in nine data sets. Seven pair RNA with surface proteins, chromatin accessibility or the electrical recordings of the same neurons (Patch-seq), and two pair the brain and nerve-cord partners of fly neurons (Fig.~\ref{fig:savings}c; Methods). In these tests, the unpaired cells were cells of the same data set that kept only one of their two modalities, and the paired cells were centred on the unpaired cells' averages (Methods). In the \SyNInf{} of these \SyNTest{} data sets in which any method recovered dependence, paired-only estimation needed at least \SyGeo{} times as many paired cells as the estimator to reach each data set's target (descriptive geometric mean, exploratory; 95\% CI, \SyGeoLo--\SyGeoHi; range, \SyMin{} to at least \SyMax; Fig.~\ref{fig:savings}c). Over the eight benchmark data sets with recovered dependence, SemiCCA, which also uses unpaired cells, needed at least \GenGTwo{} times as many (geometric mean; 95\% CI, \GenGTwoLo--\GenGTwoHi; Supplementary Note~\Sref{sec:gen}).

\begin{figure*}[tp]
\centering
% Generated by extension/synthesis/make_assets.py. Do not edit.
\definecolor{propcol}{HTML}{217D91}
\definecolor{pairedcol}{HTML}{5D6670}
\definecolor{semicol}{HTML}{C9803A}
\definecolor{champcol}{HTML}{8E4B8C}
\definecolor{ridgecol}{HTML}{9AA3AB}
\definecolor{meancol}{HTML}{A64442}
\definecolor{sitecol}{HTML}{6A8FB3}
\begin{tikzpicture}[font=\sffamily]
\begin{axis}[name=a,scale only axis,width=.3\textwidth,height=3.9cm,xmode=log,log basis x=2,
 xmin=40,xmax=3600,ymin=-20,ymax=100,xtick={50,100,200,400,800,1600,3200},xticklabels={50,100,200,400,800,1600,3200},ytick={-20,0,20,40,60,80,100},
 xlabel={Paired cells},ylabel={Dependence recovered (\%)},
 title={\textbf{a}\enspace Melanoma screen (development)},
 axis x line*=bottom,axis y line*=left,tick align=outside,ymajorgrids=true,grid style={gray!18},axis line style={gray!55},tick label style={font=\sffamily\scriptsize},label style={font=\sffamily\small},clip=false,title style={font=\sffamily\small,at={(0,1)},anchor=base west,yshift=5pt}]
\draw[dashed,black!45] (axis cs:40,30) -- (axis cs:3600,30);
\addplot[color=propcol,line width=1.1pt,mark=*,mark size=1.5pt] coordinates { (50,37.41) (100,48.06) (200,62.94) (400,72.39) (800,78.95) (1600,83.90) (3200,86.91) };
\addplot[color=pairedcol,line width=.8pt,mark=square*,mark size=1.3pt] coordinates { (50,9.10) (100,15.60) (200,28.06) (400,43.11) (800,57.81) (1600,67.01) (3200,74.36) };
\addplot[color=semicol,line width=.8pt,mark=triangle*,mark size=1.6pt] coordinates { (50,16.36) (100,31.06) (200,44.72) (400,53.32) (800,65.78) (1600,69.93) (3200,79.27) };
\addplot[color=ridgecol,line width=.8pt,dashed,mark=o,mark size=1.3pt] coordinates { (50,13.74) (100,27.21) (200,42.70) (400,57.30) (800,65.54) (1600,72.46) (3200,82.17) };
\addplot[color=champcol,line width=.8pt,mark=diamond*,mark size=1.7pt] coordinates { (50,19.44) (100,40.35) (200,54.39) (400,69.73) };
\end{axis}
\begin{axis}[at={(a.outer north east)},anchor=outer north west,xshift=.3cm,name=b,scale only axis,width=.3\textwidth,height=3.9cm,xmode=log,log basis x=2,
 xmin=20,xmax=480,ymin=-20,ymax=80,xtick={25,50,100,200,400},xticklabels={25,50,100,200,400},ytick={-20,0,20,40,60,80},
 xlabel={Paired cells},ylabel={Dependence recovered (\%)},
 title={\textbf{b}\enspace External test (prespecified)},
 axis x line*=bottom,axis y line*=left,tick align=outside,ymajorgrids=true,grid style={gray!18},axis line style={gray!55},tick label style={font=\sffamily\scriptsize},label style={font=\sffamily\small},clip=false,title style={font=\sffamily\small,at={(0,1)},anchor=base west,yshift=5pt}]
\draw[dashed,black!45] (axis cs:20,30) -- (axis cs:480,30);
\fill[propcol,opacity=.15] plot coordinates { (25,-6.62) (50,15.48) (100,27.06) (200,37.85) (400,46.54) (400,68.48) (200,56.28) (100,44.14) (50,27.81) (25,10.92) } -- cycle;
\addplot[color=propcol,line width=1.1pt,mark=*,mark size=1.5pt] coordinates { (25,5.15) (50,22.61) (100,36.03) (200,47.02) (400,56.93) };
\addplot[color=pairedcol,line width=.8pt,mark=square*,mark size=1.3pt] coordinates { (25,2.58) (50,5.70) (100,12.15) (200,20.22) (400,31.92) };
\addplot[color=semicol,line width=.8pt,mark=triangle*,mark size=1.6pt] coordinates { (25,11.67) (50,16.14) (100,26.91) (200,37.87) (400,46.71) };
\addplot[color=champcol,line width=.8pt,mark=diamond*,mark size=1.7pt] coordinates { (25,0.86) (50,8.60) (100,22.78) (200,41.38) (400,54.50) };
\addplot[color=ridgecol,line width=.8pt,dashed,mark=o,mark size=1.3pt] coordinates { (25,11.54) (50,16.97) (100,27.13) (200,35.33) (400,47.22) };
\end{axis}
\path (a.outer south west) -- (b.outer south east) coordinate[pos=.5] (legab);
\begin{axis}[hide axis,scale only axis,width=1pt,height=1pt,xmin=0,xmax=1,ymin=0,ymax=1,at={(legab)},anchor=north,yshift=-2pt,legend style={draw=none,font=\sffamily\scriptsize,at={(0.5,0.5)},anchor=north,legend columns=3,column sep=6pt,cells={anchor=west}}]
\addlegendimage{color=propcol,line width=1.1pt,mark=*,mark size=1.5pt}
\addlegendentry{Estimator}
\addlegendimage{color=pairedcol,line width=.8pt,mark=square*,mark size=1.3pt}
\addlegendentry{Paired cells only}
\addlegendimage{color=semicol,line width=.8pt,mark=triangle*,mark size=1.6pt}
\addlegendentry{SemiCCA}
\addlegendimage{color=champcol,line width=.8pt,mark=diamond*,mark size=1.7pt}
\addlegendentry{Champollion}
\addlegendimage{color=ridgecol,line width=.8pt,dashed,mark=o,mark size=1.3pt}
\addlegendentry{Reference regression}
\end{axis}
\begin{axis}[at={(a.outer south west)},anchor=outer north west,yshift=-1.25cm,name=c,scale only axis,width=.23\textwidth,height=4.6cm,xmode=log,log basis x=2,
 xmin=0.5,xmax=12,ymin=-.6,ymax=10.6,xtick={0.5,1,2,4,8},xticklabels={0.5,1,2,4,8},
 ytick={0,1,2,3,4,5,6,7,8,9,10},yticklabels={{\textbf{Geometric mean}, descriptive},{Fly CNS, male},{Fly CNS, female},{Visual cortex Patch-seq},{Motor cortex Patch-seq},{Bone marrow multiome},{Colon},{Bone marrow},{Blood, COVID-19},{Blood, 8 donors},{T cells (external test)}},
 xlabel={Saving over paired cells only},
 title={\textbf{c}\enspace Saving in the test data sets},
 axis x line*=bottom,axis y line*=left,tick align=outside,xmajorgrids=true,grid style={gray!18},axis line style={gray!55},tick label style={font=\sffamily\scriptsize},label style={font=\sffamily\small},clip=false,title style={font=\sffamily\small,at={(0,1)},anchor=base west,yshift=5pt},y tick style={draw=none},yticklabel style={font=\sffamily\scriptsize}]
\draw[black!50] (axis cs:1,-.6) -- (axis cs:1,10.6);
\draw[color=propcol,line width=.7pt] (axis cs:2.631,1) -- (axis cs:3.133,1);
\addplot[only marks,color=propcol,mark=o,mark size=1.7pt] coordinates { (2.932,1) };
\draw[color=propcol,line width=.7pt] (axis cs:3.144,2) -- (axis cs:3.849,2);
\addplot[only marks,color=propcol,mark=*,mark size=1.7pt] coordinates { (3.590,2) };
\draw[color=propcol,line width=.7pt] (axis cs:1.198,3) -- (axis cs:1.587,3);
\addplot[only marks,color=propcol,mark=*,mark size=1.7pt] coordinates { (1.392,3) };
\node[font=\sffamily\scriptsize,color=black!70,anchor=west] at (axis cs:1.05,4) {no dependence recovered};
\draw[color=propcol,line width=.7pt] (axis cs:1.713,5) -- (axis cs:3.019,5);
\addplot[only marks,color=propcol,mark=o,mark size=1.7pt] coordinates { (2.812,5) };
\draw[color=propcol,line width=.7pt] (axis cs:3.304,6) -- (axis cs:7.701,6);
\addplot[only marks,color=propcol,mark=o,mark size=1.7pt] coordinates { (6.361,6) };
\draw[color=propcol,line width=.7pt] (axis cs:1.000,7) -- (axis cs:5.988,7);
\addplot[only marks,color=propcol,mark=o,mark size=1.7pt] coordinates { (4.750,7) };
\draw[color=propcol,line width=.7pt] (axis cs:1.000,8) -- (axis cs:2.604,8);
\addplot[only marks,color=propcol,mark=o,mark size=1.7pt] coordinates { (2.049,8) };
\draw[color=propcol,line width=.7pt] (axis cs:1.000,9) -- (axis cs:2.075,9);
\addplot[only marks,color=propcol,mark=o,mark size=1.7pt] coordinates { (1.962,9) };
\draw[color=propcol,line width=.7pt] (axis cs:3.222,10) -- (axis cs:6.034,10);
\addplot[only marks,color=propcol,mark=*,mark size=1.7pt] coordinates { (4.875,10) };
\draw[color=black,line width=.9pt] (axis cs:2.442,0) -- (axis cs:3.208,0);
\addplot[only marks,color=black,mark=diamond*,mark size=2.6pt] coordinates { (3.072,0) };
\end{axis}
\begin{axis}[at={(c.outer north east)},anchor=outer north west,xshift=.2cm,name=d,scale only axis,width=.14\textwidth,height=4.6cm,xmin=-5,xmax=60,
 ymin=.4,ymax=7.6,ytick={1,2,3,4,5,6,7},
 yticklabels={{Paired cells only},{500 unpaired cells},{Original coordinates},{Axes from paired cells},{RNA axes only},{Without cell-type map},{Estimator}},
 xtick={0,20,40,60},xlabel={At 100 paired cells (\%)},
 title={\textbf{d}\enspace Source of the saving},
 axis x line*=bottom,axis y line*=left,tick align=outside,xmajorgrids=true,grid style={gray!18},axis line style={gray!55},tick label style={font=\sffamily\scriptsize},label style={font=\sffamily\small},clip=false,title style={font=\sffamily\small,at={(0,1)},anchor=base west,yshift=5pt},y tick style={draw=none}]
\addplot[only marks,mark=*,mark size=1.7pt,color=propcol] coordinates { (15.56,1.15) (41.88,2.15) (14.70,3.15) (11.19,4.15) (38.51,5.15) (44.51,6.15) (48.06,7.15) };
\addplot[only marks,mark=o,mark size=1.7pt,color=pairedcol] coordinates { (20.54,0.85) (34.51,1.85) (18.94,2.85) (15.93,3.85) (37.09,4.85) (35.69,5.85) (37.09,6.85) };
\end{axis}
\end{tikzpicture}

\caption{\textbf{A threefold saving of paired cells across tissues, species and pairs of modalities.}
\textbf{a}, Development: fraction of the dependence within culture conditions recovered in held-out perturbations of the melanoma screen (\SpFrGroups{} groups from \SpFrTargets{} knocked-out genes) against the number of paired cells. All methods used the same paired cells, and those that use unpaired cells the same unpaired cells. Paired cells only, paired-only estimation, the best at each budget of four methods that use paired cells alone; SemiCCA and reference regression, the better of pooled and per-condition fits. Dashed line, the external test's target.
\textbf{b}, Prespecified external test in OverCITE-seq (\SpExOrfs{} held-out overexpressed genes); band, 95\% interval of the estimator from resampling held-out genes.
\textbf{c}, Number of paired cells needed by paired-only estimation divided by the number the estimator needed, at each test's target (30\% recovery in the external test; half of the best recovery in the benchmark), with 95\% intervals from resampling held-out units. Open symbols, lower bounds: paired-only estimation did not reach the target within the largest budget, or the estimator reached it at the smallest. CNS, central nervous system. Diamond, descriptive geometric mean over the \SyNInf{} data sets with recovered dependence, including the lower bounds.
\textbf{d}, Recovery with 100 paired cells by the estimator, by versions with one part removed or replaced or with only 500 unpaired cells per modality, and by paired-only estimation, in the melanoma (filled) and monocyte (open) screens.}
\label{fig:savings}
\end{figure*}

\subsection{Where the dependence lies predicts how much the axes save}
To isolate what the unpaired axes save, we compared block shrinkage with single-block shrinkage, which shrinks all entries together as one block; neither used cell-type maps. We call the number of paired cells that single-block shrinkage needs to reach a given accuracy, divided by the number block shrinkage needs, the axes' saving. Two sets of shares describe the blocks: the dependence shares, the fraction of the dependence in each block, and the noise shares, the fraction of the per-cell noise. With oracle-linear shrinkage, the axes' saving follows from these shares in closed form (Proposition~\Sref{prop:law}; Fig.~\ref{fig:overview}c). It is one when the two sets of shares coincide, as they nearly do for random axes. It grows as the dependence concentrates in blocks that carry little of the noise, and at low accuracy it approaches $1+\chi^2(w\,\|\,\pi)$, where the chi-squared divergence $\chi^2$ measures how far the dependence shares $w$ depart from the noise shares $\pi$. We call this relation the law.

The law's inputs can be measured. Unpaired cells give the noise shares that the modalities would have if they were independent. A small pilot of paired cells gives unbiased raw estimates of each block's squared dependence; we form plug-in estimates of the dependence shares by setting negative estimates to zero and normalizing (Proposition~\Sref{prop:reveal}).

We developed the protocol for applying the law on the nine benchmark data sets after they had been scored (Fig.~\ref{fig:pred}a) and then tested it prospectively, fixing the protocol before extracting \PredNData{} more data sets. Analysed with panels of 100, 200 and 400 features, these gave \PredNProblems{} problems. In \PredNEstimable{} of them, from \PredNEstimableData{} data sets in two studies, the axes' saving could be measured within the budgets tested (Fig.~\ref{fig:pred}c; Supplementary Table~\Sref{tab:pred}). All six prespecified hypotheses were met (Supplementary Table~\Sref{tab:registry}). The law's median error was \PredLawErr\% (95\% CI, \PredLawErrLo--\PredLawErrHi\%; Fig.~\ref{fig:pred}a), and it ranked the problems almost perfectly (Spearman $\rho=\PredRho$). Averaging the problems within each data set preserved their ordering (follow-up; Supplementary Table~\Sref{tab:lawunits}).

In a follow-up analysis, a random pilot of 50 to 200 paired cells was used to estimate the law's inputs and to set the total budget for block shrinkage, pilot included. At a target of 50\% recovery, block shrinkage with these budgets reached the target in \RvPilotHalfFiftyReached\% to \RvPilotHalfTwoHundredReached\% of random pilots (Fig.~\ref{fig:pred}d; Supplementary Table~\Sref{tab:pilotdesign}).

\begin{figure*}[tp]
\centering
% Generated by extension/synthesis/make_assets.py. Do not edit.
\definecolor{protcol}{HTML}{217D91}
\definecolor{chromcol}{HTML}{8C5BA8}
\definecolor{wirecol}{HTML}{C9803A}
\definecolor{ephyscol}{HTML}{6BAF5A}
\definecolor{meancol}{HTML}{A64442}
\definecolor{pairedcol}{HTML}{5D6670}
\begin{tikzpicture}[font=\sffamily]
\begin{axis}[name=a,scale only axis,width=.2\textwidth,height=.2\textwidth,xmode=log,ymode=log,log basis x=2,log basis y=2,xmin=.5,xmax=48,ymin=.5,ymax=48,xtick={0.5,1,2,4,8,16,32},ytick={0.5,1,2,4,8,16,32},xticklabels={.5,1,2,4,8,16,32},yticklabels={.5,1,2,4,8,16,32},xlabel={Predicted saving},ylabel={Observed saving},title={\textbf{a}\enspace 50\%, 10 of 29 measurable}, axis x line*=bottom,axis y line*=left,tick align=outside,ymajorgrids=true,xmajorgrids=true,grid style={gray!18},axis line style={gray!55},tick label style={font=\sffamily\scriptsize},label style={font=\sffamily\small},clip=false,title style={font=\sffamily\small,at={(0,1)},anchor=base west,yshift=5pt}]
\addplot[black!45,domain=.5:48,samples=2] {x};
\addplot[black!25,dashed,domain=.5:38.4,samples=2] {1.25*x};
\addplot[black!25,dashed,domain=.625:48,samples=2] {x/1.25};
\addplot[only marks,mark=o,mark size=1.5pt,color=black!35] coordinates {(2.639,2.207) (1.925,1.703) (2.791,2.302) (5.799,5.194) (2.982,2.769) (12.265,10.477) (6.293,6.445) (2.450,2.466) (17.978,14.723) (2.176,2.118) (1.476,1.416) (5.680,5.600) (2.345,2.223) (8.817,8.813)};
\addplot[only marks,mark=*,mark size=1.8pt,color=protcol] coordinates {(6.763,5.380) (9.697,8.327) (14.549,11.567)};
\addplot[only marks,mark=square*,mark size=1.8pt,color=protcol] coordinates {(10.025,7.059) (14.716,11.430)};
\addplot[only marks,mark=triangle*,mark size=1.8pt,color=wirecol] coordinates {(2.340,1.874) (2.930,2.494) (3.479,3.088)};
\addplot[only marks,mark=diamond*,mark size=1.8pt,color=wirecol] coordinates {(3.545,3.259) (4.271,3.846)};
\end{axis}
\begin{axis}[name=b,at={(a.outer north east)},anchor=outer north west,xshift=.2cm,scale only axis,width=.2\textwidth,height=.2\textwidth,xmode=log,ymode=log,log basis x=2,log basis y=2,xmin=.5,xmax=48,ymin=.5,ymax=48,xtick={0.5,1,2,4,8,16,32},ytick={0.5,1,2,4,8,16,32},xticklabels={.5,1,2,4,8,16,32},yticklabels={.5,1,2,4,8,16,32},xlabel={Predicted saving},title={\textbf{b}\enspace 30\% and 70\%}, axis x line*=bottom,axis y line*=left,tick align=outside,ymajorgrids=true,xmajorgrids=true,grid style={gray!18},axis line style={gray!55},tick label style={font=\sffamily\scriptsize},label style={font=\sffamily\small},clip=false,title style={font=\sffamily\small,at={(0,1)},anchor=base west,yshift=5pt}]
\addplot[black!45,domain=.5:48,samples=2] {x};
\addplot[black!25,dashed,domain=.5:38.4,samples=2] {1.25*x};
\addplot[black!25,dashed,domain=.625:48,samples=2] {x/1.25};
\addplot[only marks,mark=o,mark size=1.8pt,color=protcol] coordinates {(10.071,7.411) (15.527,7.819) (25.469,15.171)};
\addplot[only marks,mark=square,mark size=1.8pt,color=protcol] coordinates {(14.673,8.120) (22.234,14.430) (34.392,23.010)};
\addplot[only marks,mark=o,mark size=1.8pt,color=ephyscol] coordinates {(2.173,0.818) (2.068,0.920)};
\addplot[only marks,mark=triangle*,mark size=1.8pt,color=wirecol] coordinates {(1.741,1.652) (2.056,1.970) (2.387,2.255)};
\addplot[only marks,mark=diamond*,mark size=1.8pt,color=wirecol] coordinates {(2.458,2.476) (2.866,2.954)};
\node[font=\sffamily\scriptsize,anchor=north west,align=left] at (axis cs:.55,44.16) {open, 30\%\\ filled, 70\%};
\end{axis}
\begin{axis}[name=c,at={(b.outer north east)},anchor=outer north west,xshift=.2cm,scale only axis,width=.2\textwidth,height=.2\textwidth,xmode=log,log basis x=2,xmin=1,xmax=64,ymin=.4,ymax=3.6,xtick={1,2,4,8,16,32,64},xticklabels={1,2,4,8,16,32,64},ytick={1,2,3},yticklabels={{Not reached\\within budgets},{Above target\\at 25 cells},{Measured}},yticklabel style={align=right},xlabel={Predicted saving},title={\textbf{c}\enspace All 29 problems}, axis x line*=bottom,axis y line*=left,tick align=outside,ymajorgrids=true,xmajorgrids=true,grid style={gray!18},axis line style={gray!55},tick label style={font=\sffamily\scriptsize},label style={font=\sffamily\small},clip=false,title style={font=\sffamily\small,at={(0,1)},anchor=base west,yshift=5pt}]
\addplot[only marks,mark=*,mark size=1.4pt,color=protcol] coordinates {(20.165,1.000)};
\addplot[only marks,mark=*,mark size=1.4pt,color=chromcol] coordinates {(3.288,0.820) (4.463,1.000) (5.814,1.180)};
\addplot[only marks,mark=*,mark size=1.4pt,color=chromcol] coordinates {(28.251,0.820) (40.710,1.000) (51.909,1.180)};
\addplot[only marks,mark=*,mark size=1.4pt,color=ephyscol] coordinates {(1.770,0.820) (1.786,1.000) (1.887,1.180)};
\addplot[only marks,mark=*,mark size=1.4pt,color=protcol] coordinates {(2.392,1.820) (3.348,2.000) (4.881,2.180)};
\addplot[only marks,mark=*,mark size=1.4pt,color=protcol] coordinates {(4.993,1.820) (5.904,2.000) (8.100,2.180)};
\addplot[only marks,mark=*,mark size=1.4pt,color=protcol] coordinates {(4.768,1.820) (6.887,2.000) (10.362,2.180)};
\addplot[only marks,mark=*,mark size=1.4pt,color=protcol] coordinates {(6.763,2.820) (9.697,3.000) (14.549,3.180)};
\addplot[only marks,mark=*,mark size=1.4pt,color=protcol] coordinates {(10.025,2.820) (14.716,3.180)};
\addplot[only marks,mark=*,mark size=1.4pt,color=wirecol] coordinates {(2.340,2.820) (2.930,3.000) (3.479,3.180)};
\addplot[only marks,mark=*,mark size=1.4pt,color=wirecol] coordinates {(3.545,2.820) (4.271,3.180)};
\end{axis}
\begin{axis}[hide axis,scale only axis,width=1pt,height=1pt,xmin=0,xmax=1,ymin=0,ymax=1,at={(c.outer south east)},anchor=north east,yshift=-2pt,legend style={draw=none,font=\sffamily\scriptsize,at={(1,0.5)},anchor=north east,legend columns=2,column sep=5pt,cells={anchor=west}}]
\addlegendimage{only marks,mark=*,mark size=1.4pt,color=protcol}
\addlegendentry{RNA--protein}
\addlegendimage{only marks,mark=*,mark size=1.4pt,color=chromcol}
\addlegendentry{RNA--chromatin}
\addlegendimage{only marks,mark=*,mark size=1.4pt,color=wirecol}
\addlegendentry{Fly wiring}
\addlegendimage{only marks,mark=*,mark size=1.4pt,color=ephyscol}
\addlegendentry{Patch-seq}
\end{axis}
\begin{axis}[name=d,at={(b.outer south west)},anchor=outer north west,yshift=-1cm,scale only axis,width=.2\textwidth,height=.2\textwidth,xmin=.5,xmax=3.5,ymin=0,ymax=100,xtick={1,2,3},xticklabels={30\%,50\%,70\%},ytick={0,20,40,60,80,100},xlabel={Target},ylabel={Achieved (\%)},title={\textbf{d}\enspace Budgets set from a pilot}, axis x line*=bottom,axis y line*=left,tick align=outside,ymajorgrids=true,xmajorgrids=true,grid style={gray!18},axis line style={gray!55},tick label style={font=\sffamily\scriptsize},label style={font=\sffamily\small},clip=false,title style={font=\sffamily\small,at={(0,1)},anchor=base west,yshift=5pt}]
\draw[dashed,black!45] (axis cs:0.6,30) -- (axis cs:1.4,30);
\draw[color=protcol!45,line width=.7pt] (axis cs:0.76,34.1) -- (axis cs:0.76,90.4);
\addplot[only marks,mark=triangle*,mark size=1.8pt,color=protcol!45] coordinates { (0.76,41.3) };
\draw[color=protcol!75,line width=.7pt] (axis cs:0.92,35.4) -- (axis cs:0.92,94.2);
\addplot[only marks,mark=triangle*,mark size=1.8pt,color=protcol!75] coordinates { (0.92,41.1) };
\draw[color=protcol,line width=.7pt] (axis cs:1.08,34.4) -- (axis cs:1.08,96.7);
\addplot[only marks,mark=triangle*,mark size=1.8pt,color=protcol] coordinates { (1.08,48.1) };
\draw[color=black,line width=.7pt] (axis cs:1.24,11.6) -- (axis cs:1.24,25.5);
\addplot[only marks,mark=diamond,mark size=2.1pt,color=black] coordinates { (1.24,22.6) };
\draw[dashed,black!45] (axis cs:1.6,50) -- (axis cs:2.4,50);
\draw[color=protcol!45,line width=.7pt] (axis cs:1.76,48.5) -- (axis cs:1.76,89.8);
\addplot[only marks,mark=triangle*,mark size=1.8pt,color=protcol!45] coordinates { (1.76,54.2) };
\draw[color=protcol!75,line width=.7pt] (axis cs:1.92,48.4) -- (axis cs:1.92,94.4);
\addplot[only marks,mark=triangle*,mark size=1.8pt,color=protcol!75] coordinates { (1.92,52.9) };
\draw[color=protcol,line width=.7pt] (axis cs:2.08,48.4) -- (axis cs:2.08,96.7);
\addplot[only marks,mark=triangle*,mark size=1.8pt,color=protcol] coordinates { (2.08,52.5) };
\draw[color=black,line width=.7pt] (axis cs:2.24,40.5) -- (axis cs:2.24,48.0);
\addplot[only marks,mark=diamond,mark size=2.1pt,color=black] coordinates { (2.24,46.0) };
\draw[dashed,black!45] (axis cs:2.6,70) -- (axis cs:3.4,70);
\draw[color=protcol!45,line width=.7pt] (axis cs:2.76,60.3) -- (axis cs:2.76,89.8);
\addplot[only marks,mark=triangle*,mark size=1.8pt,color=protcol!45] coordinates { (2.76,67.6) };
\draw[color=protcol!75,line width=.7pt] (axis cs:2.92,60.5) -- (axis cs:2.92,94.3);
\addplot[only marks,mark=triangle*,mark size=1.8pt,color=protcol!75] coordinates { (2.92,69.1) };
\draw[color=protcol,line width=.7pt] (axis cs:3.08,60.5) -- (axis cs:3.08,96.7);
\addplot[only marks,mark=triangle*,mark size=1.8pt,color=protcol] coordinates { (3.08,69.8) };
\draw[color=black,line width=.7pt] (axis cs:3.24,60.5) -- (axis cs:3.24,69.0);
\addplot[only marks,mark=diamond,mark size=2.1pt,color=black] coordinates { (3.24,67.0) };
\end{axis}
\begin{axis}[hide axis,scale only axis,width=1pt,height=1pt,xmin=0,xmax=1,ymin=0,ymax=1,at={(d.outer east)},anchor=west,xshift=.3cm,legend style={draw=none,font=\sffamily\scriptsize,at={(0,0.5)},anchor=west,legend columns=1,cells={anchor=west}}]
\addlegendimage{only marks,mark=triangle*,mark size=1.8pt,color=protcol!45}
\addlegendentry{Pilot of 50 cells}
\addlegendimage{only marks,mark=triangle*,mark size=1.8pt,color=protcol!75}
\addlegendentry{Pilot of 100 cells}
\addlegendimage{only marks,mark=triangle*,mark size=1.8pt,color=protcol}
\addlegendentry{Pilot of 200 cells}
\addlegendimage{only marks,mark=diamond,mark size=2.1pt,color=black}
\addlegendentry{Law with every reservoir cell}
\addlegendimage{dashed,black!45}
\addlegendentry{Target}
\end{axis}
\begin{axis}[hide axis,scale only axis,width=1pt,height=1pt,xmin=0,xmax=1,ymin=0,ymax=1,at={(a.outer south west)},anchor=north west,yshift=-2pt,legend style={draw=none,font=\sffamily\scriptsize,at={(0,0.5)},anchor=north west,legend columns=1,cells={anchor=west}}]
\addlegendimage{only marks,mark=o,mark size=1.5pt,color=black!35}\addlegendentry{Benchmark (development)}
\addlegendimage{only marks,mark=*,mark size=1.8pt,color=protcol}\addlegendentry{Mouse lymphoid organs, 111 proteins}
\addlegendimage{only marks,mark=square*,mark size=1.8pt,color=protcol}\addlegendentry{Mouse lymphoid organs, 206 proteins}
\addlegendimage{only marks,mark=triangle*,mark size=1.8pt,color=wirecol}\addlegendentry{Fly visual neurons, FlyWire}
\addlegendimage{only marks,mark=diamond*,mark size=1.8pt,color=wirecol}\addlegendentry{Fly visual neurons, BANC}
\end{axis}
\end{tikzpicture}

\caption{\textbf{Where the dependence lies predicts how much the axes save.} Saving: the axes' saving, the number of paired cells needed by single-block James--Stein shrinkage divided by the number needed by block shrinkage along the unpaired axes, both without cell-type maps.
\textbf{a}, At 50\% recovery, observed saving against the saving predicted by Proposition~\Sref{prop:law} from each block's dependence and noise in the reservoir (the cells from which paired cells are drawn), in the prospective test (coloured symbols; the \PredNEstimable{} of \PredNProblems{} problems with a measurable saving; FlyWire and BANC, two fly connectomes) and on the nine benchmark data sets used to develop the law's protocol (grey circles). Solid line, equality; dashed lines, $\pm25\%$.
\textbf{b}, The same at 30\% (open symbols) and 70\% (filled) recovery.
\textbf{c}, Every prospective problem at 50\% recovery, placed by the law's prediction: measured; not measurable because a curve already exceeded the target with 25 paired cells; or not measurable because a curve did not reach it within the budgets tested. Colours, pairs of modalities.
\textbf{d}, Budgets set from a pilot (follow-up): recovery achieved when a random pilot of 50, 100 or 200 reservoir cells set the total budget at a target of 30\%, 50\% or 70\%, the pilot counted in the total. Points, median over problems of each problem's median over 100 replicates; bars, interquartile range over problems; diamonds, the budget set from every reservoir cell; dashed lines, the targets.}
\label{fig:pred}
\end{figure*}

\subsection{Population averages identify the dependence only where populations differ}
Differences between populations offer a route to the dependence that needs no paired cells. Suppose, as the cell-type maps do, that every population shares one channel, and that each population is profiled in both assays, even if in different cells. The populations' average RNA and protein then determine the channel along the directions in which their average RNA profiles differ. Population averages alone leave it undetermined along all other directions (Proposition~\Sref{prop:means}; Fig.~\ref{fig:overview}d).

Blood donors differ mainly in their mix of cell types. Under a prespecified protocol, we learned a channel from the \XsPatients{} patients of the COVID-19 study that supplied the atlas above\cite{stephenson2021} by regressing each patient's average protein profile on their average RNA profile across patients, a form of ecological regression\cite{goodman1953,hoerl1970}. Each patient's RNA and protein averages came from different cells. We then predicted the gene--protein correlations of eight donors of a second cohort\cite{hao2021}. Relative error was the Frobenius norm of the prediction's error divided by that of the held-out correlation matrix, where the Frobenius norm is the square root of the sum of squared entries; predicting zero correlation gives 1. The channel's relative error was \XsHaoErrPfMeans, against \XsHaoNullErr{} when the training patients' averages were paired at random (\textit{P}~=~\XsHaoPermP, the smallest attainable\cite{phipson2010}; Fig.~\ref{fig:pairfree}a). Its improvement over predicting zero correlation was \XsHaoShare\% (95\% CI, \XsHaoShareLo--\XsHaoShareHi\%) of that achieved by the same regression fitted to the training patients' paired cells. As expected when donors differ only in their mix of cell types (Proposition~\Sref{prop:mixture}), the dependence the channel captured lay between cell types: the part of the correlation due to differences between cell types was nearly as large as the whole (\XsMixHaoBetweenShare\% of its Frobenius norm; exploratory; Supplementary Fig.~\Sref{sfig:ident}).

Perturbations create variation within a cell type, but along few directions. In the melanoma screen, under a prespecified protocol, channels learned from the averages of \PtTrainPops{} populations, one per knocked-out gene and culture condition, recovered \PtRfPrimary\% of the dependence in held-out perturbations (95\% CI, \PtRfPrimaryLo--\PtRfPrimaryHi\%). The same model fitted to the paired cells of other knockouts recovered \PtRfPaired\% (Fig.~\ref{fig:pairfree}b). For the genes and proteins of the interferon response, which the knockouts affected most, the channels recovered \PgIfnPf\% of the dependence (exploratory). This recovery rested on the perturbations with the largest effects: without the five training perturbations that shifted RNA most, it fell to \DsFiveIfnStr\%, against \DsFiveIfnRand\% without five random ones (exploratory; Fig.~\ref{fig:pairfree}c). Population averages thus identified the dependence along the directions in which populations differed, as Proposition~\Sref{prop:means} predicts.

\begin{figure*}[tp]
\centering
% Generated by extension/synthesis/make_assets.py. Do not edit.
\definecolor{propcol}{HTML}{217D91}
\definecolor{pairedcol}{HTML}{5D6670}
\definecolor{semicol}{HTML}{C9803A}
\definecolor{champcol}{HTML}{8E4B8C}
\definecolor{ridgecol}{HTML}{9AA3AB}
\definecolor{meancol}{HTML}{A64442}
\definecolor{sitecol}{HTML}{6A8FB3}
\begin{tikzpicture}[font=\sffamily]
\begin{axis}[name=a,scale only axis,width=.25\textwidth,height=3.6cm,xmin=0,xmax=1.65,
 ymin=.4,ymax=3.6,xtick={0,.4,.8,1.2,1.6},ytick={1,2,3},
 yticklabels={{Own paired cells},{Paired cells, other donors},{No pairs: population averages}},
 xlabel={Relative error},
 title={\textbf{a}\enspace Blood, second cohort},
 axis x line*=bottom,axis y line*=left,tick align=outside,xmajorgrids=true,grid style={gray!18},axis line style={gray!55},tick label style={font=\sffamily\scriptsize},label style={font=\sffamily\small},clip=false,title style={font=\sffamily\small,at={(0,1)},anchor=base west,yshift=5pt},y tick style={draw=none}]
\draw[dashed,black!45] (axis cs:1,.4) -- (axis cs:1,3.6);
\addplot[only marks,mark=*,mark size=.6pt,color=propcol!45] coordinates { (0.1282,1.090) (0.1256,1.107) (0.1467,1.124) (0.1235,1.141) (0.1057,1.159) (0.1105,1.176) (0.1095,1.193) (0.1112,1.210) };
\addplot[only marks,mark=*,mark size=.6pt,color=propcol!45] coordinates { (0.1424,2.090) (0.1377,2.107) (0.1383,2.124) (0.1401,2.141) (0.1385,2.159) (0.1271,2.176) (0.1203,2.193) (0.1295,2.210) };
\addplot[only marks,mark=*,mark size=.6pt,color=propcol!45] coordinates { (0.3142,3.090) (0.3108,3.107) (0.3052,3.124) (0.3030,3.141) (0.2952,3.159) (0.3175,3.176) (0.2438,3.193) (0.2433,3.210) };
\addplot[only marks,mark=*,mark size=1.8pt,color=propcol] coordinates { (0.1201,1.15) (0.1342,2.15) (0.2916,3.15) };
\addplot[only marks,mark=o,mark size=1.8pt,color=meancol] coordinates { (0.3485,0.85) (0.5280,1.85) (1.0820,2.85) };
\end{axis}
\begin{axis}[hide axis,scale only axis,width=1pt,height=1pt,xmin=0,xmax=1,ymin=0,ymax=1,at={(a.outer south)},anchor=north,yshift=-2pt,legend style={draw=none,font=\sffamily\scriptsize,at={(0.5,0.5)},anchor=north,legend columns=2,column sep=6pt,cells={anchor=west}}]
\addlegendimage{only marks,mark=*,mark size=1.8pt,color=propcol}
\addlegendentry{All cells}
\addlegendimage{only marks,mark=o,mark size=1.8pt,color=meancol}
\addlegendentry{Within cell types}
\end{axis}
\begin{axis}[at={(a.outer north east)},anchor=outer north west,xshift=.1cm,name=b,scale only axis,width=.18\textwidth,height=3.6cm,xmin=-50,xmax=105,
 ymin=.4,ymax=3.6,xtick={-50,-25,0,25,50,75,100},xticklabels={$-$50,,0,,50,,100},
 ytick={1,2,3},
 yticklabels={{Own paired cells},{Paired cells, other knockouts},{No pairs: perturbations}},
 xlabel={Dependence recovered (\%)},
 title={\textbf{b}\enspace Melanoma screen},
 axis x line*=bottom,axis y line*=left,tick align=outside,xmajorgrids=true,grid style={gray!18},axis line style={gray!55},tick label style={font=\sffamily\scriptsize},label style={font=\sffamily\small},clip=false,title style={font=\sffamily\small,at={(0,1)},anchor=base west,yshift=5pt},y tick style={draw=none}]
\draw[dashed,black!45] (axis cs:0,.4) -- (axis cs:0,3.6);
\draw[propcol,line width=.8pt] (axis cs:43.96,1.22) -- (axis cs:53.20,1.22);
\draw[black,line width=.8pt] (axis cs:19.80,0.78) -- (axis cs:27.75,0.78);
\draw[propcol,line width=.8pt] (axis cs:58.30,2.22) -- (axis cs:69.92,2.22);
\draw[black,line width=.8pt] (axis cs:31.65,1.78) -- (axis cs:42.22,1.78);
\draw[propcol,line width=.8pt] (axis cs:32.37,3.22) -- (axis cs:51.44,3.22);
\draw[black,line width=.8pt] (axis cs:4.05,2.78) -- (axis cs:9.42,2.78);
\addplot[only marks,mark=*,mark size=1.7pt,color=propcol] coordinates { (48.33,1.22) (64.17,2.22) (42.52,3.22) };
\addplot[only marks,mark=o,mark size=1.7pt,color=pairedcol] coordinates { (18.66,1) (36.17,2) (5.66,3) };
\addplot[only marks,mark=diamond*,mark size=2pt,color=black] coordinates { (24.00,0.78) (37.27,1.78) (6.79,2.78) };
\end{axis}
\begin{axis}[hide axis,scale only axis,width=1pt,height=1pt,xmin=0,xmax=1,ymin=0,ymax=1,at={(b.outer south east)},anchor=north east,yshift=-2pt,legend style={draw=none,font=\sffamily\scriptsize,at={(1,0.5)},anchor=north east,legend columns=3,column sep=4pt,cells={anchor=west}}]
\addlegendimage{only marks,mark=*,mark size=1.7pt,color=propcol}
\addlegendentry{Interferon response}
\addlegendimage{only marks,mark=o,mark size=1.7pt,color=pairedcol}
\addlegendentry{Other}
\addlegendimage{only marks,mark=diamond*,mark size=2pt,color=black}
\addlegendentry{All genes and proteins}
\end{axis}
\begin{axis}[at={(a.outer south west)},anchor=outer north west,yshift=-1.3cm,name=c,scale only axis,width=.2\textwidth,height=3.6cm,
 xmin=.5,xmax=3.5,ymin=-5,ymax=65,xtick={1,2,3},xticklabels={5,10,20},ytick={0,20,40,60},
 xlabel={Perturbations removed},ylabel={Interferon response (\%)},
 title={\textbf{c}\enspace Removing training perturbations},
 axis x line*=bottom,axis y line*=left,tick align=outside,ymajorgrids=true,grid style={gray!18},axis line style={gray!55},tick label style={font=\sffamily\scriptsize},label style={font=\sffamily\small},clip=false,title style={font=\sffamily\small,at={(0,1)},anchor=base west,yshift=5pt}]
\draw[dashed,black!45] (axis cs:.5,42.52) -- (axis cs:3.5,42.52);
\addplot[only marks,mark=*,mark size=.8pt,color=pairedcol!55] coordinates { (0.780,41.89) (0.789,41.97) (0.798,43.91) (0.807,42.70) (0.816,43.18) (0.825,44.17) (0.834,23.78) (0.843,8.97) (0.852,42.56) (0.861,43.22) (0.870,41.81) (0.879,42.07) (0.888,41.99) (0.897,42.11) (0.906,42.49) (0.915,5.84) (0.924,42.77) (0.933,43.87) (0.942,40.91) (0.951,41.93) (0.960,42.50) (0.969,42.40) (0.978,42.35) (0.987,42.30) (0.996,41.15) (1.004,42.92) (1.013,41.20) (1.022,41.98) (1.031,41.86) (1.040,42.51) (1.049,42.01) (1.058,42.12) (1.067,41.98) (1.076,9.05) (1.085,43.70) (1.094,5.48) (1.103,44.03) (1.112,42.17) (1.121,42.57) (1.130,5.98) (1.139,41.68) (1.148,42.23) (1.157,9.62) (1.166,43.60) (1.175,42.34) (1.184,42.23) (1.193,42.38) (1.202,8.63) (1.211,6.11) (1.220,42.05) (1.780,42.38) (1.789,41.52) (1.798,41.64) (1.807,43.21) (1.816,42.41) (1.825,41.84) (1.834,41.98) (1.843,44.16) (1.852,41.51) (1.861,42.83) (1.870,42.08) (1.879,42.27) (1.888,8.77) (1.897,43.45) (1.906,41.91) (1.915,42.65) (1.924,41.88) (1.933,42.99) (1.942,41.42) (1.951,44.77) (1.960,41.97) (1.969,42.42) (1.978,41.73) (1.987,8.21) (1.996,41.31) (2.004,42.09) (2.013,41.34) (2.022,41.44) (2.031,7.18) (2.040,3.57) (2.049,42.33) (2.058,3.99) (2.067,42.17) (2.076,42.25) (2.085,26.39) (2.094,44.25) (2.103,41.65) (2.112,41.97) (2.121,10.66) (2.130,42.42) (2.139,43.26) (2.148,44.02) (2.157,42.19) (2.166,44.51) (2.175,42.28) (2.184,40.88) (2.193,42.28) (2.202,42.78) (2.211,42.47) (2.220,41.68) (2.780,40.67) (2.789,41.63) (2.798,9.59) (2.807,47.01) (2.816,41.84) (2.825,47.38) (2.834,42.20) (2.843,8.12) (2.852,46.95) (2.861,40.91) (2.870,43.13) (2.879,42.04) (2.888,42.50) (2.897,42.04) (2.906,43.63) (2.915,42.27) (2.924,41.18) (2.933,42.23) (2.942,41.50) (2.951,44.49) (2.960,40.94) (2.969,41.28) (2.978,40.70) (2.987,41.40) (2.996,41.42) (3.004,42.79) (3.013,41.52) (3.022,41.06) (3.031,43.96) (3.040,41.85) (3.049,5.08) (3.058,40.88) (3.067,40.05) (3.076,43.17) (3.085,41.32) (3.094,41.54) (3.103,40.86) (3.112,42.68) (3.121,42.38) (3.130,43.27) (3.139,39.61) (3.148,43.79) (3.157,5.36) (3.166,2.27) (3.175,40.83) (3.184,5.06) (3.193,41.53) (3.202,41.88) (3.211,46.10) (3.220,39.56) };
\addplot[only marks,mark=triangle*,mark size=2.1pt,color=propcol] coordinates { (0.88,40.56) (1.88,7.98) (2.88,8.42) };
\addplot[only marks,mark=square*,mark size=1.9pt,color=meancol] coordinates { (1.12,7.90) (2.12,4.66) (3.12,5.12) };
\end{axis}
\begin{axis}[hide axis,scale only axis,width=1pt,height=1pt,xmin=0,xmax=1,ymin=0,ymax=1,at={(c.outer south west)},anchor=north west,yshift=-2pt,legend style={draw=none,font=\sffamily\scriptsize,at={(0,0.5)},anchor=north west,legend columns=3,column sep=4pt,cells={anchor=west}}]
\addlegendimage{only marks,mark=triangle*,mark size=2.1pt,color=propcol}
\addlegendentry{Exposure (planned)}
\addlegendimage{only marks,mark=square*,mark size=1.9pt,color=meancol}
\addlegendentry{Effect strength}
\addlegendimage{only marks,mark=*,mark size=1pt,color=pairedcol!55}
\addlegendentry{Random}
\end{axis}
\begin{axis}[at={(c.outer north east)},anchor=outer north west,xshift=.1cm,name=d,scale only axis,width=.2\textwidth,height=3.6cm,xmin=-300,xmax=105,
 ymin=.4,ymax=4.6,xtick={-300,-200,-100,0,100},xticklabels={$-$300,$-$200,$-$100,0,100},
 ytick={1,2,3,4},
 yticklabels={{Paired cells, other knockouts},{Control: shuffled pairs},{Control: pseudo-populations},{No pairs: perturbations}},
 xlabel={Dependence recovered (\%)},
 title={\textbf{d}\enspace Monocyte screen (replication)},
 axis x line*=bottom,axis y line*=left,tick align=outside,xmajorgrids=true,grid style={gray!18},axis line style={gray!55},tick label style={font=\sffamily\scriptsize},label style={font=\sffamily\small},clip=false,title style={font=\sffamily\small,at={(0,1)},anchor=base west,yshift=5pt},y tick style={draw=none}]
\draw[dashed,black!45] (axis cs:0,.4) -- (axis cs:0,4.6);
\draw[black,line width=.8pt] (axis cs:20.70,0.85) -- (axis cs:43.08,0.85);
\draw[propcol,line width=.8pt] (axis cs:72.71,1.15) -- (axis cs:99.69,1.15);
\draw[black,line width=.8pt] (axis cs:-287.47,3.85) -- (axis cs:-157.29,3.85);
\draw[propcol,line width=.8pt] (axis cs:-35.65,4.15) -- (axis cs:19.48,4.15);
\addplot[only marks,mark=*,mark size=1.7pt,color=propcol] coordinates { (83.69,1.15) (-12.44,2.15) (-3.23,3.15) (-2.95,4.15) };
\addplot[only marks,mark=diamond*,mark size=2pt,color=black] coordinates { (32.37,0.85) (-45.59,1.85) (-10.24,2.85) (-212.41,3.85) };
\end{axis}
\end{tikzpicture}

\caption{\textbf{Without paired cells, population averages identify the dependence only where populations differ.}
\textbf{a}, Relative error (the Frobenius norm of the error divided by that of the held-out correlations; 1 for predicting zero correlation) of predicted correlations between \XsGenes{} genes and \XsProteins{} proteins in eight donors of a blood cohort not used for training (small points, donors; large points, averages), for all cells and within matched cell types; dashed line, no dependence. No pairs, a channel learned from the population averages of \XsPatients{} patients of another study; paired cells, other donors, the same regression fitted on those patients' paired cells; own paired cells, the held-out donors' own paired cells, which learning without pairs cannot use.
\textbf{b}, Melanoma screen: fraction of the dependence within conditions recovered in \PtGroups{} held-out groups of knocked-out gene and condition, for the genes and proteins of the interferon response (\PgIfnGenes{} Hallmark genes and their antibodies; circles), for all other genes and proteins (open circles) and for all of them (diamonds, the primary endpoint), with 95\% intervals from resampling held-out perturbations. No pairs, channels learned from the averages of \PtTrainPops{} training populations; paired cells, other knockouts, the same model fitted to their paired cells.
\textbf{c}, Interferon-response recovery without pairs after removing the training perturbations ranked highest by exposure (how far a perturbation moves RNA along directions in which cells already vary), as planned (triangles), by effect strength (the size of its shift in average RNA; squares) or at random (points; \DsDraws{} draws); dashed line, all training perturbations.
\textbf{d}, Prespecified replication in the monocyte screen (\RpGroups{} held-out groups): recovery over all genes and proteins (diamonds) and for the interferon-response genes and proteins (circles), with 95\% intervals; controls learn the channel from populations relabelled at random within condition (pseudo-populations) or from averages paired at random.}
\label{fig:pairfree}
\end{figure*}

\section{Discussion}
Only paired cells show directly how two modalities vary together, yet a few sufficed, because unpaired cells supplied two complementary kinds of structure. Their principal axes separate the dependence from the noise, so that a few paired cells estimate the dependence pooled over cell types, and how much the axes save can be predicted from where the dependence lies relative to the noise. Their covariance within each cell type then turns the pooled estimate into one for each type. This mapping step determined whether a reference from another study transferred: references of the same assay transferred, and reweighting to the paired cells' cell-type proportions let skewed references reach the target again. Correlation structure measured with the same assay carried over between sites, studies and animals, whereas averages did not; batch correction assumes the same when it removes shifts in average and scale between batches\cite{johnson2007,korsunsky2019}.

For a laboratory planning a paired study, these results suggest using same-assay unpaired profiles from the study when available, or an atlas reweighted to the paired sample's cell-type proportions. A paired pilot can help assess the potential saving from the unpaired axes. Reference regression, which likewise passes the paired estimate through the atlas's covariance, came closest to the estimator. Without paired cells, population averages determine the dependence only along the directions in which they differ, as when bulk tissue profiles are separated into cell types\cite{newman2019}. Along the other directions, other information or assumptions are needed, such as covariance shared between populations (Supplementary Note~\Sref{sec:ident}). These directions show which parts of a prediction made without pairs rest on population averages and which rest on further assumptions, such as those of alignment methods\cite{welch2019,cao2022,demetci2022}.

\subsection*{Limitations}
\textit{Scope.} Our estimator targets linear dependence within populations, so nonlinear relations are outside its scope. Because unpaired data do not determine the dependence (Proposition~\Sref{prop:pairing}), the saving rests on the dependence concentrating along the leading unpaired axes, as it did here. The exact covariance relation underlying the cell-type maps assumes a shared channel (Proposition~\Sref{prop:map}); the implemented maps regularize this relation and may retain bias (Supplementary Note~\Sref{sec:semi}). The risk bound for the estimated shrinkage factors (Proposition~\Sref{prop:oracle}) is absolute and can be loose for small blocks or heavy-tailed cross-products. It assumes independent cross-products of paired cells, which contrasts break; after contrasts, the noise estimate is calibrated only empirically (Proposition~\Sref{prop:contrast}).

\textit{Centring and transfer.} The validation test examined a workflow revised after the bone-marrow test had been scored. As prespecified, the bone-marrow test centred each population's paired cells on their own average, which biased the noise estimate downwards and, with few paired cells, gave estimates worse than predicting zero correlation (Fig.~\ref{fig:atlas}a; Supplementary Fig.~\Sref{sfig:transfer}a). With nuclear RNA as its unpaired cells, neither the estimator nor paired-only estimation reached the target, so its second hypothesis was not met. Using nuclear RNA for both the axes and the maps, the estimator recovered only \MpAssayMapped\% of the dependence (exploratory). Centring on the averages of other sites or animals made the estimator worse than predicting zero correlation at every budget (Supplementary Fig.~\Sref{sfig:transfer}); the prespecified fly analysis, which centred on the other animals' averages, recovered \FxPoolHundred\% with 100 paired neurons. With unweighted maps, the estimator did not reach the target with either skewed atlas within 800 paired cells (Fig.~\ref{fig:reference}b), and with 25 paired cells reweighted maps were worse than predicting zero correlation.

\textit{Validation test.} The validation covers one blood study and five broad cell types, and its readout and the comparison with other methods were specified after it had been scored. The test assumes common cell-type annotations supplied by the original studies; it does not evaluate labels inferred separately from each retained modality. As in the other tests, its intervals are conditional on the fitted predictions, reference data, panels, paired samples and fold design (Supplementary Note~\Sref{sec:review} examines single draws and reference resampling). Using the atlas required \VlAtlasOverOwn{} times as many paired cells as using the paired samples' own unpaired cells (95\% CI, \VlAtlasOverOwnLo--\VlAtlasOverOwnHi), so the prespecified criterion for calling the atlas no worse was not met. Given the atlas, the estimator's advantage over reference regression was small and unstable (ratio \AltRatioRidge; 95\% CI, \AltRatioRidgeLo--\AltRatioRidgeHi; Supplementary Table~\Sref{tab:onedraw}). Its advantage over Champollion was not significant. A larger share of reference regression's top-ranked pairs were reproducible associations, and the estimator predicted more directions correctly than reference regression in only \RvDonorBetterReg{} of \RvDonorN{} donors. The readout's associations are reproducible between halves of a held-out sample, not across cohorts, and most lay between the finer cell states annotated by the study's authors, within which \RvFineReps{} rather than \RvPrimaryReps{} qualified. Against the dependence within finer states, predictions for broad cell types kept an advantage of \RvFineOneMinusPo{} percentage points over paired-only estimation in direction but recovered \RvFineRfAtlas\%, against \RvFineRfPo\% for paired-only estimation and \RvFineRfReg\% for reference regression. For the weak association between cytotoxic genes and CD56 in NK cells, the estimate was worse than predicting zero correlation.

\textit{Benchmark.} The \SyGeo-fold saving is a descriptive geometric mean at each data set's own target, computed after all tests had been scored, not the saving to expect in a new study. The benchmark's target and its rule for setting aside data sets without recovered dependence were amended after some data sets had been scored; at the registered target, the geometric-mean saving was \GenFrozenGOne{} (95\% CI, \GenFrozenGOneLo--\GenFrozenGOneHi). Six of the nine savings are lower bounds, and at a fixed 20\% recovery the saving could be measured in only \RvFixedTwentyN{} of \RvFixedSets{} data sets (Supplementary Table~\Sref{tab:fixed}). The external test's registered scoring code scored a one-sided variant of the estimator, not the planned two-sided one; both met all three hypotheses (Supplementary Note~\Sref{sec:semi}).

\textit{The law.} The law approximates the axes' saving, not the estimator's saving over paired-only estimation. An earlier form, fixed before the benchmark to predict the latter, did not predict the ranking of the benchmark's savings at the registered target (Spearman $\rho=\GenLawFrozenRho$) or the amended one ($\rho=\GenLawRho$; \textit{P}~=~\GenLawP), and the law was developed on those data sets afterwards. That unpaired noise shares approach the true ones as the dependence weakens is shown only for jointly Gaussian measurements (Proposition~\Sref{prop:reveal}). The prospective test rests on \PredNEstimable{} of \PredNProblems{} problems, from \PredNEstimableData{} data sets of two studies; in them the law overpredicted the axes' saving by about a fifth at 50\% recovery and by three quarters at 30\%. It did not hold for sparse chromatin counts. Treating the four data sets as exchangeable gives a rank-test \textit{P} value of \RvLawHalfMeansP{}, the smallest attainable; the two study families provide limited evidence for prediction across studies. A pilot of up to 200 paired cells predicted the axes' saving with a median error of \PredPilotErr\% (95\% CI, \PredPilotErrLo--\PredPilotErrHi\%). Because the law assumes the best factor in every block, the budget it set from all available paired cells was a median \RvPilotHalfResRatio{} times what block shrinkage needed and reached 50\% recovery in only \RvPilotHalfResReached\% of replicates, and pilots of 100 or 200 cells usually exceeded the budget that the law set.

\textit{Learning without pairs.} In blood, within cell types, where donors hardly differ, the channel's predictions were worse than predicting zero correlation (relative error \XsHaoWErrPfMeans; 95\% CI, \XsHaoWErrPfMeansLo--\XsHaoWErrPfMeansHi; Fig.~\ref{fig:pairfree}a), so the prespecified hypothesis on recovery within cell types was not met. In colon, a channel learned in blood was barely better than predicting zero correlation (relative error \XsColonErrPfMeans), achieved \XsColonShare\% (95\% CI, \XsColonShareLo{} to \XsColonShareHi\%) of the improvement that paired cells gave and was worse than predicting zero correlation within cell types (relative error \XsColonWErrPfMeans), so neither colon hypothesis on recovery was met. In the melanoma screen, learning without pairs recovered \PtRfShare\% (95\% CI, \PtRfShareLo--\PtRfShareHi\%) of what the same model recovered from paired cells, short of the prespecified half. Removing the five training perturbations ranked highest by exposure, the planned ranking (how far a perturbation moves RNA along directions in which cells already vary), left the interferon recovery at \DsFiveIfnExp\%, against \DsFiveIfnRand\% after random removals, so the planned criterion of a larger loss was not met. In the monocyte screen, learning without pairs did not replicate under the same prespecified rules: its predictions were worse than predicting zero correlation (\RpPrimary\%; 95\% CI, \RpPrimaryLo{} to \RpPrimaryHi\%) and than both controls (Fig.~\ref{fig:pairfree}d), and none of its five hypotheses was met.

\subsection*{Conclusions}
Tens to a few hundred paired cells, centred within their populations, became far more informative when combined with unpaired profiles of the same assay, from the study itself or from an atlas reweighted to the paired cells' cell-type proportions. Where the dependence lies along the unpaired axes indicates how much they save, and without paired cells, population averages identify the dependence only along the directions in which populations differ.

\section{Methods}

\subsection{Estimator}
In development and the benchmark, the training cells of each population were assigned pseudo-randomly to a reservoir (25\%), from which paired cells were drawn, and an unpaired pool (75\%). Half of the pool's cells contributed only RNA and half only the second modality. Paired cells were centred on the unpaired averages of their populations and scaled by the pooled within-population standard deviations in development, the external test, the benchmark and the law's test; centred on their own populations' averages in the prespecified bone-marrow test; and centred by contrasts (below) in the validation test and the reanalyses of the bone-marrow and fly tests. The paired cells' average cross-product $\hat\theta$, a genes-by-proteins matrix, was rotated onto the axes of the unpaired data, $A=U^\top\hat\theta V$. The columns of $U$ are the eigenvectors of the pool's latent RNA correlation (below), and those of $V$ the eigenvectors of its protein correlation. Blocks were formed by RNA axes 1--5, 6--20, 21--60 and the rest, crossed with protein axes 1--3 and the rest. Each block $A_b$ was multiplied by $(1-\hat t_b/\|A_b\|^2)_+$, where $\hat t_b$, the block's noise, estimates the trace of the covariance of the block's average from the per-cell products. This estimate is unbiased when the paired cells are independent and their scaling is fixed. The shrunk matrix was rotated back, $\tilde\theta=U\tilde AV^\top$. With $B$ paired cells, the best fixed multiplier for block $b$ leaves an expected squared error of $r_b^2\tau_b/(Br_b^2+\tau_b)$ (Proposition~\Sref{prop:oracle}). Here $r_b^2$ is the sum of squares of the block's true dependence and $\tau_b$ the total variance of a single cell's cross-products in the block.

A cell-type map is built for each condition $c$, a cell type or, in the perturbation screens, a culture condition (the Supplementary Notes call it a condition map). It is $G_c=\Sigma_cU_k\Lambda_k^{-1}U_k^\top+(I-U_kU_k^\top)$, where $U_k$ and $\Lambda_k$ hold the leading $k=120$ eigenvectors and eigenvalues of the pooled latent RNA correlation. Here $\Sigma_c$ is the condition's latent RNA correlation rescaled by the ratios $D_{x,c}$ of its RNA standard deviations to the pooled ones. The map and these ratios were shrunk towards the identity by the split-half reliability $s_c$ of the difference between the condition's covariance and the pooled one, $\hat G_c=I+s_c(G_c-I)$, and the condition's estimate was $\hat D_{x,c}^{-1}\hat G_c\tilde\theta\hat D_{y,c}^{-1}$. Reweighted maps replace the pooled latent RNA correlation $S$ by $S+\sum_c(w_c-w^0_c)\Sigma_c$, rescaled to standard deviations pooled with the same weights. Here $w^0_c$ is the reference's share of condition $c$ and $w_c$ the paired draw's share of contrasts of condition $c$.

Contrasts replace a population's $n\ge2$ drawn paired cells, with RNA or protein vectors $u_1,\dots,u_n$, by the $n-1$ orthonormal Helmert contrasts $h_j=(u_1+\dots+u_j-j\,u_{j+1})/\sqrt{j(j+1)}$, $j=1,\dots,n-1$; a single drawn cell contributes nothing. Each contrast has the population's covariance. Together, the contrasts give the same sum of cross-products as the cells centred on their own average, but unlike those cells they are uncorrelated with one another (and independent for Gaussian cells). Contrasts were scaled by their pooled standard deviations. In independent resamples of paired cells, the noise estimate after contrasts was slightly conservative, at \RvCalBmMin{} to \RvCalBmMax{} times the variance it estimates, as its expectation after contrasts implies (Proposition~\Sref{prop:contrast}; Supplementary Tables~\Sref{tab:noisecal} and~\Sref{tab:noiseblocks}).

\subsection{Latent RNA correlation}
Each RNA count was split into two halves by binomial thinning with probability 50\% (count splitting\cite{batson2019,neufeld2024count}), and both halves were log-normalized in the same way as the full data. A gene's reliability is the correlation across cells between its two halves, converted to full depth by the Spearman--Brown formula\cite{spearman1910,brown1910}; one minus the reliability estimates the fraction of its variance due to counting noise. The latent correlation, the RNA correlation with counting noise removed, was the measured RNA correlation with off-diagonal entries multiplied by 0.9, each diagonal entry reduced by the gene's noise fraction\cite{spearman1904} and eigenvalues floored at 0.001 times the mean diagonal (Proposition~\Sref{prop:noise}; Supplementary Note~\Sref{sec:noise}).

\subsection{Evaluation}
The cells of each held-out group $g$, such as one cell type in one held-out sample, were split into two halves. With $T_g$ the matrix of correlations between genes and proteins in the group's held-out cells and $T_{A,g}$ and $T_{B,g}$ its estimates from the two halves, the fraction of the dependence recovered by predictions $C_g$ was $\sum_g(2\langle C_g,T_g\rangle-\|C_g\|^2)/\sum_g\langle T_{A,g},T_{B,g}\rangle$. To first order, the numerator estimates without bias the reduction in squared error relative to predicting zero correlation, and the denominator, because the halves are independent, the squared Frobenius norm of the true dependence.

All methods in a comparison received the same draws of paired cells. Each draw was scored on its own, and the scores were averaged over draws, so that a curve gives the accuracy a study with one draw can expect. Rerunning the prediction code of every test (for the benchmark, the first of three folds) confirmed that every curve is the mean of single-draw scores (Supplementary Note~\Sref{sec:review}). The paired cells a method needed to reach an accuracy target were read from the running maximum of its curve, by log-linear interpolation between the two budgets around the target. A curve that never reached the target was capped at the largest budget. When the comparator failed to reach the target and the estimator reached it, this gave a lower bound on the saving; other censored comparisons were unresolved, as specified for each test in the Supplementary Information. Paired-only estimation was represented at each budget by the best of four methods, each pooled over conditions or fitted per condition: James--Stein shrinkage in the original coordinates, the simple and flexible covariance shrinkage estimators SCOSE and FCOSE\cite{muandet2014,ramdas2015}, and cross-validated low-rank approximation. Where the paired cells were centred by contrasts, so were those of these methods.

\subsection{Test designs}
\textit{Development and external test.} The estimator was developed on the Perturb-CITE-seq melanoma screen\cite{frangieh2021} and the ECCITE-seq THP-1 screen\cite{papalexi2021}, as released by scPerturb\cite{peidli2024}. Dependence was scored within groups of cells sharing a knocked-out gene and a culture condition, for knocked-out genes held out from training, with 10 (melanoma) or 4 (THP-1) draws per budget. Comparator settings are in Supplementary Note~\Sref{sec:semi}. The external test used OverCITE-seq of primary T cells\cite{legut2022}. Every third overexpressed gene was held out with all its cells, the panel comprised 200 genes, including those encoding the proteins targeted by the 14 detected antibodies, and budgets were 25 to 400 paired cells with 20 draws. Its three hypotheses compared the paired cells that the estimator needed to reach 30\% recovery with those that paired-only estimation, SemiCCA and Champollion needed.

\textit{Benchmark.} Nine data sets were extracted into a common format (Supplementary Note~\Sref{sec:gen}). Four paired RNA with surface proteins, two in blood\cite{hao2021,stephenson2021} and one each in bone marrow\cite{luecken2021sandbox} and colon\cite{mennillo2024}. One paired RNA with chromatin accessibility in bone marrow\cite{luecken2021sandbox}, two were Patch-seq\cite{cadwell2016} data sets of mouse motor\cite{scala2021} and visual cortex\cite{gouwens2020,gala2021}, and two were the connectomes of a female\cite{bates2026} and a male\cite{berg2025} fly. In the connectomes, the features were synapse counts with groups of partner neurons on the brain and nerve-cord sides of the neck. Units were held out in three folds. Within each fold, 25\% of the training cells formed the reservoir. Features were chosen from training cells: 200 variable features of the first modality, and the surface proteins, 200 variable peaks, groups of partners or all electrophysiological features of the second. The methods of the development comparison other than Champollion were run unchanged, with five draws per budget from 25 to 800 paired cells, and the target was half of the best recovery of any method.

\textit{Bone-marrow test.} The NeurIPS 2021 bone-marrow CITE-seq and multiome data\cite{luecken2021sandbox} were analysed in eight cell types present in both assays. Each of the 12 CITE-seq batches, one donor at one site, was held out once, with paired cells drawn from 25\% of the cells of the other two batches of its site. Unpaired cells came from the rest of those two batches, from the CITE-seq batches of other sites whose donors did not appear at the held-out site, or, for RNA, from the multiome nuclei (RNA and chromatin accessibility measured in the same nuclei) of those sites and donors. Panels of 200 genes and the detected antibodies were chosen from the other sites' cells. As prespecified, paired cells were centred on their own populations' averages and rescaled by $\sqrt{n/(n-1)}$; the reanalysis centred them by contrasts. Budgets were 25 to 800 paired cells with five draws (Supplementary Note~\Sref{sec:deploy}).

\textit{Validation test.} The data were the CITE-seq samples of a study of healthy blood\cite{babdor2026}, profiled with 10x 5$'$ chemistry and 140 antibodies. We analysed its \VlNFolds{} samples of \VlNDonors{} donors in \VlNPools{} processing pools, in the five cell types that the atlas holds with enough cells per patient (B cells, CD4 and CD8 T cells, natural killer (NK) cells and CD14 monocytes). Each sample was held out once, with paired cells drawn from 25\% of the cells of the next two samples of its pool. Unpaired cells came from the atlas of Stephenson et al.\cite{stephenson2021} (600 cells per patient), from the other pools' samples whose donors did not appear in the held-out sample's pool, or from the paired samples' remaining cells. A pool holds at most one sample of a donor, and a donor's two samples lie in different pools, so the held-out donor contributed none of the cells used for estimation in its fold (Supplementary Table~\Sref{tab:allocation}). Panels were fixed from the atlas alone: 200 genes and the \VlNAntibodies{} antibodies shared by the two panels and detected in at least 1\% of atlas cells. Paired cells were centred by contrasts, and budgets were 25 to 800 paired cells with five draws (Supplementary Note~\Sref{sec:val}). All methods received the original studies' cell-type labels. The validation labels were transferred from a reference annotated using RNA and protein, and the atlas labels used RNA clusters supported by protein markers\cite{babdor2026,stephenson2021}.

In the readout, a gene--protein pair was a reproducible association in a held-out sample and cell type if its correlation had the same sign in both halves, with two-sided $P<0.01$ in each by Fisher's $z$; only samples and cell types with at least 20 cells per half were scored. For every draw, we computed three quantities. The first was the fraction of reproducible associations whose direction the estimate predicted correctly, with an estimate of zero counting as half correct. The second was the fraction of the 100 pairs with the largest absolute estimate in each sample and cell type that were reproducible associations with the predicted sign. The third was the recovered fraction over the \BoNCognate{} pairs of a panel gene and the panel protein it encodes. Paired-only estimation was represented, for each quantity, by the best of the eight paired-only fits (four methods, each pooled or per cell type). A relationship's value was the average, over the genes of its program, of their within-type correlation with its protein, and its recovery the recovered fraction of this value over held-out samples. In the comparison with other methods, SemiCCA, reference regression (ridge regression of protein on RNA) and Champollion received the estimator's paired contrasts, panels and atlas, with tuning from these alone. SemiCCA was fitted in three forms, pooled over cell types, pooled and then mapped to each cell type through the atlas's cell-type maps, or refitted in each cell type, and reference regression in the same three forms with the Gram matrix of the paired cells or of the atlas. The best form at each budget represented each method. Champollion was run on one draw of five.

\textit{Predicting the saving.} The law was tested on ten data sets (Supplementary Note~\Sref{sec:pred}). Five were CITE-seq data sets: of mouse lymphoid organs with 111 and 206 antibodies and of human blood and lymphoma from 10x Genomics\cite{totalvi2021}, and of human bone marrow\cite{stuart2019}. Two paired RNA with chromatin accessibility, in human fetal and adult mouse cortex\cite{trevino2021,chen2019snare}, one was human cortical Patch-seq\cite{lee2023gaba}, and two paired the optic-lobe and central-brain partners of fly visual neurons in two connectomes\cite{dorkenwald2024,bates2026}. A problem was a data set with a panel of 100, 200 or 400 features of the first modality chosen from unpaired cells. Its cells were split into scoring cells (25\%), a reservoir (30\%) and an unpaired pool (45\%), and the quantity estimated was the pooled within-population covariance in units of the pool's standard deviations. Block and single-block James--Stein shrinkage, and block shrinkage along random orthonormal axes, were fitted to 20 draws per budget ($25\times2^k$ cells, up to 12,800). The axes' saving at 50\% recovery was measurable when both curves crossed that target from below within the budgets tested. The law used unbiased raw estimates of each block's squared dependence and noise from all reservoir cells; squared-dependence estimates were truncated at zero before totals and normalized shares were formed. The pilot prediction used the same rule on the first $\min(200,0.2N)$ of the $N$ reservoir cells. In the pilot design, a random pilot of 50, 100 or 200 reservoir cells set the total budget for block shrinkage at a target of 30\%, 50\% or 70\%. Block shrinkage was then fitted to that many cells, pilot included (100 replicates per problem, target and pilot size).

\subsection{Learning without pairs}
We assumed that populations share a Gaussian channel,
\begin{equation}
 \begin{gathered}
 y\mid x\sim N(Wx+b,\Psi),\\
 \mu_{y,s}=W\mu_{x,s}+b,
 \end{gathered}
 \label{eq:channel}
\end{equation}
so that the populations' protein averages are linear in their RNA averages. Channels were estimated by weighted ridge regression\cite{hoerl1970} of population protein averages on population RNA averages, with numbers of cells as weights and each training population's RNA and protein moments taken from disjoint halves of its cells. A channel was transferred to a recipient through the Gaussian coupling of the recipient's RNA and protein covariances with interaction $B=W^\top\Psi^{-1}$,
\begin{equation}
 \begin{gathered}
 C=S_x^{1/2}U\operatorname{diag}\{g(\beta)\}V^\top S_y^{1/2},\\
 S_x^{1/2}BS_y^{1/2}=U\operatorname{diag}(\beta)V^\top,\\
 g(\beta)=\frac{(1+4\beta^2)^{1/2}-1}{2\beta},
 \end{gathered}
 \label{eq:closed}
\end{equation}
the entropic optimal-transport coupling of Gaussian measures\cite{janati2020,mallasto2022} (Proposition~\Sref{prop:gaussian}).

In blood, channels were learned from \XsPatients{} patients of Stephenson et al. (\XsTrainCells{} cells) and scored in the eight donors of Hao et al. (\XsHaoCells{} cells) and in 12 colon donors\cite{mennillo2024}. For each recipient donor, the relative error of a prediction was the Frobenius norm of its difference from the held-out correlation matrix divided by the norm of that matrix, so that predicting zero correlation gives 1; relative errors were averaged over donors. The share of the improvement achieved with paired cells was each donor's $(1-e)/(1-e_{\mathrm{paired}})$, averaged over donors, with $e$ and $e_{\mathrm{paired}}$ the relative errors without and with pairs. The permutation null paired the training patients' RNA and protein averages at random within site (Supplementary Note~\Sref{sec:cross}). In the melanoma screen, every fourth of the \PtTargets{} perturbed genes was held out, and channels were learned from \PtTrainPops{} populations of one knocked-out gene in one culture condition, centred within condition. The interferon response was the union of the Hallmark interferon-$\gamma$ and interferon-$\alpha$ response sets\cite{liberzon2015}. A perturbation's exposure measured how far it moves RNA along directions that vary within conditions, and its effect strength the size of its shift in average RNA. The monocyte screen, analysed with the same rules, served as a replication (Supplementary Note~\Sref{sec:rep}).

\subsection{Statistics}
Intervals are percentile intervals from bootstrap resamples of the independent units of each analysis\cite{efron1979}: held-out genes, batches (stratified by site) or units, donors, or data sets with their problems (2,000 resamples, except 1,000 in development and 10,000 in the blood test of learning without pairs). Repeated draws of paired cells were averaged, never counted as replicates. The interval of the summary across test data sets combines the $r$-th bootstrap resample of every data set. Permutation \textit{P} values were one plus the number of permutations with error at most the observed error, divided by one plus the number of permutations\cite{phipson2010}. A hypothesis of the validation readout was met only if its criterion held at every one of its budgets, an intersection--union test\cite{berger1982}. Its eight planned comparisons had Bonferroni-simultaneous intervals from 20,000 resamples of donors, and the \textit{P} values of its three hypotheses were adjusted by Holm's procedure\cite{holm1979}.

The prespecified tests, whose protocol, code and data were registered before any held-out result was computed, were the external test, the benchmark, the bone-marrow test, the validation test, the prospective test of the law and the tests of learning without pairs. The follow-up analyses were the readout of the validation test, its comparison with other methods and the analyses of Supplementary Note~\Sref{sec:review}. All other analyses are exploratory. Supplementary Table~\Sref{tab:registry} lists every registered hypothesis with its outcome, every amendment and deviation, and every exploratory analysis. Analyses were run in Python 3.11 with NumPy 2.0.2 (ref.~\citenum{harris2020}), SciPy 1.13.1 (ref.~\citenum{virtanen2020}) and scikit-learn 1.8.0 (ref.~\citenum{pedregosa2011}).

\subsection{Use of large language models}
OpenAI Codex and Anthropic Claude were used to develop and run analysis code, retrieve literature, edit the manuscript and prepare the \LaTeX{} source. The authors directed this work and are accountable for all content; the tools are not authors.

\section*{Data availability}
Source data are available from ArrayExpress (E-MTAB-10026), the Gene Expression Omnibus (GSE164378, GSE194122, GSE193736, GSE150599, GSE128639, GSE162170, GSE126074 and GSE314416) and Zenodo (record 7041849, the scPerturb releases of the Frangieh et al. and Papalexi et al. screens); MSigDB Hallmark gene sets are from gsea-msigdb.org. The colon data are the public count matrices in Figshare article 21919356, version 3 (ref.~\citenum{mennillo2024}). The 10x Genomics blood and MALT lymphoma CITE-seq data and the mouse lymphoid-organ data are from \url{github.com/YosefLab/totalVI_reproducibility}. Patch-seq data are from \url{github.com/berenslab/mini-atlas} (motor cortex), \url{github.com/AllenInstitute/coupledAE-patchseq} (visual cortex) and \url{github.com/AllenInstitute/human_patchseq_gaba} (human cortex). Fly connectome tables (BANC v888, male CNS v0.9, FAFB v783 and MANC v1.2.1) are from the public compilation of the fly connectome tutorial (\url{gs://lee-lab_brain-and-nerve-cord-fly-connectome/compiled_data}). The reviewer archive contains derived result records and prediction manifests; raw count matrices and large prediction arrays are not redistributed in it. Download, extraction, prediction and scoring scripts are supplied for full reconstruction from the public data.

\section*{Code availability}
All code, with the registered protocols, executable checks and the scripts that generate every number, figure and table, is provided to reviewers and will be deposited in a public repository with a persistent identifier on publication.

\ifdefined\ReaderVersion\else
  \setlength{\bibsep}{1.5pt}
\fi
\begingroup
\ifdefined\ReaderVersion\balance\else\singlespacing\fi
\renewcommand{\bibfont}{\normalfont\small}
\sloppy
\bibliography{references}
\endgroup

\end{document}
